% 本文件由 LaTeX 排版 Agent 根据有效 translation.json 自动生成。
% author=Tertullian; work_id=0312; title=驳斥马西昂

\chapter{\WorkTitle{驳斥马西昂}}

% structure: p0001 source=h1 target=omitted reason=page-title-already-rendered-by-parent

卷第一\newline{}卷第二\newline{}卷第三\newline{}卷第四\newline{}卷第五\par

\SourceRecord{Translated by Peter Holmes.；Ante-Nicene Fathers；Vol. 3.；Edited by Alexander Roberts, James Donaldson, and A. Cleveland Coxe.；Buffalo, NY: Christian Literature Publishing Co.,；1885.；<http://www.newadvent.org/fathers/0312.htm>.}

% structure: p0001 source=h1 target=section reason=page-title

\section{\Emph{驳斥马西昂，卷一}}

\Emph{其中描述了马西昂的神。表明他完全缺乏\OriginalTerm{真天主}{true God}的一切属性。}\par

% structure: p0003 source=h2 target=subsection reason=relative-heading-rank; base=h2

\subsection{\Emph{第一章 序言 新著作的理由。} \Place{本都} \Emph{将其粗野性格传给本地\OriginalTerm{异端者}{heretic}} \Person{马西昂} \Emph{。以简短抨击描述其\OriginalTerm{异端}{heresy}。}}

无论我们过去曾经撰写过什么反对\Person{马西昂}的作品，从此刻起都不再算数。我们正在从事一部新著作以取代旧作。我最初的论文，由于过于仓促写成，后来被一部更充实的论著所取代。后者在我完全出版之前，被一个当时还是兄弟、后来却成为背教者的人以欺骗方式丢失了。他碰巧抄录了其中一部分，错误百出，然后出版了它。因此，需要修订的作品；新版的机缘促使我对论著进行了相当大的增补。因此，我这部作品的当前文本——它是第三版，取代第二版，但从此以后应被视为第一版而非第三版——需要为这版著作本身加一篇序言，以免读者若偶然遇到其分散的各种版本时感到困惑。\par

所谓的\Place{尤克辛海}，其本性自相矛盾，其名称具有欺骗性。正如你不会因其位置而认为它好客，它也同样因附着于其野蛮特性的某种污名而与我们所处的更文明水域隔绝。最凶悍的民族居住在那里，如果生活是在马车上度过的，确实可以称之为\Emph{居住}。他们没有固定居所；他们的生活没有文明的萌芽；他们不加节制地放纵淫欲，并且大多赤身裸体。再者，当他们满足隐秘的欲望时，他们将箭袋挂在车轭上，以警告好奇和鲁莽的观察者。这样，他们毫不脸红地亵渎他们的战争武器。他们将父母尸体与羊一同切碎，并在宴会上吃掉。那些没有这样死去以成为他人食物的人，被认为死得可诅咒。他们的女性并未因性别而变得柔和谦逊。她们露出胸部，将战斧挂在胸前，并且更喜欢战争而非婚姻。她们的气候也同样粗野。白天从不晴朗，太阳从不愉快；天空始终多云；全年都是冬季；唯一的风是愤怒的北风。水只有通过火才能融化；河流因冰而不再流淌；群山覆盖着层层积雪。万物僵冷，一切因寒冷而僵硬。那里没有任何生命的热力，除了那种凶猛，正是这种凶猛为戏剧提供了陶里人的献祭、科尔基斯人的爱情和高加索的折磨的故事。然而，在\Place{本都}，没有什么比\Person{马西昂}出生在那里更野蛮和悲哀的了，他比任何\Place{斯基泰人}更肮脏，比\Place{萨尔马特人}的马车生活更漂泊，比\Place{马萨格泰人}更不人道，比\Place{阿马宗人}更鲁莽，比（本都的）云更黑暗，比它的冬天更寒冷，比它的冰更易碎，比\Place{伊斯特河}更阴险，比\Place{高加索}更崎岖。不仅如此，真正的普罗米修斯，\OriginalTerm{全能天主}{Almighty God}，被\Person{马西昂}的亵渎摧残。\Person{马西昂}甚至比那个野蛮地区的野兽更残忍。因为有什么河狸曾比他更阉割——他废除了婚姻之纽带？有什么本都老鼠曾有这样的啃咬能力，像他那样将福音啃成碎片？真的，啊，尤克辛海，你产生了一个对哲学家来说比对基督徒更可信的怪物。因为昔尼克派的\Person{狄奥格内斯}曾手提灯笼在正午找一个人；而\Person{马西昂}却熄灭了他的\OriginalTerm{信德}{faith}之光，因此失去了他已找到的天主。他的门徒不会否认他最初的信德与我们相同；他自己的一封信证明了这一点；因此，将来\OriginalTerm{异端者}{heretic}可以从他的案例被指定为那样一种人：他弃绝了先前的东西，后来选择了过去不存在的东西。因为\Emph{正如}过去和起初传下来的东西将被视为真理，\Emph{同样}后来引入的东西将被视为\OriginalTerm{异端}{heresy}。但另一篇简短的论著将坚持这个反对异端者的立场，他们甚至无需考虑其教义就应被驳斥，理由是他们的观点具有新颖性而被视为异端。现在，只要允许任何争议，我将暂时（免得我们简洁的新颖性原则被随时用来帮助我们而被归咎于缺乏自信）开始陈述我们对手的\OriginalTerm{信仰规则}{rule of belief}，以免没有人知道我们的主要争论点是什么。\par

% structure: p0006 source=h2 target=subsection reason=relative-heading-rank; base=h2

\subsection{\Emph{第二章。} \Person{马西昂} \Emph{在} \Person{克尔冬} \Emph{的协助下教导两位神祇；他如何构建了关于恶神与善神的\OriginalTerm{异端}{heresy}。}}

那来自本都的异端者引入了两位神，如同他自己船难中的孪生碰撞岩：一位是无法否认的，即我们的造物主；另一位是他永远无法证明的，即他自己的\Quote{神}。这个不幸者从其谬见的最初概念得自于我们的主的一句话，那句话是针对人而非神的，其中祂以好树和坏树为例作出安排：\BibleQuote{好树不结坏果子，坏树也不结好果子。}意思是说，诚实的心和正直的信德不能产生恶行，正如邪恶的性情不能产生善行。如今（如同现今许多人，特别是那些有异端倾向的人一样），当他病态地沉思恶的起源问题时，他的感知因研究的异常而变得迟钝；当他发现造物主宣告：\BibleQuote{是我创造了灾祸，}\BibleRef{依 45:7}——因为他已经从其他论据中得出结论（这些论据令每个扭曲的心满意足），即天主是恶的创造者——于是他将造物主比作结坏果子的坏树，即道德上的恶，然后假定应当有另一位神，如同好树结好果子一样。因此，当他发现基督有一种不同的性情，仿佛是一种单纯而纯粹的仁慈，与造物主不同，他便轻易地论证说，在基督身上启示了一个新而陌生的神性；然后用一点酵母将整个信仰的面团发酵，以他自己异端的酸味调味。\par

此外，他还有一个名叫\Person{Cerdon}的怂恿者来支持这种亵渎——这使得他们更加轻易地以为自己最清楚地看见了两位神，尽管他们是瞎眼的；因为事实上，他们没有以健全的信德看见那一位神。对于视力有病的人，即使一盏灯也看似许多盏。因此，他的其中一位神（他不得不承认的那一位），他通过将恶归咎于祂的属性而加以毁损；另一位神（他费力构想的那一位），他则建立在善的原则上。他以什么条目来安排这些性质，我们通过自己的反驳来表明。\par

% structure: p0009 source=h2 target=subsection reason=relative-heading-rank; base=h2

\subsection{第三章 天主的合一性。祂是至高存在，不可能有第二个至高者。}

主要的，事实上全部的争执在于\Emph{数目}这一点：是否允许两位神，或出于诗的许可（若必须如此），或出于绘画的想象，或通过第三种方式——我们现在必须加上——出于异端的乖谬。但基督信仰的真理已明确宣告了这一原则：\Quote{天主若不是一个，就不是天主；}因为我们更恰当地相信，那不像其当有的样子的东西是不存在的。然而，为使你知道天主是唯一的，请问天主是什么，你会发现祂不是别的，正是唯一的。就人能够对天主下定义而言，我提出一个定义，所有人的良心也会承认——即天主是伟大的至高者，永恒存在，无生、无造、无始、无终。因为这样的状态必须归于那使天主成为伟大至高者的永恒性，因为正是为此目的，天主才拥有这一属性；并以此类推其他性质：所以天主在形式、理性、权能和力量上是伟大的至高者。既然所有人都同意这一点（因为没有人会否认天主在某种意义上乃是伟大的至高者，除非那人能够说出相反的意见，即天主是某种较低的存在，以便通过剥夺祂的一个属性来否认天主），那么伟大的至高者本身的状态必然是什么？必然是没有什么与祂同等，即没有另一位伟大的至高者；因为，如果有，祂就会有一个同等者；如果祂有同等者，祂就不再是伟大的至高者，因为那使伟大的至高者成立的条件和（可以这样说）法则——即不允许任何事物与伟大的至高者相等——就被推翻了。因此，那伟大的至高者必须是一个\Emph{独一无二者}，因没有同等者，从而不停止为伟大的至高者。所以，祂只能以祂存在的条件来存在，即祂的绝对独一无二性。既然天主是伟大的至高者，我们\Emph{基督信仰}的真理已正确地宣告：\Quote{天主若不是一个，就不是天主。}这并非我们怀疑祂是天主，而说：祂若不是唯一的，就不是；而是因为我们定义那我们所深信的存在之所是，即没有它天主就不是天主，也就是说，伟大的至高者。但伟大的至高者必然是独一无二的。因此，这独一无二者将是天主——并非以其他方式为天主，而是作为伟大的至高者；也并非以其他方式为伟大的至高者，而是作为没有同等者；也并非以其他方式为没有同等者，而是作为独一无二者。那么，无论你引入什么别的神，你至少无法以任何其他方式来维持其神性，除非也将神性的属性——永恒和至高无上——归于他。因此，两位伟大的至高者如何能共存，既然至高存在的属性是没有同等者——这属性只属于独一者，绝不可能存在于两者中？\par

% structure: p0011 source=h2 target=subsection reason=relative-heading-rank; base=h2

\subsection{第四章 为神圣合一性辩护，驳斥异议。人间权力与天主主权之间没有类比。否则异议不成立，因为为何只停在两位神？}

但是有人或会争辩说，可能存在两个伟大的至高者，各自在自己的领域内独立而分离；甚至可能援引世上的王国为例，它们数量虽多，但在各自的区域内却是至高无上的。这样的人将会假定人的境况总是与神的境况可比。如果这种推理方式尚可容忍，那有什么能阻止我们引进——我不说第三位神或第四位——而是与地上君王一样多的神呢？现在所讨论的是天主，其主要特性是祂不容许与自己相比。因此，自然本身，如果不是\Person{依撒意亚}，或者不如说是天主藉\Person{依撒意亚}说话，将会恳切地问：\BibleQuote{你们拿谁来与我相比？}人的境况或许可与神的境况相比，但不能与天主相比。天主是一回事，属于天主的是另一回事。再者，你援引君王作为伟大至高者的例子，务要谨慎，以便能恰当使用。因为君王虽然在祂的宝座上是仅次于天主的至高者，但仍低于天主；当他与天主相比时，他将会被移离那伟大的至高地位，而归于天主。既然如此，你怎能用一个在比较的一切目的上都失败的例子来与天主比较呢？为什么，当君王中的至高权力显然不能是多重的，而只是独特且单一的，却唯独在祂——万王之王，并因祂权能的极大超越和一切其他等秩对祂的臣服，可以说是统治的顶峰——身上破例呢？但即使是在另一种统治形式中——他们个别地联合掌权——若将他们各自的王权特权（姑且这样说）在各方面相互比较，以显明其中谁在权能的基本特征和力量上更优越，那么至高威严必然归于唯一的一位——所有其他都因比较的结果逐渐被排除和移除出最高权威。因此，虽然分散在多人手中时，最高权威看似多重，但在其自身的权能、性质和状况中，它是独特的。所以，如果两位神被比较，如同两位君王和两位最高权威，那么按照比较的意义，权威的集中必然要归于二者之一；因为从祂的优越性已清楚看出祂是至高者，而对手已被击败，并被证明不是更大的，无论其有多大。现在，由于对手的失败，另一位在权能上是独特的，仿佛在祂独特的优越中拥有某种孤单。因此，我们在此点得出的必然结论是：要么我们必须否认天主是伟大的至高者——这是任何智者都不会允许自己做的；要么就说天主没有其他任何人与祂共享权能。\par

% structure: p0013 source=h2 target=subsection reason=relative-heading-rank; base=h2

\subsection{第五章 双重原则瓦解；无论数量多少，多位神祇更一致。从马西昂的二元论产生的荒谬和对虔敬的损害。}

但\Person{马西昂}是根据什么原则将他的至高权能限定为\Emph{两位}呢？我要先问：如果有两位，为何不是更多？因为如果\Emph{数量}能与神性的本体相容，那么数量越多越好。\Person{华伦提努}更一致也更慷慨；他既想象出两位神祇，即\OriginalTerm{深渊}{\GreekText{Βυθός}}和\OriginalTerm{沉默}{\GreekText{Σιγή}}，便倾泻出一群神圣本质，不少于三十个\OriginalTerm{永世}{\GreekText{Αἰών}}的后裔，如同\Person{埃涅阿斯}的母猪。然而，任何拒绝承认多位至高存在的原则，也必定拒绝两位，因为在\Quote{一}之后最低的数字中已有复数。在\Quote{一}之后，\Emph{数量}开始了。同样，能够容纳两位的原则也能够容纳更多。在\Quote{二}之后，当\Quote{一}被超越时，\Emph{众多}就开始了。简而言之，我们觉得理性本身明令禁止相信多于一位神，因为同一规则规定只有一位天主而非两位，这一规则宣称天主必须是那个作为伟大至高者、无与伦比的存在；而无与伦比的存在也必然是独特的。但是，进一步说，假设两位至高存在、两种同等的权能有什么用途或好处呢？当两位相等者彼此没有差异时，数字上的区别又有什么意义？因为同一个事物在两者中就是同一个。即使有几个相等者，它们也只是同一个，因为作为相等者，它们彼此没有区别。因此，如果两个存在者彼此没有区别，既然两者都假设为至高，两者都是神，那么其中任何一个都不比另一个更卓越；因此，没有优越性，它们数字上的区别就没有理由。此外，神性中的数字应当与最高的理性一致，否则对祂的崇拜就会受到质疑。试想，如果我看到两位神在我面前（两者都是至高存在，彼此相等），我若崇拜两者，我是在做什么？我会非常担心我崇拜的丰富会被视为迷信而非虔敬。因为两者如此平等，且每一者都包含在对方中，我只需崇拜一位就能恰当地侍奉两者；正因如此，我恰恰证明了它们的平等和统一，只要我互相在对方中崇拜它们，因为这两者在一位中都\Emph{临在}于我。如果我崇拜两者中的一位，我同样会意识到似乎是在轻视那数字区别的无用性，这区别是多余的，因为它不表示任何差异；换句话说，我认为更安全的做法是根本不崇拜这两位神中的任何一位，而不是怀着某种良心上的顾虑崇拜其中一位，或者徒劳地崇拜两者。\par

% structure: p0015 source=h2 target=subsection reason=relative-heading-rank; base=h2

\subsection{第六章 \Person{马西昂}违背自己的理论。他假装他的众神是平等的，但实际上使它们不同。然后，承认它们的神性，又否认这种差异。}

到目前为止，我们的讨论似乎暗示马西昂使他的两位神相等。因为当我们主张天主应被相信为唯一至高的天主，排除任何与祂相等的可能性时，我们是在假设两位相等之神的情况下讨论这些主题；但尽管如此，通过教导根据至高者的律法不可能存在相等者，我们已充分确认两位相等者不可能存在。然而，我们十分清楚马西昂使他的神不相等：一位审判、严厉、在战争中强大；另一位温和、平静、单纯善美而卓越。让我们同样仔细考虑问题的这一方面，即（神性中的）区别是否无论如何能容纳两位，既然其中的相等未能做到。在此，关于至高者的同样规则将保护我们，因为它决定了神性的全部状况。现在，挑战并在某种意义上阻止我们对手的意思——他不否认造物主是天主——我极公平地向他提出反对：他没有余地在他的神中放入任何区别，因为一旦承认他们是同等的，就不能再说他们不同；确实不是因为人类在同一名称下可能非常不同，而是因为神圣存在不能被说或被认为天主，除非是至高的天主。因此，由于他必须承认他所不否认的天主是至高的天主，他不能允许对至高者作出这样的贬低以至于使他服从另一位至高者。因为如果他服从于任何一位，他就不再是至高者。此外，天主不会失去祂神性中的任何属性——例如，祂的至高性。因为按照这种说法，甚至马西昂更强大的神的至高性也会受到威胁，如果它可能在造物主那里贬值。因此，当两位神被宣称是两位至高的天主时，必然得出他们谁也不比谁更大或更小，谁也不比谁更高或更低。如果你称较低者为神而否认他是神，你就否认了这个较低者的至高性。但当你承认两位神都是神圣时，你就承认他们都是至高的。你不能从他们中任何一位夺走什么，也不能添加什么。通过允许他们的神性，你否定了他们的区别。\par

% structure: p0017 source=h2 target=subsection reason=relative-heading-rank; base=h2

\subsection{第七章 除了天主之外，其他存在在圣经中也被称为神。这一反对意见轻浮，因为这不是名称问题。神圣本质是争议所在。异端，就其一般术语而言，至此已论。}

但你会用天主之名为这一论点辩解，声称这个名称是模糊的，也适用于其他存在；正如经上所载：\BibleQuote{天主立于大能者的会中，在众神中施行审判。}又说：\BibleQuote{我曾说：你们是神。}因此，正如至高的属性对这些存在不合适——虽然他们被称为神——对造物主也不合适。这是一个愚蠢的反对意见；我的回答是，提出者未能考虑到，同样强烈的反对意见也可以针对马西昂的（更优越的）神：他也被称为神，但并不因此被证明是神圣的，正如天使和人——造物主的\Emph{手工}——也不是一样。如果名称的相同性提供了条件相等的推定，那么有多少下贱的仆人会傲慢地摆出君王的名字——你们的亚历山大、凯撒和庞培！然而这一事实并不减损王族人物的真正属性。不仅如此，外邦人的偶像也被称为神。然而没有一个是神圣的，因为被称为神。因此，我并非为了天主之名，为了它的声音或书写形式，而在造物主中要求至高性，而是为了名字所属的本质；当我发现只有那个本质是非受生、非受造、唯一永恒、万物的创造者时，我不是把至高性归于其名，而是归于其实，不是归于其称谓，而是归于其状况。因此，因为我所归与之的本质被称为神，你便以为我归之于名，因为我必须用一个名称来表达那个本质，而那位被称为天主的存在正是由这个本质构成的，并且因其本质而非其名而被视为至高的天主。简言之，马西昂本人将这一特性归于他的神时，是归于本性，而非词汇。因此，我们基于本质而非基于名而归予天主的至高性，应该平等地存在于由该实体构成的两位存在中——该实体正是天主的名称所给的；因为但凡他们被称为神（即至高的存在，当然基于他们非受生、永恒、因此伟大而至高的本质），那么作为至高的天主这一属性就不能被认为在一个至高天主中比在另一个中更少或更差。如果至高的天主的幸福、崇高和圆满在马西昂的神那里成立，那么在我们的天主那里也一样成立；如果不在我们的天主那里，那么也\Emph{不}会在马西昂的神那里成立。因此，两位至高的天主将既不会相等，也不会不相等：不会相等，因为我们刚才阐述的原理——至高者不容许与自身比较——禁止相等；不会不相等，因为另一个关于至高者的原则——祂不能有任何减少——阻止不相等。所以，马西昂，你被自己的本都海潮所困。真理的波浪从四面淹没你。你既不能设立相等的神，也不能设立不相等的神。因为就\Emph{数量}的问题而言，根本不存在两位。虽然两位神的问题全部悬而未决，我们已将讨论限制在一定范围内，现在我们将就各自的特殊性进行辩论。\par

% structure: p0019 source=h2 target=subsection reason=relative-heading-rank; base=h2

\subsection{第八章 具体观点。马西昂的神的新颖性对其主张致命。天主是永恒的，祂绝不可能在任何方面是新。}

首先，马西翁派何其傲慢地建立他们那愚蠢的系统，提出一位新神，仿佛我们对旧神感到羞愧！正如学童以新鞋为傲，但他们的旧老师却打掉了他们趾高气扬的虚荣。如今，当我听说一位新神，在旧世界、旧时代、旧神之下无人知晓、无人听闻；这位神（在过去如此漫长的世纪中\Emph{被视为}无，而人们对其无知本身就是古老）是由一位特定的\Quote{耶稣基督}启示的，别无他人；基督启示了祂，\Emph{他们说}——基督本身是新的，\Emph{按他们的说法，甚至}在古代名号中也是新的——我对其这种自负感到庆幸。因为借助它，我就能立刻证明他们关于新神祇之主张的异端性质。这种新奇之事，竟如同异教徒通过某种新的、再三更新的、各神化独有的称号，制造出诸神一般。哪有什么新神，除了假神？连萨图尔努斯凭其古老的声名也无法证明是神，因为那也曾是某个时候首次赋予他神性时的新的托词。相反，活生生的、完美的天主既不源于新奇也不源于古老，而源于其自身的真实本性。永恒没有时间。它本身就是一切时间。它行动，因而不能受苦。它不能被生，因而没有年龄。天主若是旧的，就失去了未来的永恒；若是新的，就失去了过去的永恒。新奇见证开端，古老威胁终结。再者，天主既独立于开端和终结，也独立于时间，时间只是开端和终结的仲裁者和度量者。\par

% structure: p0021 source=h2 target=subsection reason=relative-heading-rank; base=h2

\subsection{第九章 马西翁的诺斯替式虚妄徒然，因为真天主既非未知也非不确定。他所承认是天主的创造主，独自提供了判断真天主的归纳。}

如今我十分清楚他们凭借何种感知能力夸耀他们的新神；即是他们的认知。然而，正是这种对新事物的发现——对常人如此惊人——以及新事物所固有的自然满足，我想要驳斥，并进而挑战这未知之神的证据。因为他们凭其认知呈现给我们为新神的那一位，他们证明在那认知之前是未知的。让我们严格保持在论证的范围和尺度内。请说服我可能有一位未知之神。我发现，无疑，祭台被大量献给未知之神；但那不过是雅典的偶像崇拜。还有献给不确定之神；但那也只是罗马人的迷信。此外，不确定之神不为人熟知，因为关于它们没有确定性；由于这种不确定性，它们因此是未知的。那么，我们应该为马西翁的神刻上这两个称号中的哪一个？我想两者兼有，作为一个仍然是\Emph{不确定}、并且从前是\Emph{未知}的存在。因为既然造物主是已知的天主，便使他成为未知；同样，作为确定的天主，便使他成为不确定。但我不会如此偏离正轨，以至于说：如果天主是未知且隐藏的，祂被笼罩在如此黑暗的区域中，这区域本身必然是新的且未知的，并且如今也同样不确定——确实是一块巨大的区域，无疑比它所隐藏的天主更大。但我将简要陈述我的主题，然后充分探讨，断言天主既\Emph{不可能}、也\Emph{不应该}曾是未知的。不可能，因为祂的伟大；不应该，因为祂的美善，尤其是祂（按马西翁的假设）在这两方面的属性比我们的造物主更卓越。然而，既然我观察到，在某些点上，每一个新的、迄今未知之神的证据，为了检验，应当与造物主的形象相比较，那么我的职责首先是表明，我正是按照一个确定的计划采用这一方法，以便我能更有信心地用于支持我的论证。在所有其他考虑之前，请问：你们承认造物主是天主，并且凭你们的认知承认祂先在，为何不知道那另一位神应当由你们用同样的探究方法去考察，这方法曾教导你们如何首先找到一位神？每一在先的事物都为后者提供了规则。在目前的问题中，提出了两位神：未知者和已知者。关于已知者，没有问题。祂显然存在，否则就不会被认识。争论是关于未知之神。他可能不存在；因为，如果他存在，他早已被认识了。凡是长久未知之物，便是被质疑的对象；只要它仍然如此可疑，便是不确定的；而它一直处于这种不确定状态时，可能根本不存在。你们有一位神，就他为人所知而言，是确定的；就未知而言，是不确定的。既然如此，你们认为将不确定之事提交给确定之事的规则、形式和标准来证明，是合理的吗？现在，如果对我们目前这个本身充满不确定性的话题，还应用来自不确定性的论证，我们将陷入由处理这些不确定论证而产生的一系列问题中，这些问题因其不确定性而危及信仰，我们将漂入那些宗徒所不喜爱的难解问题中。再者，如果在那些条件不同的事物中，他们无疑会以规则中确定、无疑、清晰的一面来预判不确定、可疑、复杂之处，那么很可能这些点将不会被提交给确定性的标准来决定，因为它们本质条件的差异使它们在所有其他方面免于适用这样的标准。因此，既然我们命题的主题是两位神，那么他们的本质条件必须相同。因为，就他们的神性而言，两者都是无生、非受造、永恒的。这将是他们的本质条件。所有其他点，马西翁本人似乎轻视了，因为他将它们放在不同的类别中。它们在处理顺序上是次要的；实际上，它们不必被引入讨论，因为在本质条件上没有争议。现在我们的争论没有涉及这一点，因为两者都是神。因此，那些条件共同性明显的事物，当被放在那共同条件的基础上检验时，尽管它们不确定，也必须提交给那些确定性的标准，这些确定性与它们在其本质条件的共同性中被归为一类，从而使它们也分享其证明方式。因此，我将以最大的信心主张：那今日不确定的那一位不是天主，因为他迄今为止是未知的；因为凡是显然为天主者，从这一事实本身也显然他从未是未知的，因此也从未是不确定的。\par

% structure: p0023 source=h2 target=subsection reason=relative-heading-rank; base=h2

\subsection{第十章 造物主从起初就因祂的创造被认作真天主。在梅瑟启示之前，已被人的灵魂和良心承认。}

因为确实，作为万物的创造者，祂从起初就与他们同时被发现，他们自身显现出来，为的是使祂能被认识为天主。尽管许久之后，梅瑟似乎率先在他著作的圣殿中引入对宇宙天主的认识，但那种认识的诞生不可因此从 \WorkTitle{梅瑟五书} 算起。因为梅瑟的卷册绝非开创了对创造者的认识，而是从一开始就指出这认识要追溯到乐园和亚当，而非埃及和梅瑟。因此，人类的大部分，虽然不知道梅瑟的名字，更不用说他的著作，却知道梅瑟的天主；甚至当偶像崇拜以其极端盛行笼罩世界时，人们仍以祂自己的名分别称祂为天主、万神之神，并说：\Quote{若天主允准，}\Quote{如天主所愿，}\Quote{我把你托付给天主。}那么，请想一想，他们是否认识这位他们见证能行万事的祂？他们并不把这归于梅瑟的任何著作。灵魂先于预言。从一开始，对天主的认识就是灵魂的嫁妆，在埃及人、叙利亚人和本都各族中都是一样的。因为他们的灵魂称犹太人的天主为他们的天主。哦，野蛮的异端者，不要把亚巴郎置于世界之前。即使创造者曾是一个家族的天主，祂也并不晚于你的神；甚至在本都，祂在以前就被人认识。那么，从那先来的取你的标准：从确定的（必须判断）不确定的；从已知的到未知的。天主永不会被隐藏，天主永不会缺少。祂将永远被理解，永远被听见，甚至被看见，以祂所愿意的任何方式。天主有我们的整个存有和这我们居住的宇宙作祂的见证。因此，因为并非未知，祂被证明既是天主，又是唯一的天主，尽管另有人仍竭力证明他的主张。\par

% structure: p0025 source=h2 target=subsection reason=relative-heading-rank; base=h2

\subsection{第十一章 天主外在于祂的证据；但产生这证据的外在受造物实际上并非外在于祂，因为万物都属于天主。马西昂的天主，无物可显，根本就不是天主。马西昂的方案荒谬地有缺陷，不能为他的新天主的存在提供证据，而这证据至少应能与创造者的充分证据竞争。}

他们说得对。因为人岂不正是凭其自身（固有）特质比凭异质更少被认识？无人如此。那么，我坚持此说：对天主而言，如何会有异质？倘若祂亲自存在，便无物为异质。因为这是天主的属性：万物属于祂，万物归于祂；否则，我们不会如此轻易听到这问题：祂与异质者有何相干？——这一点将在适当处更充分地论及。现在只需指出：凡无物证明属于他者，便无人证明其存在。因为正如造物主因万物属祂、无物为异而被证明为天主，毫无疑问的天主；同样，对立之神从无物属他、万物因此为异的事实中，被视为非神。既然宇宙属于造物主，我看不出有其他任何神的余地。万物充满其作者，被他占据。若受造物中任何处有虚空无神，那虚空明确是假神的虚空。藉虚假，真理被显明。为何无数假神不能在某处为马西翁的神找到一席之地？因此，我坚持：从造物主的特性看，天主必已从某个专属祂自己的世界的作品中被认识，既包含其人类成分，也包含其余有机生命；甚至世间的错误也敢于称那些它有时承认的人为神，理由是每在这种情况下，总看到某种提供生活用途与益处的东西。相应地，从天主的特性看，这也被相信是一种神圣功能：即教导或指出人类事务中方便与必需之事。如此彻底地，赋予假神以影响力的权威，正是从先前流向真神的同一源头借来的。马西翁的神至少应当产出一样属于自己的蔬菜；这样他才能被当作一个新\Person{Triptolemus}来传扬。否则，请提出一个配得上神的理由，解释他为何（假设他存在）没有创造任何东西；因为他（假设他存在）必须是创造者，正是基于我们清楚我们自己的天主之所以存在，恰恰因为祂是我们这个宇宙的创造者这一原则。因为，至此一次，这规则将成立：他们不能既承认造物主是天主，又证明他们所愿同样被信为神的那位是神圣的，除非他们将他调整为祂的标准——祂是所有人共认的天主；即：既然无人怀疑造物主是天主，明确因祂创造了宇宙，那么基于同样理由，无人应相信那创造无物者也是神——除非有些好的理由出现。而此理由必限于二：他要么\Emph{不愿}创造，要么\Emph{不能}。没有第三种理由。现在，说他不能，是配不上神的理由。至于不愿，我欲探究其是否配得上神。告诉我，马西翁，你的神曾否愿自己被认识？若不，他为何从天堂降下、传道、受苦并从死者中复活？岂非为被认识？无疑，既然他被认识，他便愿如此。因为若他不愿，便不会有任何事临到他。什么比他在肉身的卑辱中显现更促进对他的认识呢？——这屈辱若肉身仅为幻影，则更低下。因为这更为可耻：他，自悬木而惹来造物主的诅咒，却只假装取了肉身。他本可在一些自己创造的证据中为认识自己奠定远更高贵的基础，尤其在他要对立于那位在其领域内从起初便无任何作品而不被认识者时。因为，如何发生：造物主虽不知（按马西翁派所言）有位高于自己的神，且甚至用誓言宣告祂独自存在，却通过如此伟大的作品守护着对自己的认识——若假设祂独一而无对手，祂本可对之毫不关心；而那位更高之神，明明深知其下级对手在能力上装备精良，却全然未作预备以使自己被承认？实则祂本应产出更光辉、更崇高的作品，以便祂能按照造物主的标准，从其作品中被认识为神，甚至以更高贵的作为显示自己比造物主更有能力、更为仁慈。\par

% structure: p0027 source=h2 target=subsection reason=relative-heading-rank; base=h2

\subsection{第十二章。无需此神存在的外在证据便不可能承认神。马西翁拒绝其神的此类证据，实属无耻与恶毒。}

但即使我们能容许祂存在，我们仍应论证祂是没有原因的。因为那没有什么（可以证明自己存在的东西），便是没有原因的，因为（这种）证明正是某个证明所属之人存在的全部原因。然而，\Emph{正如}没有东西应当没有原因，即没有证明（因为若没有原因，就等于不存在，它缺乏那使事物存在的原因的证明），\Emph{同样}，我更应当相信天主不存在，而非相信祂没有原因地存在。因为没有原因者，即因缺乏证明而没有原因。然而，天主不应当没有原因，也就是说，没有证明。因此，每当我指出祂没有原因地存在时，虽然（我承认）祂存在，但我实际上已断定祂不存在；因为，若祂存在，祂就不可能完全没有原因地存在。同样，甚至就信德本身而言，我说祂试图从人那里无原因地获得信德，而人通常是根据祂工程的见证所形成的对天主的观念而相信祂；（我重申，无原因地）因为祂并没有提供那种使人得以认识天主的证明。因为尽管多数人相信祂，但他们并非单凭理性立即相信，而没有在配得天主的工程中拥有某种神性的标记。因此，基于这种不作为和缺乏工程的理由，祂既冒昧又恶毒：冒昧，在于追求一种不属于祂、且未为之奠定基础的信德；恶毒，在于使许多人因未获得信德的基础而陷入不信的罪责。\par

% structure: p0029 source=h2 target=subsection reason=relative-heading-rank; base=h2

\subsection{第十三章 \Person{马西昂}派贬低受造界，然而受造界是天主相称的见证。这一价值藉异教哲学家的例证得以阐明，他们常将受造的各部分赋予神性。}

当我们正将一位没有任何相称于天主且配得天主之尊的见证（如造物主的见证）可供证明自己的神从这一（神性）等级中驱逐出去时，\Person{马西昂}那些最无耻的追随者却以傲慢无礼的态度攻击造物主的工程，企图加以毁灭。的确，他们说，世界是一件伟大的工程，配得一位神。那么，造物主难道根本不是神吗？祂无疑是神。因此，世界并非不配天主，因为天主没有造过任何不配祂自己的东西；尽管世界是为人类而非为祂自己而造的，并且每个工程都小于其制造者。然而，若作为我们如此这般受造的作者不配天主，那么全然什么也不造——甚至一件虽然不配却仍可激发对更好尝试之希望的产物——岂非更不配祂！那么，关于这所谓不配的世界结构——在\Place{希腊}人中它也被赋予一个装饰与恩宠的名字，而非卑贱之名——那些智慧之士（一切异端的精神皆源自他们的天才）曾称上述不配的元素为神：如\Person{泰勒斯}称水，\Person{赫拉克利特}称火，\Person{阿那克西美尼}称气，\Person{阿那克西曼德}称天体，\Person{斯特拉托}称天与地，\Person{芝诺}称气与以太，\Person{柏拉图}称众星（他称之为火神族）；至于世界，当他们思量其大小、力量、能力、尊荣与荣耀——以及那些促成万物产生、滋养、成熟与再生的各个元素之丰盛、有序与规律——大多数哲学家便犹豫不定，不敢为该世界设定开端与终结，唯恐其构成元素（尽管确实伟大）不再被视为神性之物，而它们正是\Place{波斯}\OriginalTerm{祆僧}{magi}、\Place{埃及}\OriginalTerm{圣师}{hierophants}与\Place{印度}\OriginalTerm{裸行智士}{gymnosophists}所崇拜的对象。民众的迷信本身，受普通偶像崇拜的激励，当为其偶像所承载的古代死者的名字与传说感到羞愧时，便诉诸对自然物的解释，从而极其巧妙地掩饰其耻辱，将\Person{朱庇特}喻为炽热之物，将\Person{朱诺}喻为气态之物（根据\Place{希腊}词的字面意义）；同样，\Person{维斯塔}喻为火，\Person{缪斯}喻为水，\Person{大母神}喻为大地（其庄稼被收割，被强壮的手臂耕种，被沐浴浇灌）。同样，\Person{奥西里斯}每当被埋葬、被期待复活并喜悦地重新寻回时，便是大地果实复归、元素复苏、年轮更替之秩序的象征；\Person{密特拉}的狮子亦是干旱与焦枯自然的哲学圣事。对我来说，自然元素（在位置与状态上最为卓越）被视为神圣而非不配天主，这已足够。然而，我将降格至更卑微之物。一处篱笆上的小花（我不是说草地上的），任何海洋中的一个小贝壳（我不是说\Place{红海}的），一只水禽的迷途羽毛（我不说孔雀的），我想，将会向你证明，造物主不过是个拙劣的工匠！\par

% structure: p0031 source=h2 target=subsection reason=relative-heading-rank; base=h2

\subsection{第十四章 受造界的各部分均证明造物主的美善，而\Person{马西昂}却诋毁祂。由此暴露其自相矛盾。\Person{马西昂}自己的神在建立其宗教时，竟毫不犹豫地使用造物主的工程。}

如今，当你嘲笑那些微小的动物时——它们的荣耀造物主特意赋予了它们丰富的本能和资源，以此教导我们：伟大在卑微中得着证明，正如宗徒所说，软弱中也显出能力（\BibleRef{格后 12:5}）——你且试着模仿蜜蜂的蜂巢、蚂蚁的土堆、蜘蛛的网罗、蚕的丝线；你若懂得，也请忍受那些侵扰你床榻与房屋的虫子，斑蝥的有毒喷液、苍蝇的尖刺、蚊虫的鞘与螫针。既然这些小动物都能如此影响你的喜乐或痛苦，以至于你连在它们身上也无法轻看它们的创造主，那么对于更大的动物呢？最后，请环绕你自己，审察人内内外外。我们天主亲手所造的这作品，也必使你喜悦，因为你的主，那个更好的神，如此爱它，甚至为你之故不惜从第三层天降下到这些贫乏的元素之中，并为此缘故，竟在造物主这可怜的居所里被钉在十字架上。事实上，直到如今，他并未藐视造物主所造的水——他用这水洗涤他的子民；也未藐视油——他用油敷抹他们；也未藐视蜜与奶的调和——他用此赐予他们孩提般的滋养；也未藐视饼——他用饼来代表他自己真正的身体，因此在他的圣事中，他竟需要造物主的\Quote{卑劣元素}。然而你却作门徒高于师傅，作仆人高于主人；你的辨别力比他还高；你毁弃他所需要的。我要察看，你至少在这点上是否诚实，以致对那些你所毁弃之物毫无渴慕。你是天空的敌人，却喜欢在屋中享受它的清新。你诋毁大地，虽然它是你肉身的原始母亲，仿佛它是你确凿的敌人，却从它那里提取一切肥甘作为你的食物。你也谴责海洋，却不断使用它的产物，并视之为更神圣的食物。若我献你一朵玫瑰，你必不会嫌恶它的造作者。你这伪善者！无论你以多少刻苦来表露你身为马西昂派，即你的造物主的弃绝者（因为若这世界不中你的意，你本当视这种刻苦为殉道），但你终必与造物主的物质之作结合，无论你将溶解于何种元素。你的顽固何其难改！你诋毁那些你赖以生存和死亡的事物。\par

% structure: p0033 source=h2 target=subsection reason=relative-heading-rank; base=h2

\subsection{第十五章 马西昂之神启示的迟延。敌对神祇所处的空间问题。马西昂实际上（虽然似乎无意识地）在其体系中拥有九位神，而非两位。}

总而言之，或者，如你愿意，在一切之先，你曾说他有自己的受造物，自己的世界，自己的天空；诚然，当我们论及你本身的那位宗徒时，我们确实会看到那第三重天的事。与此同时，无论那（受造的）实体是什么，它无论如何都应当与它自己的神一同出现。但如今，为何主自\Person{提庇留·凯撒}十二年就被启示，而祂的受造物直到\Person{塞维鲁}皇帝十五年都未被发现呢？既然祂的受造物比造物主的卑微作品更卓越，当它的主和作者不再隐藏时，它当然应该停止隐藏自己。因此我问：如果它不能在这个世界显现，它的主又如何在这个世界出现呢？如果这个世界接受了它的主，为什么不能接受那受造的实体呢？除非它比它的主更大。但现在出现了一个关于空间的问题，涉及那上面的世界及其神。因为，看，如果他在自己之下、造物主之上拥有自己的世界，那么他必定把这世界放在一个位置上，这位置的空间在他自己的脚与造物主的头之间是空的。因此神自己既占据了\OriginalTerm{局部空间}{local space}，也使世界占据了\OriginalTerm{局部空间}{local space}；而这\OriginalTerm{局部空间}{local space}也将比神和世界加起来还要大。因为绝无可能，那容纳者不大于那被容纳者。事实上，我们必须小心，免得这里那里留下小块空地，让某个第三位神也带着自己的世界混进来。现在，开始数算你的众神吧。\OriginalTerm{局部空间}{local space}就是一个神，它不仅大于神，而且还是\OriginalTerm{无始}{unbegotten}、\OriginalTerm{未受造}{unmade}，因此是\OriginalTerm{永恒的}{eternal}，与神同等，神曾一直在其中。再者，既然他也从某种\OriginalTerm{无始}{unbegotten}、\OriginalTerm{未受造}{unmade}且与神同时的底层材料中制造了一个世界，正如\Person{马西昂}对造物主所主张的那样，你也同样把这个世界提升到那包围了两个神（即神与\OriginalTerm{质料}{matter}）的局部空间的尊严。因为按照神性的尺度，\OriginalTerm{质料}{matter}也是一个神，它确实是无始、未受造、永恒的。但如果他是从无中造了他的世界，那么（我们的异端者）也将不得不对造物主作同样的断言，因为他在世界的实体中将质料从属于造物主。但他也应当从质料中造他的世界，才是合理的，因为作为神，他的经历与同样作为神的造物主面前的经历相同。因此，如果你愿意，你可以这样数算\Person{马西昂}的三个神：造物主、局部空间和质料。此外，他同样使造物主成为局部空间中的一个神，而局部空间本身要根据完全相同的尊严尺度来评价；他将质料从属于作为其主的造物主，而质料却是无始、未受造、并因此永恒的。与这质料，他进一步关联了邪恶——一个无始的本原与无始的对象，未受造与未受造，永恒与永恒；因此他造了第四位神。因此，在较高的情况中你有神性的三个实体，在较低的情况中有四个。当加上各自的基督——一个在提庇留时代显现，另一个是造物主所应许的——\Person{马西昂}便明显地受到那些假定他认为只有两个神的人的误解，而实际上他暗示了至少九个，尽管他自己并不知道。\par

% structure: p0035 source=h2 target=subsection reason=relative-heading-rank; base=h2

\subsection{第十六章 \Person{马西昂}从可见与不可见之物的对照中假定两位神的存在。这个相对原则事实上是造物主——独一的天主——所有可见与不可见之作品的特征。}

由于，既然那另一个世界并未显现，其神也未曾显现，他们唯一的出路便是将事物分为可见与不可见两类，并指派两位神祇为它们的作者，从而将不可见者归给他们自己的（至高）天主。然而，除了一个异端者之外，还有谁能让自己的心智相信，创造中的不可见部分属于那位先前从未显示任何可见事物者，而不属于那位通过他在可见世界中的运作而使人也相信不可见者呢？因为，在有了某些（工作的）样品之后才表示赞同，远比毫无样品就表示赞同合理得多。我们将看到，甚至（你们喜爱的）宗徒也将不可见的创造\BibleRef{哥 1:16}归于哪一位作者——当我们来查考他时。目前（我们暂不引用他的见证），因为我们大多正借助常识和公正的论据为将来相信圣经的支持铺路。因此，我们断言，这种可见与不可见事物之间的差异，必须基于此理由归于造物主，甚至因为他的整个工作都由多样性构成——有形与无形；有生命与无生命；有声与无声；能动与不动；能产与不育；干与湿；冷与热。人本身也同样被多样性所调和，无论在其身体还是在其感知中。他的一些肢体强壮，另一些软弱；一些美观，另一些不雅；一些成双，一些单一；一些相似，另一些不似。同样，在他的感知中也存在多样性：时而喜乐，时而焦虑；时而爱，时而恨；时而愤怒，时而平静。既然情况如此，因为我们的整个受造界是由其各部分之间的相互竞争而形成的，那么不可见者应当归于可见者，而不应归于任何其他作者，只应归于那位被归于其对应物的作者，它们标志着造物主本身的多样性，他命令他所禁止的，也禁止他所命令的；他也击打，也医治。为什么他们认为他只在一类事物中是划一的，只作为可见事物（且仅为可见事物）的造物主？而实际上他应当被相信创造了可见与不可见两者，正如他创造了生命与死亡，或恶事与和平一样。诚然，如果不可见的受造物比可见的伟大（而在它们自身的范围内是伟大的），那么更为伟大的也应属于那拥有伟大者的；因为无论是伟大的，还是更为伟大的，都不适合作为那位似乎连最小事物也不拥有者的财产。\par

% structure: p0037 source=h2 target=subsection reason=relative-heading-rank; base=h2

\subsection{第十七章：正如马尔西翁派所假装的，至高天主拯救人是不够的；他也必须创造人。天主的存有由其创造得到证明，这是对他的属性的优先考虑。}

被这些论点所迫，他们喊道：\Quote{一项工作就足以证明我们的神；他借着他至高且最卓越的良善拯救了人，这胜过（创造）所有的蝗虫。}这是怎样的一个优越的神，竟无法找到任何工作比那较小的神所造的\Emph{人}更伟大！如今毫无疑问，你首先要做的是证明他存在，与平常证明天主存在的方式相同——借着他的工作；然后才借着他的善行。因为首要的问题是：他是否存在？然后是他的属性是什么？前者要由他的工作来检验，后者则由这些工作的恩惠来检验。不能简单地因为他被说成为人带来了救赎就断定他存在；只有在确定他存在之后，才有空间说他实施了这一解放。甚至这一点也必须有其自身的证据，因为可能他既存在，却又没有施行所声称的救赎。如今，在我们工作中涉及未知神问题的那部分，两点已经足够清楚——他什么也没有创造；而且他应当是一位造物主，以便借着他的工作而被认识；因为，如果他存在，他就应当被认识，而且从万物起始就被认识；因为天主隐藏起来是不合适的。我必须回到那未知神问题的本身主干，以便也可以探讨它的其他旁支。因为首先应当追问的是：为什么他后来才引人注意——为什么如此迟延，而不是在最初？从受造物来看，作为天主他确实是如此紧密相连的（联系越紧密，他的良善就越大），他本应从未被隐藏。因为不能假装没有任何达到认识天主的方法或好的理由，当从起初人就在世界中，救赎如今是为人而来的；正如造物主的恶意一样，善良的天主正是为了反对它而施行了救赎。因此，他要么是不知道他自己必要显现的好的理由和方法，要么是怀疑它们；或者要么是不能，要么是不愿面对它们。所有这些选择都不配于天主，尤其是至善至高的天主。然而，这一主题我们将在以后更充分地论述，并谴责那迟来的显现；目前我们只是指出它。\par

% structure: p0039 source=h2 target=subsection reason=relative-heading-rank; base=h2

\subsection{第十八章：尽管他们自以为是，马尔西翁派的天主在创造的证据和充分的启示两方面都缺乏凭证。}

那么，他现今已经进入人们的视野，就在他愿意的时候，就在他能够的时候，就在命定时辰来到的时候。因为或许他先前受到了他的主星、或某种厄运星体、或处于四分相位的土星、或处于三分相位的火星的阻碍。马西昂派非常热衷占星术，他们也不以借由造物主（他们所贬抑的）所造的星辰来谋生为耻。我们在此也必须讨论这（新）启示的性质：马西昂的最高天主是否以\Emph{与他相称的方式}显现自己，从而确保他存在的证明；并且以\Emph{真理的方式}显现，以致人相信他就是那已经被证明曾以与他本性相称的方式启示出来的那一位。因为与天主相称的事物将证明天主的临在。我们主张，天主首先必须从\Emph{本性}被认识，然后通过\Emph{教诲}得到确证：从本性是通过他的工程，通过教诲则是借由他启示的宣告。现在，在一个排除本性的情况下，没有提供任何本性的（认识）方法。因此，他应当仔细地通过宣告来启示他自己，尤其是因为他必须在相反于一位已经借由如此众多伟大工程（包括创造和启示宣告）才勉强满足人类信德的那一位的情况下被启示。这启示是如何发生的呢？如果借由人的猜测，不要说天主可能通过他自己以外的任何方式被认识，而且不仅要以造物主的标准为依归，还要以天主的伟大和人类的渺小为条件；这样看来，人不可能比天主更大，因为人已经通过某种方式将天主拖入公众认知，而天主自己却不愿意借由自己的力量被认识，尽管人类的渺小，根据全世界多方面的经验，已经更容易为自己塑造诸神，而不是追随人们现今借由本性所认识的真正天主。至于其余，如果人能如此造神——正如\Person{Romulus}造了\Person{Consus}，\Person{Tatius}造了\Person{Cloacina}，\Person{Hostilius}造了\Person{Fear}，\Person{Metellus}造了\Person{Alburnus}，以及某位权威不久前造了\Person{Antinous}——那么同样的事也可以被他人允许。至于我们，我们已经在\Person{Marcion}那里找到了我们的领航员，尽管他既不是国王也不是皇帝。\par

% structure: p0041 source=h2 target=subsection reason=relative-heading-rank; base=h2

\subsection{第十九章 耶稣基督，造物主的启示者，不可能是马西昂的天主；后者是在基督后约一百一十五年才由异端者得知，而且其原则（即律法与福音的对立）完全与耶稣基督的教训不合。}

然而，马西昂派说，我们的天主，虽然他从起初并借着受造物没有显现自己，却已在基督耶稣内启示了自己。将有一卷书专门论述基督，处理他的全部状态；因为最好将这些主题彼此区分开来，以便得到更充分、更有条理的处理。同时，在此阶段，我只要——且简要地——证明基督耶稣不是任何其他天主的启示者，而只是造物主，就足够了。在提庇留十五年，基督耶稣愿意从天上降下，作为拯救之灵。我确实没有兴趣去探究老安东尼努斯当政的哪一年。那位怀有如此恩慈目的的主，更像是从本都吹来的一股瘟疫般的热风，呼出这种健康或救恩，而这是马西昂从他的本都所教导的。这位导师无疑是在安东尼努斯时期的异端者，在不虔敬者中的虔敬者。现在，从提庇留到安托尼努斯·比乌斯，大约有一百一十五年又六个半月。他们将基督与马西昂之间的间隔正好放在这里。那么，既然马西昂，如我们所示，在安东尼努斯时代首次将这位天主引入人们的注意，事情就立刻明了了——如果你是一个敏锐的观察者。日期已经决定了这个案子：那位在安东尼努斯统治下首次公开出现的主，并没有在提庇留时代显现；换句话说，安东尼努斯时代的天主并不是提庇留时代的天主；因此，马西昂所清楚宣讲的主，并不是基督（他早在提庇留时代就宣布了他的启示）所启示的。现在，为了清楚证明这个论点的其余部分，我将从我真正的对手那里获得材料。马西昂特殊且主要的工作是将律法与福音分开；他的门徒不会否认，在这一点上他们拥有进入并巩固于他的异端的最好借口。这就是马西昂的\WorkTitle{对反命题}，或相互矛盾的命题，旨在使福音与律法相冲突，以便从包含它们的两个文件之不同而争辩两个不同天主的存在。因此，既然正是律法与福音之间的这种对立暗示了福音的天主不同于律法的天主，那么，在上述分离之前，那位从分离本身论证中得知的天主就不可能已被认识。因此，他不可能被基督（他在分离之前来到）所启示，而必然是马西昂所杜撰的，他是打破福音与律法之间和平的作者。现在，这和平，从基督显现到马西昂的狂妄教义时期一直保持完好无损，毫无疑问是由那种思想方式所维持的，即坚信律法和福音的天主不是别的，正是造物主，而在如此长时间之后，被这位本都的异端者引入了一个分离。\par

% structure: p0043 source=h2 target=subsection reason=relative-heading-rank; base=h2

\subsection{第二十章 马西昂以圣保禄与圣伯多禄的争执来为他律法与福音之间的对立辩护，被证明误解了圣保禄的立场和论点。马西昂的教义被圣保禄的教训所驳斥，这教训完全与造物主的法令一致。}

这个极其明显的结论需要我们去辩护，以对抗对方的喧嚷。因为他们声称，\Person{马西昂}将法律与福音分开，并非在信德准则上有所创新，而是将其在先前被掺杂之后加以恢复。啊，\Person{基督}，最忍耐的主，祢容忍了多年这种对祢启示的干扰，直到\Person{马西昂}果真前来拯救祢！现在他们引证\Person{伯多禄}本人和其他宗徒柱石的例子，说他们被\Person{保禄}指责为没有按福音的真理正直行事——正是这位\Person{保禄}，当时尚在恩宠的初步阶段，总之，战战兢兢，惟恐他先前奔跑或现今奔跑是徒然的，那时他首次与那些先于他作宗徒的人交往。因此，因为他作为一个新皈依者，在反对犹太教的热情中，认为他们的行为有可指责之处——甚至他们交往的混杂性——但后来他自己在实践中成了对一切人成为一切，为要赢得一切——对犹太人，就如犹太人；对在法律之下的人，就如在法律之下——你会认为他的指责，仅仅是针对那种注定连指控者也会接受的行为，被怀疑在公共教义的一点上对天主说了诡诈的话。然而，关于他们的公共教义，正如我们已说过的，他们完全和谐地携手合作，并在福音的共融中同意了职责的分工，正如他们在其他一切方面一样：\BibleQuote{无论是我，或是他们，我们都这样传。}\BibleRef{格前 15:11}再者，当他提到\Quote{某些假弟兄偷偷混进来}，想要把迦拉达人拐到另一个福音中时，他自己表明福音的这种掺杂并非要使他们转向另一位神和基督的信仰，而是要永久延续法律的教导；因为他责备他们坚持割损，并遵守时期、日子、月份和年份，按照这些犹太礼仪，他们本该知道如今已被废除，按照造物主自己意图的新安排，祂早在古时藉着祂的先知预言了此事。因此祂藉着依撒意亚说：旧事已过。\BibleQuote{看，我要行一件新事。}\BibleRef{依 43:19}在另一处：\Quote{我要订立一个新盟约，不像我牵着他们祖先的手领他们出埃及时所订立的盟约。}同样藉着耶肋米亚说：\BibleQuote{你们要订立一个新盟约，自行割损归向上主，除掉你们心中的包皮。}\BibleRef{耶 4:4}因此，宗徒所坚持的正是这种割损和这种更新，当他禁止那些古代礼仪时——这些礼仪的创立者自己曾宣布它们有一天要废止；因此藉着欧瑟亚说：\BibleQuote{我也要使她的一切欢乐、她的庆节、她的初一、她的安息日以及她的一切盛会终止。}\BibleRef{欧 2:11}同样藉着依撒意亚说：\Quote{初一、安息日、召集集会，我无法忍受；你们的圣日、斋戒和庆节，我的灵魂憎恨。}如今，如果连造物主在如此久远之前就已丢弃了这一切，而宗徒如今宣告它们应当被弃绝，那么宗徒的意思与造物主的法令完全一致，这就证明宗徒所传的正是那位他如今希望人承认其旨意的天主，此外并无别的天主；他指责为假的那些宗徒和弟兄，其明确原因是他们正在将造物主基督的福音从造物主所预言的新境况拉回到祂所丢弃的旧境况。\par

% structure: p0045 source=h2 target=subsection reason=relative-heading-rank; base=h2

\subsection{第21章。圣\Person{保禄}在宣布废除天主的一些古老法令时，并未传讲一位新天主。在相信造物主就是\Person{基督}所启示的天主这一点上，从未有过任何犹豫，直到\Person{马西昂}的异端出现。}

如果他是为了宣扬一位新天主而急于废除旧天主的法律，那么他为何没有关于新天主的任何规诫，而只关于旧法律？除非是对创造主的信仰仍要继续，而祂的法律却要终结——正如圣咏作者所宣告的：\BibleQuote{让我们挣开他们的枷锁，抛开他们的绳索！列国为什么嚣张，万民为什么妄想？世上的君王起来，宰辅聚集，一同反对上主，反对祂的受傅者。}的确，如果保禄所宣扬的是另一位天主，那么关于法律是否应遵守的问题就毫无疑义了，因为法律当然不属于新主——法律的敌人。天主的崭新与差异不仅会消除关于旧法律（异己的法律）的所有争论，甚至会消除对它的所有提及。但当时整个问题在于：虽然法律的天主与在基督内所宣扬的天主是同一的，但祂的法律却受到贬损。因此，对创造主及其基督的信仰仍持恒，唯有生活方式和纪律才发生动摇。有人争论关于吃祭偶像之物，有人关于妇女蒙头，有人关于婚姻与离婚，甚至有人关于复活的希望；但关于天主，却无人争论。如果这个问题也曾进入争论，那么它必然会在宗徒的作品中找到依据，而且是一个重大而关键的问题。毫无疑问，在宗徒时代之后，关于天主信仰的真理受到了败坏；但同样确定的是，在宗徒在世时，他们关于这一重大问题的教导并未受到任何败坏；因此，任何其他教导都没有权利被视为宗徒的，唯有今日在宗徒建立的教会中所宣布的教导。然而，你将找不到任何宗徒来源的教会不将其基督信仰建立在创造主之上。但若教会从一开始就败坏了，那么纯正的教会在哪里？难道在创造主的敌人那里吗？那么，请给我们指出一个由宗徒传承而来的你们的教会，你们就赢了。因此，既然从各方面看，从基督直到马西翁的时代，在神圣真理的规范中，除了创造主之外没有其他天主，那么我们的论点就有了充分的确证。我们已经表明，这位异端者的天主起初是通过分离福音和法律而被认识的。因此，我们先前的立场得以成立：任何人凭自己臆想所发明的天主，都不应被相信；除非那人是先知，但那样他自己的臆想就与此无关了。然而，如果马西翁能宣称具有这种灵感特质，那就必须证明这一点。不能有怀疑或搪塞。因为所有异端都被这一个真理的楔子驱除了：基督被证明只启示了创造主，而不是别的天主。\par

% structure: p0047 source=h2 target=subsection reason=relative-heading-rank; base=h2

\subsection{第二十二章 天主的善性被视为本性；马西翁的天主在此有缺陷。它在最初需要时没有来拯救人类。}

但是，如果不放松仅凭成规所作的辩护，让我们有余地驳斥他的一切其他攻击，这敌基督又如何能彻底被推翻呢？因此，让我们接下来思考天主自身，或者不如说，我们借着基督所拥有的祂的影子和幻象，并以使祂超越造物主的条件来审查祂。无疑，将有明确的准则来审查天主的善。然而，我的首要问题是发现并把握这一属性，然后将其引申为规则。当我从时间方面审视这一主题时，我从未在物质存在的开端，或那些应该与它一同被发现的原因的开端，看到它从那里开始做一切必须做的事。因为那时已经有死亡，和有死亡之刺的罪，还有造物主的恶意，而另一位神的善本应准备好带来解救；如果它要证明自己是一种自然的\Emph{作用}，就应遵循这一首要规则，即当原因开始时立即施以援手。因为在天主内，一切事物都应是自然和内在的，就像祂自己的状态一样，以便它们是永恒的，从而不被视为偶然和外来，因而暂时且缺乏永恒。因此，在天主内，善必须是持久而不断的，就像储存在祂本性特质的宝库中，随时备用，应先于其原因和物质发展；如果这样先存，它就会\Emph{作为基础}支撑每一个原始物质原因，而不是远远观望、置身事外。简言之，在这里我也必须追问：祂的善为什么从起初就不运作？就像我们追问祂本身：为什么祂从一开始就没有被启示出来？那么，为什么没有呢？因为如果祂存在，就必须通过祂的善被启示出来。绝不能认为天主会缺乏能力，更不能认为祂不会履行祂的一切自然功能；因为这些功能如果被阻止运行，它们就不再是自然的。此外，天主自身的本性不知何为不活动。因此，如果祂的善运作，就可以认为它有一个开端。由此显然，祂因本性在任何时候都不愿行使祂的善。实际上，祂不可能因本性而不愿意，因为本性指引自身，如果停止行动，它也就不再存在。然而，在\Person{马尔基昂}的神那里，善在某时停止了运作。因此，一个能够在任何时候停止其行动的善，不是自然的，因为自然属性与这种停止不相容。如果它无法证明是自然的，就不再被视为永恒，也不配属于神性；因为只要它不自然，它就无法从过去提供任何自我延续的手段，也无法保证将来会延续。事实上，它从起初就不存在，而且无疑也不会持续到终。因为它在将来的某个时间可能不再存在，就像它在过去某个时期已经不存在一样。既然\Person{马尔基昂}之神的善在起初就失败了（因为他从起初就没有拯救人），这一失败必定是出于意志而非软弱。现在，故意压制善将被发现怀有邪恶的目的。因为什么邪恶能比有能力行善却不愿行善、阻挠有益之事、或容许伤害更甚呢？因此，对\Person{马尔基昂}的造物主的所有描述，都必须转移到他的新神身上，这位新神通过延迟自己的善，助长了前者的残忍行为。因为任何人若有能力阻止某事发生，若该事发生，便被认为负有责任。人因尝了一棵可怜树上的果子而被判处死刑，由此产生罪及其惩罚；现在所有人都在灭亡，他们甚至从未见过乐园的一寸土地。而这一切，你的那位更好的神要么不知道，要么容忍着。难道他因此而被视为更好，而造物主却被视为更坏吗？即使这是他的目的，他也足够恶毒，因为他既容许（邪恶）的工作进行以加重对手的坏名声，又让世界遭受错误的困扰。你会怎么看待一位医生，他通过扣留药物助长疾病，通过延迟处方延长危险，以便他的治疗更加昂贵和著名？对\Person{马尔基昂}之神必须作出这样的判决：容忍邪恶，鼓励错误，甜言蜜语谈论他的恩宠，在他的善上支吾其词，这善他并非单纯出于自身而显示，如果他本性为善而非习得，如果他在属性上至尊至善而非训练所得，如果他是从永恒而非从\Person{提庇留}开始为神，不（更真实地说），从\Person{切尔东}和\Person{马尔基昂}才开始。然而，就目前情况而言，我们所考虑的这种神更适合于\Person{提庇留}，以便神性的善在他的帝国统治下被引入世界！\par

% structure: p0049 source=h2 target=subsection reason=relative-heading-rank; base=h2

\subsection{第二十三章 天主的善的属性被视为理性。\Person{马尔基昂}之神在此也有缺陷；他的善是非理性的且应用不当。}

这是另一个适用于他的规则。天主的一切属性应当既合乎理性又合乎自然。我要求在其良善中有理性，因为除了那理性地善的事物外，没有其他能恰当地被视为善；良善本身更不能在任何非理性中被察觉。一件带有某种理性成分的恶事，比一件丧失一切理性品质的善事不被视为恶，更容易被视为善。现在我否认\Person{马尔西昂}之神的良善是理性的，首先基于这一理由：因为它着手进行一个与他疏远的人类受造物的救恩。我知道他们将会引用的辩解：那毋宁是一种首要而完美的良善，自愿地、自由地倾注在陌生人身上，没有任何友谊的义务，根据我们奉命要爱仇敌的原则，而仇敌也正是因此对我们来说是陌生人。既然从一开始他就没有关心过人类（从一开始就对他陌生），他通过这种疏忽预先确定了他与一个异类受造物毫无关系。此外，爱陌生人或仇敌的规则是以爱邻人如同自己的诫命为前提的；而这一诫命，虽然来自造物主的法律，连你也应当接受，因为它非但没有被基督废除，反而被祂所确认。因为你被命令爱仇敌和陌生人，其目的是为了你能更好地爱邻人。对不应得的善意的要求是当得的善意的增加。但当得的先于不应得的，如同主要品质，而且比其他品质更值得作为它们的陪伴和同伴。因此，既然神圣良善的合理性的第一步是使自己在正义中显示于其正当的对象上，而只有在第二阶段才通过一种超出经师和法利塞人正义的额外正义显示于一个异类对象上，那么，第二步如何能归于那在第一步上失败者（他没有将人作为其正当对象，并因此使他的良善有缺陷）呢？此外，一种有缺陷的仁慈，没有合适的对象可施予，如何能满溢到一个异类的对象上？澄清第一步，然后再证明下一步。没有秩序，任何事都不能声称是理性的，理性本身更不能在任何情况下放弃秩序。假设现在\Emph{神圣的}良善从其理性运作的第二阶段开始，也就是说，作用于陌生人，那么如果它在任何其他方面受损，这第二阶段将在理性上不一致。因为只有当良善的第二阶段（即对着陌生人显示的阶段）在操作时不对有优先权利者造成损害时，它才被视为理性的。正是正义使一切良善在万事万物之前成为理性的。因此，当在其正当对象上显示时，如果它是正义的，它在主要阶段就是理性的。同样，当向着陌生人显示时，如果它不是不正义的，它就能显得理性。但那在错误中、且是\Emph{那}代表一个异类受造物所显示的良善，是什么样的良善呢？因为一种仁慈，即使在有害时，如果是为自家人施行的，或许在某种程度上可被视为理性。然而，按照什么规则，一种不正义的仁慈（是代表一个陌生人，连诚实的仁慈都合法地不应给予他）可以被辩护为理性的呢？因为有什么比如此善待一个异奴——将他从他的主人那里带走，宣称他是另一个人的财产，并教唆他背叛主人的性命；而这一切，为了更显不义，是在他仍住在他主人的家中，靠他主人的谷仓生活，仍在他主人的鞭笞下战栗时——更不义、更不公正、更不诚实的呢？这样的解救者——我几乎要说绑匪——甚至在世俗中也会受到谴责。现在，\Person{马尔西昂}的神的性格正是如此：扑向一个异类世界，从他的天主那里抢夺人，从父亲那里抢夺儿子，从导师那里抢夺学生，从主人那里抢夺仆人——使他对他的天主不虔，对其父亲不孝，对其导师忘恩，对其主人毫无价值。如果理性的仁慈使人如此，请问非理性的仁慈会把他变成什么样的人呢？我想没有比以下更无耻的了：他用属于另一个人的水受洗归于他的神，向着属于另一个人的天向他的神伸开双手，在属于另一个人的地上向他的神跪拜，用属于另一个人的饼向他的神献上感恩，并以施舍和爱德的方式，为了他的神分发属于另一位天主的恩赐。那么，他们那如此良善的神是谁，以致人通过他变成了恶；又如此有恩宠，以致激起了那另一位天主（他事实上是他自己的真正主）对人的愤怒呢？\par

% structure: p0051 source=h2 target=subsection reason=relative-heading-rank; base=h2

\subsection{第二十四章 \Person{马尔西昂}之神的良善仅不完美地显现；它拯救的人数稀少，且仅仅拯救他们的灵魂。\Person{马尔西昂}对身体的轻视是荒谬的。}

但既然天主是永恒而理性的，我认为祂在所有方面都是完美的。\Quote{你们应当是成全的，如同你们的天父是成全的一样。} \BibleRef{玛 5:48} 那么，请证明你的神的良善也是完美的。我们已经充分表明它实际上是\Emph{不完美的}，因为它被发现既非自然的也非理性的。然而，现在要用另一种方法阐明同样的结论：它不仅仅是简单的完美，而且是软弱的、无力的、枯竭的，未能涵盖其对象的全部数目，也并未在它们所有之中显明自身。因为并非所有人都因它而得救；而创造者的子民，无论是犹太人还是基督徒，都被排除在外。如今，既然大部分人都这样丧亡，那种在大多数情况下不起作用、仅在少数中稍有效果、在许多中毫无作用、屈服于丧亡并与之共谋的良善，如何能被捍卫为完美的呢？如果这么多人都将丧失救恩，那么更完美的将不是良善，而是恶意。因为正如良善的运作带来救恩，阻止它的则是恶意。然而，既然这种良善只救少数人，从而更倾向于不救的选项，那么它通过不提供帮助比通过帮助更能显出其完美。现在，你将无法在涉及创造者的方面把良善归于你的神，如果它面对所有人时都失败的话。因为无论你召谁来评判这个问题，你必须把神圣性情提交裁决，视之为良善的分配者（如果这个称号能够成立的话），而非你声称你的神是良善的挥霍者。因此，既然你仅凭良善就偏爱你的神胜过创造者，并且他宣称唯独拥有这属性，他就不应该对任何人缺乏它。然而，我现在并不想证明马西昂的神在良善上不完美，因为大多数人丧亡。我满足于通过以下事实来说明这种不完美：即使被他拯救的人也只获得不完美的救恩——也就是说，他们仅在灵魂方面得救，而身体方面却失丧了，因为按照他的说法，身体并不复活。那么，这种对救恩的割裂从何而来，如果不是来自良善的缺失呢？还有什么能比使整个人恢复救恩更能证明完美的良善呢？既然被创造者完全定罪，他就应该被最仁慈的神完全恢复。我反而认为，按照马西昂的规矩，身体受了洗，被剥夺了婚姻，在告解中受到残酷的折磨。虽然罪归于身体，但它们却先有灵魂有罪的贪欲；甚至，罪的第一个动机必须归于灵魂，肉体是作为仆役来行动的。不久之后，当脱离灵魂时，肉体就不再犯罪。因此，在这件事上，良善是不义的，也是不完美的，因为它将更无害的本性——它并非出于意志而是出于顺从才犯罪——留给毁灭。尽管基督并未披上肉体的真实性，正如你们的异端所乐于假设的那样，祂仍然屈尊接受了肉体的外表。因此，祂肯定会因这假装的接受而对它有所顾及。此外，人若非肉体，又是什么呢？因为无疑正是从身体而非精神元素，人之本性的作者给人起了名字。\Quote{上主天主用地上的尘土造了人，} 并非用精神的本质；这后来才来自天主的嘘气：\Quote{人就成了一个有灵的生物。} 那么，人是什么？无疑是用尘土造的；天主把他放在乐园里，因为祂是塑造了他，而不是吹气使他成为存在——一个肉体的结构，而非精神的结构。既然如此，你有什么脸面去主张那种良善的完美呢？它不仅在人的得救的某一点上失败，而是在整体能力上失败。如果那是一种完整的恩宠和实质的怜悯，只给灵魂带来救恩，那么我们现在完整享受的生命岂不更好？而只复活一部分将是一种惩罚，而非解放。完美良善的证明是，人在获救后，应从恶神的住所和权势下被解救出来，得到最良善、最仁慈的天主的保护。可怜上你当的马西昂，你正发着高烧；你疼痛的肉体长满了各种荆棘和蒺藜。你不只是暴露于创造者的雷霆之下，或战争、瘟疫和祂其他更重的打击之下，甚至还暴露于祂的爬虫之下。你认为自己在哪方面脱离了祂的国度，当祂的苍蝇仍在你的脸上爬行时？如果你的解救在于将来，为什么不在现在呢，好让它得以完美地完成？在祂面前，我们的境况远非如此，祂是我们族类的作者、审判者和受伤的元首！你展示祂只是一个良善的天主；但你不能证明祂是完美良善的，因为你没有被祂完美地解救。\par

% structure: p0053 source=h2 target=subsection reason=relative-heading-rank; base=h2

\subsection{天主并非单纯美善的存在；其他属性也属于祂。\Person{马西昂}在描绘其单纯美善且无情的天主时显出自相矛盾。}

关于这善的问题，我们在这些论证纲要中已表明它决不可能与神性相容——因为它既不自然、不合理，也不完美，反而是错误的、不义的、不配称为善的——因为就神性本性的相符性而言，确实不适合将那被认为拥有\Emph{这样的}善且不是以修改的方式，而是单纯且唯一的方式者视为天主。此外，在这点上完全可以讨论：天主是否应被视为具有单纯之善，而排除所有那些\Person{马西昂}派从他们的神那里转移到创造主身上的其他属性、感觉和情感，而我们也承认这些也是创造主值得拥有的特性，但只是因为我们视他为天主。那么，基于此，我们将否认他是天主，因为在他身上找不到一切适合天主之存有的东西。如果（\Person{马西昂}）选择采纳\Person{伊壁鸠鲁}学派的任何一位，以\Person{基督}之名称他为天主，理由是幸福且不朽坏者不会给自己或他人带来任何麻烦（因为\Person{马西昂}在研究这种神圣漠然的观点时，从他身上移除了所有审判特质的严厉与能量），那么他有责任将他的概念发展成一个不受干扰且懒散的天主（那么他与\Person{基督}有什么共同之处呢？\Person{基督}既通过教训给犹太人带来麻烦，也通过感受给自己带来麻烦？），或者承认他拥有与他人相同的情感（在这种情况下，他与\Person{伊壁鸠鲁}有什么关系呢？\Person{伊壁鸠鲁}既不是他的朋友，也不是基督徒的朋友？）。因为一个在过去时代处于静止状态，不愿意通过任何工作传达任何关于自己的知识的存在，在如此漫长的时间之后，竟然开始关心人的救恩（当然是由他自己的意志），——难道他因此不就变得容易接受一个新意志的冲动，以至于明显地向所有其他情感敞开吗？但什么样的意志没有欲望的激励呢？谁渴望他不想要的东西呢？此外，关心将是意志的另一个伴侣。因为谁想要某个对象并渴望拥有它，而不也\Emph{关心}得到它呢？因此，当（\Person{马西昂}的天主）对人类的救恩既有了意志又有了欲望时，他肯定给自己和他人带来了一些关心和麻烦。这\Person{马西昂}的理论暗示，尽管\Person{伊壁鸠鲁}反对。因为他正是在他的意志、欲望和关心所对抗的事物——无论那罪或死亡——尤其是它们的暴君与主，即人的创造主那里，为自己树立了一个对手。再者，没有敌对竞争，任何事物都不会运行其进程，而竞争本身也不会没有敌对性。事实上，当（\Person{马西昂}的天主）意愿、渴望和关心拯救人时，他在行动中就已经遇到了对手：既在拯救所来自的那一位那里（因为他当然意味着自由来与他对抗），也在拯救所产生的事物那里（因为意图的自由有利于其他事物）。因为竞争必然伴随其附庸的情欲，无论其竞争针对什么对象：愤怒、纷争、仇恨、蔑视、愤怒、怨毒、厌恶、不快。既然所有这些情感都存在于竞争中，而且拯救人的竞争又激发了它们，并且人类的拯救又是善的运作，那么这种善若无其附属物——即那些他借以对抗创造主的感觉和情感——就毫无用处；因此甚至不能认为它是无理性的，仿佛它缺乏适当的感觉和情感。当我们来为创造主辩护时，我们将更充分地坚持这些观点，在那里它们也将招致我们的谴责。\par

% structure: p0055 source=h2 target=subsection reason=relative-heading-rank; base=h2

\subsection{在正义的属性上，\Person{马西昂}的天主软弱得毫无盼望且不像神。他厌恶邪恶，却不惩罚行恶者。}

但在此处，仅通过对其孤立之善的阐述就足以证明他们的神极端乖张，因为他们拒绝赋予他他们在创造主身上所谴责的种种情感。如今，如果他不易产生竞争、愤怒、损害或伤害的情感，如同一个克制自己不使用审判权的人，我无法知道任何纪律体系——而且是一个完整的体系——如何能与他一致。因为他怎能发布命令，如果不打算执行它们？或者禁止罪，如果不打算惩罚它们，反而拒绝法官的职能，仿佛对严厉和审判惩罚的所有概念都格格不入？因为他为何禁止那犯下后他并不惩罚的行为？如果他没有禁止他不打算惩罚的事，而他却惩罚他没有禁止的事，那就合理得多。不仅如此，他甚至有义务允许他以如此不合理的方式禁止，以致不为犯罪行为附加任何惩罚。因为即使是现在，那是被默许的，即被禁止却未施任何报复的行为。此外，他只禁止犯下他不喜欢发生的行为。因此，他是最懒散的，因为他对做他厌恶的事不感到愤怒，尽管不悦本应是他受冒犯的意志的伴侣。如今，如果他受到冒犯，他应当发怒；如果发怒，他应当施加惩罚。因为这样的施加是怒气正当的果实，而怒气是不悦的债务，不悦（如我所说）是受冒犯的意志的伴侣。然而，他不施加惩罚；所以他不受冒犯。\par

他不感到受冒犯，因此他的意志未受伤害，虽然那事是他所不愿意发生的；而这过犯如今是在他意志的默许下犯下的，因为凡不冒犯意志的事，就不是违背意志而犯的。现在，若这神圣德性或美善的原则是：确实不愿意某事发生并加以禁止，但对其实施却不被触动，那么我们就断言，当他宣告他的不愿意时，他已经是被触动了；而他在意愿它不被做成之后，既然已被其可能性触动，却不受其成就的触动，那是徒劳的。因为他因不愿而禁止了它。那么，当他宣告不愿意并随之禁止时，他岂不就是在行使审判的行为？因为他判断它不应被做，并慎重地宣告它应被禁止。因此，到此时他甚至履行了审判者的角色。若天主履行审判职能是不合适的，或者至少只到某种程度合适，即祂仅仅可以宣告不愿意并宣布禁止，那么当过犯发生后，祂甚至不能惩罚它。如今，再没有比不执行祂所不喜悦并禁止之事的报应，更不配天主性体的了。\Emph{第一}，祂有义务对祂所颁布的任何判决或法律施加惩罚，以维护祂的权威并保持对其的服从；\Emph{第二}，因为对祂所不喜悦之事，并因那不喜悦而禁止之事，敌对性的反对是不可避免的。此外，天主宽恕恶人比惩罚恶人更不合适，尤其是在至善至圣的天主身上，祂并非别样地完全良善，而是作为邪恶的仇敌，并且\Emph{那样}到如此程度，以致祂藉着对邪恶的仇恨显示对良善的爱，并藉着根除后者来实现对前者的保卫。\par

% structure: p0058 source=h2 target=subsection reason=relative-heading-rank; base=h2

\subsection{第27章 如此软弱的天主之教义对宗教与道德的危害性影响}

再者，他藉着不愿而审判邪恶，藉着禁止而谴责它；而另一方面，他藉着不报复而宣告它无罪，藉着不惩罚而释放它。这样的神是何等真理的支吾者！何等与自己决定不一的伪善者！害怕谴责他所真正谴责的，害怕憎恨他所不爱的，允许那他所不允许的事，选择表明他所不喜悦的而非深深检视它！这将成为一种想象的良善，纪律的幻影，对职责敷衍了事，对罪过掉以轻心。听着，你们这些罪人；以及你们这些尚未到此地步的人，听着，好使你们也达到如此境地！一个更好的神已被发现，他从不感到冒犯，从不发怒，从不施加惩罚，他没有预备地狱之火，没有外在黑暗中的咬牙切齿！他纯粹而简单地良善。他确实禁止一切不法行为，但只在言语上。他在你们内，只要你们愿意为了外表而向他致敬，好使你们显得尊敬天主；他不需要你们的恐惧。马西昂派如此满足于这些伪装，以致他们毫无对他们的神的恐惧。他们说，只有恶人才会被怕，好人会被爱。愚昧的人，你说你所称为主的那位不应被畏惧，而你所给他的这个头衔本身就指明一种必须被畏惧的能力？然而，你怎能爱，却没有一些畏惧而你不爱呢？确实，（这样的神）既不是你的父，你对祂的出于义务的爱应与因祂的能力而有的畏惧一致；也不是你真正的主，你应因祂的人性而爱祂，并因祂是导师而畏惧祂。绑架者确实以这种方式被爱，但他们不被畏惧。因为能力不会被畏惧，除非它是公正而规范的，尽管它甚至在腐败时也可能被爱：因为它靠的是引诱，而非权威；靠的是谄媚，而非正当影响。还有什么比不惩罚罪恶更直接的谄媚呢？那么，来吧，如果你不害怕天主为良善，为何你不爆发出各种情欲，并因此认识到，我相信，那是对所有不敬畏天主之人生活的主要享受？为何你不常去那疯狂的竞技场、血腥的斗兽场和淫荡的剧院的惯常娱乐？为何在迫害中，当香炉献上时，你不立即藉着否认信仰来赎回你的性命？你说\Quote{天主禁止}并加倍强调。所以你的确害怕罪恶，并藉着你的恐惧证明，那禁止罪恶的主是可畏惧的对象。这与你们向那你们不害怕的神所献的谄媚敬意完全不同，那敬意实际上与它自身的行为同样悖逆：禁止某事却不附加惩罚的制裁。他们行事更加徒劳，当被问及：在那大日子里，每个罪人将如何？他们回答，他将被丢到看不见的地方。这岂非也是一个司法判断的问题？他被判定应受摒弃，且是经过谴责的判决；除非罪人被抛弃倒也为了他的救恩，好使如此宽容也符合你们那至善至美之神的性格！而被抛弃又是什么？不过是失去那人若非被摒弃本可得着的东西——即他的救恩。因此他的被抛弃将导致救恩的丧失；除了那发怒并被冒犯的权威，即那惩罚罪恶的审判者，不可能对任何人颁布这个判决。\par

% structure: p0060 source=h2 target=subsection reason=relative-heading-rank; base=h2

\subsection{第28章 这乖谬的教义剥夺了圣洗圣事的一切恩宠。若马西昂是对的，这圣事就不会赋予罪的赦免、重生、或圣神的恩赐。}

那么，他被抛弃之后会怎样？他们说，他将被投入造物主的火中。那么，他们的神难道没有制定任何补救措施，以便惩罚那些得罪他的人，而不必诉诸残忍的手段，将他们交给造物主？而造物主又会做什么呢？我想他会为他们预备一个加倍硫磺的地狱，如同对待亵渎他自己的人一样；除非他们的神出于热心，也许会对他的对手的叛臣显示仁慈。哦，这是何等的神！处处乖谬，无处理性，事事虚空，因此是虚无！——在他的状态、状况、本性以及每项安排中，我都看不到连贯和一致；不，甚至在他信仰的圣事中也看不到！因为照他看来，圣洗有何用处？如果是罪的赦免，他如何证明他赦免了罪，既然他丝毫没有显示他保留了罪？因为他若行使审判者的职能，他就会保留罪。如果是脱离死亡，他自己既然从未\Emph{向}死亡交付，又怎能\Emph{从}死亡中解救？因为他若一开始就定罪了罪，就必须将罪人交付死亡。如果是人的重生，他从未生育，又如何重生？因为从未做过任何事的人不可能重复做。如果是赐予圣神，他最初没有赋予生命，又如何赐予圣神？因为生命在某种意义上是对圣神的补充。因此，他给那从未被他盖印的人盖印；他洗那从未被他玷污的人；他将那本在救恩之外的血肉完全浸入救恩的圣事中！没有农夫会灌溉不能给他回报果实的土地，除非他和\Person{马西昂}的神一样愚蠢。那么，为什么要把圣德强加给我们最软弱、最不配的血肉，作为负担或荣耀？我又该说什么呢？一种纪律使那已经圣洁的更加圣洁，这是无用的。为什么要给软弱者加重负担，或给不配者加增荣耀？为什么不给所负担或荣耀的以救恩作为报酬？为什么不给这项事工其应得的赏报，不赏赐血肉以救恩？为什么甚至允许其中的圣德荣誉死去？\par

% structure: p0062 source=h2 target=subsection reason=relative-heading-rank; base=h2

\subsection{第二十九章 \Person{马西昂}禁止婚姻。\Person{戴尔都良}雄辩地辩护婚姻为神圣，并谨慎辨别\Person{马西昂}的教义与他自己\OriginalTerm{蒙丹主义}{Montanism}之间的区别。}

按马西昂的说法，除非肉身处于童贞、寡居或独身状态，或藉离婚购得圣洗的资格，否则它不被浸入圣事的水中——仿佛连不能生育者也不是从婚姻结合中获得其肉身似的。无疑，这样的方案必然涉及对婚姻的禁止。那么，让我们看看这是否正当：我们并非像某些尼苛劳党人那样，以维护情欲和奢靡为目的而企图摧毁圣德的幸福，而是作为已经认识圣德并追求它、偏好它的人，同时又不损害婚姻；我们不是以善代替恶，而是以更好的代替好的。因为我们不拒绝婚姻，只是节制而已。我们也不将圣德定为规则，而只是推荐它，视之为善，甚至更好的状态，只要各人按自己的能力审慎地持守它；但同时，每当有人诋毁婚姻为污秽之物以贬低造物主时，我们便热切地为之辩护。因为祂也赐福于婚姻，视之为尊贵的身分，以促进人类的繁衍；祂确曾如此对待祂的整个创造，使之成为健全而有益的用途。饮食不因此而被谴责，因为若以过于精致的口味享用，便会导致贪饕；衣裳也不该被指责，因为若过于奢华地装饰，就会陷入虚荣和骄傲。同样，婚姻本身也不该被拒绝，因为若不加节制地享受，便会燃起欲火。原因与过失之间，状态与其过度之间，有着巨大区别。因此，受责备的不是这类制度本身，而是对它的过度使用；这符合其创立者自己的判断，祂不仅说了\BibleQuote{你们要生育繁殖}（\BibleRef{创 1:28}），也说了\Quote{不可奸淫}和\Quote{不可贪恋你邻人的妻子}；并且用死亡来威胁那些不贞洁的、亵渎的以及可憎的奸淫和与男人、兽类行逆性之事的罪。如果对婚姻设定了限制——例如灵性的规则，在基督的服从下、藉护慰者的权威所维持的只许一次婚姻——那么，划定这限制的权柄应属于那曾经广泛允许者；属于那曾散布、如今收集者；属于那曾栽种、如今砍伐者；属于那曾撒种、如今收割者；属于那曾说过\BibleQuote{有妻子的要像没有妻子一样}（\BibleRef{格前 7:29}），而又曾说过\Quote{你们要生育繁殖}的那一位；终点属于那起点之主。然而，树被砍伐并非因为它应受责备，谷物被收割也并非因为它该被定罪——而仅仅是因为它们的时期到了。同样，婚姻状态并不需要圣德的钩镰，仿佛它是邪恶的；而是因为它已成熟到可被解脱，并准备好迎接那将在长久中藉收割带来丰硕收成的圣德。这使我注意到马西昂的天主：他斥责婚姻为邪恶和不洁之事，实则损害了那看似事奉的圣德本身。因为他破坏了圣德赖以存在的材料；如果没有婚姻，就没有圣德。当过度成为不可能时，节制的所有证明便丧失了；因为许多事物正是由其对立面得以显明。正如\BibleQuote{软弱中成全了刚强}（\BibleRef{格后 12:9}），同样，节制因婚姻的许可而得以彰显。如果那给予人追求节制生活的机会被夺走，谁还能被称为节制呢？饥荒能给人留下什么节制食欲的余地？贫穷能提供什么放弃野心的事？阉人配得什么抑制情欲的功劳？然而，完全阻止人类的繁衍，据我所知，可能与马西昂那至善至美的天主颇为一致。因为他禁止人生育，又夺走那产生人生命的制度，又如何能渴望人的得救呢？他不容许人存在，又如何能找到一个人来施予他的善？他憎恶人的起源，又怎能爱他？也许他害怕人口过剩，免得他因解放太多人而疲惫；免得他不得不制造许多异端；免得马西昂派的父母生出太多马西昂的高贵门徒。法郎的残忍，在婴儿出生时杀害他们，与之相比也将显得不那么不人道。因为法郎毁灭生命，而这位异端者的天主却拒绝给予生命：一个将人从生命移除，另一个则不接纳任何人进入生命。两者在杀人方面没有区别——人都被二者杀害：前者在出生后，后者在未出生时。我们本该感谢你，我们的异端者的天主，只要你禁止了造物主在男女结合上的安排；因为这样的结合，你的马西昂才得以出生！然而，关于马西昂的天主已经够了——从我们关于唯一神性的定义及其属性的条件来看，他显然完全不存在。但这整篇小作品的进程直接指向这个结论。因此，如果有人认为我们迄今取得的成果甚微，请他继续期待，直到我们审查马西昂所引用的圣经本身。\par

\SourceRecord{Translated by Peter Holmes.；Ante-Nicene Fathers；Vol. 3.；Edited by Alexander Roberts, James Donaldson, and A. Cleveland Coxe.；Buffalo, NY: Christian Literature Publishing Co.,；1885.；<http://www.newadvent.org/fathers/03121.htm>.}

% structure: p0001 source=h1 target=section reason=page-title

\section{\WorkTitle{驳斥玛尔基翁，第二卷}}

\Emph{戴尔都良在此表明，被玛尔基翁诬蔑的造物主，或称德谬格，乃是真实而良善的天主。}\par

% structure: p0003 source=h2 target=subsection reason=relative-heading-rank; base=h2

\subsection{第一章 玛尔基翁论证的方法错误而荒谬。论证的正确进路。}

再版这部小作品的机会（我们在第一部书的序言中已提及它的经历）使我们得以在驳斥玛尔基翁关于两位天主的论题时，对每位都按其主题的划分给予各自的描述和章节，界定一位天主根本不存在，并主张另一位确实是天主；至此与这位本都的异端者保持同步，他随意接纳一位而排除另一位。因为他若不摧毁真理的体系，就无法构建他谬误的学说。他发现必须拆除某些东西，才能构建他所希望的理论。然而，这个过程如同建造房屋却没有准备合适的材料。讨论本应只集中于这一点：那取代造物主的天主并非天主。然后，当假神被某些预设地界定独一完美神性的规则排除后，关于真天主就不会再有疑问。祂存在的证据将会清晰，而在支持任何其他神的一切证据都缺失的情况下，这证据更加明显；关于祂应毫无争议地受到尊荣这一点也会更加清楚：祂应被敬拜而非判断，被虔诚地事奉而非挑剔地处理，甚至因其严厉而畏惧。因为人更需要的难道不是对真天主的审慎评估吗？他可以说是偶然碰上这位天主的，因为没有别的神。\par

% structure: p0005 source=h2 target=subsection reason=relative-heading-rank; base=h2

\subsection{第二章 造物主天主的真道。异端者自称认识神性存有，其实与启示相悖并颠覆之。天主的本性及道路超越人类发现。亚当的异端。}

我们现在已经扫清了道路，得以默观全能的天主，即宇宙的主宰和造物主。我认为，祂的伟大正彰显于此：从起初祂就使自己被人认识；祂从未隐藏自己，而总是光辉照耀，甚至在\Person{罗慕路斯}以前，更不用说\Person{提庇留}的时代了；只有异端者，唯独他们，不认识祂，尽管他们如此费力地探讨祂。他们因此认为必须假定另一位神存在，因为他们更能够指责而非否认那位存在如此明显的天主，他们从感官的推论中汲取关于天主的一切思想；正如某个盲人或视力不全的人，因看不见视界的对象，便假定存在另一个日头，有着更温和、更健康的射线。人啊，只有一个日头统治这世界；即使你对它另有所想，它仍是最好的和有益的；虽然在你看来它可能过于猛烈有害，或者可能过于污秽腐败，但它忠实于自身存在的规律。既然你看不透这些规律，那么即使存在另一个无论多么伟大和良善的日头，你也同样无法承受它的光芒。如今，你视力有缺陷，无法辨识较低之神，对于那更高的那一位，你又作何观感？你实在太纵容自己的软弱了；你不致力于求证事物，一旦发现天主存在，就确信、无疑、因而充分认识祂，尽管你认识祂只限于祂旨意放置证据的那一面。但你甚至不是明智地否认天主，而是无知地谈论祂；不，你以貌似聪明指责祂，而如果你认识祂，你绝不会指责，更不会谈论祂。你确实称呼祂的名，却否认那名的本质真理，即那被称为天主的伟大；不承认它是那样的伟大，以至于倘若人有可能在各方面认识它，那就不再是伟大了。依撒意亚如此早地以宗徒般明晰，预见到异端心中的思想，问道：\BibleQuote{谁曾知道主的心？谁作过祂的谋士？祂与谁商议？……谁教导祂知识，将明智的道路指示给祂？}宗徒同意他的话，大声说：\BibleQuote{啊！天主的富饶、智慧与知识多么深奥！祂的判断是多么不可测量，祂的道路是多么不可探测！}\BibleRef{罗 11:33}\BibleQuote{祂的判断不可测量，}因为那是审判者天主的判断；\BibleQuote{祂的道路不可探测，}因为包含一种智慧和知识，无人曾指示给祂，除非是那些批评天主的人，他们说：天主不应是这样，而应是这样；好像除了天主圣神，还有谁知道天主的事呢？\BibleRef{格前 2:11}再者，他们拥有世界的精神，\BibleQuote{凭天主的智慧，未能因智慧认识天主，}\BibleRef{格前 1:21}他们自以为比天主更聪明；因为，正如世界的智慧在天主面前是愚妄，同样，天主的智慧在世界的眼中是愚拙。然而，我们知道\BibleQuote{天主的愚拙比人更智慧，天主的软弱比人更刚强。}\BibleRef{格前 1:25}因此，天主尤其伟大，当祂在人眼中显得渺小时；尤其良善，当人在判断中不认为祂良善时；尤其独一，当人在祂身上看到两位或更多时。如今，如果从起初\BibleQuote{属血气的人不接受天主圣神的事，}\BibleRef{格前 2:14}便将天主的法律视为愚拙，因而忽略遵守；并且作为进一步的后果，由于他没有信德，\BibleQuote{连他以为有的也要被夺去}——例如乐园的恩宠和天主的友谊，如果他继续服从，本可借此认识天主一切的事——那么，他被贬为物质本性，被放逐到耕种的劳苦中，在他那低沉而恋土的劳动中，将这种源自土地的世俗精神传递给他的全族，这全族是纯粹属血气和异端的，不接受属于天主的事，这有什么奇怪呢？或者谁还会犹豫，把亚当的重大之罪称为异端？当他凭自己的意志而非天主的意愿作出选择时犯下此罪。但是亚当从未对他的无花果树说：你为什么这样造成我？他承认自己受了迷惑；他没有隐瞒那诱惑者。他是个非常粗鲁的异端者。他是不服从的；但他并没有亵渎他的造物主，也没有责怪那造物主，他从生命之初就发现造物主如此善良卓越，也许他从起初就立祂为自己的审判者。\par

% structure: p0007 source=h2 target=subsection reason=relative-heading-rank; base=h2

\subsection{第三章 天主由祂的化工而被认识。祂的良善在祂的创造能力中显示；但本性永远常在；内在于天主，在它所有展示之前。这良善的第一阶段先于人。}

因此，当我们开始审查已知的天主，而当问题出现：祂如何为我们所知时，我们理应从祂的工程入手，这些工程先于人存在；这样，祂的美善，既与祂自身一同被立即发现，又被规约并确定下来，就能向我们提示某种意义，使我们得以理解随后事物的秩序是如何产生的。此外，马西昂的门徒们在承认我们的天主的美善时，或许也能学习到，同样地，这一美善如何堪当天主性，依据的正是那些我们据以证明它在他们的神那里不值得的理由。然而，正是这一点——在他们体系中至关重要的一点——\Emph{\Person{马西昂}}并未在任何其他神祇身上发现，而是从自己的神那里推演出来。那么，最初的美善是造物主的美善，借此天主不愿永远隐藏；换言之，不愿没有某种东西使天主得以被认识。因为，实际上，有什么比认识天主并享有祂更美善呢？虽然这一美善并未显现出来，因为尚不存在任何事物使之显现，但天主预知最终会显现出什么美善，因此祂着手发展自己圆满的美善，以成就那将要显现的美善；这美善并非突然的，不是出自某种偶然的恩赐或激动的冲动，好像只能从它开始运作的时刻算起。因为如果美善在开始运作时产生了自己的开端，那么它在行动时实际上本身并没有开端。然而，当美善完成了一个初始行动，时间季节的体系就开始了，为了辨别和标明这些季节，天上的星辰和光体被安置在各自秩序中。天主说：\BibleQuote{让它们存在，作为季节、日子和年份} \BibleRef{创 1:14}那么，在这时间进程之前，创造时间的美善本身没有时间；同样，在那由同一美善发端的开端之前，它也没有开端。因此，由于它完全没有开端的秩序和时间的方式，它将被视为拥有一个无垠且无尽的时代；也不可能视之为突发的、偶然的或冲动的情绪，因为它没有任何东西能引起这样的判断；换言之，没有任何时间序列。因此，它必须被视为一个永恒属性，植根于天主，存至永远，并因此堪当天主性，永远羞辱马西昂的神的仁慈，因为他的神不仅落后于（我不说）一切开端和时间，而且落后于造物主本身的恶意——如果美善中真的可能发现恶意的话。\par

% structure: p0009 source=h2 target=subsection reason=relative-heading-rank; base=h2

\subsection{第四章 藉永恒圣言创造人的下一步：灵性与物质恩赐，以及人自由意志的祝福。}

天主的良善既为此为人预备了寻求认识祂的道路，便在最初的启示之外，又增添了这一点：祂首先为人预备了一个居所——首先是无垠的宇宙结构，然后是更广阔的（更高世界的）结构——好使人在宏大与微小的舞台上操练并进步于他的试炼中，从而从天主赐予他的\Emph{善}（即他的高位）被提升至天主的\Emph{至善}，即某种更高的居所。在这善工中，\Emph{天主}使用了一位最卓越的臣仆，即祂自己的圣言。\Quote{我的心，}祂说，\Quote{涌出了我最优美的圣言。}让\Person{马尔西翁}由此学到这棵真正最优美之树的佳果的最初一课。但他却像一个最笨拙的小丑，将一根好枝嫁接在坏砧木上。然而，他那亵渎的嫩芽绝不会强壮；它将与栽种者一同枯萎，这样好树的本质便会显明。看那最终的结果：圣言何其丰盛！天主发出祂的\Emph{\OriginalTerm{命令}{fiat}}，事就成了；天主看了，认为好。\EditorialNote{创世纪第一章}并非祂不知道那好直到看见，而是因为它是好的，所以祂看见了，尊重了，并盖上了印记；并且藉着允准它们接受这一默观，完成了祂工程的良善。如此，天主祝福了祂所造的好，为将自己作为完整而完美的、言语和行为皆善的那位，传达给你。此时，圣言尚不知诅咒，因为祂与恶行无涉。我们将看到什么原因也要求天主有此\Emph{事}。与此同时，世界由一切美善构成，清楚地预示了为谁预备这一切，就有多少美善在为他预备。谁能比祂自己的肖像和模样的人更有资格居住在天主的工程中呢？那肖像是藉着比平时更具行动力的良善雕琢而成的，没有用颐指气使的话语，而是用友善的手，并以近乎可亲的话语为先导：\BibleQuote{让我们按我们的肖像，按我们的模样造人。}\BibleRef{创 1:26}良善说出了这话；良善用地上的尘土塑造了人，成为如此伟大的肉体实质，由同一材料建造出众多特质；良善向他吹入一个灵魂，不是死的，而是活的。良善赐给他统管万物的权柄，使他享受并管理它们，甚至给它们命名。此外，良善还为人附加了愉悦，使他作为整个世界的主人，可以逗留在更高的喜乐中，被迁移到乐园，从世界进入教会。这同一位良善也为他预备了一个相称的助手，好使他的命运中没有一样不是好的。因为祂说：那人独居不好。祂深知\Person{玛利亚}的性别，以及教会的性别，对他将是何等的祝福。然而，你们所指责并曲解为争端课题的法律，是良善加给人的，旨在使人幸福，使他紧贴天主，以免他显出卑贱而非自由，降低到其它受他统治的动物水平——那些动物不受天主约束，免于一切烦人的管束——而是作为唯一的人类，夸耀他独自堪当从天主领受法律；并作为有理性的存在，能够理解和知识，被约束在理性自由的界限内，服从于那位将万有置于他权下的主。为确保这法律的遵守，良善也藉助这样的制裁来谋划：\BibleQuote{你吃的日子必定死。}\BibleRef{创 2:17}因为祂如此指明过犯的后果，是极其仁爱的行为，免得对危险的不知鼓励了对服从的忽视。既然这制裁在颁布法律之前就已给出理由，它也构成了后来遵守法律的一个动机，即对违犯附加了惩罚；的确，提出惩罚的那位仍不愿惩罚被实施。所以，在这些事物中和此刻，学习我们的天主的良善吧；从祂卓越的工程、祂慈爱的祝福、祂宽厚的恩赐、祂恩慈的眷顾、祂如此良善而仁慈的法律和警告中学习。\par

% structure: p0011 source=h2 target=subsection reason=relative-heading-rank; base=h2

\subsection{第五章 马尔西翁之挑剔的考量。驳斥其异议，即人的堕落显示天主的失败。人之存在的完美在于其自由，天主特意赋予他自由。堕落归咎于人自己的选择。}

那么现在，你们这些狗，宗徒将你们置于外（见：\BibleRef{默 22:15}），并向真理的天主狂吠，让我们来应付你们的各种问题。这些就是你们不断啃咬的争论骨头！如果天主是良善的，且预知未来，并能避免邪恶，祂为何允许人——祂自己的肖像和模样，并且因灵魂的起源，也是祂的本体——被魔鬼欺骗，而因不服从法律而堕入死亡？因为如果祂是良善的，因此不愿意发生这样的灾难，且预知，因此不会不知道将要发生的事，且有足够的能力阻止其发生，那么那个结果就绝不会发生，这在神圣的伟大三个条件下本是不可能的。然而，既然它发生了，相反的主张就肯定是真实的，即天主既不是良善的，也不是预知的，也不是全能的。因为\Emph{正如}如果天主是祂被传说的那样——良善、预知、全能——就不可能发生这样的结果，\Emph{同样}这个结果实际发生了，因为祂不是这样的天主。作为回应，我们必须首先为造物主中受到质疑的属性辩护——即祂的良善、预知和全能。但我不会在这个问题上久留，因为基督自己的定义（见：\BibleRef{若 10:25}）立刻来帮助我们。必须从工作中获得证据。造物主的工程立刻证明祂的良善，因为它们是良善的，正如我们已经指出的，也证明祂的全能，因为它们是强大的，并且确实是从无中生有的。即使它们是由某种（先存的）物质造成的，如有些人所主张的，它们也同样是从无中生有的，因为它们\Emph{原先}不是现在\Emph{这个样子}。简而言之，它们既伟大因为它们是良善的；天主也同样全能，因为万物都属于祂，因此祂是无所不能的。但我要如何说祂的预知呢？预知的证人就像它所默启的先知一样多。毕竟，我们期望在宇宙的创造者身上看到什么预知的称号呢？因为正是藉着这个属性，祂在指定万物位置时预知了它们，并在预知它们时指定了它们的位置。还有罪本身。如果祂没有预知这一点，祂就不会在死亡刑罚下警告大家谨慎提防它。如果天主具有这些属性，这些属性必然使人不可能也不应该遭受任何邪恶，然而邪恶确实发生了，那么也让我们考虑人的状况——是否是\Emph{它}（人的状况）才是导致那本不能通过天主发生的事情发生的真正原因？然后我发现，人被天主造成自由的，是他自己意志和力量的主人；祂的神圣肖像和模样在人身上没有任何东西比这种本性构造更能体现。因为不是藉着他的面容和身体的轮廓，尽管它们在人性中如此多样，他表达了他对天主形式的相似；而是藉着他从天主自身所获得的那种本质（即属神的，对应天主的形状），以及他意志的自由和力量。他的这种状态甚至被天主随后施加给他的法律所证实。因为法律不会被施加给一个没有能力履行法律所要求的服从的人；同样，死亡的刑罚也不会被用来威胁罪，如果人在他意志的自由中不可能蔑视法律。在造物主后来的法律中，你也会发现，当祂在人面前设立善恶、生死时，整个纪律的进程都安排在天主的诫命中，祂呼召人远离罪，并警告和劝诫他们；而这没有其他理由，只是因为人是自由的，有服从或反抗的意志。\par

% structure: p0013 source=h2 target=subsection reason=relative-heading-rank; base=h2

\subsection{第六章 这种自由就其原始创造而得到辩护；也适于展示天主的良善与目的。如果人是出于必然而非选择而成为善或恶，则赏罚是不可能的。}

但是，尽管从我们的论证可以理解，我们只是如此肯定人对其意志拥有不受束缚的控制，以致发生在他身上的事应归咎于他自己，而非归咎于天主；然而，为了你现在仍不反对，说他不应当被如此构成，因为他的自由和意志能力可能变成有害的，我将首先主张他如此构成是正确的，以便我更有信心地称赞他实际的本性构成，以及这构成本身堪当神性存在这一额外事实；导致人按此构成被创造的原因表明是更好的原因。此外，如此构成的人将受到天主的良善和祂的旨意这两者的保护，这两者在我们的天主身上总是和谐一致。因为祂的旨意没有良善就不是旨意；祂的良善没有旨意也不是良善，除非在\Person{马尔基昂}的神的情形中，祂毫无旨意地良善，正如我们已经表明的。那么，合适的是：天主应该被认识；这无疑是良善而合理的事情。同样合适的是：应该有某种堪当认识天主的事物。有什么比天主的肖像和模样更堪当呢？这也无疑是良善而合理的。因此，合适的是：那位是天主肖像和模样的，应当被造成拥有自由意志和自主；以至于这同一件事——即意志的自由和自制——可以被算作他身上的天主的肖像和模样。为此，一种与此性格相适应的本质被赋予人，即天主的气息，祂自身是自由而不受约束的。但如果你对此另有看法，那么，当人拥有整个世界时，他为何没有首先在自制中统治——作他人的主人，却作自己的奴隶？天主的良善，你可以从祂赐予人的恩赐中学到，而祂的旨意则从祂对万事的安排中学到。目前，让我们只关注天主的\Emph{良善}，它赐予人如此大的恩赐，即他意志的自由。天主的\Emph{旨意}需要另有机会处理，因为它提供同样意义的教导。现在，只有天主在本性上是良善的。因为那位拥有无起始之物的，不是通过受造，而是通过本性拥有它。然而，人完全通过受造而存在，有一个开始，并随着那个开始获得了其存在的形式；因此他并非在本性上倾向于善，而是通过受造，不把良善当作自己的属性，因为（正如我们所说）他倾向于善不是由于本性，而是由于受造，是按照其良善造物主，即一切美善的创作者的规定。因此，为了使人拥有自己从天主得来的良善，并在人身上从此有一个属性，在某种意义上一种良善的本性特质，在他本性的构成中，作为天主赐予他的良善的形式见证，被赋予他意志的自由和能力，这种意志的自由和能力应当使善被人自发地施行，作为他自己的属性，理由是：在将要被自愿行使的良善的事情上，所需的要求恰是如此，也就是说，通过他意志的自由，既无偏袒亦无奴役于其本性的构成；因此，人如果确实按照其本性构成，但仍然作为其意志的结果，作为其本性的属性，来展示他的良善，他就恰好在这一点上为善；并且通过类似的意志活动，他也表明自己在抵御恶方面过于坚强（因为甚至这一点，天主当然也预见了），他是自由的，是自己的主人；因为，如果他缺少这\Emph{自制}的特权，以至于甚至行善也是出于必然而非意志，那么他在其奴役的无助中，就会遭受恶的篡夺，成为恶的奴隶如同善的奴隶一样。因此，完整的意志自由在两种倾向上都被赐予了他；以至于作为自己的主人，他可以通过自发遵守善而不断遇到善，并通过自发回避恶来回避恶；因为，即使人处于其他环境，在天主的审判中，他仍有当然的责任根据其意志的动向（当然，被视为自由的）来行正义。但善或恶的赏报都不能支付给那个被发现出于必然而非选择而行善或行恶的人。真正的法律就在于此，它并不排除，反而通过自发地服从或自发地犯罪来证明\Emph{人的}自由；人的意志对于任一结局的自由是如此明显。因此，既然天主的良善和旨意都在赐予人意志的自由中被发现，那么在忽略了在任何讨论之前必须确定的良善和旨意的原始定义之后，我们就不应该基于后续的事实，竟敢说天主不应当以这种方式造人，因为结果与假定对天主合适的不同。我们更应当在适当考虑了天主应当如此创造\Emph{人}之后，保持这一考虑不受损害，并审视事情的其他方面。无疑，对于为人的堕落感到愤慨但未检查他被造事实的人，很容易在未审查祂的旨意的情况下，将所发生的事归咎于造物主。总之：天主的良善，从祂的工程开始就被完全考虑，足以使我们相信天主不可能产生任何恶；而人的自由，经过再次思考，会向我们表明，只有它才应当为它自己所犯的过错负责。\par

% structure: p0015 source=h2 target=subsection reason=relative-heading-rank; base=h2

\subsection{如果天主以任何方式限制了人的自由，\Person{马尔基昂}本会准备好另一个相反的刁难。人的堕落为天主所预知。天主补救并与其真理和良善一致地做了预备。}

通过这样的结论，一切都完整地保留给天主：祂本性的美善、祂治理与预知的目的，以及祂能力的丰盈。然而，倘若你不再追问是否有什么事可能违背天主的意愿，你就应当从天主的属性中减去祂在全部受造中至高目标的坚定与最卓绝的真实。因为，当你持守这善天主的坚定与真实——它们确实能从理性受造物得到证明——时，你就不会对以下事实感到惊奇：天主没有干预阻止祂不愿发生的事，以便使祂所愿的免受伤害。因为，既然祂一次而永远地（并且正如我们所证明的，相称地）允许了人自由意志和自主权，那么祂凭着自己创造中的权威，确实允许人享受\Emph{这些恩赐}；就祂自己而言，是依照祂作为天主的本性，即为了善（因为谁会允许任何敌对自身的事呢？）；就人而言，是依照他自由的冲动（因为当一个人给予别人某物去享受时，谁不是伴随着许可，让他全心全意、心甘情愿地享受呢？）。因此，必然的结果是，天主必须将祂的预知和能力与祂一次而永远地赐予人的自由分开（换言之，保持在自己内），而祂本可用这些预知和能力阻止人因错误地行使自由而陷入危险。因为，如果祂曾干预，祂就会收回祂有目的、善意地允许的人意志的自由。但是，假设天主干预了；假设祂取消了人的自由，警告人远离那棵树，并阻止那狡猾的蛇与女人交谈；马西昂岂不会高呼：多么轻浮、无常、不忠实的主，竟取消祂已赐予的恩赐！祂为何允许自由意志，若后来又收回？为何允许后又收回？让祂选择在何处给自己烙上错误的印记：是在祂起初造人时，还是在后来取消时？如果祂曾限制人的自由，岂不更显得祂因缺乏对未来的预见而被欺骗了吗？然而，当祂任由自由充分发展时，谁不会说祂是在对事情结局无知的情况下这样做的？天主确实预知人会滥用祂所造的禀赋；然而，有什么能比天主的坚定和祂受造物的真实——无论它们如何——更相称于天主呢？人必须看清，若他未能善用所领受的美好恩赐，他自身如何因不愿遵守的法律而有罪，而不是立法者对祂自己的法律行骗，不允许其诫命得以履行。每当你倾向于对造物主发出这样的指责（这对你最合适），请温柔地在心中为祂想起祂的坚定、忍耐和真实，因为祂使祂的受造物既理性又完善。\par

% structure: p0017 source=h2 target=subsection reason=relative-heading-rank; base=h2

\subsection{第八章  人禀赋自由，超越天使，甚至胜过诱惑他堕落的天使，当他悔改并重新服从天主时}

因为天主造人，不仅是为了让他过自然的生活，而是为了让他过德性的生活，即与天主及其法律相关的生活。因此，当人被塑造成有生命的灵魂时，天主赐予他\Emph{生命}；但当他被要求遵守法律时，天主命令他\Emph{善度生活}。同样，天主表明人并非为死亡而被造，因为祂如今愿意人恢复生命，宁要罪人悔改而非死亡。\BibleRef{则 18:23}因此，正如天主为人安排了生命的状态，人也为自己带来了死亡的状态；而这既非由于软弱，也非由于无知，以致不能归咎于造物主。诚然，诱惑者是一个天使；但被诱惑者是自由的，是自主的；作为天主的肖像和模样，他比任何天使都强大；作为天主的\OriginalTerm{气息}{afflatus}，他也比天使所属的物质性的灵更高贵。经上说：\Quote{祂使祂的天使成为风，使祂的仆役成为火焰。}如果人因过于软弱而不能统治，且低于天使（天主并未将如此臣服者交给天使），天主就不会使万物都服于人；如果人不能承受如此重大的负担，天主也不会将法律的重担加于他；再者，如果天主知道人因无助而无辜，祂也不会用死亡的刑罚来威胁他；简言之，如果祂造人是软弱的，那祂就不会通过意志的自由和独立来造他，而是通过不赋予这些禀赋。因此，即使现在，同一个人、同样的灵魂实体、同样的亚当境况，凭着相同的自由和意志的能力——当意志顺从天主的法律时——也能战胜同一个魔鬼。\par

% structure: p0019 source=h2 target=subsection reason=relative-heading-rank; base=h2

\subsection{第九章  回答另一异议，即堕落可归咎于天主，因为人的灵魂是造物主神性本质的一部分。在人的罪中，天主的\OriginalTerm{气息}{Divine Afflatus}没有过错，而是附加其上的人意志有过错。}

但是，你说，无论造物主的本质以何种方式被认为有瑕疵，当天主的\Emph{气息}，也就是灵魂，在人内犯罪时，部分的过错必可追溯至原始的全体。如今，为应对此异议，我们必须解释灵魂的本性。我们必须从一开始就把握希腊圣经的意义，其中用的是\Emph{气息}，而非\Emph{神}。一些希腊文的诠释者，未思考词语的差异，且不关心其确切含义，以\Emph{神}代替\Emph{气息}；他们以此给异端者提供了用瑕疵玷污天主之神，也就是天主本身的机会。现在问题来了。请注意，\Emph{气息}小于\Emph{神}，虽然它来自神；它是神的微风，但它不是神。微风比风更稀薄；虽然它出于风，但微风不是风。人可称微风为神的肖像。同样，人是天主的肖像，也就是神的肖像；因为天主是神。因此，\Emph{气息}是神的肖像。肖像绝不等同于实物本身。肖似本体是一回事，而本体本身是另一回事。所以，虽然\Emph{气息}是神的肖像，却不能如此比较天主的肖像，即因为本体——也就是神，或换言之，神圣的存在——是无瑕疵的，因此\Emph{气息}，也就是肖像，也绝不可能犯错。在这方面，肖像将小于本体，而\Emph{气息}次于神，因为虽然它无疑拥有神性的真特征，如不朽的灵魂、自由与自主、很大程度的预知、理性、理解与知识的能力，但即使在那些方面，它仍然是肖像，永不及神性的实际能力，也达不到绝对免除瑕疵——这一特性只赐予天主，即本体，而与肖像完全不相容。肖像，虽然可能表达本体的所有特征，但缺乏其内在能力；它没有能动性。同样，灵魂，作为神的肖像，不能表达神的单纯能力，即其幸福的免于犯罪。若非如此，它就不是灵魂，而是神；不是接受了灵魂的人，而是天主。此外，从另一角度看，并非一切属于天主的都被视为天主，因此你不可主张祂的\Emph{气息}就是天主，即免于瑕疵，因为它是天主的嘘气。在你自己的一举动中，比如向一支笛子吹气，你不会因此使笛子成为人，尽管是你自己人的气吹入其中，正如天主吹了祂自己的神。事实上，圣经借着明确说\BibleQuote{天主将生命的气息吹入人的鼻孔，人就成了一个有灵的生物，而不是一个赋予生命的神}\BibleRef{创 2:7}，将那个\Emph{灵魂}与造物主的状况区分开来。作品必然与工师不同，且低于他。陶瓶不会是陶匠，尽管是陶匠所造；同样，\Emph{气息}，因为是神所造，也因此不会是神。灵魂常与气息同名。你也要注意，不要从气息下降到更低的质量。所以，你（你说）承认灵魂的软弱，你之前否认了！无疑，当你要求它与天主同等，即免于瑕疵时，我主张它是软弱的。但是当与天使比较时，我不得不主张那掌管万有的比二者更强，天使是服役的\BibleRef{希 1:14}，他必将审判天使\BibleRef{格前 6:3}，如果他坚守天主的法律——一种他起初拒绝的服从。现在，这不服之罪是天主的\Emph{气息}可能犯的：它是可能的，但并不适当。这种\Emph{可能性}在于它本性的薄弱，因为是气息而非神；而\Emph{不适当性}则出于它的意志能力，因为是自由的，而非奴隶。它还得到了警告的帮助，警告人不可犯罪，以免死亡，这警告旨在支持其薄弱的本性，并指引其自由选择。因此，灵魂不能再显得是因其与天主的亲缘关系而犯罪，即借着气息，而是借着那加于其本性的东西，即它的自由意志，这确实由天主按祂的目的和理性赐予它，但被人类按其选择鲁莽地运用了。既然如此，天主行动的整个过程就摆脱了一切致恶的指控。因为意志的自由不会将自己的过错归咎于赐予者，而是归咎于滥用者。那么，你想要归咎于造物主的邪恶是什么？如果是人的罪，那就不是天主的过错，因为那是人的作为；那位成为罪的禁止者，甚至谴责者，不应被视为罪的作者。如果死亡是邪恶，死亡不会责备威胁它者，而是责备藐视它者。因为人因藐视而引入了死亡，它必定不会出现，倘若人未曾藐视它。\par

% structure: p0021 source=h2 target=subsection reason=relative-heading-rank; base=h2

\subsection{第十章 另一异议被驳斥，即那诱惑人犯罪的魔鬼本身是天主的受造物。不，原始的革鲁宾只是天主的工作。魔鬼的本性是因任性而附加的。在人的恢复中，魔鬼在其自己的战场上被击败。}

然而，你若选择将恶的记述从人转移到那作为罪恶煽动者的\OriginalTerm{魔鬼}{devil}身上，借此也将责任归咎于\OriginalTerm{造物主}{Creator}，因为祂造了魔鬼——因为祂造了那些精神体，即天使——，那么由此推出：受造之物，即天使，将属于造它的那一位；而那并非由天主所造者，那\OriginalTerm{魔鬼}{devil}或\OriginalTerm{控告者}{accuser}，不能不说是由自己造成的；这乃是通过对天主的虚假诬蔑：首先，天主如何禁止他们吃所有树上的果子；然后，假称如果他们吃了就不会死；第三，仿佛天主嫉妒他们具有神性的所有权。那么，这种对人类的谎言和欺骗以及对天主的毁谤的恶意是从哪里来的呢？当然不是来自天主，祂使天使善，正如祂的善工一样。事实上，在他成为魔鬼之前，他显为受造物中最有智慧的；智慧并非恶。你若翻阅\Person{厄则克耳}的预言，你即刻会看出，这位天使受造时本是善的，而因选择而败坏。因为论到\Place{提洛}的君王，是就魔鬼而说的：\Quote{主的话又传给我说：\InnerQuote{人子，你要对提洛的君王作一首哀歌，对他说：主上主这样说：你完美无比，充满智慧，美丽绝伦}}（这属于他，作为最高天使、\OriginalTerm{总领天使}{archangel}、最智慧者）；\Quote{你在你天主的\OriginalTerm{乐园}{paradise}中、在欢乐中出生的，你原来佩带的宝石是：红玛瑙、黄玉、金刚石、水苍玉、条纹玛瑙、玉髓、青金石、绿宝石和红玉；用黄金装满你的仓库和宝库。从你受造之日，在你被安置在\OriginalTerm{革鲁宾}{cherub}中，在天主的圣山上，你在火石中行走，你的日子无可指摘，从你受造之日，直到在你身上发现了不义。由于你生意繁荣，你充满了仓库，你犯了罪，}等等。显然，这段描述正属于天使的悖逆，而非君主：因为没有任何人生在天主的乐园里，连\Person{亚当}本人也不例外，他是被移植到那里的；也没有人与一个\OriginalTerm{革鲁宾}{cherub}一同被安置在天主的圣山上，即高天之上，主曾见证\OriginalTerm{撒殚}{Satan}从那里坠落；也没有被拘留在火石和燃烧星座的闪光之中，\OriginalTerm{撒殚}{Satan}如闪电般从那里被击落。\BibleRef{路 10:18}不，那正是罪恶本身的作者，以一个有罪之人的形象被指称：他一度无可指摘，在他受造之时，被天主塑造为善，如同善的\OriginalTerm{造物主}{Creator}创造无可指摘的受造物，并饰以一切天使的荣耀，与天主联合，善与善同在；但后来他自愿转向恶。\Emph{从你的不义被发现的那一天}，他说，\Emph{它们被发现}，——将那些伤害归咎于他，就是当他被逐出对天主的忠诚时对人造成的伤害——从那时起他犯了罪，当他传播罪恶时，并由此运用\Quote{他生意繁荣}，即他的邪恶，甚至他的过犯的数目，因为他自己作为一个精神体，不比（人）逊色，被造时拥有\OriginalTerm{自由意志}{free-will}。因为天主绝不会不赋予一个将要与祂相近的存在以这种自由。然而，通过预先定他的罪，天主证明他因自行渴望内里自发产生的邪恶而偏离了受造本性的状态；同时，通过允许他的计划运作，祂一贯地行使自己的美善，推迟魔鬼的毁灭，与推迟人的恢复同样理由。因为祂为一场战斗留出了空间：在这场战斗中，人可以用使他屈服于仇敌的同样自由意志来粉碎仇敌（证明过错全在自己，而非天主），从而通过胜利相称地恢复救恩；在这场战斗中，魔鬼也可能受到更苦的惩罚，因为他被他先前伤害的人所征服；在这场战斗中，天主显为更加美善，因为祂等候人从现世生活返回一个更荣耀的乐园，有权去摘生命树上的果子。\par

% structure: p0023 source=h2 target=subsection reason=relative-heading-rank; base=h2

\subsection{第十一章 若人在犯罪后天主行使祂的公义与审判的属性，这符合祂的美善，并增强对天主性格完美的真实认识。}

因此，从起初直到人堕落之前，天主仅是良善的；此后祂成为一位严厉的法官，并且如马西昂派所愿意的，残忍的。女人立即被判在痛苦中生产，并服侍她的丈夫（\BibleRef{创 3:16}），尽管她先前曾听到祝福中毫无痛苦地宣告种族繁衍：\Quote{\Emph{你们要生育繁殖}}，并且她本应成为男性的助手而非奴隶。大地也立即受了诅咒（\BibleRef{创 3:18}），而它先前是蒙祝福的。立即长出荆棘和蒺藜，而那里曾经生长青草、菜蔬和结果实的树木。立即为了面包而有了汗水和劳动，而先前每棵树上都提供着自发的食物和不需劳作的滋养。从此以后，人\Quote{人\Emph{归于}尘土}，而不像先前\Quote{\Emph{出于}尘土}；从此以后\Quote{\Emph{归于死亡}}，而先前是\Quote{\Emph{归于生命}}；从此以后穿着皮衣，而先前赤身露体而无羞耻。因此，天主的先前良善出于本性，祂后来的严厉则由一个原因引起。一个是内在的，另一个是偶然的；一个是祂自己的，另一个是适应的；一个从祂发出，另一个被祂接受。但是，\Emph{本性}不能合理地允许祂的良善一直不运作，\Emph{原因}也不能允许祂的严厉以伪装或隐藏的方式逃脱。天主将前者为自己预备，将后者为时机预备。现在你们应当着手证明法官的角色与邪恶结盟——你们一直梦想着另一位纯粹良善的神，仅仅因为你们无法\Emph{理解天主是一位法官}；尽管我们已经证明天主也是一位法官。或者，如果祂不是法官，那么祂就是一个反常且无用的纪律的创始者，这纪律不应被辩护——换句话说，不应被审判。然而，你们没有证明天主是法官，你们也没有证据表明祂是法官。你们无疑将不得不控告正义本身——它提供了法官——或者把她归入邪恶的种类，也就是说，把不义加到良善的称号上。但那样，正义就是邪恶，如果不义是良善的话。然而你们被迫宣称不义是最坏的事物之一，同理也被迫将正义归入最卓越的事物。因为没有什么与邪恶为敌的不是良善，也没有什么与良善为敌的不是邪恶。由此，正如不义是邪恶，同样地正义是良善。它不应被仅仅视为良善的一个种类，而是良善的实际遵守，因为良善（除非正义被如此控制以致成为公正的）如果不公正就不会是良善。因为没有什么不义的是良善；而另一方面，凡是公正的都是良善。\par

% structure: p0025 source=h2 target=subsection reason=relative-heading-rank; base=h2

\subsection{第十二章 良善与正义的属性不应被分离。它们在真天主内是相容的。正义在天主性中的功能描述。}

既然良善与正义之间有这种联合与一致，你们不能规定它们的分离。你们有何面目决定将你们两位神分开，认为在他们各自的状态中，一位是独特的良善神，另一位是独特的正义神？哪里存在正义，那里也存在良善。简言之，从起初，创造主既是良善的，也是正义的。祂的两种属性一同前进。祂的良善创造了世界，祂的正义安排了世界；在这个过程中它甚至当时就裁定世界应由好的材料形成，因为它与良善商议。正义的工作是显而易见的，即在光与暗、昼与夜、天与地、上面的水与下面的水、海的聚积与干地的块体、较大的光体与较小的光体、昼的光体与夜的光体、男与女、知识树与死亡和生命树、世界与乐园、水生动物与陆生动物之间所宣告的分别。正如良善构思了一切，同样正义区分了一切。随着后者的决定，一切都被安排和整理有序。每个元素的地点、性质、效果、运动、状态，以及每个的升起和落下，都是创造主的正义决定。不要以为祂作为法官的职能必须被定义为从邪恶开始时才开始，从而用邪恶的原因玷污祂的正义。通过这样的考虑，我们表明这一属性与良善（万物的创造者）一同前进——正义本身也应被视为内在和本性的，而非偶然地归于天主，因为她被发现存在于祂——她的主——那里，作为祂工作的仲裁者。\par

% structure: p0027 source=h2 target=subsection reason=relative-heading-rank; base=h2

\subsection{第十三章 对天主正义的进一步描述；自人堕落后它调节了天主的良善。天主对我们爱的要求与对我们敬畏的要求得以调和。}

然而，当恶后来爆发，天主的美善开始有对手与之争战时，天主的正义也获得了另一项功能，即按照人对其美善的祈求来引导它。结果是：天主的美善原本自由地流淌，天主自发地行善，现在却被打断，转而按照每个人的功德来施予；它赐予配得的人，拒绝不配的人，从忘恩者那里收回，并且报复所有它的敌人。因此，正义的全部职能在这方面成为美善的媒介：凡它藉审判所谴责的，凡它藉谴责所惩罚的，凡它（用你的话说）无情追逼的，实际上是以善相待而非加害。因为对审判的恐惧有助于善，而不是恶。现在，善与敌人争战，单凭自身不足以推荐自己。无论如何，如果它能做到\Emph{如此之多}，它也无法站稳脚跟；因为敌人使它失去了坚固性，除非有某种恐惧的力量介入，迫使那些极不情愿的人也追求善并加以珍惜。但是，当如此多的恶的诱惑攻击人时，谁会渴望那可以藐视而不受惩罚的善呢？又有谁会珍惜那可以失去而无危险的东西呢？你读过通向恶的道路是多么宽阔，\BibleRef{玛 7:13}相比之下走向善的道路是多么狭窄；难道若不是因为那里有可怕的东西，所有的人不都会滑下那条路吗？我们畏惧造物主那令人战栗的威胁，却仍然难以离开恶。假如祂不威胁呢？你将这正义称为恶吗？因为它完全不利于恶？你会否认它是善吗？因为它着眼于善？你希望天主成为什么样的存在？是否应该更愿意祂成为那样的，以致罪恶在祂之下猖獗，而魔鬼嘲弄祂？你会认为祂是一位善的天主，祂却能让人的罪因安逸而更糟吗？谁是善的作者，除非是那也要求善的？同样，谁是恶的陌生人，除非是那与恶为敌的？谁是与恶为敌的，除非是那战胜恶的？谁是战胜恶的，除非是那惩罚恶的？因此，天主完全是善的，因为祂在一切事上都站在善的一边。事实上，祂是全能者，因为既能帮助也能伤害。仅仅施益是比较小的事，因为它不能做别的，只能行善。对于这样的行为，我怎能满怀信心地期待善，如果它只有这种能力？我怎能追随无罪的奖赏，如果我不顾念恶行的报应？我不得不怀疑，那资源不足以同时应对两种情况的，是否可能无法酬报其中一种或另一种。至此，正义正是天主自身圆满的丰满，显明天主既是完美的父亲，又是完美的主：作为父亲，祂是仁慈的；作为主，祂是训导的；父亲在祂权能的温和中，主在祂权能的严厉中；父亲必须以孝爱之情来爱，主则必须敬畏；爱祂，因为祂喜欢仁爱胜过祭献，\BibleRef{欧 6:6}敬畏祂，因为祂厌恶罪恶；爱祂，因为祂喜欢罪人的悔改胜过他的死亡，\BibleRef{则 33:11}敬畏祂，因为祂厌恶不悔改的罪人。因此，神圣的法律就这两方面规定了职责：\Emph{你当爱天主}，和\Emph{你当敬畏天主}。它为顺从的人提出了一个，为悖逆的人提出了另一个。\par

% structure: p0029 source=h2 target=subsection reason=relative-heading-rank; base=h2

\subsection{第十四章 两种恶：刑罚性的与罪孽性的。天主不是后一种恶的作者，而只是前一种（刑罚性的）恶的作者，并且这恶包含在祂的正义之内。}

在所有场合，天主都与你会面：祂击伤，也治愈；祂使人死，也使人活；祂使人卑微，也使人高举；祂\Quote{制造灾祸}，也\Quote{缔造和平}；——这样，从祂眷顾的这些（对比）中，我就能得到对异端者的回应。看，他们说，祂在经文中承认自己是灾祸的创造者：\Quote{是我制造灾祸}。他们取了一个词，其单一形式将两种灾祸（因为罪和惩罚都称为\Emph{灾祸}）混淆和模糊，并且要我们在每处经文都理解为祂是所有恶事的创造者，以便将祂指定为恶的作者。相反，我们区分该词的两层含义，将罪的恶与罚的恶分开，即\OriginalTerm{罪之恶}{mala culpæ}与\OriginalTerm{罚之恶}{mala pœnæ}，将每一类归于各自的作者——魔鬼是罪的恶\OriginalTerm{罪}{culpæ}的作者，天主是罚的恶\OriginalTerm{罚}{pœnæ}的创造者；这样，一类被视为道德上的恶，另一类则被归类为正义针对罪的恶执行惩罚判决的行动。对于后一类与正义相容的恶，天主因而公开承认是创造者。它们对于忍受它们的人来说无疑是恶，但就其本身而言是好的，因为是正义的、维护善的、敌视罪的。在这方面，它们也配得上天主。否则，就证明它们是不义的，从而表明它们属于罪恶的一类，即不义的恶；因为如果它们属于正义，它们就不再是恶事，而是好事——只对坏人来说是恶，就连直接的好事也被他们谴责为恶。在这种情况下，你必须断定：人，虽然故意藐视神圣律法，却是不公正地承受了他本想逃脱的判决；那时代的邪恶被洪水、后来被（\Place{索多玛}的）火不公正地击打；\Place{埃及}，尽管极其堕落和迷信，而且更糟的是，它骚扰寄居的百姓，却被十灾的惩罚不公正地击打。\Emph{天主}使\Person{法郎}的心硬。然而，他理应被引导走向毁灭，因为他已经否认了天主，已经多次骄傲地拒绝了祂的使者，给祂的百姓加重了重担，（总之）作为一个埃及人，他早已在外邦偶像崇拜上得罪了天主，崇拜朱鹭和鳄鱼胜过永生的天主。甚至祂自己的百姓，天主也因他们的忘恩负义而惩罚了他们。祂还派熊攻击少年，因为他们对先知不敬。\par

% structure: p0031 source=h2 target=subsection reason=relative-heading-rank; base=h2

\subsection{第十五章 天主的严厉与理智和正义相容。施行时，并非任意，而是为了纠正。}

那么，在一切事上，首先要仔细考虑审判者的正义；如果其目的明确，那么它的严厉以及严厉在过程中的运作，就显得与理智和正义相容。现在，为了不在此点上拖延太久，（我要求你）也提出其他理由，以便你能谴责\Emph{审判者的}判决；减轻罪人的过犯，以便你能指责对他的司法定罪。不要理会谴责审判者，而要证明祂是不义的。好吧，即使祂要求父亲们的罪由子女承担，人民的硬心使得这种纠正措施对他们成为必要，以便他们顾及后代而遵守神圣律法。因为谁不是更关心自己的子女而不是自己呢？再者，如果父亲的祝福也注定归于后裔，远在这些后裔有任何功德之前，那么父亲的罪为何不能也归于子女呢？正如恩宠，过犯也是如此；恩宠和过犯同样传递整个种族，但确实保留了一项后来的法令，据此可以避免说：\Quote{父亲吃了酸葡萄，孩子的牙就发酸} \BibleRef{耶 31:29} 换句话说，父亲不应担当儿子的罪孽，儿子也不应担当父亲的罪孽，而每个人都应承担自己的罪；这样，律法的严苛在人民的硬心之后有所缓和，正义不再审判\Emph{整个民族}，而是审判个人。然而，如果你接受真理的福音，你就会发现，当审判者将父亲的罪报应在儿子身上时，这判决落在了谁身上——就是那些（足够刚硬）以致自动诅咒自己承担这种定罪的人：\Quote{祂的血归在我们和我们的子孙身上} \BibleRef{玛 27:25} \Emph{这}，因此，天主的眷顾在整个过程中安排了这事，正如它曾听见的那样。\par

% structure: p0033 source=h2 target=subsection reason=relative-heading-rank; base=h2

\subsection{第十六章 天主的严厉拥有附属的品质，与正义相容。若将人的情感归于天主，则不可用人的不完美尺度来衡量。}

那么，就连祂的\OriginalTerm{严厉}{severity}也是善的，因为它是正义的：当\OriginalTerm{审判者}{judge}是善的时，那就是正义的。其他属性同样也是善的，藉由这些属性，善的严厉的善工得以完成其历程，无论是\OriginalTerm{忿怒}{wrath}、\OriginalTerm{妒忌}{jealousy}还是严峻。因为所有这些对严厉而言如同严厉对\OriginalTerm{正义}{justice}一样不可或缺。一个本该敬畏的时代，竟不知羞耻，必须受到惩罚。因此，那些属于审判者的属性，当它们实际上无可责备时，如同审判者本身一样，绝不能被当作过失加诸于祂。假若你认为医生是必要的，却指责他的工具，因为它们切割、烧灼、截肢或收紧；然而没有专业工具，任何有值的医生都不存在——那会怎么说呢？你若要指责，不妨指责那位切割不当、截肢笨拙、烧灼鲁莽的行医者；甚至指责他的器械是粗糙的工具。你的行为同样不合理，因为你确实承认天主是审判者，但同时却摧毁了祂执行审判职能的那些运作和性情。我们是由先知和基督教导认识天主，而不是由哲学家或\Person{伊壁鸠鲁}。我们相信天主确实曾在世上生活，并为了人的救恩而取了卑微的人性形态，我们远非像那些拒绝相信天主关心任何事物的人那样思考。这种论调是怎样传到异端者中的呢：如果天主发怒、妒忌、激动、忧伤，那么祂必定会败坏，因而必定死去。然而幸运的是，相信天主确实死了，却又永远活着，甚至也是基督徒信经的一部分。他们的愚妄至极，竟从人的事物来预断神圣事物；以致因为人在败坏的状态中发现这类情欲，所以就必须在神中也存在同样的感受。要分辨本性，并给它们各自的感受，这些感受如同它们的本性所要求的那样多样，尽管它们似乎有共同的名称。的确，我们读到天主的右手、眼睛和脚：但这些不能与人的相比较，因为它们联系在同一个名称上。现在，神性的身体与人的身体之间的差异有多大，尽管它们的肢体使用相同的名称，神性灵魂与人的灵魂之间的差异也将同样大，尽管它们的感受用相同的名称来指示。在人性中，这些感受因人的实体之可朽坏而变得腐败；在天主内，它们因神圣本质的不朽坏而成为不朽坏的。你真的相信\OriginalTerm{创造者}{Creator}是天主吗？当然，你回答。那么，你怎能认为在天主内有任何属于人的东西，而不是全部是神圣的呢？你既然不否认祂是天主，你就承认祂不是人；因为你承认祂是天主时，实际上已经断定祂无疑不同于任何一种人的状况。此外，虽然你与他人一样承认人是天主嘘气而成为有生命的灵魂，而不是天主由人而来，但你却是明显荒谬地将人的特征置于天主内，而不是将神圣的特征置于人内，并将天主穿上了人的肖像，而不是把人放在天主的肖像中。因此，这应被视为人在天主内的相似：人的灵魂具有与天主相同的情感和感受，尽管它们不是同一种类；根据它们的本性，在状况和结果上都有不同。然后，关于相反的感受——我指的是温良、忍耐、慈悲，以及它们之母——良善——你为什么根据人的品质来形成对这些神圣表现的意见呢？因为我们确实不完全拥有它们，因为只有天主才是完美的。同样，对于其他感受——即忿怒和恼怒——我们并没有以那么幸福的方式被它们影响，因为只有天主才是真正幸福的，由于祂不朽坏的特性。祂可能会发怒，但不会恼怒，也不会受到危险的诱惑；祂会被触动，但不会被颠覆。祂必须使用所有手段，因为所有偶然情况；有多少原因就有多少感受：为恶人而发怒，为忘恩负义者而义愤，为骄傲者而妒忌，以及其他一切阻碍邪恶的东西。同样，为迷途者而慈悲，为不悔改者而忍耐，为有功德者而卓越的恩赐，以及其他一切对善者必需的东西。祂以祂自己特有的方式被所有这些情感所感动，这方式与祂应当被感动的方式极为相宜；而且，正是由于祂，人也以同样属于他自己的方式被相似地感动。\par

% structure: p0035 source=h2 target=subsection reason=relative-heading-rank; base=h2

\subsection{第十七章 在历史中追寻天主的治理与祂的诫命，你将发现它充满祂的良善。}

这些考量表明，天主作为审判者的整个秩序是有所作为的，并且（容我用更相称的言辞表达）是保守祂大公至善之善的——这善既远离司法情绪，又在其自身状态中纯然无染——但\OriginalTerm{马西翁派}{Marcionitae}却拒绝承认它存在于同一天主之内。\Quote{降雨给义人，也给不义的人；使祂的太阳上升，光照恶人，也光照善人}——这是任何别的神都未曾施行的恩惠。诚然，马西翁胆敢从福音中删去基督对\OriginalTerm{造物主}{Creator}的这一证言；然而世界本身刻有\Emph{造物主的良善}，这铭文被每个人的良心读得出来。不仅如此，造物主这长久的忍耐将导致马西翁被定罪；我说的是那忍耐，它等待罪人悔改而非死亡，喜爱仁慈胜过祭献\BibleRef{欧 6:6}，它使尼尼微人免于已向他们宣布的毁灭\BibleRef{纳 3:10}，因希则克雅的眼泪而恩赐他延年益寿\EditorialNote{列王纪下 20:1}，并在巴比伦王彻底悔改后恢复他的王位\BibleRef{达 4:33}；还有那仁慈，它应百姓的恳求救回了即将被处死的撒乌耳之子\BibleRef{撒上 14:45}，并在达味承认自己因乌黎雅家所犯的罪时赐予他无条件的宽恕\BibleRef{撒下 12:13}；它一再定罪以色列家，也一再复兴它，责备与安慰同样频繁。因此，不要仅仅视天主为审判者，也要注意祂作为\OriginalTerm{至善者}{the Most Good}的行为典范。当你注意祂施行报复时，也当考虑祂显示仁慈时。在天平上，以祂的温柔反对祂的严厉。当你发现两种品质在造物主中共存时，你就会在祂身上找到那使你想到有另一位神的情形。最后，来考察祂的道理、训导、戒律和劝谕。你也许会说不难从人类法律中找到同样美好的规定。但梅瑟和天主存在于你们所有的\OriginalTerm{来库古}{Lycurgus}和\OriginalTerm{梭伦}{Solon}之前。没有一个后来时代不是源自原始源泉。无论如何，我的造物主并非从你们的神那里学来颁布这样的诫命：不可杀人；不可奸淫；不可偷盗；不可作假见证；不可贪恋你邻人的财物；孝敬你的父亲和你的母亲；以及，你应当爱邻人如同自己。除了这些关于纯洁、贞洁、正义和虔敬的首要大诫之外，还加上了人道的条例，例如每第七年释放奴隶得自由；同一时期土地休耕；也给予贫民一份；并且从踹谷的牛口上除去笼头，让它在劳苦之前享受劳动果实，好叫先在动物身上显示的仁慈能从这些初级操练提升到对人的安慰。\par

% structure: p0037 source=h2 target=subsection reason=relative-heading-rank; base=h2

\subsection{第十八章 辩护天主的某些法律是良善的——这些法律曾被\OriginalTerm{马西翁派}{Marcionitae}攻击，例如\OriginalTerm{同态复仇法}{Lex Talionis}。从社会与道德观点看此法及其他各类条例的有益目的。}

然而，我能以更大信心辩护法律中哪些部分为善，胜过异端所如此渴望的部分呢？——就是那报复法令，要求以眼还眼，以牙还牙，以伤还伤。\BibleRef{出 21:24} 现在这里没有任何许可互相伤害的意味；相反，总体上是遏制暴力的规定。对于一个非常顽固、对天主缺乏信德的民族而言，期待从天主那里得到后来由先知所宣告的报复，似乎令人厌烦，甚至难以置信：\Quote{报仇是我的事；我必报复，主说。} 因此，暂时间，错误的行为要因对即将来临的报应的恐惧而被遏制；所以这报应的许可乃是对挑衅的禁止，从而可以制止一切冲动的伤害，而通过第二种的许可，第一种因恐惧而被阻止，并通过这第一种的威慑，第二种便不会犯下。通过同一法律也达到另一个结果，即由于其中所含的激情滋味，更容易激发对报复的恐惧。没有比承受你加于他人的痛苦本身更痛苦的事了。再者，当法律从人的食物中有所剥夺，宣布某些曾受祝福的动物为不洁时，你应当明白这是鼓励节制的措施，并从中看出一个缰绳，加在了那虽吃着天使的食物，却贪求埃及的黄瓜和甜瓜的食欲上。也要在其中认出一种预防措施，针对那些食欲的伴侣，即色欲和奢侈，它们通常因食欲的节制而冷却。因为\Quote{百姓坐下吃喝，起来玩耍。}\BibleRef{出 32:6} 此外，为了抑制对金钱的热切渴望（只要它是由食物的需要所引起），昂贵的肉食和饮品的欲望被从他们的能力中除去。最后，为了使人更容易被天主教育以守斋，他习惯了那些既不丰富也不奢侈、并且不太可能纵容奢侈者食欲的食物。当然，造物主更该受责备，因为祂是从自己的百姓那里拿走食物，而不是从更忘恩负义的马西翁派那里。至于那些沉重的祭献，以及他们仪式和供物的繁琐苛责，无人应当指责它们，仿佛天主特别为自己要求这些：因为祂明确地问：\Quote{你们许多的祭献于我有什么益处？} 以及\Quote{谁向你们手中要求这些？}\BibleRef{依 1:11} 但他应在此看出天主方面的审慎安排，这表明祂愿意用那时代的迷信所包含的那种服务，将倾向于偶像崇拜和过犯的百姓绑在祂自己的宗教上；祂可以借此将他们从那里召回来，同时要求他们将这服务归给祂自己，仿佛祂愿意在制造偶像上不犯任何罪。\par

% structure: p0039 source=h2 target=subsection reason=relative-heading-rank; base=h2

\subsection{第十九章. 法律的细致规定意在使百姓依赖天主. 天主为彰显其美善而派遣先知. 引用许多优美经文以说明此属性.}

但即使在日常生活的交往中，在家庭和公共的人际来往中，甚至对最小器皿的照料，祂也以各种可能的方式做出了明确的安排；这样，当他们随处遇到这些法律指示时，就绝不会在任何时刻脱离天主的视线。因为有什么比\BibleQuote{以主的法律为乐？}更能使人幸福呢？\BibleQuote{他日夜默思这法律。}这法律的颁布者并非出于严厉而颁布它，而是为了最高仁爱的利益，旨在克服这民族的硬心，并通过劳苦的服事，凿出一份（尚未）经历过顺服的忠诚：因为我刻意不去触及法律在灵性和预言关系中的神秘意义，以及其中充满几乎各种各类的预表。目前，只需知道它简单地将人系于天主，因此除了不愿事奉天主的人之外，没有人应当指责它。为了促进法律这仁慈而非沉重的目的，先知也由同一的天主的良善所设立，教导配得上天主的诫命，即人应当\BibleQuote{停止作恶，学习行善，寻求正义，为孤儿伸冤，为寡妇辩护：}\BibleRef{依 1:16}喜爱天主的规劝；避免与恶人接触；\BibleQuote{释放受压迫者：}\BibleRef{依 58:6}撤销不义的判决，\BibleQuote{将你的饼分给饥饿的人，将无家可归的穷人接到你家中，见赤身露体的人就给他衣服，不要隐藏自己不看你的骨肉。}\BibleQuote{禁止舌头不出恶言，嘴唇不说诡诈的话；要离恶行善，寻求和睦，一心追赶：}生气却不要犯罪；就是说，不要停留在愤怒中，也不要发怒：\BibleQuote{不从恶人的计谋，不站罪人的道路，不坐亵慢人的座位。}那么在哪里呢？\BibleQuote{看，弟兄们同居共处，是多么美好，多么愉快！}他们（正如他们所行）日夜默思主的法律，因为\BibleQuote{依赖主，胜过依赖人；投靠主，胜过投靠王侯。}因为人将从天主得到什么报偿呢？\BibleQuote{他要像一棵树栽在溪水旁，按时候结果子，叶子也不枯干；凡他所做的尽都顺利。}\BibleQuote{手洁心清，不妄称天主的名，不对邻人起假誓的人，他必从主得蒙福佑，从他救恩的天主得蒙怜悯。}\BibleQuote{主的眼目看顾敬畏祂的人，看顾仰望祂慈爱的人，要救他们的灵魂脱离死亡，}即永死，\BibleQuote{并在他们饥饿时养他们，}即渴望永生。\BibleQuote{义人多有苦难，但主救他脱离这一切。}\BibleQuote{主的圣人之死，在主眼中极为宝贵。}\BibleQuote{主保全他们一切的骨头，连一根也不折断。}主必救赎祂仆人的灵魂。我们从造物主的大量圣经中引用了这少数引文；我想无须更多，足以证明祂是至善的天主，因为它们充分表明了祂良善的诫命及其初果。\par

% structure: p0041 source=h2 target=subsection reason=relative-heading-rank; base=h2

\subsection{第二十章 马西昂派指控天主煽动希伯来人掠夺埃及人。论该神圣安排中的辩护。}

但是，这些（异端者的）\Quote{\OriginalTerm{无礼的墨鱼}{saucy cuttles}}，以它们为形象，关于可食之物的法律（\EditorialNote{申命纪 14}）禁止了这种鱼类食物，一旦他们发现自己被驳倒，便喷出他们亵渎的黑色毒液，并四处散布他们所各自怀有的目标（如今已显而易见），当他们发出这样的断言和抗议时，意图遮蔽并玷污造物主恩惠重新点燃的光明。然而，我们将追踪他们邪恶的计谋，甚至穿越这些黑云，将他们的暗中诽谤的伎俩拖到光天化日之下，特别强调地将欺诈和偷窃金银的罪名加给造物主——希伯来人奉祂命令向埃及人施行此事。来吧，不幸的异端者，我甚至召你作证；先看看这两个民族的情况，然后你将对命令的发出者形成判断。埃及人对希伯来人提出对这些金银器皿的要求。希伯来人则提出反诉，声称根据他们各自祖先的契约，并经双方书面协议的证明，他们劳苦奴役的欠薪应当偿还给他们，为了他们辛苦制造的砖块，以及他们所建造的城市和宫殿。你将如何裁决呢？你这发现\OriginalTerm{至善的天主}{most good God}的发现者？是希伯来人必须承认欺诈，还是埃及人必须给予补偿？因为他们坚持说，双方辩护律师已经这样解决了问题：埃及人要求归还器皿，希伯来人要求回报他们的劳动。但尽管他们这样说，埃及人当时当地正当地放弃了他们要求归还的主张；而希伯来人直到今日，不顾马西翁派，仍然重新主张要求更大的赔偿，坚持说，无论他们借出的金银数量有多大，也不足以补偿，即使六十万人的劳力每人每天只值\Quote{一个法寻}。然而，哪一方数量更多——是要求器皿的人，还是居住在宫殿和城市中的人？再者，哪一样更大——埃及人对希伯来人的怨恨，还是他们向他们显示的\Quote{恩惠}？难道自由人被贬为奴役劳动，是为了希伯来人仅仅得以通过法律诉讼向埃及人追讨伤害；还是为了让他们的长官坐在长凳上，展示他们被残忍鞭打而可耻撕裂的背部和肩膀？并不是凭着几个盘子和杯子——无疑总是属于更少的富人的财产——任何人会断言补偿本应判给希伯来人，而是凭着这些人的全部资源和全体人民的贡献。因此，如果希伯来人的案子是正当的，造物主的案子也必然是正当的；也就是说，祂的命令，当祂既使埃及人无意识地感恩，又在他们迁移时通过一种默示的回报——\Emph{他们长期奴役的}微薄安慰——给予祂的子民全然的解放。祂命令索取的显然少于他们所应得的。埃及人还应当把他们的男婴也归还给希伯来人。\par

% structure: p0043 source=h2 target=subsection reason=relative-heading-rank; base=h2

\subsection{第21章 安息日之法律解释。围绕耶里哥的八天游行。拾柴的违犯。}

同样，在其他方面，你也责备祂反复无常、不稳定，因为祂的诫命自相矛盾，例如祂禁止在安息日做工，却在围攻耶里哥时命令约柜围绕城墙行走八天；换句话说，当然，实际上是在安息日。然而，你不考虑安息日的法律：它禁止的是人的工作，而非天主的工作。（\BibleRef{出 20:9}）因为经上说：\BibleQuote{六天你要劳作，做你一切的工作；但第七天是主你天主的安息日：在这一天你不可做任何工作。}什么工作？当然是你的工作。结论是，祂从安息日移除了那些祂此前在六天中所命令的工作，也就是你自己的工作；换言之，日常生活的属人工作。现在，抬约柜绕行显然不是一项普通的日常职责，也不是属人的工作，而是一项罕见的神圣工作，并且，由于当时是由天主的直接命令所吩咐的，是属天主的工作。我本可以充分解释这预表什么，若不是要详细展开造物主所有证据的形式会是一个冗长的过程，而且你很可能拒绝允许。更切题的是，如果你在明显的事情上被真理的单纯而非好奇的推理所驳倒。因此，在目前的例子中，关于安息日禁止人的劳动而非天主的劳动，有一个清楚的区别。相应地，那个在安息日出去捡柴的人被处以死刑。因为他做的是他自己的工作，而这是法律所禁止的。然而，那些在安息日抬约柜绕行耶里哥的人，却安然无恙。因为他们做的不是他们自己的工作，而是天主的工作，并且是奉祂明确的命令。\par

% structure: p0045 source=h2 target=subsection reason=relative-heading-rank; base=h2

\subsection{第22章 铜蛇和金革鲁宾并未违反第二条诫命。它们的意义。}

同樣，當祂禁止製造天上、地上及水中萬物的形象時，也說明了理由，作為禁止一切隱藏偶像崇拜的物質展示。因為祂補充說：\Quote{你不可向他們叩拜，也不可事奉他們。}然而，主後來命令梅瑟製造的銅蛇形象，並未給偶像崇拜提供藉口，而是為了醫治被火蛇咬傷的人。\BibleRef{戶 21:8}我對這醫治所預表的事緘默不言。同樣，金革魯賓和色辣芬也純粹是約櫃雕飾中的一種裝飾；它們的裝飾出於完全脫離偶像崇拜情況的理由，而正是因著偶像崇拜，製造形象才被禁止；它們顯然不與這禁令相抵觸，因為它們並不在禁令所指的那種形象形式中。我們已經談論過祭獻的合理設立，作為使他們的敬拜從偶像轉向天主；若祂後來拒絕了這敬拜，說：\Quote{你們獻那麼多祭品，對我有什麼用處？}\BibleRef{依 1:11}——祂無非是要我們明白，祂從未真正為自己要求這種敬拜。因為祂說：\Quote{我不吃公牛的血肉}；在另一處：\Quote{永恆的天主既不饑也不渴。}雖然祂曾垂顧亞伯爾的祭品，聞到諾厄全燔祭的馨香，但祂又能從羊肉或焚燒犧牲的氣味中得到什麼喜悅呢？然而，那些獻上從天主領受的、無論是食物還是馨香之氣的人，其純樸而敬畏天主的心，蒙天主悅納，作為對天主表示敬意，天主並不那麼需要所獻之物，而是需要那促使奉獻的心。現在設想，某個依附者向一無所缺的富人國王獻上微不足道的禮物，禮物的數量和品質會使富人和國王蒙羞，還是敬意的考量會使他們歡喜？然而，若這個依附者，無論是自願還是奉命，向他獻上適合其身份的禮物，並遵守對國王應有的禮儀，卻沒有信德和純潔的心，也不準備履行其他順服的行為，那位國王或富人豈不會因而喊道：\Quote{你們獻那麼多祭品，對我有什麼用處？我厭煩了你們的慶節、你們的節日和安息日。}祂稱它們為\Emph{你們的}，因為是奉獻者按自己的意願，而非按天主的宗教而行（因為他把它們當作自己的，而非天主的），\Emph{全能者在此處}表明，祂拒絕那些祂曾命令向祂獻上的祭品，是多麼合乎情況，又多麼合理。\par

% structure: p0047 source=h2 target=subsection reason=relative-heading-rank; base=h2

\subsection{第二十三章 \Person{天主}在揀選和棄絕同一人（如\Person{撒烏耳}王）上的目的，答覆馬爾基翁派的挑剔}

現在，雖然你堅持認為祂在對待人的事情上變化無常，有時在該贊許的地方不贊許；或者缺少先見，將贊許賜給本該被棄絕的人，好像祂不是責備自己過去的判斷，就是無法預見將來的判斷；然而，即使是一個好法官，也沒有什麼比根據當下的功過而既棄絕又揀選更一致的了。\Person{撒烏耳}被揀選，\EditorialNote{撒慕爾紀上第九章}但他尚未輕視先知\Person{撒慕爾}。\EditorialNote{撒慕爾紀上第十三章}\Person{撒羅滿}被棄絕；但他已成為外方婦女的獵物，摩阿布和漆冬偶像的奴隸。造物主必須做什麼，才能逃脫馬爾基翁派的指責？祂必須提前定罪那些目前行為端正的人，因為將來的過犯嗎？但預先定罪尚未該受定罪的人，不是好天主的標記。那麼，祂必須拒絕驅逐罪人，因為他們先前的好行為嗎？但一個公義的法官，不會因不再實踐的先前德行而赦免罪惡。如今，人中有誰如此無可指責，以致天主能一直揀選他，而永遠無法棄絕他？另一方面，誰又如此缺乏任何善行，以致天主能永遠棄絕他，而永遠無法揀選他？那麼，請給我一個總是行善的人，他就不會被棄絕；也請給我一個總是行惡的人，他就不會被揀選。然而，若同一個人，在不同場合追求善惡兩者，被\Person{天主}——既是善又是公義的存有——在兩個方向上報應，祂並不因反覆無常或缺乏先見而改變判斷，而是以最堅定不移且預見性的決定，按每種情況的功過施行賞報。\par

% structure: p0049 source=h2 target=subsection reason=relative-heading-rank; base=h2

\subsection{第二十四章 \Person{天主}後悔的事例，特別是在尼尼微人身上，加以解釋和辯護}

此外，关于祂行事中所出现的悔改，你同样以乖僻的方式解释，就好像祂是出于反复无常和缺乏远见而悔改，或是因想起某种恶行而悔改；因为祂确实说过：\BibleQuote{我后悔立了撒乌耳为王}\BibleRef{撒上 15:11}，非常像是在说祂的悔改带有承认某种恶事或错误的意味。但事实并非总是如此。因为即使在善工中也会出现悔改的承认，作为对那些对恩惠忘恩负义之人的责备和谴责。例如，在撒乌耳这件事上，造物主在拣选他作王并赐予他祂的圣神这件事上并未犯错，祂就他的美好形体作了声明，说祂最合宜地拣选了他，因为那时他是最优秀的人，所以（如祂所说）在以色列子民中没有比他更好的。\BibleRef{撒上 9:2}祂也并非不知道他后来会变成怎样。因为没有人会支持你将缺乏远见归于那位天主，既然你不否认祂是天主，你也承认祂是有预见性的；因为这一神圣属性存在于祂内。然而，正如我所说，祂用对自己悔改的承认加重了撒乌耳的罪债；但既然在拣选撒乌耳这件事上没有任何错误或过失，那么这种悔改就被理解为责备他人而非自责。你说：那么，在关于尼尼微人的事上，我发现了一个自责的案例，因为\WorkTitle{约纳书}记载：\BibleQuote{天主就后悔了祂所说的要降于他们的灾祸，祂没有实行。}\BibleRef{纳 3:10}据此，约纳自己对主说：\BibleQuote{因此我先前逃往塔尔史士；因为我知道祢是有恩惠、有怜悯、缓于发怒、富有慈爱的天主，且后悔降灾。}\BibleRef{纳 4:2}因此，他很好地预先提出了至善天主的属性，即对恶人极其忍耐，对承认和哀叹自己罪过的人（正如当时尼尼微人所做的）极其丰盛地施予仁慈和良善。因为如果具有这属性的是至善者，那么你必须首先放弃你的立场，即与恶接触是与这样的存在者——即至善的天主——不相容的。而且因为马西昂也主张，好树不应结坏果子，但他却提到了\Quote{恶}（在讨论的段落中），而至善的天主是不能与恶相关的，那么有没有任何对\Quote{恶}的解释，能使它们甚至与至善者相容呢？有。简而言之，我们说，这里的恶并非指那些可归于造物主作为恶者的本性，而是指那些可归于祂作为审判者的权能。根据此，祂宣告：\BibleQuote{我造灾祸}\BibleRef{依 45:7}，和\BibleQuote{我计划降灾祸给你们}\BibleRef{耶 18:11}，这里指的是惩罚性的灾祸，而非有罪的恶事。这些灾祸带有何种污名，它们与天主的审判性角色相符，我们已充分解释。尽管它们被称为\Quote{恶}，但在一位审判者身上并不可谴责；也并非因为这一名称就表明审判者是恶的：同样，这具体的恶将被理解为属于这一类司法性的灾祸，并与它们一起与（作为审判者的）天主相容。希腊人有时也用\Quote{恶}一词指麻烦和伤害（非恶意的），正如在你们这段引文中也意指如此。因此，如果造物主因这样的恶而悔改，表明受造者应受谴责并应为其罪恶受罚，那么在此情况下，没有可归咎于造物主的罪责，因为祂有理有据地判定了如此充满不义之城的毁灭。因此，祂公正地判决，不怀任何恶意，而是出于正义原则，而非出于恶意。然而，祂称之为\Quote{恶}，是因为痛苦本身所涉及的恶和应得。那么，你会说，如果你以正义为名原谅了恶，理由是祂公正地判定了尼尼微人的毁灭，那么即便按照这一论点，祂也应受责备，因为祂对一正义行为后悔了，而正义行为本不应后悔。当然不会，我答道；天主永不会对正义行为后悔。现在我们必须理解天主的悔改意味着什么。因为虽然人最常因想起罪过而后悔，有时甚至因某项善行的不愉快而后悔，但在天主身上决非如此。因为天主既不犯罪，也不谴责善行，因此祂内没有对善行或恶行后悔的空间。这一点甚至在我们所引用的经文中已为你确定。撒慕尔对撒乌耳说：\BibleQuote{上主今日从以色列撕裂了你的王国，赐给了比你更好的近人}\BibleRef{撒上 15:28}；以色列将分为两部分：\BibleQuote{因为祂不转意，也不后悔；因为祂不像人那样后悔。}\BibleRef{撒上 15:29}因此，按照这一定义，天主的悔改在任何情况下都与人的悔改形式不同，因为它从不被视为缺乏远见、反复无常或对善功或恶行的任何谴责的结果。那么，天主悔改的方式是什么呢？如果你避免将其与人的情况相联系，那么它已经很清楚了。因为它的意义别无其他，只是先前意愿的简单改变；这在人身上也是可接受的，更不用说在天主身上了，因为祂的每个意愿都是无瑕的。在希腊语中，悔改一词（\OriginalTerm{悔改}{\GreekText{μετάνοια}}）的形成并非来自对罪的承认，而是来自心思的改变，我们在天主身上已经证明，这种改变是由不同情况的发生所调节的。\par

% structure: p0051 source=h2 target=subsection reason=relative-heading-rank; base=h2

\subsection{第二十五章 天主对亚当在堕落时和加音在犯罪后的处理：令人钦佩的解释和辩护。}

现在是时候了，为了应对所有这类异议，我应当继续解释和澄清你们视为琐碎、薄弱和矛盾的其他点。天主呼唤亚当：\BibleRef{创 3:9}\Quote{你在哪里？}好像不知道他在哪里；当亚当声称他赤身的羞耻是（他躲藏的原因）时，天主询问他是否吃了树上的果子，好像祂有所怀疑。绝不是；天主既未对罪的犯下不确定，也未对亚当的下落无知。固然，应当召唤那因意识到自己的罪而躲藏的罪人，并带他到上主面前，不仅藉着呼唤他的名字，而且要以直击他当时所犯之罪的一击。因为这语气不应仅仅以疑问的语调来读：\Quote{亚当，你在哪里？}而是以沉重而恳切的声音，带着责备的意味：\Quote{哦，亚当，你\Emph{在哪里}？}——意即：你已不在这里，你处于沉沦中——如此，这声音是那既责备又悲伤之上主的表达。但当然，乐园的某部分逃脱了那将宇宙握在手中如鸟巢、以天为宝座、以地为脚凳之主的眼目，以致祂在召唤亚当出来之前，看不到他躲藏和吃禁果时在哪里！狼或小偷逃不过你葡萄园或花园看守人的注意！我想，天主以祂更敏锐的视觉，从高处无法看到祂脚下任何东西？愚昧的异端者，你们如此轻蔑地对待天主伟大和人受教训的如此精妙的论证！天主以看似不确定的方式发问，正是为了在此证明人拥有自由意志，可以选择否认或承认，并给人自由承认自己过犯的机会，从而减轻它。同样，祂询问加音他的弟弟在哪里，好像尚未听见亚伯尔的血从地上喊冤，为的是让加音也有机会藉同样的意志力量自愿否认，从而加重他的罪行；这样，就为我们提供了认罪而非否认罪的榜样：如此，甚至在那时就开始了福音的教导：\Quote{凭你的话，你将被成义；凭你的话，你将被定罪。}\BibleRef{玛 12:37}虽然亚当因律法下的处境注定死亡，但藉着上主的话仍给他保留了希望：\Quote{看，亚当成了我们中的一个。}即，因将来人进入天主性体。接着是什么？\Quote{现在，恐怕他伸手，也摘取生命树的果子，吃了，就永远活着。}这样插入表示现在时的小词\Quote{现在}，表明祂暂时、目前延长了人的生命。因此祂并没有真正诅咒亚当和厄娃，因为他们是要被恢复的，并且已经藉着认罪得到了缓解。然而，加音，祂不仅诅咒；而且当他希望藉死亡赎罪时，甚至禁止他死，以致他除了自己的罪，还要背负这禁令的重担。这样，这将证明我们天主的\Quote{无知}，祂因这缘故而假装如此，为的是犯罪的人不应不知道他该做什么。当论及索多玛和哈摩辣时，祂说：\Quote{我现在要下去，看看他们所行的，是否全像那达到我面前的呼声一样；如果不像，我也要知道。}那么，祂在这事上也是因无知而不确定，并想要知道吗？还是这是必要的语气，表达威胁而非怀疑，以调查为表象？如果你嘲笑天主的\Quote{下去}，好像祂不降下就不能完成审判，小心你不要以同样沉重的打击攻击你自己的神。因为祂也下来完成祂所愿意的。\par

% structure: p0053 source=h2 target=subsection reason=relative-heading-rank; base=h2

\subsection{第26章. 天主的宣誓：其意义。梅瑟恳求天主息怒不降罚以色列，是基督的预像。}

但天主也会起誓。我想，他是凭着马尔基昂的天主起誓吗？不，不，马尔基昂说，是一个更虚浮的誓言——凭他自己起誓！当他明知（\BibleRef{依 44:8}）没有别的神，尤其当他起誓正是为了这一点，即除了他自己以外根本没有别的神，那他还能做什么呢？那么，你是指控他发假誓，还是发虚誓？但既然他像你所说的，不知道有另一个神，他就不可能显得发了假誓。因为他凭他所知道的那位起誓，确实没有发假誓。而他起誓说没有别的神，这也不是虚誓。如果没有那些相信有别的神的人，比如当时拜偶像的人和现今的异端者，那才真是虚誓。因此，他凭自己起誓，好让你相信天主，即使他起誓说除他自己以外根本没有别的神。但是马尔基昂，是你自己迫使天主这样做的。因为甚至在那么早的时候，你就已被预见。因此，如果他在应许和警告中都起誓，并以此迫使人相信那起初难以相信的事，那么凡是能使人相信天主的，就没有什么配不上天主的。但（你说）天主那时甚至在他的烈怒中也非常卑劣，当他因百姓拜金牛而发怒时，他向他的仆人梅瑟请求说：\BibleQuote{容我发怒攻击他们，消灭他们；我要使你成为一个大民族。}\BibleRef{出 32:10} 因此，你坚持认为梅瑟比他的天主更好，因为他祈求平息，甚至阻止了他的怒气。梅瑟说：\Quote{你不可这样做；否则连我也一起消灭。} 你同百姓一样可怜，因为你不认识基督，他预表在梅瑟身上，作为圣父的求情者，并为百姓的得救而献出自己的生命。不过，那时百姓确实被赐给了梅瑟，这已经够了。他作为仆人所能向主求的，主也对他自己作了要求。正为此，他对他的仆人说：\Quote{容我消灭他们}，好让梅瑟通过恳求和奉献自己来阻止（所威胁的审判），并让你通过这个例子知道，天主赐予一个忠信的人和先知多大的特权。\par

% structure: p0055 source=h2 target=subsection reason=relative-heading-rank; base=h2

\subsection{第二十七章 其他异议之考量。天主在降生成人中的屈尊。在此救恩计划中无一物有损于天主性体。全能圣父堪当地维持其天主尊威，绝不为人所见。马尔基昂派妄论之乖谬。}

现在，为了简要回顾你迄今所收集的、用以拆毁造物主的其他论据（你认为这些论据卑鄙、软弱、不相称），我将以简单明确的陈述提出：天主若不承担人的情感和情爱，就无法与人进行任何交往；祂藉此使祂威严的力量得以缓和——这种威严对于人有限的能力来说，无疑是不可承受的——而祂所采用的这种屈尊降卑，对祂自己而言固然是降低身份，但对人却是必要的，正因如此，它才堪当成为天主的作为：因为没有什么比人的得救更堪当天主。如果我是与外教人辩论，我会更详细地讨论这一点；尽管与异端者的争论也并无太大不同。既然你们现在自己也相信，天主曾以人的本性的形状和其他一切状况行动，那么你们当然不再需要被说服天主使自己适应人性，而是以你们自己的信仰感到义务。因为如果你们所信的天主，甚至从祂更高的状态，将祂威严的至高尊荣降低到如此卑微，以至于受死，甚至死于十字架，那么为什么你们不能认为一些屈辱也适合于我们的天主呢？——只不过比犹太人的凌辱、十字架和坟墓更能容忍罢了。这些屈辱，难道从此以后就要对基督（祂是情感的主体）形成偏见，认为祂分享了那你们以分享人性品质为辱的神性吗？现在，我们相信基督确实经常以天主圣父的名行事；祂实际上从起初就与人交往；实际上与族长和先知交流；祂是造物主的圣子；祂是祂的圣言；天主通过从自身发出祂而使其成为圣子，并从此将祂置于每一安排和祂旨意的管理之上，使祂稍微低于天使，正如在达味中记载的。在祂这种降低身份的状态中，祂从圣父那里接受了正被你们斥为人的那些方面的安排；从起初，甚至那时，祂就学习祂最终将要成为的人的状态。是祂下降，是祂询问，是祂要求，是祂发誓。然而，关于圣父，我们共有的福音本身将作证，祂从未被看见，正如基督的话所说：\Quote{除了圣子，没有人认识圣父。}\BibleRef{玛 11:27} 因为甚至在旧约中，祂已宣告：\Quote{没有人能看见我，并且活着。}\BibleRef{出 33:20} 祂的意思是说圣父是不可见的，在祂的权柄和祂的名下，那位作为圣子显现的天主就是天主。但在我们这里，基督是以基督的身份被接纳的，因为即使这样，祂也是我们的\Emph{天主}。因此，任何你们认为堪当天主的属性，都必须在圣父那里找到——祂是不可见、不可接近、平和，并且可以这样说，哲学家的天主；而那些你们斥为不相称的品质，则必须认为存在于圣子身上——祂曾被看见、被听见、被遇见，是圣父的见证者和仆人，在自己内结合了人和天主，天主有大能的行为，人则有软弱的品质，以便祂能向人给予多少，就从天主那里取走多少。在你们眼中，我的天主的一切耻辱，实际上就是人类救恩的圣事。天主与人交谈，好使人学会像天主一样行事。天主与人平等交往，好使人能与天主平等交往。天主被显为微小，好使人能变得极其伟大。你们这样鄙视一位天主，我几乎不知道你们是否\Emph{\OriginalTerm{出于信德}{ex fide}}相信天主曾被钉在十字架上。那么，你们对造物主这两种品格的颠倒该有多大！你们称祂为\Emph{审判者}，却谴责审判者那仅按案情是非行事的严厉为残忍。你们要求天主\Emph{极其良善}，却又鄙视祂那与祂仁慈相符的、根据人境遇的卑微而低声下气交往的温柔为卑贱。祂无论伟大还是渺小都不讨你们喜欢，既不做你们的审判者，也不做你们的朋友！如果同样的特征也在你们的神身上被发现呢？\Emph{祂}也是审判者，我们已在适当章节中表明：因为是审判者，祂必然严厉；因为是严厉，祂也必然残忍——如果真是残忍的话。\par

% structure: p0057 source=h2 target=subsection reason=relative-heading-rank; base=h2

\subsection{第二十八章 通过对比，局势翻转有利于真天主，反对马尔西昂}

现在，关于软弱和恶毒，以及造物主的其他（所谓的）特性，我也将提出一些\OriginalTerm{反命题}{antitheses}，与\Person{马西昂}所提出的相抗衡。如果我的天主不知道有谁比祂更高，你的神也完全不知道有谁比自己更低。正如\Person{赫拉克利特}\Quote{晦涩者}所说的：无论向上还是向下，结果都一样。如果祂确实并非不知（自己的位置），那么祂从一开始就必定已经想到这一点。罪和死亡，以及罪的作者——魔鬼——以及我的天主允许存在的一切邪恶，你的神也同样允许；因为他容许祂允许它们存在。我们的天主改变了祂的旨意；同样，你的神也改变了。因为那位如此晚才向人类投下目光者，改变了那长期拒绝投下目光的旨意。我们的天主在某一事例上后悔降灾；你的神也是如此。因为藉着他最终顾念人的救恩这一事实，他表现出对先前忽视的后悔，这种后悔是恶行所当受的。但忽视人的救恩将被视为恶行，仅仅因为它已被你的神的后悔所补救。你们说，我们的天主命令了一件欺诈行为，但涉及金银之事。然而，既然人比金银更宝贵，你的神就更加欺诈，因为他剥夺了人的主和造物主。我们的天主要求以眼还眼；但你的神甚至造成更大的伤害（按你们的想法），因为他阻止报复行为。有哪一个人不等挨第二下就会还手呢？（你们说）我们的天主不知道祂该选择谁。你的神也不知道，因为如果他预知结果，就不会选择叛徒\Person{犹达斯}。如果你们声称造物主在某个事例中行了欺骗，那么你们的基督就有更大的虚假，因为他的身体本身是不真实的。许多人被我的天主的严厉所消灭。那些不被你的神拯救的人，实际上被他安排去毁灭。我的天主下令杀死一个人。你的神愿意自己被处死；他对自己同样是一个凶手，正如对他想被其杀死的那个人一样。此外，我要向\Person{马西昂}证明，被他的神杀死的人很多；因为他使每个人都成为杀人犯：换言之，他注定他们灭亡，除非人们对基督尽了一切义务。但真理的正直之德满足于很少的资源；谬误则需要许多东西。\par

% structure: p0059 source=h2 target=subsection reason=relative-heading-rank; base=h2

\subsection{第29章 \Person{马西昂}自己的\WorkTitle{反论集}，若只除去该作品的名目和目的，便为真天主的恒常属性提供了证据。}

但是，如果需要更详尽地拆解它们，以便为造物主维护良善天主和审判者的特质——我们在两方面的例子之后，已经表明这些特质如此堪为天主——那么我本应更密切、更充分地攻击\Person{马西昂}自己的\WorkTitle{反论集}。然而，既然良善与公义这两个属性共同构成了作为全能者的天主性体之恰当圆满，我就能满足于现在已经简要地驳斥了他的\WorkTitle{反论集}——该作品旨在从（造物主的）作为、或祂的律法、或祂的伟大工程的性质中引出区别，从而将基督与造物主分开，如同至善者与审判者分开，如同慈悲者与冷酷者分开，以及带来救恩者与造成毁灭者分开。事实上，它们反而联合了那两位存在，它们将两位放在那些（属性的）差异中，而这些差异在天主内却是相容的。因为只要除去\Person{马西昂}的书名以及作品本身的意图和目的，你就无法得到更好的证明，表明同一位天主既是至善又是审判者，因为这两种特质唯独在天主内才恰当地被找到。诚然，在所选例子中为反对基督与造物主而做的努力，反而更有助于他们的合一。因为神性者的本性完全相同，既良善又严厉，正如通过相同的例子和相似的证据所表明的，以至于祂愿意向那些祂首先施加严厉的人显示祂的良善。时间上的差异不足为奇，因为同一位天主在恶被制服后显出仁慈，而当初恶尚未被制服时祂曾如此严厉。因此，藉助\WorkTitle{反论集}，造物主的管理计划可以更容易地被表明是基督所\Emph{改革}的，而非\Emph{摧毁}；是\Emph{恢复}的，而非\Emph{废除}的；尤其因为你们使你们的神远离一切如尖刻行为等，甚至远离与造物主的任何竞争。现在既已如此，那\WorkTitle{反论集}如何表明祂在每一件争论的事上都是造物主的对手呢？好吧，即使在这里，我也承认：在这些事上，我的天主一直是忌邪的天主，祂凭自己的权柄特别注意，使祂所做的一切在其开端都有更强健的生长；而且这是以一种良善的、理性的竞争方式，趋向成熟。从这个意义上说，世界本身也会承认祂的\Quote{反命题}，即从其自身元素的对立中看出，尽管世界是以最高的理性被规范。因此，最轻率的\Person{马西昂}，你所当做的本是表明（你所教导的两位神中）一位是光明之神，另一位是黑暗之神；这样你就会发现更容易说服我们：一位是良善之神，另一位是严厉之神。然而，这个\Quote{反命题}（或管理形式的变换）将恰当地属于那一位，因为它在世界（的管理）中真正属于祂。\par

\SourceRecord{Translated by Peter Holmes.；Ante-Nicene Fathers；Vol. 3.；Edited by Alexander Roberts, James Donaldson, and A. Cleveland Coxe.；Buffalo, NY: Christian Literature Publishing Co.,；1885.；<http://www.newadvent.org/fathers/03122.htm>.}

% structure: p0001 source=h1 target=section reason=page-title

\section{\WorkTitle{驳斥马尔基昂，第三卷}}

\Emph{此处证明基督是创造世界之天主子，由先知所预言，取了与我们相同的人性，藉着真实的降生成人。}\par

% structure: p0003 source=h2 target=subsection reason=relative-heading-rank; base=h2

\subsection{第一章 引言；结合本书主题对前述论证的简要陈述}

追随我最初论文的轨迹——我们正在稳步恢复那已失落的文字——我们现在按主题顺序来讨论基督，尽管在证明了只有唯一的天主之后，这似乎是余事。因为无疑已经足够清楚地规定了：基督必须被视为只属于造物主，而非其他天主，当已经确定除了造物主之外没有其他天主应成为我们信仰的对象。基督如此明确地宣讲祂，而宗徒们也一个接一个地同样清楚地肯定：基督只属于祂自己所宣讲的那一位——即造物主——因此在马尔基昂的丑闻之前，从未搅动起任何关于第二位天主（因此也没有第二位基督）的提及。这一点最容易通过检视宗徒的教会和异端的教会来证明，从这些教会我们被迫宣称：\Emph{那里}无疑出现了对（信仰）准则的颠覆，凡发现有后来产生的意见之处。——这一点我已经在我第一本书中提及了。对这一点的讨论即使在现在也无疑有价值，因为我们即将对基督（的主题）进行单独检视；因为，在证明基督是造物主的\Emph{圣子}时，我们实际上将马尔基昂的天主排除在外。真理应当运用她所有可用的资源，而不是跛行。然而，在我们简明的信仰准则中，她自有其方式。但我已下定决心，像一个热忱的人，以各种方式和在任何地方与我的对手在其异端的疯狂中交锋，这疯狂是如此之大，以至于他发现更容易假定那从未被听说过的基督已经来了，而不是那一直被预言的基督。\par

% structure: p0005 source=h2 target=subsection reason=relative-heading-rank; base=h2

\subsection{第二章 基督的来临为何应被预先宣告}

于是立即切入主题，我必须面对这个问题：基督是否应该如此突然地降临？（我回答：不。）首先，因为祂是祂圣父的天主子。这是秩序之理：圣父应在圣子之前宣告圣子，圣子应在圣父之前为圣父作证；圣父应为圣子作证，圣子应为圣父作证。其次，因为除了圣子的名号外，祂还是被派遣者。因此，派遣者的权威必须首先在被派遣者的见证中出现；因为以他人权威而来的人，不会自己凭己意提出它，反而会寻求它的保护，因为首先到来的是赋予他权威之人的支持。如果圣父从未提名祂，基督就不会被承认为圣子；如果没有派遣者给祂使命，祂就不会被相信为被派遣者：圣父若有，则特意提名祂；派遣者若有，则特意派遣祂。凡违反规则之事都会引发怀疑。万物的首要秩序不会允许在认知上圣父后于圣子，或派遣者后于被派遣者，或天主后于基督。没有什么能在被认知上优先于其本源，在秩序上亦然。突然一位圣子，突然被派遣，突然基督！相反，我认为从天主而来没有什么是突然的，因为没有什么不是由天主所安排和制定的。如果被安排，为何不也被预先宣告，以便通过预言表明它是被安排的，并通过安排表明它是神圣的？确实，如此伟大的工程——我们可以肯定它需要预备，因为它是为了人的救恩——不能因此成为一件突然的事，因为它是透过信德才能生效的。既然它必须被相信才能有用，那么为了确保这信德，它就要求一种建立在预先安排和预先宣告之基础上的预备。当信德藉着这样的过程得到告知后，天主可正当地向人要求信德，而人也可将信德寄于天主；既然知识已使之成为可能，人便有责任相信那些他确实从预先宣告中学到要相信的事。\par

% structure: p0007 source=h2 target=subsection reason=relative-heading-rank; base=h2

\subsection{第三章 单有奇迹而无预言，不足以证明基督的使命}

你说，这样的程序并不必要，因为祂要立即通过奇妙作为的证据证明自己就是圣子、被派遣者、以及实实在在的天主的基督。然而，我不得不否认单凭这类证据就足以作为对祂的见证。后来祂自己剥夺了这类证据的权威，因为当祂宣称有许多人要来，\Quote{显大奇事和异能}，以致连被选者都要被引诱，但却不可接待他们时，祂表明了相信奇事和异能是何等轻率——这些连假基督都能轻易做到。不然，如果祂意欲通过某种证据——我指的是奇迹——来被认可、被理解、被接受，那么祂为何禁止承认那些拥有同样证据、且其来临同样突然而无任何权威宣告的其他假基督呢？如果祂因先于他们到来，并先于他们显示大能作为的迹象，便抢先占据了人信仰的权利——正如先到者占据浴场首位——从而在权利上抢在所有后来者之前，那么要小心，祂自己可能也会落入后来者的状况，如果祂被发现落后于创造主——创造主早已被认识，且已行过类似祂的奇迹，并同样警告过人们不要相信后来者。因此，如果先到并发出警告是为了限制和设定信仰的界限，那么祂自己就必须被定罪，因为祂在被人认识方面较晚；而制定关于后来者的这一规则的权威，只能属于创造主，因为祂不可能居后于任何人。如今，我要证明创造主有时通过祂古时的仆人显示，并在其他情况下保留给祂的基督显示的奇迹，正是你们声称唯独通过你们的基督的信仰才有的那些奇迹，那么我甚至可以从这一点合理主张，正因为如此，基督更不应仅仅因奇迹而被相信，因为这些奇迹只会显示祂属于创造主，而非任何别的天主，因为祂的奇迹与创造主的大能作为相符，无论是通过祂的仆人施行还是保留给祂的基督。尽管如果在你们的基督身上发现其他证据——即是新的——我们也会更愿意相信它们同样属于同一位天主，而不属于那位只有新证据、缺乏赢得信仰同意的古老证据的天主。因此，即使从这一理由出发，祂也应当既通过建立信仰的预言，也通过奇迹来被宣布，尤其是当祂要与创造主的基督相对抗时——创造主的基督将要带着自己的征兆和先知到来，以便通过不同类型的证据作为祂的对手而显耀。但是，一个从未被预言的基督如何能被从未被预言的天主所预言呢？因此，不可避免的结论是：你们的天主和你们的基督都不是信仰的对象，因为天主不应被未知，而基督应通过天主被认识。\par

% structure: p0009 source=h2 target=subsection reason=relative-heading-rank; base=h2

\subsection{第四章 马西昂的基督并非预言的对象。异端者这一理论的荒谬后果。}

我想，他不屑于仿效我们天主的秩序，因为他是令天主不悦的，且必当被击败。他愿以新的方式作为新的存在而来——在圣父的宣告之前为子，在派遣者的权威之前为被派遣者；如此他便能亲自传播一种最怪异的信仰，即人们应相信基督在祂的存在被知晓之前已经来临。在此，我适宜处理另一个问题：他为何不在基督之后来？因为我观察到，在如此漫长的时期内，他的主以极大的忍耐容忍了那位极其冷酷的造物主，而造物主却一直向人们宣告祂的基督。我认为，不论什么理由促使他如此行，而推迟了他自己的启示和干预，正是同一个理由加给他容忍造物主的责任（造物主在\Emph{祂的}基督内也有自己的安排要实行），以便在对立的天主和对立的基督的整个计划完成和成就之后，他再引入他自己适当的安排。但他厌倦了如此长久的忍耐，因而未能等到造物主的历程结束。他容忍他的基督被预言，却拒绝允许他被显现，这是无益的。要么他打断对手的时间进程是没有正当理由的，要么他长久以来没有打断它也是没有正当理由的。是什么在\Emph{起初}阻止了他？又是什么在\Emph{最后}扰乱了他？然而，就现今的情况而言，他在两方面都自相矛盾：他在造物主之后如此迟缓地启示自己，又在祂的基督之前如此仓促地显现；而他本该早就以辩驳来对抗前一位，对后一位则暂缓对抗——不该容忍前一位如此长久的无情敌意，也不该惊扰后一位，后者当时尚在静止中！在二者的情形中，当他剥夺了他们被视为至善天主的称号时，他至少表明自己是反复无常和不确定的：对造物主（在愤怒上）冷淡，却对祂的基督炽热，而且对二者都无能为力！因为他既未能约束造物主，也未能抵挡祂的基督。造物主仍如祂的真实所是。祂的基督也必如经上论及祂所记载的那样来临。既然他不能通过惩罚来纠正造物主，他又为何在造物主之后来？既然他无法阻止基督显现，他又为何在基督之前启示自己？相反，如果他确实惩罚了造物主，那么（我想）他在造物主之后启示自己，是为了让需要纠正的事物先出现。因此（当然），他应当等待基督先显现，然后以同样的方式惩罚基督；这样他就成为在基督之后的惩罚者，正如他在造物主的情形中所做的那样。另有考虑：由于他在第二次来临时要跟在基督之后，既然他在他的第一次来临时对造物主采取了敌对行动，摧毁了属于造物主的法律和先知，那么他必定也会在第二次来临时反对基督，推翻祂的国。那时，毫无疑问，他将结束他的历程，然后（如果可能）才值得相信；否则，如果他的工作已经完成，他的来临就是徒劳的，因为他确实不能再完成任何事了。\par

% structure: p0011 source=h2 target=subsection reason=relative-heading-rank; base=h2

\subsection{先知风格的若干特征：其解释原则}

这些初步的评论，我斗胆在讨论的第一步，当冲突尚处于，可以说，远距离时做出。但是，既然从现在起，我必须在明确的问题上与我的对手近身搏斗，我察觉到我甚至必须在此处推进一些战线，战斗将在此处发动；它们就是创造主的圣经。因为我将必须证明，基督出自创造主，按照这些圣经，后来在创造主的基督身上应验了；我发现有必要阐述圣经本身的形式和，可以这样说，本质，免得当它们被应用于讨论的主题时，因被卷入争议而分散读者的注意力，并且使它们的证明与主题本身的证明相混淆。现在我引出预告预言的两个条件，它们要求我们的对手在讨论的未来阶段表示同意。其一，未来的事件有时被宣布为仿佛已经过去。因为对于天主本性而言，将它所决定的一切视为已成事实是合宜的，因为在他内，永恒自身支配着时间的统一状态。事实上，对于先知的预卜而言，更自然的是，将所预见的事物表现为已被看见并已实现，即使它还在预见之中；换句话说，就是无论如何都会成为未来的事物。例如，在\BibleBook{}{依撒依亚}中：\BibleQuote{我把我的背转给打击我的人，把我的脸颊（我暴露）给他们的手。我没有因羞耻和唾沫遮掩我的脸。}因为无论那时说话的是基督，如我们所持，还是先知，如犹太人所言，说的是关于他自己，无论哪种情况，那尚未发生的事听起来仿佛已经成就了。另一个特性是，非常多的事件是通过谜语、比喻和寓言来比喻性地预告的，并且必须从与字面描述不同的意义来理解。因为我们读到\BibleQuote{山岳滴下新酒}，\BibleRef{岳 3:18}但不是好像有人期望从石头中流出\Quote{\Emph{葡萄汁}}，或从岩石中流出它的煮液；也听到\BibleQuote{流奶和蜜的地方}，但不是好像你设想你会从地上收集萨莫斯糕饼；天主也并非，真的，以水闸管理员或农夫的角色提供服务，当他说：\BibleQuote{我要使河流在旱地涌出；我要在旷野中栽植香柏和黄杨木。}同样，当预言外邦人的归化时，他说：\BibleQuote{田野的走兽要尊敬我，巨龙和猫头鹰，}他肯定从不打算从鸟和狐的幼崽，以及从奇谈和寓言中的歌手那里获取他的吉兆。但是何必在此主题上赘述？当我们的异端者所采纳的那位宗徒亲自解释那条允许为打谷的牛犊笼上嘴的法律，并非关乎牛，而是关乎我们自己；\BibleRef{格前 9:9}并且还声称那跟随（以色列子民）并供给他们饮水的磐石就是基督；此外，教导迦拉达人，关于亚巴郎两个儿子的叙述在他们的进程中具有比喻意义；又向厄弗所人暗示，当起初宣告男人应离开父亲和母亲，与妻子成为一体时，他将其应用于基督和教会。\BibleRef{弗 5:31}\par

% structure: p0013 source=h2 target=subsection reason=relative-heading-rank; base=h2

\subsection{第六章 马西昂派和犹太人的错误在某些点上的共同性。关于基督被拒斥的预言的考察。}

既然在犹太先知文献中显然存在这两个特征，那么读者当记得，每当我们从中引用任何证据时，双方都同意讨论的要点不是圣经的形式，而是它被召来证明的主题。因此，当我们的异端者狂妄地断言那从未被预言过的基督已经来了时，按照他们的假设，那始终被预言的基督就尚未出现；这样他们不得不与犹太人的错误联手，并借助它来构建他们的论点，借口说犹太人自己相当确信来的是另外一位：因此他们不仅把祂当作陌生人拒绝，而且作为敌人杀害了祂，尽管如果祂只是他们中的一员，他们无疑会承认祂并以虔诚的信仰追随祂。我们的船长当然不是从罗得斯法律得到他的航海智慧，而是从本都法律，它告诫他不要相信犹太人没有权利得罪他们的基督；然而（即使他们的行为没有任何类似的事被预言过），单是容易犯错的人性就可能使他推测，犹太人作为人，完全可能犯下这样的罪，而无须对他们的情感抱有任何不公的偏见，他们的罪在先前是如此可信。然而，既然确实预言了他们将不承认基督，因此甚至要置祂于死地，那么就可以推知，祂被他们忽视并杀害，而他们事先已被指出将会对祂犯下这样的恶行。如果你要求对此提供证明，那么与其推出那些声称基督能受苦而死的圣经段落（因为若祂未被拒绝，祂就不能真正受苦），并将它们留作受苦主题之用，我目前只满足于引用那些仅仅显示基督被拒绝的可能性经文。这很容易做到，因为经文表明认识的全部能力被创造主从这民族那里取走。祂说：\BibleQuote{我要除掉智慧人的智慧；并隐藏聪明人的聪明；}\BibleRef{依 29:14}又说：\Quote{你们用耳朵听，却不明白；用眼睛看，却不领悟；因为这百姓的心变迟钝，耳朵沉重，眼睛紧闭；免得他们用耳朵听见，用眼睛看见，用心领悟而回转，我就治好他们。}这种对健全官能的迟钝是他们自己造成的，他们用嘴唇爱天主，心却远离祂。既然基督是由创造主所宣告的，\Quote{祂形成闪电，创造风，并向人宣告祂的基督，}如\Person{岳厄尔}先知所说，既然犹太人的全部希望，更不用说外邦人了，都集中在基督的显现上——那么显然，他们被剥夺了认识与理解的能力——即智慧与知识——将不能认识和理解所预言的，即基督；那时他们的智慧的首领将误解祂——即经师和有识之士或法利塞人——而百姓像他们一样，用耳朵听却不明白施教的基督，用眼睛看却不领悟基督，尽管祂给了他们征兆。同样，别处说：\Quote{谁眼瞎，难道是我的仆人？谁耳聋，难道是统治他们的那位？}祂又通过同一\Person{依撒意亚}责备他们说：\BibleQuote{我养育子女，将他们养大，他们却背叛了我。牛认识主人，驴认识主人的槽，以色列却不认识，我的子民不思考。}\BibleRef{依 1:2}我们确实确切知道，基督作为创造主的圣神常在先知们内说话（因为先知说：\Quote{我们圣神的位格，主基督，}祂从起初就如父的代权者以天主之名被听见和看见），我们很清楚，祂实际上责备以色列的话，与预言祂将要对他斥责的话是一样的：\BibleQuote{你们离弃了主，激起了以色列圣者的怒气。}\BibleRef{依 1:4}然而，若你宁愿将犹太人的无知从起初就归于天主本身，而不归于基督，因不愿承认甚至古时创造主的圣言和圣神——即祂的基督——就被他们轻视和不认识，那么即使在这诡计中你也会被击败。因为当你不否认创造主之子、圣神和实体也是祂的基督时，你必须承认那些不认识父的人，同样因他们自然实体的同一性而不认识子；因为如果圆满的实体已超出了人的理解，那么它的部分，特别是在分享圆满时，更是如此。如今，当仔细考虑这些事后，便可清楚看出犹太人如何既拒绝基督又杀害祂；不是因为他们视祂为陌生的基督，而是因为虽然祂是他们的，他们却不认识祂。因为他们如何能理解那从未被宣告过的陌生者，既然他们未能理解那关于祂曾有持续预言的那一位？拥有预言实质基础的事物，可能被理解或不理解，从而有知识或错误的素材；而缺乏这样素材的事物则不允许智慧的结果。因此，他们厌恶基督并迫害祂，并非好似祂属于另一个神，而只是作为一个他们视为施行奇迹的术士和教义上的敌人的人。因此，他们只把祂当作一个普通人，且是他们中的一员——即一个犹太人（只不过是一个背弃者和犹太教的破坏者）——来审判，并按照他们的法律惩罚了祂。如果祂确实是陌生人，他们就不会审判祂。他们远非表现出理解祂是陌生的基督，以至于他们甚至没有判断祂对他们的本性是陌生的。\par

% structure: p0015 source=h2 target=subsection reason=relative-heading-rank; base=h2

\subsection{第七章 预言指出基督两种不同的状态：一是卑微的，一是尊威的。这一事实指向基督的两次来临。}

我们的异端者现在将与犹太人一起有最充分的机会认清他错误的线索，他正是从犹太人那里借来了这场讨论的指引。然而，瞎子领瞎子，两人一起掉进坑里。我们断言：既然先知们指明了基督的两种不同状态，那么这些状态也同样预兆了相同次数的来临；第一次来临将是在卑微中，那时祂必须像羊一样被牵去宰杀作祭品，像羔羊在剪毛者前缄默无声，不开口，外貌也不俊美。因为（先知）说：我们论祂宣告过：\BibleQuote{祂像一棵嫩苗，像从干地生出的根；祂无威仪，也无美貌；我们看见祂时，祂没有可羡慕之处：祂的形象被毁损；} \BibleQuote{比世人子更憔悴；是个受苦的人，深知如何承担我们的病弱；} \BibleQuote{被圣父立为绊脚石和使人失足的磐石；} \BibleRef{依 8:14} \BibleQuote{使祂暂时比天使微小一点；} 祂自称是\BibleQuote{虫，不是人，是众人羞辱，是百姓所藐视的。} 这些卑微的标志完全适合祂的第一次来临，正如祂尊威的标志适合祂的第二次来临，那时祂将不再是\BibleQuote{绊脚石和使人失足的磐石}，而是在被弃绝之后成为\BibleQuote{房屋的角石}，被接纳并被高举到圣殿的顶端——即祂的教会——正是达尼尔书中从山岩凿出的那块石头，要打碎并压垮世俗国度的像。关于这次来临，同一位先知说：\BibleQuote{看，有一位像人子者乘着天上的云彩而来，来到万古常存者面前；有人引祂到祂跟前，赐给祂统治权、尊荣和国度，使各民族、各邦国、各语言的人都侍奉祂。祂的统治是永远的统治，永不消逝；祂的国度也不会被毁灭。} \BibleRef{达 7:13} 那时，祂确实将有荣耀的形体和超越世人子的无瑕美貌。\BibleQuote{你比世人子更俊美，}（圣咏作者）说，\BibleQuote{恩宠倾注在你唇上，因此天主永远祝福了你。大能者啊，愿你腰间佩剑，带着你的荣耀和威严。} 因为圣父使祂暂时比天使微小一点之后，\BibleQuote{将赐祂尊荣和荣耀为冠冕，将万有置于祂脚下。} \BibleQuote{那时他们仰望他们所刺透的那位，并要哀悼祂，各族各族；} 无疑，因为他们曾经拒绝在祂人性卑微中承认祂。祂是人，耶肋米亚说，谁能认识祂？因此，依撒意亚问：\BibleQuote{谁能述说祂的世代？} \BibleRef{依 53:8} 同样，在匝加利亚书中，基督耶稣——圣父的真大司祭——以若苏厄的身份，且以祂名字的奥迹，被描绘成两种服装，关联祂的两次来临。起初祂穿着污秽的衣服，即处于受苦和必死肉身的卑微中；然后魔鬼敌对祂，作为叛徒犹达斯的鼓动者，更不用说在祂受洗后诱惑祂；之后祂脱去最初的污秽衣服，穿上司祭袍、冠冕和纯洁的礼冠，即第二次来临的荣耀和尊荣。此外，若我还能对\BibleQuote{赎罪大节}上献的两只公山羊作出解释，它们不也预表基督的两种本性吗？它们大小相同，外貌非常相似，因主的容貌同一；因为祂不会以别的形貌来临，祂必须被那些也曾刺伤和刺透祂的人所认出。其中一只公山羊被绑上朱红线，被百姓赶出营外进入旷野，伴随着咒骂、唾弃、拉扯和刺伤，用一切标志标明主的苦难；而另一只作为赎罪祭被献，并给圣殿的司祭作食物，则证明了祂第二次显现的痕迹——那时（一切罪过已被赎清）属灵圣殿的司祭，即教会，将享用主恩宠的肉，而其余的人则未品尝就离开救恩。既然第一次来临在预像中极其模糊、充满各种侮辱，第二次却荣耀并完全合于天主，那么他们只关注自己易于理解和相信的第二次来临，从而被更隐晦、更卑微的第一次来临所欺骗，也并非不当。因此，直到今日他们否认他们的基督已经到来，因为祂未在威严中显现，而他们忽视了祂也应在卑微中来临的事实。\par

% structure: p0017 source=h2 target=subsection reason=relative-heading-rank; base=h2

\subsection{第八章 马西昂幻象论的荒谬；基督降生成人的真实性。}

我们的异端者现在必须停止从犹太人那里借用毒物——正如谚语所说，\Quote{毒蛇}从\Quote{蝮蛇}那里借毒——并从此吐出他自己本性的恶毒，就如他断言基督是幻影时那样。确实，他的这个意见必然会在他的早熟且有些夭折的马吉安派信徒中得到维护，宗徒若望称这些人为假基督，因为他们否认基督曾降生成人；他们这样做并非为了确立另一个神的权利（因为在这一方面他们也已被同一宗徒所谴责），而是因为他们从一开始就假设了一位降生成人的天主是不可信的。现在，假基督马吉安越是坚定地抓住这个假设，他当然就越准备好拒绝基督的肉身实体，因为他已经使他的神被我们认识为既非肉身的创造者，也非肉身的恢复者；正因如此，他当然被认为是至善的，且远离造物主的欺骗和谬误。因此，他的基督，为了避免所有这些欺骗和谬误，以及可能的话，避免被归入造物主，就成了他外表所不是的，并伪装成他所不是的——降生成人却没有肉身，是人却不是人，同样地，是天主的基督却不是天主！但他为何不也传播天主的幻影？我怎能相信一个在外部实体上完全错误的人，关于内在本质的事？当他在一个明显事实上被发现如此虚假时，怎么可能相信他在奥迹上是真实的？此外，当他混淆了精神的真理与肉身的错误时，他怎能在他自己内结合光与黑暗、或真理与错误的相通，而宗徒说这些不能共存？\BibleRef{格后 6:14} 然而，既然基督是肉身这一事实已被发现是谎言，那么凡由基督肉身所做的一切事都是虚假的——每一项接触、触摸、饮食的行为，是的，甚至祂的奇迹。如果祂通过触摸或被人触摸而治愈某人，任何由身体行为所做的事，在祂自己身体完全缺乏实在性的情况下，都不能相信是真实地做成的。没有任何实质性的东西能被允许由非实质性的东西产生；没有任何丰满由空虚产生。如果外表是假定的，行为也是假定的；如果工作者是想象的，工作也是想象的。按照这个原则，基督的苦难也将被发现不足以证明我们对祂的信德。因为那并非真实受苦的人，什么也没有遭受；而且一个幻影不能真实地受苦。因此，天主的全部工程都被推翻了。基督的死亡——其中蕴含着基督徒名称的全部重量和果实——被否认了，尽管宗徒如此明确地断言它无疑是真实的，使它成为福音、我们的救恩和他自己宣讲的基础。他说：\Quote{我首先传给你们的是：基督为我们的罪死了，而且被埋葬了，并在第三天复活了。} 此外，如果祂的肉身被否认，祂的死亡又如何能被断言？因为死亡是肉身固有的苦难，肉身按造物主的律法通过死亡回归它所出自的尘土。现在，如果祂的死亡因否认祂的肉身而被否认，那么祂的复活就没有确定性。因为祂没有复活，正因祂没有死，甚至因为祂没有拥有肉身的实在性，而死亡与复活都是属于肉身的。同样地，如果基督的复活被取消，我们的复活也被摧毁了。如果基督的\Emph{复活}没有实现，那么基督为此而来的目的也不会实现。因为正如那些说没有死人复活的人被宗徒从基督的复活所驳倒一样，如果基督的复活落空，死人的复活也被扫除了。如此，我们的信德便是空的，宗徒们的宣讲也是空的。而且，他们甚至显明自己是天主的假见证，因为他们见证天主复活了基督，而天主并没有复活祂。我们仍留在我们的罪中。\BibleRef{格前 15:13} 而那些在基督内安眠的人已经丧亡了；确实，他们命定要复活，但也许像基督一样以幻影状态复活。\par

% structure: p0019 source=h2 target=subsection reason=relative-heading-rank; base=h2

\subsection{第九章. 驳斥马吉安从天使及天主之子降生成人前的显现所引出的异议。}

如今，在你这个讨论中，当你设想我们必会遇到创造主的天使的例子时，就好像他们以幻象状态（即仅仅表象的肉身）与\Person{亚巴郎}和\Person{罗特}交往，却果真交谈、饮食和工作，正如他们受命所做的那样——但你将不被允许以你所摧毁的那位天主的行为作为例子。因为你的神越是更好、更完美的存在，他就越是与那位全然不同的神的一切例子不相称；没有这种差异，他根本不会更好或更完美。此外，其次，你必须知道，你不会被认可，认为在天使中只有表象的肉身，而非真实而坚固的人性实体。因为（按你的说法）在表象的肉身中彰显真实的知觉和行动对他并不困难，那么将这些真实的知觉和行动赋予真实的肉身实体（作为其正当的创造者和形成者）就更容易多了。但你的神，也许因为他根本没有造出任何肉身，所以引入了他无法产生实体的那个东西的单纯幻影，这倒很合理。然而，我的天主，祂用地上的尘土造出了那在真正的肉身品质中的（虽非来自婚配之种），也同样有能力将任何质料构成的肉身赋予天使，因为祂甚至从无中建造了世界，化作如此众多不同的形体，而且只凭一句话！实际上，如果你的神应许人有一天会拥有天使的真正本性（因为他曾说过：\BibleQuote{他们好像天使一样}\BibleRef{路 20:36}），那么我的神为什么不也从任何来源将人的真正实体赋予天使呢？因为你也不会告诉我，你那边天使的本性是从哪里来的；所以我只需断定，这对于神是合适且适当的——即那呈现于三种感官（视觉、触觉、听觉）的事物的真实性。对神而言，行骗比从任何质料（即使没有生育）产生真实肉身更难。然而，对于其他异端者，他们坚持认为天使中的肉身如果是真人的，就必须从肉身而生，我们有一个确切的答案：那确实是\Emph{人的}肉身，却\Emph{非受生}的。它确实是人的，因为神是真实的，既不能说谎也不能欺骗；并且（天使）除非以人的实体，否则无法以人的方式与人交往；它又是非受生的，因为除了基督，没有人能因肉身而生成为人，以便藉祂的诞生重造我们的出生，并进一步藉祂的死解除我们的死亡，并藉着在祂曾为了死而生的那肉身的复活。因此，那时祂自己与天使一起以真实的肉身显现给\Person{亚巴郎}，这肉身尚未经历出生，因为还不准备死，尽管它那时已经在学习与人交往。天使就更合适了，他们从未接受为我们而死的使命，也没有藉出生而获得肉身的短暂经验，因为他们不打算藉死再放下它；但无论他们从哪里获得肉身，无论后来怎样完全脱去它，他们从未假装它不是真实的肉身。因为创造主\BibleQuote{使祂的天使成为风，使祂的仆役成为火焰}——既有真实的灵，也有真实的火——同样也真实地使他们成为肉身；因此，我们现在可以回想并提醒异端者，祂曾应许有一天会将人形成为天使，而祂曾将天使形成为人。\par

% structure: p0021 source=h2 target=subsection reason=relative-heading-rank; base=h2

\subsection{第十章 真正的降生比\Person{马西昂}幻想的肉身更相称于天主。}

因此，既然你不得求助于创造主的任何例子（因为这些例子与主题无关，且有它们自身特殊的原因），我希望你自己陈述你的神的计划：祂展示祂的基督并非在真实的肉身中，这计划是什么。如果祂鄙视肉身为属土的（如你所说，充满粪土），为什么祂不也因此鄙视它的形像呢？因为任何本身不配受尊敬的东西，它的形像也不应受尊敬。自然状态如何，形像也如何。但祂怎能以人的实体的形像与人交往呢？既然如此，为什么不在现实（即真实实体）中，使祂的交往成为真实的？因为祂必须交往。而这一必要性本可以更好地服务于信仰而非欺诈！你所造的神是够可怜的，正因祂只能在不配的、甚至是外来的东西的肖像中显示祂的基督。在某些情况下，使用不配的东西（只要是我们自己的）是合适的，而使用即使很配但非自己的东西则很不合适。那么，为什么祂不采用其他更配的实体（尤其是祂自己的实体）而来，以免显得好像不得不使用不配和外在的实体？既然我的创造主甚至藉荆棘和火焰与人交往，后来又藉云柱和火柱，并在祂自己的显现中使用元素构成的形体，这些神力的例子充分证明，神不需要虚假或甚至真实的肉身作为工具。然而，如果我们仔细审视这个主题，实在没有一种实体配得上成为神的衣裳。无论祂乐意以什么穿上自己，祂都使之配得上自己——只是不能有虚假。那么，为什么祂会认为真实的肉身（而非虚假的）是一种耻辱呢？嗯，但祂凭对它的虚构反而荣耀了它。那么，这肉身是多么伟大，连它的幻影对于那更高级的神都是必需的！\par

% structure: p0023 source=h2 target=subsection reason=relative-heading-rank; base=h2

\subsection{第十一章 基督确实诞生；\Person{马西昂}为维护幻象诞生而进行的荒谬诡辩。}

\Person{马西昂}采纳了这些关于（他的）\Person{基督}的虚假形体的幻想，其用意在于使他的人性起源也不带有来自他人性实体的任何证据，从而使造物主的\Person{基督}可以自由地接受所有预言，这些预言把祂视为一个能够生于人世、因而具有肉身的存在。但这位本都的异端魁首\Person{马西昂}在这件事上也极为愚蠢。仿佛更容易相信的是：在神性存在中的肉身，宁可无生，也不可虚假；这种信念实际上早已被造物主的天使们所预备，他们曾以真实的、却非由出生而来的肉身说话。因为那臭名昭著的\Person{斐鲁麦娜}确实说服了\Person{阿佩莱斯}和其他脱离\Person{马西昂}的人，他们更相信\Person{基督}确实带着一个肉身的身体；但这身体并非由出生而来，而是从元素中借来的。如今，既然\Person{马西昂}担心对肉身的信仰会包含对出生的信仰，那么那看似为人者，无疑就被相信是真实而确实地出生了。因为有一个妇人曾高呼：\Quote{怀过你的胎，及你所吮吸过的乳房，是有福的！}\BibleRef{路 11:27}不然他们怎能说祂的母亲和祂的兄弟站在外面呢？\BibleRef{路 8:20}但我们在适当之处会详细讨论这一点。无疑，当祂自称为人子时，祂便已承认自己曾出生。现在我宁愿将所有这些问题提交到对福音的审查中；但尽管如此，如我已说过的，如果那看似为人者无论如何都必须被视为已出生，那么他若以为通过虚构肉身的伎俩就能完善对自己出生的信仰，便是徒劳的。因为一件被人信以为真的事，无论是他的肉身还是他的出生，不真实又有什么益处呢？或者，若你说，让人的意见算不得什么；那么你就是在欺骗的掩护下尊崇你的神，因为他知道自己与使人们对他所想的不同。那样的话，你甚至可能给他一个推定的出生，从而不把问题挂在这点上。因为愚昧的妇人有时会幻想自己怀孕，实则她们或因月经或因其他病症而肥胖。而且，无疑，既然他戴上了实体的假面具，就有责任从最初场景演完他幻想的戏，免得在肉身之初就演砸了。当然，你拒绝了出生的伪装，而拿出了真实的肉身本身。而且，无疑，即使一位神的真实出生也是一件最卑微的事。那么，来吧，将你的吹毛求疵倾泻在自然界最神圣可敬的作品上；抨击你的一切；摧毁肉身和生命的起源；称子宫为尊贵动物的下水道——换言之，即生产人的工厂；详述分娩的污秽可耻的痛苦，以及产褥本身肮脏、麻烦、可鄙的产物！然而，在你将这一切贬低到臭名昭著，以至宣称它们不堪于天主之后，出生对祂而言并不比死亡更糟，婴孩时期不比十字架更糟，刑罚不比本性更糟，定罪不比肉身更糟。如果\Person{基督}真实地承受了这一切，那么出生于祂而言是更小的事。如果\Person{基督}规避性地受难，如同幻影；那么祂也可能规避性地出生。这就是\Person{马西昂}用以杜撰另一个\Person{基督}的主要论据；我认为我们足够清楚地表明，这些论据完全不相干，因为我们教导说，与现实相比，天主在其中显示祂的\Person{基督}的那种状态的虚假，更真实地符合天主。因为祂是\Quote{真理}，祂是肉身；因为祂是肉身，祂便出生。因为当攻击的手段被摧毁时，这异端所攻击的论点反而得到了确认。因此，如果祂被视为在肉身中，因为祂出生了；而祂出生，因为祂在肉身中，并且因为祂不是幻影——那么结论就是：必须承认祂本身就是造物主的\Person{基督}，是造物主的先知们所预言的那一位，祂将带着肉身、经由人类出生过程而来。\par

% structure: p0025 source=h2 target=subsection reason=relative-heading-rank; base=h2

\subsection{第十二章。\Person{依撒意亚}关于\Person{厄玛奴耳}的预言。\Person{基督}享有此名。}

请先按你的惯常做法，挑战我们思考依撒意亚关于基督的描述，同时你争辩说它毫无吻合之处。因为，首先，你说依撒意亚的基督必须被称为\OriginalTerm{厄玛奴耳}{Emmanuel}；\BibleRef{依 7:14}然后，祂将取得大马士革的财富和撒玛黎雅的战利品，对抗亚述王。然而，那已经来临的那一位，并非以此名出生，也从未从事任何军事行动。不过，我必须提醒你，你应当审视这两段经文的上文下理。因为紧接着就给出了厄玛奴耳的解释：\Quote{天主与我们同在}；因此你不仅需要考虑那被说出的名字，还要考虑它的含义。这说法是希伯来语，\OriginalTerm{厄玛奴耳}{Emmanuel}，属于先知的本国语言；但这词的含义，\OriginalTerm{天主与我们同在}{God with us}，通过解释成了共同的财产。那么请问，这名字——与我们同在的天主——即厄玛奴耳，是否因基督光照了世界而常被用于基督的名字？我想你不会否认，因为你自己也承认祂被称为与我们同在的天主，即厄玛奴耳。否则，如果你如此愚蠢，因为在你那里祂被称呼为与我们同在的天主，而不是厄玛奴耳，因此你不愿意承认那一位已经来临，而祂的特性就是被称为厄玛奴耳，仿佛这与与我们同在的天主不是同一个名字——你将会发现，在\OriginalTerm{希伯来基督徒}{Hebrew Christians}中间，也在\OriginalTerm{马尔基昂派}{Marcionites}中间，当他们意指祂应被称为与我们同在的天主时，他们称祂为厄玛奴耳；正如每个民族，用任何他们用来表达与我们同在的天主的词语，都已称祂为厄玛奴耳，在含义上完成了声音。既然厄玛奴耳是与我们同在的天主，与我们同在的天主就是基督，祂在我们内（因为\BibleQuote{你们凡受洗归于基督的，就是穿上了基督}\BibleRef{迦 3:27}），那么基督含义上包含在这名字——与我们同在的天主——中，正如在发音上这名字是厄玛奴耳一样。因此，显然那被预言为厄玛奴耳的那一位已经来临，因为厄玛奴耳所意指的已经来临，即与我们同在的天主。\par

% structure: p0027 source=h2 target=subsection reason=relative-heading-rank; base=h2

\subsection{第十三章 依撒意亚的预言被考虑。基督母亲的童贞是一个记号。其他预言也是记号。先知多处经文中专名的比喻义。}

你同样被名字的声音引偏，当你如此理解\OriginalTerm{大马士革}{Damascus}的财富、\OriginalTerm{撒玛黎雅}{Samaria}的战利品和\OriginalTerm{亚述}{Assyria}王，仿佛它们预示创造者的基督是个战士，却不注意那经文中所含的许诺：\BibleQuote{因为孩子尚不知喊叫我父亲、我母亲以前，大马士革的财富和撒玛黎雅的战利品必被搬走，运到亚述王面前。}\BibleRef{依 8:4}你应当首先审视年龄的问题，看它是否能代表基督甚至还是婴孩，更不用说战士了。尽管，确实，祂可能以婴孩的哭叫来号召武装；可能不是用号角，而是用拨浪鼓来发出战争的警报；可能不是在马上、战车上或城墙上寻找敌人，而是在保姆的脖子上或奶妈的背上，就这样注定从祂母亲的乳房征服\OriginalTerm{大马士革}{Damascus}和\OriginalTerm{撒玛黎雅}{Samaria}！当然，当你们蛮荒\OriginalTerm{本都}{Pontus}的婴孩们冲入战阵时，情况就不同了。我想，他们是在学会撕裂之前先学会刺击；起初包裹在阳光和油膏中，后来配上袋子，以面包和黄油为口粮！既然自然当然不在任何地方容许人在生命之前学习战争，在知道父母名字之前掠夺大马士革的财富，那么这段经文必须被认为是比喻性的。好吧，但他说，自然不允许\Quote{\OriginalTerm{童贞}{virgin}女怀孕}，然而先知仍被相信。而且确实非常恰当；因为祂通过给出此事发生的原因——即它要成为记号——为令人难以置信的事被相信铺平了道路。\Quote{因此，}祂说，\BibleQuote{主要亲自给你们一个\Emph{\OriginalTerm{记号}{sign}}；看，童贞女要怀孕生子。}\BibleRef{依 7:14}来自天主的记号不会是一个记号，除非它是某种新奇而惊人的事物。然后，犹太的挑剔者们为了使我们困窘，大胆地假装圣经并非认为童贞女，而只是年轻女子要怀孕生子。然而，他们被这一考量所反驳：任何属于\Emph{\OriginalTerm{记号}{sign}}性质的事物都不可能出自日常发生的事情——年轻女子的怀孕和分娩。童贞母亲的提出被天主公正地视为\Emph{\OriginalTerm{记号}{sign}}，但好战的婴孩却没有类似的主张；因为即使在这种情况下，也不会出现记号的特征。但在那奇异而新颖的降生的记号被宣告之后，紧接着就宣布了一个记号——婴孩未来的历程，祂要吃奶油和蜂蜜。并不是说这确实是记号的性质，也不是祂\Quote{弃恶}；因为这也只是婴儿期的特征。但祂注定在亚述王面前夺取大马士革的财富和撒玛黎雅的战利品\Emph{无疑是一个奇妙的记号}。坚持祂年龄的尺度，寻求预言的旨趣，并将你因异端迟来而从那福音真理夺走的东西归还给它，预言立刻就变得可理解并宣告其自身的应验。让那些东方的\OriginalTerm{贤士}{magi}来侍奉新生的基督，在祂婴孩时期献给祂黄金和乳香的礼物；果然，一个婴孩将不战而取得大马士革的财富，且手无寸铁。\par

因为除了众所周知的事实，即东方的财富——也就是说，它的力量和资源——通常包括黄金和香料，同样确定的是，\OriginalTerm{创造者}{Creator}也使黄金成为其他民族的财富。因此，祂通过\Person{匝加利亚}说：\BibleQuote{犹大也必在耶路撒冷争战，并且要聚集四围列国的财富，就是金银。}\BibleRef{匝 14:14}此外，关于这黄金的礼物，\Person{达味}也说：\Quote{人要把阿拉伯的金子献给祂；}又说：\Quote{阿拉伯和舍巴的君王要献礼物给祂。}因为东方通常视\OriginalTerm{贤士}{magi}为君王；\Place{大马士革}在古代被认为是属于\Place{阿拉伯}的，后来在（罗马）划分叙利亚地区时划归叙利亚腓尼基。\Person{基督}随后接受了它的财富，当祂接受了它的象征——金子和香料时；而\Place{撒玛黎雅}的战利品就是贤士们本人。这些贤士发现了祂，并用礼物尊崇祂，屈膝朝拜祂作为他们的天主和君王，通过引导他们道路的星的见证，他们成为了撒玛黎雅的战利品，也就是说，偶像崇拜的战利品，因为显而易见，他们相信了\Person{基督}。祂以撒玛黎雅之名指代偶像崇拜，因为那座城市因其偶像崇拜而可耻，从\Person{雅洛贝罕}王时代起就背离了天主。这并非创造者（在祂的圣经中）比喻性地使用地名作为隐喻的罕见方式，这种隐喻源自它们罪恶的类比。例如，祂称犹太人的首领为\BibleQuote{\Place{索多玛}的统治者}，称整个民族为\BibleQuote{\Place{哈摩辣}的人民。}\BibleRef{依 1:10}在另一处祂也说：\BibleQuote{你的父亲是\Place{阿摩黎}人，你的母亲是\Place{赫特}人，}\BibleRef{则 16:3}因为他们在不义上的亲属关系；尽管祂实际上曾称他们为祂的儿子：\BibleQuote{我养育了\Emph{儿女}，并将他们抚养长大。}\BibleRef{依 1:2}同样地，在祂的意义中，\Place{埃及}有时理解为整个世界，因为它以迷信和诅咒为标记。类似地，在我们的（圣）\Person{若望}中，\Place{巴比伦}也是\Place{罗马}城的预象，如同（巴比伦）一样伟大、骄傲于王权，并迫害天主的圣徒。正是按照这种风格，祂称贤士为撒玛黎雅人，因为（正如我们所说）他们崇拜偶像，如同撒玛黎雅人一样。此外，对于短语\Quote{在\Place{亚述}王之前\Emph{或反对}他}，理解为\Quote{反对\Person{黑落德}}；贤士们当时反对黑落德，当时他们拒绝向他回报关于基督的事，因为黑落德正试图消灭基督。\par

% structure: p0030 source=h2 target=subsection reason=relative-heading-rank; base=h2

\subsection{第十四章。圣咏中某些默西亚预言的形象化风格。应用于基督的军事隐喻。}

我们这一解释将得到证实，当你假设基督在任何经文中被称为战士，因提及某种武器和此类表述，你仔细权衡这些词语其他含义的类比，并据此得出结论。\Quote{将你的剑佩在腰间}，达味说，\Quote{在你的大腿上。}但你刚才读到基督什么呢？\Quote{你比世人更美丽；恩宠倾注于你的唇间。}我觉得好笑，竟有人把娇美的容貌和优雅的嘴唇这些媚态归于一位必须佩剑作战者！同样，当接着说\Quote{愿你乘驾在威严中兴旺亨通}，理由接续道：\Quote{因着真理、温良和正义。}但谁能用刀剑产生\Emph{这些}结果，而不是相反——欺骗、严酷和伤害——必须承认，这些才是战争的本分？因此，让我们看看，是否有一种别的剑，其作用如此不同。现在，宗徒若望在默示录中，描述一把从天主口中出来的剑是\Quote{一把双刃利剑}。\BibleRef{默 1:16}这可理解为此剑是天主圣言，他因法律和福音两约而具有双刃——以智慧磨利，敌对魔鬼，装备我们对抗一切邪恶与贪欲的精神仇敌，并为了天主圣名将我们从最亲爱的事物中割舍。然而，如果你不承认若望，我们有我们共同的导师保禄，他\Quote{用真理束我们的腰，给我们穿上正义的胸甲，用和平福音的预备穿我们的脚，不是战争的；他命我们拿起信德的盾牌，藉此能熄灭魔鬼的一切火箭，戴上救恩的头盔，并拿起圣神的剑，（他说）即天主的圣言。}\BibleRef{弗 6:14}这把剑正是主亲自来送到地上的，而非和平。\BibleRef{玛 10:34}如果他是你的基督，那么他甚至是一位战士。如果他不是战士，而他挥动的剑是比喻性的，那么创造主在圣咏中的基督也可能被佩带圣言的象征之剑，没有任何军事装备。上面提到的祂美貌的\Quote{俊美}和\Quote{唇间的恩宠}完全适合这样一把剑，因为它早在达味的段落中就佩在祂的大腿上，并且有一天将被祂差遣到地上。因为这就是祂所说的：\Quote{愿你乘驾在威严中兴旺亨通}——\Emph{推进}祂的圣言进入各地，以召唤万民：注定要\Emph{兴旺}于接纳祂的信德的成功，并且\Emph{统治}，因祂藉复活战胜了死亡。\Quote{你的右手，}祂说，\Quote{将奇妙地引领你前行，}即你属神恩宠的力量，藉此基督的知识得以传播。\Quote{你的箭锐利；}你的训诫四处飞驰，你的警告和内心的谴责，刺痛并刺透每个人的良心。\Quote{万民将仆倒在你脚下，}即朝拜。因此，创造主的基督在战争中强大，并佩戴武器；同样，祂现在夺取战利品，不仅是撒玛黎雅，而且是万民。那么，承认祂的战利品是象征性的，因为你已经知道祂的武器是比喻性的。因此，既然主亲自说话，祂的宗徒也以比喻风格写下这些事，我们使用祂的解释并非鲁莽，这些记录的可靠性甚至我们的对手也承认；因此，只有\Emph{如此}程度的依撒意亚的基督已经来临，只有\Emph{如此}程度的非战士，因为依撒意亚所描述的祂并非这样的性格。\par

% structure: p0032 source=h2 target=subsection reason=relative-heading-rank; base=h2

\subsection{第十五章 基督这一称号适合作创造主之子的名称，但不适合马西昂的基督。}

关于祂肉身的讨论，以及由此而来的祂的诞生，并附带祂的名字厄玛奴耳，就此足够。至于祂的其他名字，尤其是基督之名，对方有什么可答辩的呢？倘若基督之名对你们而言如同天主之名一样普遍——以致两位神的圣子都可以恰当地称为基督，同样，每一位父也可以被称为主——那么，理性必定会反对这一论点。因为天主之名作为神性的自然称谓，可以归于所有那些声称具有神性的存在者，比如甚至偶像。宗徒说：\\BibleQuote{因为有人称为神，或在天上，或在地上。}\\BibleRef{格前 8:5} 然而，基督之名并非源于本性，而是源于安排；因此它成为祂的专有名字，因着安排而归属于祂。它也不能由任何其他神分享，尤其是作为对手的那一位，且其有自己的安排，那么祂也必须拥有与其他神不同的名字。因为，当人们为两位神设计了不同的安排后，他们如何能在这个差异的安排中承认名字的\\Emph{共同性}呢？然而，最能证明两位神是竞争对手的，莫过于发现他们的（不同）安排也伴有名字的差异。因为那不是\\Emph{不同}性质的状态，如若在称号的具体含义中没有明确标示。每当这些缺失时，就会出现希腊人所谓的\\OriginalTerm{错用}{katachresis}一词的情况，即不当用于不属于它的事物。然而，在天主内，我想，不应有缺陷，不应通过\\OriginalTerm{错用的}{katachrestic}滥用词语来建立祂的安排。这是哪位神，竟为其子从创造主那里要求名字？我说的不是不属于他的名字，而是古老且广为人知的名字，即使从这个角度看，它们对一位新奇且不为人知的神也是不合适的。再者，他告诉我们\\BibleQuote{没有人把新布补在旧衣服上}，或\\BibleQuote{没有人把新酒装在旧皮囊里}，\\BibleRef{玛 9:16} 而他自己却缝补并穿着一身旧名字呢？他如何从法律中撕裂福音，而他自己却完全披戴着法律——以基督之名，真的吗？什么阻止他用别的名字称呼自己，既然他宣讲另一个（福音），来自另一源头，且拒绝接受真实的身体，正是为了不让人以为他是创造主的基督？然而，他不愿看似是他愿意接受其名的那位，是徒劳的；因为即使他真正具有形体，若他没有接受祂的名字，他肯定更会避免被当作创造主的基督。但事实上，他拒绝了祂的名字所表明的实体真实性，即使他应通过自己的名字来证明这一真实性。因为基督的意思是\\Emph{受傅者}，而受傅无疑是身体的事。没有身体的人绝不可能受傅；绝不可能受傅的人绝不能被称为基督。如果他只是接受了名字的幻影，那又另当别论。但他问，他如何能潜入犹太人之中而被相信，除非通过一个他们惯用且熟悉的名字？那么你所描述的天主是反复无常、诡计多端的！用欺骗来推进任何计划，要么是不信任，要么是恶意的伎俩。假先知们反对创造主时，行为要坦率简单得多，他们以祂的名来作为自己的神。但我没有看到这种做法带来任何好处，因为他们更倾向于认为基督要么是他们自己的，要么是某种骗子，而不是另一位的基督；福音将证明这一点。\par

% structure: p0034 source=h2 target=subsection reason=relative-heading-rank; base=h2

\subsection{第十六章。神圣之名耶稣最适于创造主的基督。若苏厄是祂的预像。}

如果祂抓住\Emph{基督}这名，就像扒手抓住施舍篮，为何祂还愿被称为\Emph{耶稣}，这个犹太人并不那么期待的名字？因为我们藉天主的恩宠得以理解祂的奥秘，承认这名也原是为基督预备的，然而，尽管如此，犹太人却不知道此事，智慧已从他们那里被取走。简言之，直到今天，他们仍在等待基督，而非耶稣；并且他们认为\Person{厄里亚}就是基督，而非耶稣。因此，那位以基督不被期待的名字而来的，祂本可以只以那唯一为祂预备的名字而来。但因祂将两者——被期待的和不被期待的——混在一起，祂的双重计划便失败了。因为若祂是基督，恰恰是为了暗示自己是造物主的基督，那么耶稣就反对祂，因为在造物主的基督里，耶稣并不被寻找；或者若祂是耶稣，为要让自己能被归为另一位天主，那么基督就阻碍祂，因为基督不被期待属于造物主以外的任何一位。我不知道这两个名字中哪一个能站得住脚。然而，在造物主的基督里，两者都保持其位置，因为祂内也有一位耶稣被发现。你问怎么做到？那就从这里学吧，连同那些与你们同享异端的犹太人一起。当\Person{农}的儿子\Person{曷舍亚}被预定为\Person{梅瑟}的继承人时，他的旧名不是被更改，而是首次被称为\Person{若苏厄}吗？你说，确实如此。那么，我们首先注意到，这是那位要来者（基督）的预像。因为耶稣基督要引入一个新世代（因我们生在此世的旷野中）进入流奶流蜜的应许之地，即进入永生，没有比这更甘甜的了；同样，这事要成就，不是藉着\Person{梅瑟}，即不是藉着法律的规诫，而是藉着\Person{若苏厄}，藉着福音的恩宠，我们的割礼是由一把石刀完成的，即基督的割礼，因为基督是磐石（或石头）。因此那位被设立为此奥迹预像的伟大人物，实际上被祝圣，带有主自己名字的预像，被称为\Person{若苏厄}。就连那时，当基督与\Person{梅瑟}交谈时，祂亲自见证这名是祂自己的。因为与祂交谈的，除造物主的神——基督——之外，还能有谁？因此，当祂向百姓颁布这命令时：\BibleQuote{看，我派遣我的天使在你面前，在路上保护你，领你到我所为你预备的地方；你要注意他，听他的话，不要惹怒他；因为他没有回避你们，因为我的名在他身上，}\BibleRef{出 23:20}祂确实称他为\Emph{天使}，因为他将要行使的权能伟大，也因他的先知职务，宣布天主的旨意；但祂也称他为\Emph{\Person{若苏厄}}（耶稣），因为那是对祂自己未来名字的预像。祂常常确认祂这样赋予（祂仆人）的这个名；因为祂命令他今后惯常使用的名不是\Quote{天使}，也不是\Quote{曷舍亚}，而是\Quote{若苏厄}（耶稣）。因此，既然这两个名字都合于造物主的基督，它们就相应不合于非造物主的基督；同样，我们基督所预定的一切其余部分也是如此。简言之，今后在我们之间必须确立那个确定而公平的规则，对双方都必要，它应判定：另一位天主的基督与造物主的基督之间，决不应该有任何共同之处。因为你们将同样有需要维持他们的差异性，正如我们需要反对它，因为你们将同样无法证明另一位天主的基督已来临，除非你们证明祂与造物主的基督截然不同；同样我们也无法宣称（已来临的）祂为造物主的基督，除非我们证明祂正是造物主所指定的一位。现在关于他们的名字，这就是我们反对\Person{马西昂}的结论。我为自己主张基督；我为自己坚持耶稣。\par

% structure: p0036 source=h2 target=subsection reason=relative-heading-rank; base=h2

\subsection{第十七章 在\WorkTitle{依撒意亚}和\WorkTitle{圣咏}中关于基督贬抑的预言。}

让我们将祂其余的救恩计划与圣经比较。无论那可触摸、可见的卑微身体是什么，那都是我的基督，不论祂无荣耀、无高贵、无尊荣；因为祂的处境和形貌正是如此被预言的。\Person{依撒意亚}再次援助我们：\BibleQuote{我们已在祂面前宣布了（祂的道路），}他说；\BibleQuote{祂像仆人，像干旱之地的一根枝条；祂没有俊美，也没有华丽；我们看见祂，祂没有容貌，也没有美丽；但祂的容貌被轻视，比众人更受损毁。}同样，圣父稍早前对圣子说：\BibleQuote{许多人将因你而惊异，同样你的美丽在人面前也失去荣耀。}\BibleRef{依 52:14} 因为，虽然按\Person{达味}的话，祂比世人更美，但那是在属神恩宠的预表状态中，祂佩带着圣神的剑，这实在是祂的形像、美丽和光荣。然而，按照同一位先知，祂的体格\BibleQuote{是虫，不是人，是人的耻辱，百姓的弃物。}但祂并没有宣称自己有这类内在品质。在祂内住满了圣神；因此我承认祂是\BibleQuote{\Person{叶瑟}的茎干上发出的枝条。}祂绽放的花朵就是我的基督，按\Person{依撒意亚}的话，\BibleQuote{智慧和聪敏之神，超见和刚毅之神，明达和虔敬之神，以及敬畏上主之神，}安息在祂身上。\BibleRef{依 11:1} 现在，除了基督，没有人能合适地具备如此多样的属神证明。祂确实像花朵因圣神的恩宠，虽被算作\Person{叶瑟}的枝条，却从那里通过\Person{玛利亚}获得祂的世系。现在我特意问你：你是否承认祂承受这一切屈辱、苦难和安宁的命运？正是这一点使祂成为\Person{依撒意亚}所说的基督——一个受苦的人，熟悉忧患，被牵去宰杀，像羊羔在剪毛者前不出声；祂不争辩、不呼喊，街上也听不见祂的声音；祂不折断压伤的芦苇——即犹太人破碎的信德——也不熄灭冒烟的亚麻——即外邦人刚刚点燃的热忱。祂只能是那位被预言的人。应当按圣经的准则来检验祂的行为，若我没弄错，这行为以两种方式彰显：宣讲和奇迹。但我将安排这两者的处理，推迟到我决定讨论\Person{马尔西翁}实际福音的章节中去考虑祂奇妙的教义和奇迹——然而，这是为了我们当前的目的。让我们在这里以概括的方式完成我们已经开始的主题，在行文中指出基督如何被\Person{依撒意亚}预宣为宣讲者：\BibleQuote{因为你们中谁敬畏上主，}他说，\BibleQuote{听从祂圣子的声音？}\BibleRef{依 50:10} 同样作为治愈者：\BibleQuote{因为，}他说，\BibleQuote{祂承担了我们的病弱，背负了我们的痛苦。}\BibleRef{依 53:4}\par

% structure: p0038 source=h2 target=subsection reason=relative-heading-rank; base=h2

\subsection{第18章。—基督死亡的预像。\Person{依撒格}；\Person{若瑟}；\Person{雅各伯}反对\Person{西默盎}和\Person{肋未}；\Person{梅瑟}为对抗\Person{阿玛肋克}祈祷；铜蛇。}

关于祂的死亡，我想你试图引入一种分歧的看法，仅仅因为你否认创造主的基督曾预言了十字架的苦难，并且因为你争辩说，创造主不会使祂的圣子遭受祂自己曾诅咒过的那种死亡。祂说：\Quote{被诅咒的，是凡挂在树上的。} 但这个诅咒意味着什么，尽管它完全配得上简单预言的十字架，而我们目前主要探讨的正是这一点，我暂且不讨论，因为在另一处我们已经给出了理由，并且先有证据。首先，我将对预像作完整的解释。毫无疑问，这个奥迹应当通过预像以先知的方式表达出来，而且主要通过那种方法：因为按其不可信的程度，若以赤裸的预言表达出来，它就会成为绊脚石；而按其伟大程度，也有必要将其隐藏在阴影中，以致理解的困难可引导人祈求天主的恩宠。那么，首先，依撒格被他父亲献为祭品时，他自己背负了用于他死亡的木柴。这个行为甚至在那时就预示了基督的死亡，祂被圣父预定为祭献，并背起了祂受难时所负的十字架。同样，若瑟也是基督的预像，但并非在这点上（以免我拖延进程），即他因天主的事业遭受他弟兄们的迫害，如同基督按肉体遭受祂的弟兄——犹太人的迫害；而是当他被他父亲以这些话祝福时：\BibleQuote{他的荣耀有如牛犊；他的角是独角兽的角；他要以这些角撞击万民，直到地极。}\BibleRef{申 33:17}——他当然不是被描述为只有一只角的独角兽，也不是有两只角的弥诺陶罗斯；而是在他身上指出了基督——一头公牛，具有祂的两种特性：对于一些人严厉如同审判者，对于另一些人温和如同救主；祂的角就是祂十字架的末端。因为十字架的横木，其末端称为\Emph{角}；而整个结构的中间木桩则是\Emph{独角兽}。因此，藉着祂十字架的这种力量，并以此种方式\Emph{有角的}，祂现今正通过信德推动万民，将他们从大地带到天堂；将来则要通过审判推动他们，将他们从天主抛到大地。根据同一圣经的另一处，祂也将是一头公牛，当祂在属灵上被解释为雅各伯反对西默盎和肋未时，即反对经师和法利塞人；因为后者正是从前者衍生出来的。\Emph{如}西默盎和肋未，他们以他们的异端完成了他们的邪恶，并用这异端迫害基督。\Quote{让我的灵魂不要进入他们的计谋；让我的心不要与他们的集会联合：因为他们在愤怒中杀害了人，}即先知们；\Quote{并在他们的自恃中砍断了牛犊的筋，}即基督的筋。因为正是向祂他们发泄了愤怒，在杀了祂的先知之后，甚至用钉子将祂钉在十字架上。否则，在杀了人之后，又因折磨一头公牛而谴责他们，这是毫无意义的！再者，就梅瑟而言，为什么他在若苏厄与阿玛肋克作战的那个时刻，特别以伸开双手的姿势坐着祈祷，而在这样的战斗中，按理说更合适的是屈膝、捶胸、面伏于地，并以这样的俯伏姿势献上祈祷呢？这是为什么？乃是因为在奉那位有一天要与魔鬼作战的圣名进行的一场战斗中，那正是耶稣将要藉以得胜的十字架的样式是必要的。此外，为什么同一位梅瑟，在禁止一切肖像之后，却将铜蛇挂在木杆上，并当它挂在那里时，将其作为治愈的对象供人瞻仰呢？他这里不也是要显示吾主十字架的能力，藉此那条古蛇魔鬼被击败——藉此也向每一个被属灵的蛇咬伤的人，但若他以信德的目光转向它，就宣告了从罪咬伤中得治愈，并获得永远的康健吗？\par

% structure: p0040 source=h2 target=subsection reason=relative-heading-rank; base=h2

\subsection{第十九章 基督死亡的预言。}

来吧，当你在达味的言辞中读到\Quote{主从木架上为王}时，我想知道你是如何理解的。也许你以为是指犹太人的某个木制国王！——而不是指基督，祂因在十字架上受苦而战胜了死亡，并由此为王！现在，尽管死亡从亚当直到基督一直作王，为什么不能说基督从木架上为王，因为祂在祂十字架的木架上死去，封闭了死亡的国度呢？同样，依撒意亚也说：\Quote{因为有一个婴孩为我们诞生了。}\BibleRef{依 9:6} 但在这话中有什么不寻常之处，除非祂是指天主子？\Quote{有一子赐给了我们，祂肩上担负着主权。}\BibleRef{依 9:6} 现在，有哪一个国王是把统治的徽号放在肩上，而不是戴在头上作为冠冕，或拿在手中作为权杖，或作为某种王袍上的标志呢？但是，新世代的新王——耶稣基督，却把祂新光荣的权能和卓越都扛在肩上，那就是祂的十字架；这样，按照我们先前的预言，祂便可从此从木架上为主为王。耶肋米亚也向你暗示了这木架，当他向犹太人预言，他们将会说：\Quote{来吧，让我们毁坏这树连同它的果实（饼）}\BibleRef{耶 11:19} 那就是祂的身体。因为天主在你们自己的福音中也启示了这意义，当祂称祂的身体为\Emph{饼}；这样，从今以后，你便可理解，祂已将饼的形象赋予祂的身体，先知古时曾比喻性地将祂的身体变成了饼，而主自己也将要解释这奥秘。如果你还需要关于主十字架的进一步预言，第二十一篇圣咏足以提供给你，因为它包含了基督的全部受难，那时祂已预言性地宣告祂的光荣。祂说：\Quote{他们刺穿了我的手和我的脚}，这是十字架特有的残酷。再者，当祂恳求圣父援助时，祂说：\Quote{求你救我脱离狮子的口}，即死亡的牙关，\Quote{和我的卑微脱离独角兽的角}；换言之，就是脱离十字架的极端，如我们上文所示。现在，达味本人并没有遭受这十字架，犹太人的任何其他国王也没有；所以你不能认为这是关于任何其他人的受难的预言，唯独是祂，被这个民族如此显著地钉死在十字架上。如果异端者固执己见，拒绝并蔑视所有这些解释，我愿向他们承认，造物主没有给祂的基督的十字架任何记号；但他们不能从这一让步证明那位被钉者就是另一位（基督），除非他们能以某种方式表明，这死亡被他们自己的神预言为属于他的，以便从预言的差异中维持受难者的差异，从而也维持位格的差异。但由于连马西昂的基督都没有任何预言，更不用说他的十字架了，对我来说，我的基督只要有死亡的预言就够了。因为，既然死亡的\Emph{方式}没有宣告，十字架的死仍有可能被意指，那么如果预言指向另一位，这死亡就得归给另一位（基督）。此外，如果他不愿承认我的基督的死亡被预言，那么当他宣告他自己的基督确实死了，却否认他有诞生，而同时否认我的基督是可死的，却承认祂能够出生时，他的混乱就更大了。然而，我将向他显示我的基督的死亡、埋葬和复活都在依撒意亚的一句话中指明了，他说：\Quote{祂的坟墓从他们中间被移开了。}没有死亡就没有坟墓，没有复活就没有坟墓的移开。然后，他最后补充说：\Quote{因此，祂必得众多的人为产业，祂要分取众多者的战利品，因为祂倾注了祂的灵魂直到死亡。}\BibleRef{依 53:12} 因为这里陈述了这恩宠的原因，即作为祂受死之苦的补偿。同样也表明，祂将因祂的死得到这补偿，并且必在死后通过复活得到它。\par

% structure: p0042 source=h2 target=subsection reason=relative-heading-rank; base=h2

\subsection{第二十章——基督死亡在世上之后续影响之预言。达味的慈爱。这些是什么。}

至此，我在这些细节中追溯基督救恩计划的进程，已足以达到我的目的。这已证明祂就是预言所宣告的那一位，所以我们不应以预言所指定给祂的角色以外的任何其他角色来看待祂；而祂进程的事实与造物主的圣经之间的一致，其结果应使人们对圣经的信仰从那种偏见中得以恢复，这种偏见因助长意见分歧，已使其中相当一部分受到怀疑或遭否认。现在我们更进一步，从造物主的圣经中，建立起那些类似事件的上级结构，这些事件是被预言并注定在基督之后发生的。因为救恩计划若没有祂到来之后继续其进程，就不会被视为完备。看看万民，从人错误的漩涡中浮现出来，上升到神圣的造物主、神圣的基督，然后否认祂是预言的对象，如果你敢的话。圣父在圣咏中的许诺立刻会出现在你面前：\Quote{你是我的儿子，我今日生了你。你向我请求，我就将外邦人赐你作基业，将地极赐你作产业。}你将无法主张这里所指的是一些达味的儿子，而非基督；或主张地极是赐给达味的，其王国仅局限于犹太民族，而非赐给基督，祂如今在祂福音的信德中怀抱全世界。同样，祂通过依撒意亚说：\Quote{我立了你作万民的救恩计划，作外邦人的光，开启盲者的眼睛，}即那些在错误中的人，\Quote{将囚犯从监狱中领出，}即使他们从罪恶中得释放，\Quote{并从牢狱中，}即死亡的牢狱，\Quote{那些坐在黑暗中的人}——即无知的黑暗。\BibleRef{依 42:6}如果这些事是通过基督实现的，那么预言中它们就不会为任何其他者而设计，只为通过祂实现的那一位。在另一处，祂也说：\Quote{看，我已立祂为万民的见证，为万民的君王和统帅；不认识你的民族必呼求你，万民必奔向你。}\BibleRef{依 55:4}你不会将这些话解释为关于达味，因为祂先前曾说：\Quote{我要与你们订立永久的盟约，即达味确实的仁慈。}\BibleRef{依 55:3}的确，你将更必须从这些话理解到，基督因由童贞玛利亚诞生，而被视为按肉身出自达味。这关于这有关祂的许诺，圣咏中有向达味起的誓：\Quote{我必从你身体的果实而立在你宝座上。}所指的身体是什么？达味自己的身体？当然不是。因为达味不会生子。也不是他妻子的身体。因为若说\Quote{从你身体的果实}，他本应说\Quote{从你妻子身体的果实}。但是提到\Emph{他的}身体，可见祂是指向他种族中的某一位，基督的肉体要成为那身体的果实，这肉体从玛利亚胎中绽放。祂只称为身体的（子宫的）果实，因为那果实特别属于子宫，实际上只属于子宫，而不属于丈夫；祂将子宫（身体）归于达味，视为种族的首领和家族的祖先。因为将童贞女与丈夫结合不符合她的身份，所以祂将身体（子宫）归于父亲。那么，那在基督身上发现的新的救恩计划，将证明正是造物主当时以\Quote{达味确实的仁慈}之名所预许的，这仁慈属于基督，因为基督出自达味，或者说，祂的肉体本身就是达味的\Quote{确实的仁慈}，被宗教祝圣，并在复活后成为\Quote{确实}的。因此，纳堂先知在列王纪上，向达味许诺他的后裔，他说：\Quote{将从你的，}他说，\Quote{肠中生出。}现在，你若将此仅仅解释为关于撒罗满，你会让我大笑一场。因为达味显然生出了撒罗满！但此处基督岂非被称为达味的后裔，如同出自那源于达味的子宫，即玛利亚的子宫？现在，因为基督而非其他任何要建造天主的圣殿，即一个圣洁的人性，天主的圣神可以居住其中，如同在更好的圣殿中，所以应期待基督而非达味的儿子撒罗满为天主子。此外，永远的宝座和永远的王国更适合基督，而非撒罗满，一个暂时的君王。天主的仁慈也没有离开基督，而对于撒罗满，天主的愤怒却降临在他身上，在他奢侈和拜偶像之后。因为撒但激起一个厄东人作为他的敌人。因此，既然这些事无一与撒罗满相符，而只与基督相符，那么我们的解释方法必定正确；事实的结果本身也表明它们明显是为基督所预言的。因此，在祂身上，我们将拥有\Quote{达味确实的仁慈}。\Emph{祂}，而非达味，天主立为万民的见证；\Emph{祂}，为万民的君王和统帅，而非达味，他只统治以色列。正是基督，一切不认识祂的民族现在呼求祂；正是基督，一切种族现在投奔祂，他们之前不认识祂。不可能将你们（每日）看到正在实现的事说成是未来的。\par

% structure: p0044 source=h2 target=subsection reason=relative-heading-rank; base=h2

\subsection{外邦人的蒙召——在福音影响下所预言。}

因此，你无法从你这想法中为你的两个基督的区别找到根据，仿佛犹太人的基督是由造物主所定，只为使那民族从分散中复兴；而你的基督则是由至善的天主所派，为解放全人类。因为，毕竟最初的基督徒是站在造物主一边，而非\Person{玛尔西安}一边，万民都被召叫进入祂的国，这是因为天主从那（十字架的）树木建立了那国，那时还没有\Person{克尔东}出生，更不用说\Person{玛尔西安}了。然而，当你关于\Emph{万民}的召叫被驳斥时，你就转向\Emph{归化者}。你问道：万民中谁能转向造物主呢？因为先知所提及的那些人是归化者，各有不同的私人境况。依撒意亚说：\BibleQuote{看，归化者要藉着你到我这里来，} 表明那些藉着基督找到通往天主之路的人甚至就是归化者。但万民（外邦人）也像我们一样，同样被（先知）提及为信靠基督的人。他说：\BibleQuote{外邦人将因祂的名而信赖。} 此外，你在预言中以归化者替代万民，但归化者并不习惯于信靠基督的名，而是信靠梅瑟的律法，他们的教导来自梅瑟。然而，万民的拣选是从末后的日子开始的。依撒意亚正是用这些话说的：\BibleQuote{到末后的日子，上主的山，} 即天主的崇高，\BibleQuote{和天主的殿，} 即基督，天主公教会的圣殿，在其中天主受钦崇，\BibleQuote{必要矗立在群山之上，} 超越一切德能与权能的崇高；\BibleQuote{万民都要涌向它；将有许多民族前去说：来，我们攀登上主的山，到雅各伯的天主的殿里去；祂必教导我们祂的道路，我们就在祂的路上行走：因为法律将出自熙雍，上主的话将出自耶路撒冷。} \BibleRef{依 2:2} 福音将是这条\Quote{道路，} 即基督内的新法律和新话，不再是在梅瑟内。\BibleQuote{祂要在万民中施行审判，} 关于他们的错误。\BibleQuote{这些人要谴责一个大民族，} 即犹太人本身及其归化者。\BibleQuote{他们要将刀剑铸成犁头，将枪矛制成镰刀；} 换言之，他们要将伤害人的心思、敌对的口舌以及各种邪恶和亵渎，转变为温良与和平的追求。\BibleQuote{民族不再举刀攻击民族，} 不再挑起纷争。\BibleQuote{他们也不再学习战争，} \BibleRef{依 2:4} 因此你在此得知，基督被预许为并非在战争中强大，而是追求和平。现在你必定否认这些事曾被预言，尽管它们清晰可见；或者否认它们已经应验，尽管你读到它们；否则，如果你不能否认其中任何一点，那么它们必然已在祂——那位被预言者——身上应验了。因为请看祂的召叫从开始直到如今的整个过程：它是针对那些在末后的日子接近造物主天主的万民（外邦人），而不是针对归化者——归化者的拣选反倒是早期的事。确实，宗徒们已废除了你的那种信念。\par

% structure: p0046 source=h2 target=subsection reason=relative-heading-rank; base=h2

\subsection{第二十二章 宗徒的成功以及他们为福音所受的苦难，早已预言。}

宗徒的工作也早已被预言：\Quote{那传布和平福音、带来好消息的人，他们的脚是何等美丽！}, 而不是战争或坏消息。对此，圣咏相应说：\Quote{他们的声音传遍普世，他们的言语达到地极。} 就是说，那些将出自熙雍的法律和源自耶路撒冷的主的话四处传播之人的言语，为叫那所记载的得以应验：\Quote{那些远离我义德的人，已就近我的义德和真理。} 当宗徒们为此事束上腰时，他们弃绝了犹太人的长老、官长和司祭。他说：好，但难道他们首要的不是宣讲另一位神吗？恰恰相反，他们宣讲的是那位同一的天主，他们正竭尽全力实践祂的圣经！依撒意亚呼喊道：\Quote{离开，离开，}\Quote{从那里出去，不要触摸不洁之物，}即对基督的亵渎；\Quote{从她中间出来，}即从会堂中出来。\Quote{凡携带主器皿的人，务要分别出来。}\BibleRef{依 52:11} 因为主早已按照先知的前言，以祂的手臂启示了祂的圣者，即基督以大能显于万民眼前，使万民和地极都看见了来自天主的救恩。他们如此离开犹太教本身，以福音的自由交换法律的义务和重担时，正是应验了圣咏：\Quote{让我们挣断他们的捆绑，从我们身上抛弃他们的轭；}而这确实是在\Quote{外邦人咆哮，众民图谋虚妄的事；}在\Quote{地上的君王站定，官长聚集商议，敌对主和祂的基督。}之后发生的。宗徒们此后遭受了什么？你回答：各种不义的迫害，来自那些确属那位造物主的人，而这位造物主正是他们所宣讲之主的反对者。那么，为何造物主——若祂是基督的反对者——不仅预言宗徒们将遭受这些苦难，甚至对此表示不悦？因为祂既不应预言另一位神（如你所争辩，祂并不认识祂）的进程，也不应对祂曾悉心促成之事表示不悦。\Quote{看，义人怎样灭亡，没有人放在心上；虔诚的人怎样被收去，没有人思念。因为义人被接去，脱离了恶人。}\BibleRef{依 57:1} 这不是基督是谁？\Quote{他们说：来，让我们除掉义人，因为祂对我们不合宜，（且完全与我们的行为相反。）}\BibleRef{智 2:12} 因此，祂既预述又附带着基督受苦的事实，预言了祂的义人将同祂一样受苦——无论是宗徒还是后来的所有信友；祂以厄则克耳所说的同一印记标记了他们：\Quote{主对我说：你穿过城门，经过耶路撒冷中心，在那些人的额上划\OriginalTerm{陶乌}{Tau}记号。} 现在，希腊字母陶乌和我们的字母T正是十字架的形状，祂预言这将是我们额上的记号，在真正的普世耶路撒冷中，按照第二十一篇圣咏，基督的兄弟或天主的子女将把光荣归于天主圣父，借基督亲自对祂圣父所说的：\Quote{我要向我的弟兄宣扬祢的圣名，在会众中我要赞美祢。} 因为那些必须在我们的时代，以祂的名和借祂的圣神所应验的事，祂正确地预言了是出于祂。稍后祂又说：\Quote{在盛大的集会中，我的赞美源于祢。}在第六十七篇圣咏中祂又说：\Quote{在集会中赞美主天主。}因此，玛拉基亚的预言也与此一致：\Quote{万军的主说：我不喜欢你们，也不悦纳你们的供物。因为从日出之地到日落之处，我的名在外邦人中要受尊崇；在各地要献祭与我名，并奉献洁净的祭品。}\BibleRef{拉 1:10}诸如归光荣、祝福、赞美与颂歌。如今，既然这一切也在你们中间找到——额上的记号、教会的圣事、洁净祭品的奉献——你们现在应当爆发出来，宣告造物主的圣神预言了你们的基督。\par

% structure: p0048 source=h2 target=subsection reason=relative-heading-rank; base=h2

\subsection{第二十三章 犹太人因拒绝基督而遭分散与荒凉之境，早已预言。}

现在，既然你与犹太人一同否认他们的基督已经来临，也要回想他们被预言在基督之后将对自己招致的结局，因为他们不敬地拒绝并杀死了祂。这事从那天开始应验，当时，按依撒意亚所说，\BibleQuote{A man threw away his idols of gold and of silver, which they made into useless and hurtful objects of worship;}\BibleRef{依 2:20} 换言之，从基督显明真理后他抛掉偶像的时候。考虑一下先知接下来所写的是否已经应验：\BibleQuote{The Lord of hosts has taken away from Judah and from Jerusalem, among other things, both the prophet and the wise artificer;}即祂的圣神，祂建造教会，这教会确实是天主的殿、家和城。因为从此天主的恩宠在他们中间断绝了；\BibleQuote{the clouds were commanded to rain no rain upon the vineyard} of Sorech；即扣留天上的恩宠，使它们不降福于\BibleQuote{the house of Israel,} 这以色列家只产生了\BibleQuote{the thorns} 用以给主加冕，以及\BibleQuote{instead of righteousness, the cry} 用以将祂匆匆推向十字架。\BibleRef{依 5:6} 这样，法律和先知直到若翰为止，但神圣恩宠的露水从这民族中撤去了。此后，他们的疯狂仍然继续，主的名被他们亵渎，如圣经所说：\BibleQuote{Because of you my name is continually blasphemed among the nations}\BibleRef{依 52:5}（因为亵渎源自他们）；从提庇留到韦斯巴芗期间，他们没有学习悔改。因此\BibleQuote{has their land become desolate, their cities are burnt with fire, their country strangers are devouring before their own eyes; the daughter of Sion has been deserted like a cottage in a vineyard, or a lodge in a garden of cucumbers,}\BibleRef{依 1:7} 从那时起\BibleQuote{Israel acknowledged not the Lord, and the people understood Him not, but forsook Him, and provoked the Holy One of Israel unto anger.}\BibleRef{依 1:3} 同样，关于刀剑的威胁条件：\BibleQuote{If you refuse and hear me not, the sword shall devour you,}\BibleRef{依 1:20} 已经证明正是由于反叛基督，他们才灭亡。在第五十八篇圣咏中，祂向圣父要求分散他们：\BibleQuote{Scatter them in Your power.} 依撒意亚也说，在完成关于他们被火消灭的预言时说：\BibleQuote{Because of me has this happened to you; you shall lie down in sorrow.}\BibleRef{依 50:11} 但如果他们遭受这报应不是因祂的缘故，而祂曾在预言中指定他们的受苦归于自己的原因，而是为了另一位神的基督，那么这一切就毫无意义。那么，虽然你断言是另一位神的基督被造物主的掌权者和权势者（如同敌对者）钉在十字架上，我还是要说，请看造物主何等明显地保护了祂：祂得到了\BibleQuote{the wicked for His burial,} 就是那些竭力坚持祂的尸体被偷走的人，\BibleQuote{and the rich for His death,} 就是那些用金钱将祂从犹达斯的出卖以及士兵的谎言（说祂的身体被取走）中赎回的人。因此，这些事要么不是因祂的缘故临到犹太人，那样你就会因圣经的意义与事实的结果和时间的顺序相符而被驳倒；要么确实因祂的缘故发生，那么造物主除了为自己的基督，不可能施加这种报复；不仅如此，如果他们所杀的是祂主人的敌人，祂反而该给犹达斯赏报。无论如何，如果造物主的基督还没有来，预言因祂的缘故判他们遭受这样的苦难，那么当祂来临时，他们将要承受这些苦难。那时，哪里还有熙雍女子可以变成荒场？因为现在已经找不到了。哪里还有城市可以被火焚烧？因为它们已经成了废墟。哪里还有民族可以被分散？因为它已经在流亡中。让犹太回复原状，好使造物主的基督能找到它，然后你可以争辩说另一位基督已经来了。但此外，祂如何能允许主人在自己的天上漫游，而某一天却要杀死在同一地上的人？甚至祂王国的更尊贵荣耀之处已被侵犯，祂自己的宫殿和至高之处已被践踏？或者他只是表面上这样做？天主当然是忌邪的天主！然而他取得了胜利。你应该羞惭，你竟相信一个被击败的神！你对他能有什么盼望？他连自己都不能保护。因为他不是因软弱而被造物主的权势和人所粉碎，就是因恶意，想借忍受他们的恶行而在他们身上留下如此大的污点。\par

% structure: p0050 source=h2 target=subsection reason=relative-heading-rank; base=h2

\subsection{第二十五章 基督的千年与天上荣耀与祂的圣人们同在。}

是的，当然，你说，我确实从祂那里盼望那本身就证明基督的差异（不同基督）的事物，即天国，作为永远而属天的产业。此外，你的基督向犹太人应许他们最初的状态，以及他们国土的恢复；而在此生结束后，在阴府亚巴郎的怀抱中安息。哦，最卓越的天主，当祂以宽恕恢复祂在震怒中所夺取的！哦，你的天主是怎样的天主，祂既击伤又医治，创造灾祸又缔造和平！哦，何其慈悲直下至阴府的天主！关于亚巴郎的怀抱，我将在适当的地方有所论述。然而，至于犹太地的恢复，连犹太人自己，被地方和国家的名称所引导，正如所描述的那样盼望，详细陈述比喻性解释如何属灵地适用于基督和祂的教会，以及其性质和果实，将令人厌烦；此外，这主题已在另一著作中系统处理，我们题名为\WorkTitle{De Spe Fidelium}。目前，也因这个原因是多余的：我们的探讨涉及天上所应许的，而非地上的。但我们确实承认有一国度在地上应许给我们，虽然在天上之前，只是另一种存在状态；因为它将是复活后一千年，在神圣建造的耶路撒冷城中，\Quote{从天上降下来的} \BibleRef{默 21:2}，宗徒也称其为\Quote{我们天上的母亲} \BibleRef{迦 4:26}；并且，当宗徒宣告我们的\OriginalTerm{公民身份}{\GreekText{πολίτευμα}}在天上时，他预言它实质上是天上的一座城。厄则克耳知道这事，\BibleRef{则 48:30}，若望宗徒也看见了。\BibleRef{默 21:10} 而作为我们信仰一部分的新预言的话语，证实它如何预言了将有一个记号，即在显现之前这座城的图像被展示。事实上，这预言最近在一次东征中得以应验。因为即使外邦人的见证也清楚说明，在犹太地，连续四十天每天清早都有一座城悬在空中。随着时日推移，其城墙的整个轮廓逐渐消退，有时会瞬间消失。我们说，这座城是天主为迎接圣人复活而预备的，并以一切真实属灵福分的丰盛使他们焕然一新，作为对我们在世上所轻视或丧失的事物的补偿；因为祂的仆人应在他们为主的名受苦的地方获得喜乐，这是正义且合乎天主的尊荣。关于天上的国度，过程如下：在其千年之后——在此期间圣人的复活完成，他们按其功绩或早或晚复活——将接着发生审判时世界的毁灭和万物的燃烧：我们将在一瞬间被改变为天使的本质，甚至披上不朽的本性，如此被迁移到我们现在所讨论的天上国度，仿佛它从未被造物主预言过，仿佛证明基督属于另一位神，仿佛他是第一位且唯一的启示者。但现在要知道，它事实上已被造物主预言，并且即使没有预言，它对于关乎造物主的信德也有其要求。在你看来，亚巴郎的后裔，在最初的应许——如海沙般众多之后，又注定要与天上的星辰等同——这些岂不既是地上的也是天上的安排的标志吗？当依撒格祝福其子雅各布伯说：\Quote{愿天主赐给你天上的甘露，地上的肥腴} \BibleRef{创 27:28}，他的话中岂不有两类祝福的例子吗？的确，祝福的形式在此值得注意。因为就雅各布伯而言——他是那后来而更卓越之民（即我们）的预象——先是天上甘露的应许，然后是地上肥腴的应许。同样，当我们与世界分离时，先被邀请得天上的福分，然后我们才找到获得地上福分之路。你自己的福音也如此说：\Quote{你们先寻求天主的国，这些自会加给你们} \BibleRef{路 12:31}。但给厄撒乌的应许祝福是地上的，他在这地上肥腴之后补充了天上的，说：\Quote{你的居所也将得天上的甘露} \BibleRef{创 27:39}。因为犹太人的安排（他们包含在厄撒乌中，按出生是长子，但按情感是后来者）起初藉着法律充满地上的祝福，后来藉着福音由信德转向天上的祝福。当雅各布伯在梦中看见一个梯子立在地上，顶端通到天上，有天使在梯子上上去下来，主立在其上时，我们将毫不犹豫地大胆推测，主藉此梯子在审判中命定，通往天堂的道路指示给人，有些人可达，有些人则跌落。为何他一从睡梦中醒来，因惧怕那地方而战栗，就陷入对梦的诠释？他喊道：\Quote{这地方何等可怖！}然后补充说：\Quote{这决非别处，乃是天主的殿，是天堂的门！} \BibleRef{创 28:12} 因为他看见了主基督，天主的殿，也是进入天堂的门。如果创造主的安排中不能进入天堂，他当然不会提及天堂的门。但现在基督提供了一扇门，它接纳并引向\Emph{荣耀}。关于此，亚毛斯说：\Quote{祂在建造祂升入天堂的阶梯} \BibleRef{亚 9:6}；当然不仅为祂自己，也为与祂同在的祂的子民。\Quote{你要将它们束在身上}，他说，\Quote{如同新娘的装饰}。\BibleRef{依 49:18} 因此，圣神赞赏那些藉这些阶梯翱翔于天域的人，说：\Quote{他们飞翔，仿佛鸢鸟；他们如云彩，又如雏鸽，向我飞来} \BibleRef{依 60:8}——也就是说，就像鸽子一样。因为我们，按照宗徒的话，将被提到云中，迎接主（即人子，按达尼尔所说驾云而来，\BibleRef{达 7:13}），如此我们就常与主同在 \BibleRef{得前 4:17}，只要祂仍同时在天上及地上，祂对那些对两个应许都不感恩的人，呼天唤地作证：\Quote{诸天哪，请听；大地啊，请侧耳} \BibleRef{依 1:2}。事实上，就我而言，即使圣经不向我伸出天上希望之手（尽管它确实常这样做），在我目前享受地上恩赐时，我仍会对这应许有足够的假定；而且我会从那位同时也是天地之主的天主那里，盼望某种天上的事物。我因此会相信，那应许更高福分的基督，就是那也应许较低福分的那一位（之子）；祂尚且透过较小的恩赐证明更大的恩赐；祂为祂的基督独自保留了这（也许）前所未有的国度的启示，为的是，当地上的荣耀由祂的仆人宣告时，天上的荣耀可让天主亲自作其使者。然而，\Emph{你}却因祂宣告一个新国度，而主张另一个基督。你应先举出祂恩慈的某些实例，使我无充分理由怀疑那巨大应许的可信性——你说是当盼望的；不，首要之事是你应证明天堂属于祂，你声称祂是天上事物的应许者。现在，你邀请我们用膳，却不指出你的房子；你声称有国度，却不显示任何王权。难道你的基督应许天国，却没有天堂，如同祂显示为人，却没有肉身？哦，从头到尾是何等的幻影！哦，何等的巨大应许的空洞伪装！\par

\SourceRecord{Translated by Peter Holmes.；Ante-Nicene Fathers；Vol. 3.；Edited by Alexander Roberts, James Donaldson, and A. Cleveland Coxe.；Buffalo, NY: Christian Literature Publishing Co.,；1885.；<http://www.newadvent.org/fathers/03123.htm>.}

% structure: p0001 source=h1 target=section reason=page-title

\section{驳斥马西昂，卷四}

\Emph{在此，\Person{特土良}继续他的论证。\Person{耶稣}就是创造主的\Person{基督}。他从\BibleBook{路加}{福音}汲取证据；那是\Person{马西昂}部分接受的\WorkTitle{新约}唯一的历史部分。本书也可视为对\BibleBook{路加}{福音}的注释。它出色地证明了\Person{特土良}对\WorkTitle{圣经}的把握，并证明\Quote{\WorkTitle{旧约}与\WorkTitle{新约}并不相悖。}它还充满对圣经段落引人注目的阐释，包含关于启示的深刻见解，并与人的本性相关联。}\par

% structure: p0003 source=h2 target=subsection reason=relative-heading-rank; base=h2

\subsection{第一章。将\Person{马西昂}的\OriginalTerm{对立论}{Antitheses}置于他自己的福音中检验。\WorkTitle{旧约}与\WorkTitle{新约}时期中某些真正的\OriginalTerm{对立论}{Antitheses}。这些变化与颁布它们的同一真天主完全相容。}

如今我们将不虔敬且亵圣的玛尔齐盎的一切意见和全套方案，提交给那部福音本身来检验——这部福音经过他的篡改已成了他自己的。为了让人相信这部福音，他竟为它设计了一种妆奁，即一部由相互对立的陈述构成的著作，因而题名为\WorkTitle{对立命题}，其编纂目的在于将法律与福音如此割裂，以致将神性分为两位，甚至两位不同的神——每位对应一部圣约书，或更通常所说的圣约；这样他也能庇护对\Quote{按照\WorkTitle{对立命题}的福音}的信仰。然而，我本打算以特殊格斗方式，即徒手作战来攻击这些\WorkTitle{对立命题}；也就是说，我本想逐一迎战这位本都异端者的几个诡计，若不是更方便地在那部它们所支持的福音本身之中、并与之一同驳斥它们的话。尽管很容易立即以断然的异议驳回它们，但为了既使它们在论证中可被接纳，又视它们为有效的意见表达，甚至主张它们有利于我方——好让它们作者的盲目蒙受更深的羞耻——我们现在已拟定了我们自己的若干\WorkTitle{对立命题}以对抗玛尔齐盎。事实上，我确实承认，在旧约时代，在创造主之下，有一种秩序曾运行其进程；而在新时代，在基督之下，另一种秩序正在运行。我不否认，它们的文献用语、德行训诫和法律教导之间存在差异；然而，所有这些多样性都与同一位天主，即那位安排并预言这一切的天主相符。早在很久以前，依撒意亚就曾宣告：\Quote{因为法律将出自熙雍，上主的话将出自耶路撒冷}\BibleRef{依 2:3}——即另一种法律和另一句话。简言之，他说：\Quote{他将统治万民，并斥责许多民族}\BibleRef{依 2:4}；这不单指犹太民族，也指那些被福音的新法律和宗徒的新话所审判的万民，他们一旦相信，便在自己的旧错误中受到斥责。其结果是：\Quote{他们将刀剑铸成犁头，将长矛（一种打猎工具）铸成镰刀}\BibleRef{依 2:4}；也就是说，曾经凶残的心灵被他们转变为善良的品性，结出美果。又说：\Quote{你们要听我，听我，我的百姓；列王啊，要侧耳听我；因为法律将从我发出，我的正义将为万民之光}；因此祂决定并规定万民也要藉福音的法律和话得到光照。这将是那法律（按达味之言）无可指摘，因为\Quote{完美的法律能回转灵魂}，使之从偶像转向天主。这同样将是那话，关于它，同一位依撒意亚说：\Quote{因为上主将在地上的国中成就决定性的话}。因为新约简洁而短，摆脱了法律琐碎而混乱的负担。但何必详述呢？既然创造主藉同一位先知，比光本身更显明、更清晰地预言了革新？\Quote{不要追念先前的事，也不要思虑古旧的事（旧事已过，新事正在发生）。看，我要做新事，它们即将发生}\BibleRef{依 43:18}。同样藉耶肋米亚：\Quote{你们要为自己开垦新的田园，不要撒种在荆棘中；要在你们的心皮上行割损。}在另一处：\Quote{看，日子将到，上主说，我必要与以色列家和犹大家订立新约；不像我当日拉着他们祖先的手，领他们出离埃及地时与他们所立的约}。祂表明古约只是暂时的，当祂指示其变更时；也当祂应许随后将有一个永久的约时。因为藉依撒意亚祂说：\Quote{你们要听我，你们必得生存；我要与你们订立永久的盟约}，加上\Quote{达味确实的仁慈}\BibleRef{依 55:3}，为要表明那约将在基督内完成。基督出自达味的家族，按玛利亚的族谱，祂甚至藉那将要从叶瑟的树干发出的枝条，以预表的方式宣告了这事\BibleRef{依 11:1}。既然因此祂说，从创造主那里将有别的法律、别的话和新的约的安排，并藉玛拉基亚指出，连祭献本身也要在万民中领受更崇高的职务，当他说：\Quote{我没有喜悦在你们身上，上主说，也不从你们手中收纳你们的祭品。因为从日出之地到日落之地，我的名在万民中要受尊崇，在各处都要向我的名献上祭品，就是洁净的供物}\BibleRef{拉 1:10}——意指出自纯洁良心的简单祈祷——那么，凡因革新而来的每一种改变，必然在那些被改变的事物中引入多样性，从多样性又产生对立。因为没有什么东西在改变后不变得不同，而不同的东西没有不是对立的。因此，对于那个因革新而经历状态改变的事物，将可断言一种因多样性而产生的对立。带来改变的那一位，也设立了多样性；预言革新的那一位，也预告了那对立。何以在你的解释中，你把事物的状态差异归因于权能的差异？你为何将那些你从中引出\WorkTitle{对立命题}的例子曲解为对创造主不利，而实际上你可以在祂的情感和情绪中认出它们全部？祂说：\Quote{我打伤，我也医治}；又说：\Quote{我杀死，我也使人活}\BibleRef{申 32:39}——同样是那位\Quote{创造灾祸、缔造和平}者\BibleRef{依 45:7}；你甚至习惯以此指责祂喜怒无常、善变，仿佛祂禁止祂所命令的，命令祂所禁止的。那么，你为何没有列举那些也出现在创造主自然作品中的\WorkTitle{对立命题}呢？祂常与自己对立。除非我搞错了，你未能认识到，至少在你本都人中，世界是由相互敌对的多样元素构成的。因此，你的本分是首先确定光之神是一位，黑暗之神是另一位，这样你才能明确断言其中一位是法律之神，另一位是福音之神。然而，从显明的证据中，我的心灵已确信：正如祂的工程和计划以\WorkTitle{对立命题}方式存在，同样，祂宗教的奥秘也按同一规则存在。\par

% structure: p0005 source=h2 target=subsection reason=relative-heading-rank; base=h2

\subsection{第二章 \Person{圣路加}福音被\Person{马尔基翁}选为权威并被其肢解。其他福音同样具有权威。然而，接受\Person{马尔基翁}的讨论条件，并仅在\Person{圣路加}福音的基础上应对。}

你现在已经被我们简明扼要地指出了对\WorkTitle{反命题}的答复。我转而给出对福音的证明——自然不是属\Place{犹太}的，而是属\Place{本都}的——它同时被证明是败坏的；这将表明我们进展的顺序。我们首先确立：福音的新约是宗徒们所著，主亲自将宣讲福音的任务托付给他们。然而，也有宗徒式的人物，但他们并非独存，而是与宗徒们一同出现，且在宗徒们之后；因为门徒的宣讲可能遭受沽名钓誉的怀疑，如果没有师傅的权威伴随，即基督的权威伴随，因为是基督使宗徒成为他们的师傅。因此，在宗徒中，\Person{若望}和\Person{玛窦}首先向我们灌输信德；在宗徒式人物中，\Person{路加}和\Person{玛尔谷}后来重新更新它。这些人从同一信德原则出发，关系到唯一创造主天主和祂的基督，即祂由童贞女所生，来满全法律和先知。若他们叙述的顺序有些变动亦无关紧要，只要在信德要事上一致即可，而在此点上\Person{马尔基翁}是不同意的。另一方面，你要知道，\Person{马尔基翁}不为他的福音指定作者，仿佛不准他给福音标题，而撕毁（他眼中的）福音本体却不算罪行。现在我本可以在此停步，主张一部作品不应被认可，它不能昂首挺立，没有连贯性，从其标题的充份和作者的正名来看，未显出可信的希望。但我们宁愿在每一点上辩论；也不忽略任何可以公平地认为对我们有利之处。现在，在我们拥有的作者中，\Person{马尔基翁}似乎特别选中了\Person{路加}来施行其肢解过程。但\Person{路加}并非宗徒，仅是宗徒式人物；不是师傅，而是门徒，因此低于师傅——至少远远晚于他所跟随的宗徒（无疑是\Person{保禄}）晚于其他宗徒；因此，即使\Person{马尔基翁}以\Person{圣保禄}自己的名义出版他的福音，单凭该文件的权威，缺乏所有前代权威的支持，也不足以为我们的信德提供充分基础。还需要存在\Person{圣保禄}所发现的福音，他曾相信这福音，并如此热切地希望自己的福音与之相符，以至于他因此上耶路撒冷去认识并与宗徒们商议，\BibleQuote{免得他白跑，或曾经白跑了}\BibleRef{迦 2:2}；换句话说，他所学到的信德和他所宣讲的福音可能与他们的相符。于是，最后，在与（最初的）作者们商议后，并同意他们对信德准则的观点后，他们彼此握手共融，此后在宣讲福音的职务上分工，以便他们去犹太人那里，而\Person{圣保禄}去犹太人和外邦人那里。因此，既然\Person{圣路加}本人的启蒙者为他自己的信德和宣讲都渴望前辈们的权威，那么我对路加福音所要求的，对其师傅的福音是必要的，岂不更应如此？\par

% structure: p0007 source=h2 target=subsection reason=relative-heading-rank; base=h2

\subsection{第三章 \Person{马尔基翁}暗示某些宗徒不可信，这些宗徒曾受\Person{圣保禄}责备。该责备表明这并不能被视为贬低他们的权威。宗徒的福音完美可信。}

反之，在马西昂的体系中，基督宗教的奥迹始于路加的门徒身份。然而，既然它在此之前早已展开，它必然有其本身的真确材料，藉此才得以传到圣路加手中；而藉着它所承载的见证之助，路加本身才得以被接纳。可是，马西昂发现保禄致迦拉达人书（其中他责备宗徒们\Quote{没有按福音的真理正直行事}，并指控某些假宗徒歪曲基督的福音），便竭尽全力摧毁那些以宗徒名义发表的真福音书之信誉，以便为自己的福音赢得从它们那里夺去的声望。但即便如此，他责备伯多禄、若望和雅各伯——那些被视作柱石的人——也是有明显原因的。他们似乎因偏袒而改变了交往。然而，正如保禄自己\BibleQuote{成了众人所成的一切}\BibleRef{格前 9:22}，为赢得众人，伯多禄也可能采取了同样的策略，即实行与他教导稍有不同的事。同样，如果假宗徒也混了进来，他们的品性也表现在坚持割损和犹太仪式上。因此，圣保禄标记他们不是因他们的宣讲，而是因他们的行为；若他们在有关造物主天主或祂的基督方面有任何错误，他也会同样公正地加以谴责。因此，必须区分各种情况。当马西昂抱怨宗徒们（因他们的变节和虚伪）甚至败坏了福音时，他是在控告基督，因为他控告了基督所拣选的人。那么，如果宗徒们仅仅因行为不一致而受责备，却撰写了纯正的福音，而假宗徒却篡改了他们的真实记录；如果我们手中的抄本是从这些篡改的版本而来，那么宗徒著作的未经玷污的真确文本又在哪里？是什么启发了保禄，并通过他启发了路加？它要么完全被抹去，如同被洪水淹没——被伪造者的洪流冲走——那样连马西昂也不拥有真福音；要么，马西昂独自拥有的那个\Emph{版本}就是真确的，即宗徒们的版本？那么，它如何与我们手中的版本一致呢？后者据说不是宗徒之作，而是路加之作？或者，如果马西昂使用的版本不能仅仅因为与我们的一致（当然，我们的版本也在标题上遭到篡改）而归给路加，那它就是宗徒之作。因此，与我们一致的福音同样是宗徒之作，但也在标题上遭到篡改。\par

% structure: p0009 source=h2 target=subsection reason=relative-heading-rank; base=h2

\subsection{第四章。双方都声称拥有真福音。古老性是判断此类问题的真理标准。马西昂自命为福音的修订者。}

那么，我们必须跟随我们讨论的线索，以相应的精力迎接对手的每一次努力。我说，\Emph{我的}福音是真实的；马西翁说，\Emph{他的}是真实的。我断言马西翁的福音是被篡改的；马西翁断言我的福音是被篡改的。现在，有什么可以为我们解决这个问题呢？除非是\Emph{时间}的原则。这一原则规定，权威在于那些被证明更古老的事物；并且假设一个基本真理，即（教义的）败坏属于那些被证明在起源上相对较晚的一方。因为，既然错误是对真理的歪曲，那么真理必须存在于错误之前。事物必然先于它遭受任何事故而存在；对象必然先于一切对它的竞争而出现。否则，这将何等荒谬：当我们证明我们的立场是更古老的，而马西翁的立场是较晚的，我们的立场却显得是虚假的，甚至在它从真理获得客观存在之前；而马西翁的立场却被认为在我们手中经历了竞争，甚至在其发布之前；最后，那个较晚的立场反而被认为更真实——它比基督教宗教所有许多伟大事实和记录的出版晚了一个世纪，而后者当然不可能在\Emph{没有}——也就是说，\Emph{之前}——福音真理的情况下出版。那么，关于悬而未决的问题，即\BibleBook{路加}{福音}（只要它作为我们和马西翁的共同财产，能够决定真理），其我们单独接受的部分远比马西翁古老，以至于马西翁本人曾经相信过它，当时他在初信的热忱中向天主教会捐款，而该教会后来与他一同被拒绝，当他从我们的真理坠入他自己的异端时。倘若马西翁派信徒否认他曾在我们中间持有原始信仰，甚至否认他自己的书信呢？倘若他们不承认那封信呢？无论如何，他们接受他的\Emph{\WorkTitle{对立论}}；不仅如此，他们还炫耀地使用它们。从这些中取得的证据对我来说已经足够。因为如果在我们中间流传的被称为\BibleBook{路加}{福音}（我们将看看它是否也流传在马西翁中间）正是马西翁在他的\Emph{\WorkTitle{对立论}}中所争论的那一本，被犹太主义的捍卫者篡改，目的为了将法律与先知与它混杂在一起，以便他们能从中塑造他们的基督，那么他当然不可能如此争论，除非他曾发现它是（这样的形式）。没有人会在事物存在之前批评它们，当他不知道它们是否会实现。修正从不先于错误。确实，那本从提庇留时代到安敦尼时代一直混乱不堪的福音书的修正者，唯独在马西翁身上首次出现——基督如此长久地期待他，却一直后悔自己如此匆忙地派遣了他的宗徒而没有马西翁的支持！但尽管如此，异端永远在修正福音书，并在行动中败坏它们，这是人的胆大妄为，而非天主的权威；若马西翁甚至是一个门徒，他仍不是\BibleQuote{不高于他的老师；}\BibleRef{玛 10:24}若马西翁是一个宗徒，仍然如保禄所说，\BibleQuote{不拘是我，或是他们，我们都这样传；}\BibleRef{格前 15:11}若马西翁是一个先知，甚至\BibleQuote{先知的神魂将服从先知，}\BibleRef{格前 14:32}因为他们不是混乱的作者，而是和平的作者；或者若马西翁实际上是一个天使，他必更被指定为\BibleQuote{当受诅咒，而不是福音的宣讲者，}\BibleRef{迦 1:8}因为他所宣讲的是另一个福音。所以，当他修正时，他只确认了两个立场：既确认我们的福音是先前的，因为他修正了他先前遇到的那个；又确认\Emph{那个}是较晚的，他通过汇集来自我们福音的修正，从而制作了他自己的福音，一个新颖的福音。\par

% structure: p0011 source=h2 target=subsection reason=relative-heading-rank; base=h2

\subsection{第五章 依古据之准则，天主教福音书被证实为真，其中包括真正的\BibleBook{路加}{福音}。马西翁的福音仅是残缺的版本。异端者因忽略其他福音书而显其软弱与矛盾。}

因此，总而言之，如果更早的显然更真实，如果从起初而来的更早，如果由宗徒为作者的那是从起初而来的，那么，那在宗徒的教会中作为神圣宝藏被保存的，当然同样显然是从宗徒传承下来的。让我们看看格林多人从保禄饮了何种乳汁；迦拉达人被带到何种\OriginalTerm{\Emph{信德}准则}{rule of faith}下接受矫正；斐理伯人、得撒洛尼人、厄弗所人\Emph{据此}读了什么；还有罗马人，他们离宗徒如此之近，伯多禄和保禄共同将福音遗赠给他们，甚至以自己的血封印，他们有什么宣告。我们还有\Emph{圣}若望的养育教会。因为虽然马尔强拒绝了他的\BibleBook{若望}{默示录}，但（那里的）主教职序，当追溯至其起源时，仍会以若望为作者。同样，其他教会的卓越根源也被承认。因此，我说，在它们内（不仅是在那些由宗徒建立的教会内，而是在所有那些在\OriginalTerm{\Emph{基督福音的}奥秘}{mystery of the gospel of Christ}中共融的教会内），我们全力捍卫的\BibleBook{路加}{福音}从最初出版之日起就立住了脚；而马尔强福音并不为大多数人所知，并且没有任何人知道它而不同时加以谴责。当然，它也有自己的教会，但特别属于它自己——尽管它们迟来且是伪造的；如果你想了解它们的起源，你更容易在其中发现背教而非宗徒传承，因为马尔强诚然是它们的创始人，或马尔强蜂群中的某一位。连黄蜂也造蜂巢；同样，这些马尔强派也造教会。宗徒教会的同一权威也为其他福音提供了证据，我们同样通过它们并根据它们的使用而拥有这些福音——我是指\BibleBook{若望}{福音}和\BibleBook{玛窦}{福音}——而马尔谷所发表的可以确认为伯多禄的，因为马尔谷是他的翻译者。因为就连\BibleBook{路加}{福音}的形式，人们通常也归功于保禄。并且，门徒所发表的作品属于他们的师傅，这似乎很合理。那么，马尔强也应当就这些（其他福音）受到严格质问，因为他省略了它们，而偏执地坚持路加福音；好像它们也从起初就没有在教会中自由流传，如同路加福音一样。不，它们从起初就存在是更可信的；因为，作为宗徒的作品，它们更早，并且在起源上与教会本身同时代。但若宗徒们没有发表任何作品，怎么会出现他们的门徒更热衷于这样的工作呢？因为他们若没有从师傅那里得到任何教导，就不可能成为门徒。那么，如果显然这些（福音）也在教会中流传，为什么马尔强不动它们——要么如果它们被篡改就修正它们，要么如果它们未被败坏就承认它们？因为，那些歪曲福音的人，自然会对那些他们知道其权威更被普遍接受的事物更加关心其歪曲。就连假宗徒（如此称呼）也正是因为这个原因，因为他们通过其伪造来模仿宗徒。因此，\Emph{一}如他可能修正了那些需要修正之处（如果发现是败坏的），\Emph{同}样，他也坚定地暗示，凡他认为不需要修正的，都是没有败坏的。简言之，他只是修正了他认为败坏的东西；尽管，确实，连这一点也不正当，因为它并非真的败坏。因为如果宗徒的（福音）完整地传给了我们，而\BibleBook{路加}{福音}（在我们中被接受）与他们的准则如此一致，以至于在教会中接受的持久性上与它们等同，那么显然可见，\BibleBook{路加}{福音}也同样完整地传给了我们，直到马尔强亵渎性的处理。简言之，当马尔强对它下手时，它就变得与宗徒的福音不同且敌对。因此，我要劝告他的追随者，他们要么将这些福音（尽管现在为时已晚）改变为与他们自己的福音一致，从而显得与宗徒著作相符（因为他们每日都在修改自己的作品，正如他们每日被我们定罪）；要么他们为他们的师傅感到羞愧，他在两种情况下都自我定罪——一旦他良心受责而交付福音的真理，或再次通过无耻的篡改来颠覆它。这些就是我们为福音的信德拿起武器反对异端者时使用的概要论证，既坚持时间的秩序（它规定晚期是伪造者的标记），也坚持教会的权威（它支持宗徒的传统）；因为真理必然先于伪造，并且直接从那些传递它的人而来。\par

% structure: p0013 source=h2 target=subsection reason=relative-heading-rank; base=h2

\subsection{第六章。马尔强篡改福音的动机。创造者的基督与福音的基督之间没有区别。不可接受另一位基督。真正的基督与旧约安排的关系得到确认。}

但我们更进一步向前推进，并（如我们所许诺的）挑战马尔西昂自己的福音本身，意在借此证明它已被篡改。因为可以肯定，他在编写他的《\WorkTitle{Antitheses}》时竭力追求的全部目标，都集中于此：他要在旧约和新约之间建立一种差异，好使他的基督与造物主分开，作为属于这个对立的神的基督，\Emph{并}与法律和先知无关。同样可以肯定的是，出于这一观点，他删除了所有与他自己的意见相悖且有利于造物主的内容，仿佛它们是由造物主的拥护者插补的；而所有与他自己的意见一致的内容，他都保留了。我们将严格审查后一类陈述；如果它们结果更有利于我们这一边，并粉碎马尔西昂的假设，我们将接纳它们。那时就会很明显：他保留这些内容，显露出异端特有的盲目缺陷，其程度不亚于他删除前一类主题时的表现。那么，这就是我这篇小论文的趋向和形式；当然，受制于问题的两方面都已成为必要的任何条件。马尔西昂确立了一个立场：在提庇留时代，由一个先前未知的神为了救恩万民而启示的基督，与由造物主天主为了复兴犹太国而制定、且尚未来临的那一位是不同的存在。他在两者之间介入了一种巨大而绝对的差异——正如正义与善良之间的差异一样大；正如法律与福音之间的差异一样大；（简言之，）正如犹太教与基督教之间的差异一样大。由此也将产生我们的准则，由我们确定：对立之神的基督与造物主之间不应该有任何共同之处；而如果基督管理了他的安排、实现了他预言、促进了他的法律、实现了他的许诺、复兴了他的大能、重塑了他的决意、表达了他的属性、他的特性，就必须宣告基督属于造物主。我恳切请求读者始终记住这一法律和准则，然后让他开始探讨基督是属于马尔西昂的还是属于造物主的。\par

% structure: p0015 source=h2 target=subsection reason=relative-heading-rank; base=h2

\subsection{第七章：马尔西昂拒绝了\BibleBook{路加}{福音}的前面部分。因此本评论从考察\Place{葛法翁}会堂中邪魔的案例开始。邪魔所承认的那一位就是造物主的基督。}

\Person{提庇留}在位第十五年（这是\Person{马西昂}的主张），他\Quote{下到\Place{加里肋亚}的城\Place{葛法翁}}，当然是从造物主的天堂降下，因为他先前已从自己的天堂降下。那么，他经历了怎样的过程，才会被描述为先从自己的天堂降至造物主的天堂？因为，我为何要避免指摘那些不能满足普通叙述要求、却总是以谎言告终的陈述呢？当然，我们已经通过之前提出的问题一劳永逸地表达了指摘：即，当他穿越造物主的领域，并且确实与造物主为敌时，他是否可能被造物主接纳，并被造物主传送至同样属于造物主领土的地上？现在，假设他已降下，我还想知道他下降的其余过程。因为我们不可过于苛求地询问，他是否在任何地方被\Emph{看见}。进入视线是一种突然的、意外的瞥见，瞬间使目光停留在掠过眼前的物体上，但并不停留。但是，当降下发生时，它是显明的，并且进入眼睛的注意。而且，它涉及\Emph{事实}，因此迫使我们考察下降以何种状态、何种准备、何种暴力或温和程度、以及更进一步的，在白天或夜晚的何时发生；还有，谁看到了下降，谁报告了它，谁严肃地证实了这个事实——这个事实当然不容易被相信，即使经过郑重断言。简言之，这太糟糕了：\Person{罗慕路斯}尚有\Person{普罗库卢斯}作为他升天的见证者，而这位神的基督却找不到任何人宣告他从天降下；仿佛一人的升天和另一人的降下不是在同一谎言的梯子上完成的！然后，他与\Place{加里肋亚}有何关系，如果他并不属于那位预定该地区（为祂的基督）在开始传教时进入的造物主？正如\Person{依撒意亚}所说：\BibleQuote{先痛饮吧，并且迅速行动，哦，\Place{则步隆}区域和\Place{纳斐塔里}地，以及你们这些（居住）在海滨和\Place{约旦}的人，\Place{加里肋亚}的外邦人，你们坐在黑暗中的百姓，看啊，一道大光；你们这些居住在那地、坐在死荫中的人，光已升起。}然而，幸好\Person{马西昂}的神自称是外邦人的启蒙者，这样他就有更好的理由从天降下；只是，如果必须如此，他宁愿选择\Place{本都}作为降下之地，而不是\Place{加里肋亚}。但是，既然按照预言，启蒙的地点和工作都与基督相符，我们开始看出他就是预言的所指，这预言表明，在祂\Emph{传教之初}，祂来不是要废除法律和先知，而是要成全他们；\BibleRef{玛 5:17}因为\Person{马西昂}删除了这段经文，视为篡改。然而，他否认基督口中说了祂立即部分地付诸实践的话，这将是徒劳的。因为关于地点的预言祂立刻实现了。从天直接到会堂。正如谚语所说：\Quote{我们来做的事，立即去做。}\Person{马西昂}甚至必须从福音中删除\BibleQuote{我奉派遣，只为以色列家迷失的羊}\BibleRef{玛 15:24}和\BibleQuote{拿儿女的饼丢给小狗是不对的}\BibleRef{玛 15:26}——这样，基督才不至显得是以色列人。但事实将代替言语令我满足。撤回我基督的所有言语，祂的行为将说话。看，祂进入会堂；这当然是前往以色列家迷失的羊。看，祂首先将祂教导的\Quote{饼}提供给以色列人；这当然是因为他们是\Quote{儿女}，祂才显出优先。注意，祂尚未分给其他人；这当然是因为祂将这些人当做\Quote{狗}而忽略。因为，如果祂特别不属于造物主，祂还能更好地分给谁是造物主的陌生人呢？然而，祂怎能被允许进入会堂——一个如此突然出现、如此陌生的人；一个还没有人知道祂的支派、民族、家族，最后，连\Person{奥古斯都}人口普查——主诞生的最忠实见证，保存在\Place{罗马}档案中——的登记都不知道的人？他们当然会记得，如果他们不知道祂受过割礼，祂一定不能被允许进入他们最神圣的地方。即使祂有进入会堂的一般权利（像其他犹太人一样），但施教之职只允许给一个非常有名、经过考察和考验、并在其他地方经过适当经验证明后才被赋予特权的人。但是，\BibleQuote{他们都惊讶于祂的教导。}当然惊讶；因为（\Person{圣路加}说）祂的话具有权能\BibleRef{路 4:32}——不是因为祂教导反对法律和先知。毫无疑问，祂神圣的言论散发出能力和恩宠，建造而非拆毁法律和先知的实质。否则，惊讶会变成恐惧。他们不会钦佩，而是迅速而确实的厌恶，他们会给予一个毁灭法律和先知的人，并且自然随之提出敌对之神的人；因为他不可能贬低法律和先知，从而也贬低造物主，而不事先提出一个不同且敌对的神性学说。既然如此，圣经在这件事上没有其他陈述，只提到祂话语的单纯力量和权能产生惊讶，这更自然地表明祂的教导与造物主一致（通过不否认这一点），而非与造物主对立（通过不主张这一事实）。因此，祂要么被承认为属于造物主，祂是按照造物主教导的；要么将被判定为欺骗者，因为祂是按照祂来反对的那一位教导的。在同一段经文中，\Quote{一个不洁魔鬼的精神}喊道：\BibleQuote{我们与祢何干，祢耶稣？祢来毁灭我们吗？我知道祢是谁，祢是天主的圣者。}\BibleRef{路 4:33}我在此不提出问题：这个称谓是否适合一个不被称作基督的人，除非祂是由造物主派遣的？别处已经充分考虑了祂的称号。\par

我当下讨论的是，邪魔如何能知道祂被称为这一名号，因为从未有任何时候，由一位未知且迄今沉默的神祇发出过关于祂的预言，而这位神祇甚至不为祂自己的造物主所知，因此祂不可能被证明为\Quote{圣者}。那么，他能宣布一个相似的\Emph{事件}来表明那位敌对之神祇的\Quote{圣者}是什么\Emph{新}神祇呢？仅仅是他进了会堂，并且连言语上也没有做过任何反对造物主的事？因此，正如他绝不可能承认他所不认识的那位是耶稣和天主的圣者，他也确实承认了他所认识的那位（既是耶稣也是天主的圣者）。因为他记得先知曾预言天主的那位\Quote{圣者}，并且天主\Quote{耶稣}的名号在农的儿子身上。这些事他也从天使那里得知，按照我们的福音：\Quote{因此，那将要由你诞生的，将称为\Emph{圣者}，天主子；}并且\Quote{你要给他起名叫\Emph{耶稣}。}\BibleRef{玛 1:21}因此，他实际上（虽然只是一个邪魔）对主的救恩计划有一些概念，而非对任何奇异且迄今未被完全理解的神祇有概念。因为他还预先提出了这个问题：\Quote{我们与祢有什么相干？}——并非指一个陌生的耶稣，造物主的\Emph{邪魔}属于祂。他也没有说：\Quote{祢与我们有什么相干？}而是说：\Quote{我们与祢有什么相干？}仿佛在哀叹自己，并祈求免遭自己的灾祸；面对这灾祸，他接着说：\Quote{祢来毁灭我们吗？}因此，他完全在耶稣身上认出了那位有审判、复仇（可以说）严厉的天主的圣子，而非那位纯粹良善、不知毁灭也不知惩罚的天主的圣子！现在我们为何首先引用这段话？为了表明，邪魔既不承认耶稣是别的神祇，耶稣自己也没有肯定祂是别的神祇，而只是肯定祂是造物主的那一位。可是你说，耶稣斥责了他。确实如此，因为他是嫉妒的（邪灵），在他的承认中也只是无礼的，并因谄媚而邪恶——仿佛基督的最高荣耀在于来毁灭魔鬼，而不在于拯救人类；而祂的真正愿望是，祂的门徒不应因邪魔的屈服而夸耀，而应因救恩的美好而夸耀。否则祂为何斥责他？如果是因为他完全错了（在呼求中），那么祂就既不是耶稣，也不是天主的圣者；如果是因为他部分错了——因为他恰当地认为祂是耶稣和天主的圣者，但也认为祂属于造物主——那么祂斥责他就极不公平了，因为他思考了他知道应该思考的（关于祂的事），而没有设想他不知道应该设想的关于祂的事——即祂是另一位耶稣，另一位神祇的圣者。然而，如果斥责没有比我们所赋予的更合理的含义，那么就可以得出结论：邪魔没有犯错，也没有因说谎而被斥责；因为那就是耶稣自己，除了祂以外，邪魔不可能承认任何其他一位，而耶稣肯定了祂就是邪魔所承认的那一位，因为祂没有斥责他说谎。\par

% structure: p0018 source=h2 target=subsection reason=relative-heading-rank; base=h2

\subsection{第八章 来自同一章的其他证据：在纳匝肋讲道并被某些魔鬼承认为天主子基督的耶稣，是造物主的基督。适当时揭露马西昂的\OriginalTerm{幻影论的错误}{Docetic errors}。}

造物主的基督必须按预言称为\Emph{纳匝肋人}；为此，犹太人甚至因此称我们为\Emph{纳匝肋人}，追随祂。因为我们就是经上所记载的人：\Quote{她的纳齐尔人比雪更白}（\BibleRef{哀 4:7}），他们曾一度被罪的污点玷污，被无知的黑暗笼罩。但\Quote{纳匝肋人}这个称号对基督而言注定是合适的，因为祂的婴孩时期被隐藏，为此祂下去住在纳匝肋，以躲避黑落德之子阿尔赫劳。我提及此事，是因为马西昂的基督本应避免与造物主基督的\Emph{本地}地点有任何关联，当时犹太地有许多城镇并未被先知如此指定给造物主的基督。但基督将是先知们的基督，无论祂在何处被发现符合先知所言。然而，即便在纳匝肋，祂并未被记载宣讲任何新事物（\BibleRef{路 4:23}），而在另一节中，祂因一句简单的谚语而被拒绝（\BibleRef{路 4:29}；\BibleRef{路 4:24}）。在此，当我观察到他们向祂下手时，我不得不就祂的身体实质得出结论：祂不可能是幻影，因为祂能被触摸甚至被粗暴对待——当祂被抓住、带走并被带到悬崖边时。因为虽然祂从他们中间穿过去了，但祂已经经历了他们的粗暴对待，随后离开，无疑是因为人群（如常）让开了，甚至被冲开；而不是像不可触摸的伪装那样被躲避，若有这样的伪装，就不会遭受任何触摸了。\par

\Quote{\OriginalTerm{因为除非是物体，否则没有东西能触摸或被触摸}{Tangere enim et tangi, nisi corpus, nulla potest res},}\par

这句话甚至值得在世俗智慧中占有一席之地。总之，祂亲自触摸了别人，将自己可被感知的手按在他们身上，并赐予治愈的恩惠（\BibleRef{路 4:40}），这些恩惠与祂用来赐予它们的手一样真实，绝非虚幻。因此，祂正是依撒意亚所预言的基督，医治我们疾病的祂。经上说：\Quote{的确，祂承担了我们的疾苦，\Emph{背负}了我们的疼痛。}希腊人习惯于用\Quote{承担}一词，这个词也意指\Quote{除去}。一个普遍的应许对我来说足够暂时使用。无论耶稣施行了何种治愈，祂都属于我。然而，我们即将讨论治愈的种类。因此，将人从恶灵中解放出来，也是一种疾病的治愈。相应地，邪恶的灵（正如我们先前的例子）常常带着见证出去，喊叫说：\Quote{祢是天主子。}（\BibleRef{路 4:41}）——至于这是哪位天主之子，从事情本身来看已足够清楚。但祂被斥责，并被命令不要说话；这恰恰因为基督愿意自己被\Quote{人}而非不洁的灵宣告为天主子——甚至是唯独那一位基督，这事与祂相称，因为祂预先派遣了人，藉着他们祂可以为人所知，而且这些人确实是更值得的宣讲者。拒绝不洁之灵的宣告是祂的本性，因为祂的命令之下有充足的圣者。然而，那位从未被预言过的（如果祂确实希望被承认；因为如果祂不那么希望，祂的来临就是徒然的），就不会拒绝一个异类或任何物质的见证，因为祂自己并没有自己的物质，而是降生在异类物质中。而且，如今作为造物主的毁灭者，祂最希望不过被\Quote{祂的}灵所承认，并为了被畏惧而传扬；只是马尔强说他的神不被畏惧；坚持认为一个善的存在不是恐惧的对象，而只是一个审判的存在，其中存有恐惧的缘由——愤怒、严厉、审判、报复、定罪。但毫无疑问，邪恶的灵是因恐惧而被吓倒。因此，他们承认祂是一位可畏惧的天主之子，因为如果没有任何畏惧的理由，他们就会有不驯服的借口。此外，祂表明自己应当被畏惧，因为祂驱逐他们，不是像善的存在那样通过劝说，而是通过命令和斥责。否则，祂难道是因为他们使祂成为恐惧的对象而斥责他们，而祂一直都不愿被畏惧？祂希望他们以何种方式出去？他们除了带着恐惧，不能以别的方式出去。因此，祂陷入了与自己的本性相悖的困境，而祂本可以以单纯的良善立刻仁慈地对待他们。祂也陷入了另一个虚伪的境地——口是心非，当祂允许自己被魔鬼们畏惧为造物主之子时，为了驱逐他们，不是凭祂自己的权能，而是凭造物主的权威。\newline{}\Quote{祂离开，往荒野地方去了。}（\BibleRef{路 4:42}）这确实是造物主惯常的区域。圣言在那里以身体显现是合宜的，因为祂昔日曾在云中工作。那曾为法律所指定之地的状况，也适合于福音。因此，让\Quote{让旷野和干旱之地欢喜快乐}；正如依撒意亚所应许的。（\BibleRef{依 35:1}）当被群众留住时，祂说：\Quote{我也必须向别的城传天主国的福音。}（\BibleRef{路 4:42}）祂曾在任何地方显明过祂的天主吗？我想还没有。但祂是在对那些知道另一位神的人说话吗？我不这么认为。因此，如果祂未曾宣讲过，他们也未曾知道过任何其他天主，唯有造物主，那么祂是在宣告那位天主的天国，祂知道这位天主是听众所知道的唯一天主。\par

% structure: p0022 source=h2 target=subsection reason=relative-heading-rank; base=h2

\subsection{第九章 从圣路加福音第五章中找到基督属于造物主的证据，例如：召叫渔夫担任宗徒职务，以及洁净癞病人。基督与先知厄里叟相比。}

在如此众多的职业中，祂为何特别看重渔夫这一行，以致从中拣选西满和载伯德的儿子们作宗徒（因为那叙述之所以被详细铺陈，似乎不只是为了这件事本身），对那因着极多渔获而战栗的伯多禄说：\BibleQuote{不要害怕；从今以后，你要得人了。}祂这样说，便向他们启示了预言的应验，即正是祂曾通过耶肋米亚预言的那一位：\BibleQuote{看，我要派遣许多渔夫，他们要捕捉他们，}\BibleRef{耶 16:16}即捕捉人。然后他们才舍弃了船，跟随了祂，明白祂就是那开始实现祂所宣告的事的那一位。这与祂故意从船长公会中拣选，有朝一日立船长马西昂为祂的宗徒，是完全不同的情形。我们早已针对祂的\WorkTitle{对立论}指出，马西昂的立场并不能从他所假设的法律与福音之间的差异中获益，因为就连这一点也是造物主所制定的，并且在新法律、新圣言和新盟约的应许中早就预言了。然而，既然他特别仔细地引用了一位与他同样受苦的同伴（\OriginalTerm{患难中的同伴}{\GreekText{συνταλαίπωρον}}）和仇恨中的伙伴（\OriginalTerm{仇恨中的伙伴}{\GreekText{συμμισούμενον}}）作为他领域的证据，用来医治癞病，\BibleRef{路 5:12}，我便乐于与他交锋，并在一切事之前，先向他指出法律以喻意解释的力量，这法律在此癞病人的例子中（不可触摸他，反而应当将他与众人隔离）禁止与染有罪污的人有任何交往；宗徒也禁止我们甚至与这样的人一同进食，\BibleRef{格前 5:11}，因为罪的污秽会如同传染病一般，每当人与罪人混杂时便会传染。因此，主愿意更深刻地理解法律，即它通过肉体的事实来象征属灵真理——从而并非废除，反而是在建立那祂希望更认真地被承认的\Quote{法律}——祂触摸了那癞病人；尽管祂作为人可能被玷污，但祂作为神不可能被玷污，因为祂当然是不朽的。所以，这条规定不能适用于祂，即祂必须遵守法律而不触摸不洁之人，因为接触不洁之物并不会使祂受玷污。我以此教导说，这种豁免与我的基督是一致的，而当我证明它在你的基督那里是不一致的时，更是如此。现在，如果是作为法律的仇敌，祂触摸了癞病人——因藐视玷污而藐视法律的诫命——那么祂如何能被玷污呢？祂并没有一个能被玷污的身体。因为幻象是不能被玷污的。因此，那位作为幻象不能被玷污的，其免于玷污并非出于任何神圣能力，而是由于祂虚幻的空无；也不能认为祂藐视了玷污，因为祂实际上根本没有物质上的能力去接受玷污；同样，也不能认为祂废除了法律，因为祂凭藉其幻象本性而摆脱了玷污，并非出于任何德性的彰显。然而，如果造物主的先知厄里叟只洗净了叙利亚人纳阿曼，而排除了以色列那么多的癞病人，这一事实无助于区分基督，好像祂因此就更优越，因为祂洗净了这个与祂无关的以色列癞病人，而祂自己的主却未能洗净他。那叙利亚人的洗净，反而预表了全世界万民将在他们的光基督内得到洗净，因为他们曾浸染在七种大罪的污秽中：崇拜偶像、亵渎、凶杀、奸淫、私通、作假见证和欺诈。因此，他在约但河洗了七次，好像为每一种罪一次，为的是庆祝一个完美的七数所完成的赎罪；并且因为这唯一圣洗的德能和满全，就这样庄重地归给唯独基督，祂终有一天要在世上不仅建立启示，而且建立圣洗，赋予它简洁的效能。连马西昂也在这里找到一个对立：厄里叟采用了物质资源，运用水，而且七次；而基督只凭一句话，且仅一次，便立即治愈了。那么，我或许敢断言，那圣言本身也属于造物主的实体。那位原始的作者，没有什么是祂不能更有力的。的确，令人难以置信的是，造物主曾凭一句话创造了如此浩瀚的宇宙，如今却凭一句话只产生一种疾病的疗法！从什么方面可以更好地辨别造物主的基督，而不是从祂圣言的能力呢？但是基督之所以是另一位基督，是因为祂行事与厄里叟不同——事实上，主人比仆人更有能力！马西昂，你为什么制定规则，说仆人做事正像他们的主人一样？你不怕这会反而不利于你吗？如果你否认基督属于造物主，理由是祂一次比造物主的仆人更有能力——既然与厄里叟的软弱相比，祂被承认为更大者，如果真是更大的话！因为治愈是相同的，尽管运作方式有所不同。你的基督比我的厄里叟做了更多的事吗？不，你的基督的圣言究竟行了什么大事，它仅仅做了造物主的一条河流所做的事？同样原则适用于其他一切。就放弃一切人的荣耀而言，祂禁止那人宣扬治愈的事；但就法律的尊荣而言，祂要求遵循常例：\BibleQuote{你去，叫司铎检验你，并献上梅瑟所规定的献仪。}\BibleRef{路 5:14}因为祂仍要遵守法律在其预像中的象征性记号，因为它们具有预言的意义。这些预像象征：一个人一度是罪人，后来因天主的圣言洗净了罪污，就应在圣殿中向天主献上礼物，即在教会内通过耶稣基督——祂是圣父的公教司铎——献上祈祷和感谢。因此祂又加上：\BibleQuote{为给你作证据。}——无疑，这证据是要证明祂不是在废除法律，而是在成全法律；也是为了证明祂就是那预言要承担他们的疾病和软弱的那一位。马西昂，这位他自己基督的谄媚者，试图用慈悲和温柔为借口，将\Quote{证据}这个非常一致且恰当的解释排除在外。因为，祂既是良善的（这是他的话），又知道每一个从癞病中被释放的人必然会履行法律的礼仪，因此祂下了这条命令。那么，祂有坚持祂的良善（即对法律的容许）吗？还是说没有？如果祂坚持了良善，祂就永远不会成为法律的废除者；也永远不会被认为是属于另一位神的，因为不存在那种会使祂属于另一位神的对法律的废除。然而，如果祂没有坚持良善，而是后来废除了法律，那么祂在治愈癞病人时所加给他们的证据便是虚假的；因为祂在废除法律时，遗弃了祂的良善。因此，如果祂在支持法律时是良善的，那么现在作为法律的废除者，祂就成了邪恶的。但是，祂对法律的支持肯定了法律是良善的。因为没有人在支持邪恶之事时容许自己这样做。所以，如果祂容许人服从一条恶法，祂不仅邪恶，而且更坏，如果祂表现为一条善法的废除者的话。因此，如果祂命令献上礼物是因为祂知道每个治愈的癞病人都一定会献上，那么祂可能就避免命令那祂知道会自发去做的事。因此，祂降临，仿佛要废除法律，却对遵守法律的人让步，乃是徒劳的。然而，因为祂知道他们的倾向，就更应当力劝他们不要疏忽法律，既然祂为了这个目的而来。那么，祂为什么没有保持沉默，让人凭自己的简单意愿去服从法律呢？因为那样祂似乎才在一定程度上坚持了祂的忍耐。但是祂又加上自己的权威，并用这个\Quote{证据}的分量来加强。我问，这是什么证据，如果不是肯定法律的证据？当然，问题不在于祂如何肯定了法律——是作为良善，还是作为多余，还是作为忍耐，还是作为反复无常——只要马西昂，我将你从你的立场上赶走。看，祂命令法律应被成全。无论祂以何种方式命令，祂也就可以先说出那句话：\BibleQuote{我来不是要废除法律，而是要成全。}\BibleRef{玛 5:17}那么，你有什么理由从福音中删除这一完全一致的话呢？因为你已经承认，祂在行为上凭着祂的良善，做了你在言语上否认祂做过的事。我们因此有确凿的证据，证明祂说了那句话，因为祂做了那件事；并且证明你删除了主的话，而不是我们的门徒插入了它。\par

% structure: p0024 source=h2 target=subsection reason=relative-heading-rank; base=h2

\subsection{第10章　在同一章中，从治好瘫子以及耶稣自称人子进一步证明同一真理。\Person{戴尔都良}引用先知多处经文支持其论证。}

瘫子被治愈了，\newline{}\BibleRef{路 5:16}，而且是公然在众人眼前。因为依撒意亚说：\Quote{他们要看见上主的荣耀，和我们天主的威严。}\newline{}\BibleRef{依 35:2} 什么荣耀，什么威严？\Quote{你们要坚强，软弱的手，无力的膝：}这是指瘫痪。\Quote{要坚强，不要害怕。}\newline{}\BibleRef{依 35:4} \Emph{要坚强}不是徒然重复，\Emph{不要害怕}也不是多余添加；因为随着肢体的更新，按着应许，身体的力量也要恢复：\Quote{起来，拿起你的床榻；}同样地，也要有道德勇气，不怕那些说\Quote{除了天主以外，谁能赦罪呢？}的人。所以你在这里不仅看到了预言应验，应许了某种特殊的医治，也看到了治愈后伴随的症状。同样地，你也应当在同一位先知身上认出基督是赦罪者。因为他说：\Quote{他要赦免许多人的罪，并亲自除去我们的罪。}因为在较早的章节中，他以主自己的口吻说：\Quote{即使你们的罪像朱红，我要使它们白如雪；即使它们像丹红，我要使它们白如羊毛。}\newline{}\BibleRef{依 1:18} 在朱红色中他指出先知的血；在丹红色中，他指出主的血，更为明亮。关于赦罪，米该亚也说：\Quote{有哪一位天主像你？赦免罪过，宽恕你产业遗民的过犯。他不会永远怀怒，因为他喜爱仁慈。他必再转来怜悯我们，将我们的罪孽踏在脚下，将我们的一切罪恶投入深海。}\newline{}\BibleRef{米 7:18} 如果对基督没有这样的预言，我便会在创造主中找到如此仁慈的榜样，使我相信作为祂之父的那位子在同样情感上的应许。我看见尼尼微人从创造主那里获得了罪的赦免\newline{}\BibleRef{纳 3:10}——更不用说从基督那里了，因为从起初祂就以父的名行事。我也读到，当达味承认他对乌黎雅的罪时，纳堂先知对他说：\Quote{上主已赦免你的罪，你不会死；}\newline{}\BibleRef{撒下 12:13} 同样地，阿哈布王，依则贝耳的丈夫，犯有偶像崇拜和纳博特的血的罪，因他的悔改而获得了赦免；\newline{}\BibleRef{列上 21:29} 还有撒乌耳的儿子约纳堂，因他的恳求而抹去了违犯斋戒的罪。我何必一再叙述这民族本身在罪赦之后的经常复兴呢？——这是由那位喜爱仁爱胜过祭献、喜爱罪人悔改胜过死亡的天主所做的。\newline{}\BibleRef{则 33:11} 你首先必须否认创造主曾赦免过罪；然后你必须理性地表明，祂从未为祂的基督指定过这样的特权；这样你就能证明，那所谓的全新仁慈是多么新奇——当然，来自全新的基督——当你证明了它与创造主既不相符，也非创造主所预言。但是，赦罪是否属于一位据说不能保留罪的主，解罪是否属于一位甚至无权定罪的主，宽恕是否适合一位无人能向其犯罪的主——这些问题我们在别处已经遇到过，当时我们倾向于提出一些提示，而不是重新讨论。关于人子，我们的规则是双重的：基督不能说谎，以致自称是人子，如果祂实际上不是；祂也不能成为人子，除非祂由人父或人母所生。那么争论的焦点就是：祂应当被视为哪一位人父或人母的儿子？既然祂是由天主圣父所生，祂当然不是人父的儿子。如果祂不是人父的儿子，那么祂必定是人母的儿子。如果祂是人母的儿子，那么显然她必须是童贞女。因为对于没有人类父亲的人，祂的母亲就不会有丈夫；而对于没有丈夫的母亲，童贞的状态就属于她。但如果祂的母亲不是童贞女，那么祂就必须有两个父亲——一个神性的，一个人类的。因为她必须有一个丈夫，这样才不是童贞女；而有了丈夫，她就会导致祂有两个父亲——一个神性的，一个人类的——这样祂就既是天主子，也是人的儿子。这种诞生（如果可以这样称呼的话）神话故事中归于卡斯托耳或赫拉克勒斯。如果这个区别被遵守，也就是说，如果祂是人子，因为由祂母亲所生，而不是由父亲所生，并且祂的母亲是童贞女，因为祂的父亲不是人——那么祂就是依撒意亚所预言的那位由童贞女所怀的基督，\newline{}\BibleRef{依 7:14} 你马尔西昂，怎么能承认祂是人子，我实在无法看出。如果是通过人类父亲，那么你就否认祂是天主子；如果是通过一个神性父亲\Emph{也}，那么你就使基督成了神话中的赫拉克勒斯；如果只通过人类母亲，那么你就同意了我的观点；如果不是通过人类父亲\Emph{也}，那么祂就不是任何人的儿子，并且祂一定说了谎，声称自己是祂实际上不是的。只有一件事能帮助你摆脱困境：你要么大胆地将你的神实际称为基督的人类父亲，就像瓦伦提努对他的永世所做的那样；要么否认那童贞女是人，这是连瓦伦提努也没有做的。现在，如果基督在达尼尔书中正好以\Quote{人子}的称号被描述，这难道还不足以证明祂是预言的基督吗？因为如果祂自己使用了那在预言中为创造主的基督所预备的称号，那么祂无疑表明自己就是先知所指定给那称号的那一位。但也许这可以被视为名字的简单相同；然而我们已经坚持认为，如果它们具备任何差异的条件，就不应该被称为基督或耶稣。但关于\Quote{人子}这个称号，\Emph{就}它偶然出现而言，\Emph{那么}它与名字的偶然相同同时出现就有困难。因为这是纯粹的偶然，尤其是当没有出现同样的原因导致相同的时候。因此，如果马尔西昂的基督也被说成是由人所生，那么他也会得到一个相同的称号，于是就会有两个子，同样也有两个基督和两个耶稣。因此，既然这个称号是唯一属于那位有恰当理由的人，如果它被另一个人声称拥有，而那人虽然有名字的相同，却没有称号的相同，那么名字的相同甚至对那人来说也是可疑的，因为他毫无理由地声称拥有称号的相同。因此，必须相信只有一位是同一的，那位更适合接受名字和称号的人；而另一位被排除，因为他没有权利得到称号，因为他没有理由证明它。也没有任何人比那位更早、名字和称号都被分配给他的人更适合两者，即创造主的耶稣。祂曾与祂的殉道者在火窑中被巴比伦王看见：\Quote{第四位，好像人子。}\newline{}\BibleRef{达 3:25} 祂也明确地向达尼尔本人启示为\Quote{人子，乘着天上的云彩而来}作为审判者，正如圣经所示。\newline{}\BibleRef{达 7:13} 我所提出的关于人子在预言中的称号已经足够了。但圣经还给了我进一步的信息，甚至是在主自己的解释中。因为当犹太人仅仅把祂看作人，还不确定祂也是神——因为是同样天主子——他们相当正确地说，一个人不能赦罪，只有天主才能，为什么祂不接着他们关于\Emph{人}的论点回答他们说，祂有权赦罪？因为当祂提到人子时，祂也提到了一个人？除非是因为祂想借助达尼尔书中\Quote{人子}这个称号，引导他们反思，向他们表明那赦罪的是神和人——正是达尼尔预言中那唯一的人子，祂获得了审判的权能，因此当然也获得了赦罪的权能（因为审判者也会释放）；这样，当他们因记起圣经而粉碎了他们的异议时，他们就能更容易地借着祂自己的赦罪行为而承认祂就是人子本身。我还观察到一点：祂在此处之前从没有自称是天主子——而是在这段经文里第一次，祂第一次赦罪；也就是说，祂第一次使用了\Emph{审判}的职能，通过释罪。所有对立面能用来反驳这些事的论点，（我请你）仔细权衡其分量。因为它必须将自己逼到如此愚蠢的地步：一方面，坚持（他们的基督）也是人子，以免祂被指控说谎；另一方面，否认祂由女人所生，免得他们承认祂是童贞女的儿子。然而，既然神圣的权威、事物的本质和常识都不接受异端者这种疯狂的主张，我们在这里有机会以最简短的方式否决\Emph{基督身体的实体}，反对马尔西昂的幻影。既然祂由人所生，是人子，祂是来自身体的身体。我向你保证，你更容易找到一个没有心或没有脑子的人——像马尔西昂自己一样——而不是一个没有身体的基督，像马尔西昂的基督。让这成为你考察那位本都异端者的心、或至少是他的脑子的界限。\par

% structure: p0026 source=h2 target=subsection reason=relative-heading-rank; base=h2

\subsection{\Person{肋未}的蒙召。\Person{基督}与\Person{洗者若翰}的关系。\Person{基督}作新郎。新酒与旧酒的比喻。将\Person{基督}与\OriginalTerm{造物主}{Creator}联系起来之论证。}

主所拣选的那位税吏，他引以为证，表明他被一位与法律为敌者所拣选，而此人对法律是生疏的，且未受犹太教的熏陶。伯多禄的例子却被他遗忘了：虽然伯多禄是一个熟悉法律的人，却不仅被主拣选，而且得到了拥有知识的见证，这知识是由父赐给他的。\BibleRef{玛 16:17} 他（马西昂）从未在何处读到，基督被预言为外邦人的光明、希望和期望！然而，\Emph{他}（马西昂）却以赞赏的语气论及犹太人，当他说：\BibleQuote{健康的人不需要医生，而有病的人才需要。} \BibleRef{路 5:31} 因为藉着\Quote{有病的人}，他意指应理解为外邦人和税吏，就是他所拣选的；他断言犹太人是\Quote{健康的}，对他们来说医生不是必需的。既然如此，他（马西昂）在下来毁灭法律上就犯了错误，好像是为了医治有病的状态，因为生活在法律之下的人本是\Quote{健康的}，且\Quote{不需要医生}。此外，他提出\Emph{医生}的比喻，若他不去证实它，又怎么可能呢？因为，正如没有人对健康的人使用医生，同样也没有人对陌生人使用医生，如果他是马西昂的神所造的人，自有他的创造者和保护者，以及在他的基督内的一位特别好的医生。这个比较预先决定了这一点：医生通常是由病人所属的那一位提供的。再者，若翰从哪里登场？基督突然出现；同样突然地，若翰！在马西昂的体系中，万事皆如此发生。然在造物主的安排中，它们有其独特而圆满的进程。然而，关于若翰，我还要说的其他内容将在另一段中找到。对于目前摆在我们面前的几点，必须予以答复。因此，我将留心这样做——证明若翰与基督彼此相合，若翰作为造物主的先知，正如基督是造物主的基督；如此异端者将因挫败\Person{若翰}\Emph{洗者}的使命而感到羞耻，而他自己反遭挫败。因为如果若翰根本没有职务——即依撒意亚所称的\Quote{呼声}，\Quote{在旷野中呼喊的}，以及藉宣告和劝勉悔改来预备主的道路的人；如果他也没有亲自给基督施洗，连同给别人施洗，那么就没有人能够在基督的门徒又吃又喝的时候，拿他们与时常禁食祈祷的若翰门徒相比较；因为如果基督与若翰之间，以及他们的追随者之间存在任何差异，那么精确的比较就不可能，也没有一个可比点。因为没有人会感到惊奇，也没有人会困惑，尽管会出现关于不同神明的对立预言，这些预言也在行为方式上相互有别，而在此之前它们基于的权威便已不同。因此，基督属于若翰，若翰也属于基督；而两者都属于造物主，两者都来自法律和先知，是宣讲者和导师。否则，基督就会拒绝若翰的纪律，认为它属于敌对的神，并且也会为门徒辩护，认为他们走不同的道路是恰当的，因为他们献身于另一个相反的神祇的侍奉。但事实上，当他谦逊地给出理由，为什么\Quote{新郎的同伴在新郎与他们同在的时候不能禁食}，却应许\Quote{但当新郎从他们中被带走时，他们就要禁食}时，\BibleRef{路 5:34} 他既没有为门徒辩护（反而为他们开脱，好像他们并非毫无理由地受到责备），也没有拒绝若翰的纪律，反而允许它，将它归给若翰的时代，尽管为他自己时代保留。否则，他的目的就是要拒绝它，并为其反对者辩护，如果他自己当时还不属于它的话（即当时的纪律）。我也认为，圣咏中所说的新郎指的是我的基督：\BibleQuote{他如新郎走出洞房；他的出发从天的尽头，他的返回又到天的尽头。} 藉依撒意亚的口，他也欣喜地论及父说：\BibleQuote{我的灵魂要因主而喜乐，因为他给我穿上救恩的衣裳，以喜乐的外袍给我披上，如同新郎戴上冠冕，又如新娘佩戴首饰。} \BibleRef{依 61:10} 同样，他也将教会归于自己，同一圣神论及教会向他说：\BibleQuote{你要以他们全部作为饰物穿戴自己，如同新娘佩带妆饰。} \BibleRef{依 49:18} 这个配偶，基督也藉撒罗满从外邦人的蒙召中邀请她归向自己，因为经上记载：\BibleQuote{我的新娘，请与我同从黎巴嫩来。} \BibleRef{歌 4:8} 他优雅地提到黎巴嫩（当然是山），因为它在希腊人中代表乳香之名；因为他正是从偶像崇拜中聘娶了教会。现在，马西昂，否认你的极度狂妄吧（如果你能）！看哪，你甚至攻击你的神的律法。他不结合于婚姻纽带，即使缔结了，他也不允许；他只给\Emph{独身者}或阉人施洗；他将洗礼保留到死亡或离婚之后。那么，你为什么将他的基督称为新郎？这是结合男女者的称号，而非分离者的称号。你在基督的那一陈述中也犯了错误，其中他似乎在旧与新之间作出分别。你因旧皮囊而自大，因新酒而头脑昏乱；因此你把你新奇的异端补丁缝在旧的（即先前的）福音上。我想知道造物主在哪些方面与自己矛盾。当他藉耶肋米亚颁布这条诫命：\BibleQuote{你们要为自己开垦新荒地}，\BibleRef{耶 4:3} 他岂非背离旧的事物？当他藉依撒意亚宣告：\BibleQuote{旧事已过，看哪，我所造的一切都是新的}，他岂非指向新的事物？我们通常认为，旧事物的终结更由造物主应许，并由基督在现实中展示，这都是在同一且唯一的天主的权下，新旧事物都属于他。因为新酒不能放入旧皮囊，除非那人有旧皮囊；也没有人把新布补在旧衣服上，除非他手里有旧衣服。只有那些拥有材料、若要做就能做的人，才可以不做那些不当做的事。因此，既然他作此比喻的目的是要表明他正将福音的新状况与法律的旧状况分开，他便证明，他所分开的与其自身的东西不应被指责为彼此相异事物的分离；因为没有人会把自己的东西与异己之物结合，以便以后能将它们从异己之物中分离出来。分离是借助造成分离的结合而可能的。因此，他所分开的事物，他也证明它们曾是一体；若不是他分开，它们本会保持一体。但我们仍承认有一种分离，藉着革新、扩展、进步；正如果实从种子分离，尽管果实来自种子。同样，福音也从法律分离，当它从法律前进时——与法律不同，却并非异己；相异，却不相反。在基督身上，我们甚至找不到任何新颖的言论形式。无论他提出比喻还是反驳问题，都来自第七十七篇圣咏。他说：\Quote{我要开口讲比喻}（即比喻）；\Quote{我要说出\Emph{隐晦}的谜语}（即我要提出问题）。如果你想证明一个人属于另一个种族，毫无疑问你会从他的语言惯用语中取证。\par

% structure: p0028 source=h2 target=subsection reason=relative-heading-rank; base=h2

\subsection{第十二章。基督对安息日的权柄。作为安息日的主，祂使安息日回归其原本目的，即由创造主所设立，以纠正法利塞人的忽视。论及门徒在安息日掐麦穗的事例。在安息日治好枯手的人。}

论及安息日，我也有此前提：若基督没有公开宣告安息日的主，这问题就不会产生。祂废止安息日的可能性也不会引起讨论，除非祂有权废止。此外，若祂属于敌对之神，祂便有权如此行；祂做祂有权做的事，不会令任何人惊讶。因此，人们的惊讶源于他们认为基督宣告造物主为天主，却又攻击祂的安息日，这是不合适的。\newline{}现在，为能首先解决这几个问题，免得我们每遇到对手基于基督某一新奇制度的论据就不断重复它们，让我们先确立这一点：对每一项制度之新奇性的讨论之所以发生，是因为基督尚未提出任何关于新神祇的教导，因此讨论本身是不被允许的；也不能反过来论证，说从每一项制度的新奇性本身，基督就足够清楚地证明了另一位神祇，因为显然新奇性在基督身上并非值得惊奇的特征——它已被造物主预言了。当然，一位新神祇应当先被解释，然后祂的规诫才被引进；因为神祇会赋予规诫权威，而非规诫赋予神祇权威；除非（当然）事实上马西昂不是从一位导师那里获得他极其乖谬的意见，而是他的导师从他的意见中获得！\newline{}关于安息日的所有其他各点，我如此裁定：若基督干预了安息日，祂只是效法了造物主的榜样；因为在耶里哥城围攻中，约柜绕城八天，包括安息日，实际上是在造物主的命令下废止了安息日——依照那些在此《路加福音》段落中认为基督如此行的人的观点，他们不知道基督和造物主都没有违反安息日，我们稍后将证明。然而，安息日当时确实被若苏厄打破了，因此同样的指控也可能针对基督。但即使祂作为犹太人的基督，对犹太人最庄严的日子显出了仇恨，祂也只是作为祂的基督，在同样仇恨安息日的事上公开效法造物主；因为祂以依撒意亚的口呼喊：\Quote{你们的月朔和你们的安息日，我的灵魂憎恶。}\BibleRef{依 1:14}无论这些话是在何种意义上说的，我们知道，在此类问题中，必须用突然的辩护回应突然的挑战。\newline{}现在我将把讨论转移到基督的教导似乎废止安息日的那件事本身。门徒饿了；在那安息日，他们掐了些麦穗，用手搓着；如此准备食物，他们亵渎了圣日。基督原谅他们，并成了他们破坏安息日的同谋。法利塞人指控祂。马西昂诡辩地解释争论的阶段（若我可借助我主的真理来嘲弄他的技巧），无论是在圣经记载中还是基督的意图中。因为从造物主的圣经和基督的意图中，可以得出一个貌似合理的先例——如同达味的例子，他在安息日进入圣殿，勇敢地掰开供饼来提供食物。他甚至记得，这特权（我是指免于禁食的特别许可）从起初安息日本身被设立时就已赐给安息日。因为虽然造物主禁止两天收集玛纳，但祂仅在安息日前一天允许一次，以便昨天的食物储备使随后的安息日节庆免于禁食。因此，主有充分理由在废止安息日时遵循同一原则（因为这是人们会用的词）；也有充分理由在赐予安息日不禁食的特权时表达造物主的旨意。简言之，祂会当场终止安息日，甚至终止造物主本身，若祂命令门徒在安息日禁食，违背圣经和造物主旨意的意图。但祂没有直接为门徒辩护，而是原谅他们；祂提出人的需要，如同在请求宽恕；祂维护安息日的尊荣，认为这是一个免于忧愁而非免于工作的一天；祂将达味及其同伴与祂的门徒置于同一过失和减轻情节的地位；祂乐意认可造物主的宽容；祂自己按照祂的榜样是良善的——难道因此祂就与造物主疏远了吗？\newline{}于是法利塞人监视祂，看祂是否在安息日治病，\BibleRef{路 6:7}好控告祂——当然是作为安息日的破坏者，而非新神祇的提出者；因为我可能满足于坚持这一点，即别处从未宣告另一位基督。然而，法利塞人对安息日的律法完全错误，没有注意到其中条件是有限的，当它规定从劳动中安息时，对劳动做了某些区分。因为当它论及安息日说：\Quote{在这一日，你不可做你的任何工作，}\BibleRef{出 20:16}用\Quote{你的}这个词，将禁令限制在人的工作上——每人在自己的职业或事务中做的工作——而非神圣的工作。现在，医治或保存的工作不属于人，而属于天主。同样，律法中说：\Quote{你不可做任何工作，}\BibleRef{出 12:16}除了为任何灵魂可做之事，即为灵魂的救恩的事；因为天主的工作可通过人的代理为灵魂的救恩而做。但由天主来做，就是人而基督要做的事，因为祂也是天主。因此，祂愿意通过枯手复原来引导他们进入律法的这层含义，就问：\Quote{安息日行善或作恶，救生或害命，哪样是合法的？}\BibleRef{路 6:9}为了使祂在允许那为灵魂而即将做的工作的同时，提醒他们安息日的律法禁止什么——即人的工作；并吩咐什么——即神圣的工作，可为任何灵魂的益处而做，祂被称为\Quote{安息日的主}，\BibleRef{路 6:5}因为祂维护安息日作为祂自己的制度。现在，即使祂废止了安息日，祂也有权这样做，作为它的主，更是作为设立它的那位。但祂没有完全毁灭它，尽管是它的主，为的是从今以后显明，甚至在约柜绕耶里哥时，安息日也未被造物主破坏。因为那确实是天主的工作，祂亲自命令，并为了祂的仆人在战争危险中的性命而安排的。虽然祂在某处表达了憎恶安息日，称它们为\Quote{你们的安息日}，\BibleRef{依 1:13}把它们算作人的安息日，非祂自己的，因为它们被充满不义的民庆祝，没有敬畏天主，\Quote{心口不一地爱天主}，\BibleRef{依 29:13}但祂把祂自己的安息日（即那些按照祂的吩咐遵守的）放在不同的位置；因为通过同一位先知，在后来的段落中，祂宣告它们为\Quote{真实、可喜乐、不可侵犯的}。这样，基督根本没有废除安息日：祂遵守了祂的律法，在先前的事上做了有益于门徒生命的工作，因为当他们饥饿时，祂以食物缓解他们，在眼前的事上则治愈了枯手；每件事都通过事实暗示：\Quote{我来不是要废除法，而是要成全法，}\BibleRef{玛 5:17}虽然马西昂用这话封住了祂的口。因为即便在我们面前的事上，祂成全了律法，同时解释了它的条件；此外，祂清楚地展示了不同种类的工作，同时做律法从安息日的神圣性中豁免的工作，并且通过祂自己的有益行动，赋予安息日本身——从起初就由圣父的祝福祝圣——额外的圣洁。因为祂为这日提供了神圣的保护——祂的对手本会为其他一些日子追求这些保护，以逃避尊崇造物主的安息日，并恢复安息日适合的工作。同样，先知厄里叟在这日使叔能妇人的死子复活，你看，法利塞人，你也是，马西昂，那是适合的工作，从古时就在造物主的安息日上行善、救生而非害命；基督带来了什么新的东西？没有不是按照造物主的榜样、温柔、怜悯和预言的。因为在这同一事例中，祂应验了一个特定医治的预言宣告：\Quote{弱手得加强}，正如瘫痪病患中\Quote{软弱的膝}。\BibleRef{依 35:3}\par

% structure: p0030 source=h2 target=subsection reason=relative-heading-rank; base=h2

\subsection{第十三章 基督与创造主的联系得以显明。旧约中许多引文预言性地涉及耶稣生平中的某些事件——例如祂上山祈祷；祂拣选十二宗徒；祂将西满的名字改为伯多禄，以及提洛和漆冬的外邦人投奔祂。}

诚然，祂向熙雍传报喜讯，向耶路撒冷传报和平和一切祝福；祂上山去，在那里彻夜祈祷，\BibleRef{路 6:12} 并且确实蒙圣父垂听。因此，翻阅先知书，从中认识祂的全部行径。\BibleQuote{你登上高山，给熙雍传喜讯的啊，请大声呼喊；给耶路撒冷传喜讯的啊，请提高你的声音。} \BibleRef{依 40:9} \BibleQuote{他们对祂的教训十分惊讶，因为祂说话具有权柄。} \BibleRef{路 4:32} 又说：\BibleQuote{因此，我的百姓要在那一日认识我的名。} 先知\Emph{所指}的名，不就是基督吗？ \BibleQuote{是我在说话——就是我。} \BibleRef{依 52:6} 因为那曾在先知们内说话的，就是圣言，圣子。\Quote{我在此，时辰已到，在山上，如同传报和平喜讯者，如同传报美善佳音者。} 于是，十二小先知之一的\Person{鸿}说：\BibleQuote{看，那传报和平喜讯者的脚踪在山上何等快捷！} \BibleRef{鸿 1:15} 此外，关于祂夜间向圣父祈祷的声音，圣咏明确说：\BibleQuote{我的天主，我白天呼求，祢必应允；夜间呼求，也不会落空。} 在另一处论及同一声音和地点时，圣咏说：\BibleQuote{我扬声呼求上主，祂从祂的圣山上应允了我。} 你有了名的预像；你有了传福音者的行动；你有了山作为地点；夜晚作为时间；声音的回响；圣父的垂听：总而言之，你有了先知所预言的基督。但为何祂拣选了\Emph{十二}位宗徒，\BibleRef{路 6:13} 而不是别的数目？实在，我单从这一点就能断定我的基督，祂不仅由先知的话语所预言，也由事实的预兆所预示。因为我在创造主的安排中处处发现这个数目的预表：厄林的十二股水泉，\BibleRef{户 33:9} 亚郎大司祭祭衣上的十二块宝石，\BibleRef{出 28:13} 以及若苏厄命人从约旦河中取出、为约柜竖立的十二块石头。如今，同样数目的宗徒就如此被预兆：好像他们将成为泉源和河流，浇灌那先前干涸、缺乏知识的外邦世界（正如祂藉依撒意亚说：\BibleQuote{我要在无水之地放置溪流} \BibleRef{依 43:20}）；好像他们将成为宝石，照耀教会的神圣祭衣，就是基督——圣父的大司祭——所穿戴的；好像他们也将是信德坚固的石头，真正的若苏厄从约旦河的水中取出，安放在祂盟约的圣所里。马西昂的基督能为此数目拿出什么同样有力的辩护？任何事若不能显明与我的基督有关联（与先前的预表相连），就不可能显明由他无关联地做成。在谁身上发现了同一事件的预备，事件就属于谁。此外，祂将西满的名字改为伯多禄，正如创造主也曾更改亚巴郎、撒辣依和曷舍亚的名字，将后者称为若苏厄，并给前两者各加一个音节。但为什么是\Emph{伯多禄}？如果是因为他信德的活力，还有许多坚固的材料可以从其强度借来名字。难道是因为基督既是磐石又是石头？因为我们读到祂被放置为\Quote{绊脚石和使人失足的磐石}。我略去其余经文。因此祂愿意将祂最爱的门徒赋予一个名字，这名字由祂自己特殊的象征性名称所提示；因为依我看，这比可能源于祂自身任何非比喻性描述的名字更为适合。有来自提洛及其他地区甚至海外的人群到祂那里去。圣咏曾预见到这一事实：\Quote{看，外邦人的支派，提洛，以及厄提约丕雅人，都在那里。人要说：熙雍是我的母亲；在那里生了一个人}（因为神人生了），祂按圣父的旨意建造了她；好叫你知道那时外邦人如何涌向祂，因为祂就是那按圣父旨意建造教会的神人——甚至也包括其他种族。依撒意亚也这样说：\BibleQuote{看，这些从远方来；这些从北方和西方来；这些从波斯地来。} \BibleRef{依 49:12} 关于他们，祂又说：\BibleQuote{你举目四望，看，所有这些都聚集在一起。} \BibleRef{依 49:18} 且又说：\BibleQuote{你看见这些陌生奇怪的人，心里必说：这些是谁给我生的？谁给我养育了这些？这些是从哪里来的？} \BibleRef{依 49:21} 这样的基督难道不是先知所预言的基督吗？马西昂派的基督又是什么？既然真理的歪曲是他们的乐趣，他就不可能是先知所预言的基督。\par

% structure: p0032 source=h2 target=subsection reason=relative-heading-rank; base=h2

\subsection{第十四章 基督的山中圣训。其方式和内容与创造主在安排中的言行极为相似。因此它提示了这样的结论：耶稣就是创造主的基督。真福八端。}

我现在来探讨祂的那些寻常训诫，借着这些训诫，祂将祂道理的独特性适用于我所说的祂作为基督的正式宣告。\BibleQuote{贫穷的人是有福的}（因为要理解希腊文中的这个词，正是需要这样）\BibleQuote{因为天国是他们的。}\BibleRef{路 6:20}恰恰就是祂以真福开始这件事，乃造物主的特征，祂无论在起初的创造命令中，还是在宇宙最终的祝圣中，都使用了同一祝福之声；因为祂说：\BibleQuote{我的心，}祂说，\BibleQuote{已吐露出一个极好的言词。}这将是那被承认为新约创始原则的祝福之\BibleQuote{极好的言词}，仿效旧约。那么，有什么可奇怪的呢，如果祂以造物主特有的属性\Emph{开始祂的职务}，而造物主总是以同样的言词爱、安慰、保护、\Emph{并}报复那乞丐、穷人、卑微者、寡妇和孤儿？好让你相信基督这种私人的恩惠好比一条从救恩的泉源流出的溪流。事实上，面对如此丰富的这类\Emph{善}言，我几乎不知该转向何方；如同在森林、草地或苹果园中。因此，我必须寻找偶然出现的材料。\par

\BibleQuote{保护孤儿和贫困者，为谦卑者和穷人伸冤；解救穷人，从恶人手中救出贫困者。} 同样在\BibleBook{}{圣咏}第七十一篇中：\BibleQuote{祂必以公义审判民中的贫困者，拯救穷人的子民。} 并在接下来的话中论及基督说：\BibleQuote{万邦都要事奉祂。} 但\Person{达味}只统治过犹太民族，所以没有人能以为这话是指\Person{达味}说的；然而祂已经亲自承担了穷人的境况，以及那些受贫困压迫的人，\BibleQuote{因为祂要从强者的手中解救贫困者；祂必怜悯贫困者和穷人，并救赎穷人的灵魂。祂必从高利贷和不义中救赎他们的灵魂，在他们的眼中，他们的名字必被尊崇。} 又说：\BibleQuote{恶人必归于地狱，连同所有忘记天主的万邦；因为贫困者必不被永远遗忘，穷人的忍耐必不会永远消失。} 又说：\BibleQuote{谁像主我们的天主？祂居于高处，却俯视天上和地上的卑微之事！祂从灰尘中提拔贫困者，从粪堆中高举穷人，使祂与祂百姓的首领同坐，就是说，在祂自己的国度里。} 同样早先在\BibleBook{列王}{纪}中，\Person{撒慕尔}的母亲\Person{亚纳}以这些话赞美天主：\BibleQuote{祂从地上提拔穷人，从粪堆中高举乞丐，使祂与祂百姓的首领同坐（就是说，在祂自己的国度里），并坐在荣耀的宝座上（甚至君王的宝座上）。} \BibleRef{撒上 2:8} 通过\Person{依撒意亚}，祂如何痛斥压迫贫困者的人！\BibleQuote{你们为什么要焚烧我的葡萄园，穷人的掠夺物在你们家中？你们为何击碎我的百姓，磨碎贫困者的脸？} \BibleRef{依 3:14} 又说：\BibleQuote{祸哉，那些颁布不义法令的人；因为他们的法令规定了邪恶，使贫困者偏离审判，剥夺我百姓中穷人的权利。} \BibleRef{依 10:1} 这些公义的审判祂也要求为孤儿寡妇，以及为安慰贫困者本身。\BibleQuote{为孤儿伸冤，以正义对待寡妇；来，让我们和好，主说。} \BibleRef{依 1:17} 对于那在每一个卑微阶段都得到创造主如此多怜悯眷顾的人，也将获得基督所应许的国度；怜悯和同情属于祂，而且早已属于那些承受应许的人。因为即使你认为创造主的应许是属地的，而基督的是属天的，但很明显，天从来不属于任何别的神，只属于那也拥有地的神；很明显，创造主甚至赐下了较小的应许（属地的祝福），为使我更易相信祂关于更大应许（属天的祝福）的话，胜过（\Person{马西昂}的神），后者从未通过任何先前赐予的较小恩惠来证明他的慷慨。\BibleQuote{饥饿的人是有福的，因为他们将得饱足。} \BibleRef{路 6:21} 我或许可以将这句话与前一句联系起来，因为只有穷人和贫困者才会饥饿，如果不是创造主特别设计，让类似祝福的应许作为福音的预备，使人知道那是祂的应许。因为祂通过\Person{依撒意亚}论及那些祂即将从地极召来的人，即外邦人，这样说：\BibleQuote{看，他们必快速而来，毫不迟延；} \BibleRef{依 5:26} \Emph{快速}，因为他们奔向时机的圆满；\Emph{毫不迟延}，因为他们不被古法的重负所阻碍。他们既不饥饿也不口渴。因此\Emph{他们将得饱足}——这应许只给予那些饥渴的人。祂又说：\BibleQuote{看，我的仆人将得饱足，你们却要饥饿；看，我的仆人将得畅饮，你们却要口渴。} \BibleRef{依 65:13} 至于这些对比，我们将在后面看到它们是不是基督的预兆。同时，对饥饿者充满的应许是创造主天主的安排。\BibleQuote{哭泣的人是有福的，因为他们将欢笑。} \BibleRef{路 6:21} 再回到\Person{依撒意亚}的经文：\BibleQuote{看，我的仆人将欢欣踊跃，你们却要蒙羞；看，我的仆人将喜乐，你们却要因心中忧伤而哭。} \BibleRef{依 65:13} 并认出这些对比也在基督的救恩安排中。的确，喜乐和欢欣是应许给那些处于相反状态的人——即忧伤、悲哀、焦虑的人。正如\BibleBook{}{圣咏}第一二五篇所说：\BibleQuote{流泪播种的，必将欢欣收割。} 此外，欢笑与喜乐和欢欣相伴，正如哭泣与忧伤和哀恸相伴。因此，创造主在预言欢笑和哭泣之事时，是首位说哀恸的人将欢笑的那位。因此，那以安慰穷人、谦卑者、饥饿者和哭泣者开始祂（进程）的那位，立刻急切地表明自己正是祂通过\Person{依撒意亚}的口所指出的一位：\BibleQuote{主的神临于我，因为祂给我傅了油，向穷人传报喜讯。} \BibleRef{依 61:1} \BibleQuote{贫困的人是有福的，因为天国是他们的。} \BibleRef{路 6:20} \BibleQuote{祂派遣我医治破碎的心。} \BibleRef{依 61:1} \BibleQuote{饥饿的人是有福的，因为他们将得饱足。} \BibleRef{路 6:21} \BibleQuote{安慰一切哀恸的人。} \BibleRef{依 61:2} \BibleQuote{哭泣的人是有福的，因为他们将欢笑。} \BibleRef{路 6:21} \BibleQuote{给熙雍哀恸的人，赐予华冠（或荣耀）代替灰尘，喜乐的油代替哀恸，赞美的外袍代替哀伤的神。} \BibleRef{依 61:3} 既然基督一进入祂的进程就完成了这样的服务，那么祂要么就是那位预言自己到来做这一切的那位；要么如果预言这事的那位尚未到来，那么对\Person{马西昂}的基督的吩咐就是可笑的（虽然我或许应当加上必要的），这吩咐命令祂说，\BibleQuote{当人们憎恨你们，辱骂你们，并且因人子的缘故，把你们的名字当作恶的加以抛弃时，你们是有福的。} \BibleRef{路 6:22} 在这宣告中，无疑是对忍耐的劝勉。好，创造主通过\Person{依撒意亚}另说了什么？\BibleQuote{不要怕人的辱骂，也不要因他们的藐视而沮丧。} 什么辱骂？什么藐视？那是因了人子的缘故而遭受的。哪个人子？按照创造主旨意（来到）的那位。我们从哪里得到证明？从针对祂所预言的那被剪除的事；正如祂通过\Person{依撒意亚}对煽动仇恨的犹太人所说：\BibleQuote{因了你们，我的名在外邦人中受了亵渎；} \BibleRef{依 52:5} 并在另一处说：\BibleQuote{将惩罚加于那交出自己生命、被外邦人藐视（无论仆人或官长）的那位。} 现在，既然仇恨是针对那位从创造主受使命的人子而预言的，而\WorkTitle{福音书}见证说，基督徒的名号源自基督，将因人子的缘故被恨，因为祂是基督，这就决定了在仇恨的事上那人子正是按照创造主的旨意而来的那位，仇恨也是针对祂预言的。即使祂尚未到来，现今存在的对祂名字的仇恨也绝不可能先于那承受此名者。但是，祂已在我们面前受了惩罚，并为我们的缘故交出自己的生命，放下生命，且被外邦人藐视。而那位（降生于世的）将就是那人子，我们的名字也因祂而被弃绝。\par

% structure: p0035 source=h2 target=subsection reason=relative-heading-rank; base=h2

\subsection{第十五章 继续山中圣训。其祸哉与创造主的安排严格一致。许多引自\WorkTitle{旧约}的引文证明这一点。}

\BibleQuote{同样，}祂说，\BibleQuote{他们的父辈也这样对待先知。} \BibleRef{路 6:26} 这是何等善变的基督——\Emph{马尔齐安的}基督！时而摧毁先知，时而维护先知！祂摧毁他们作为对手，使他们的门徒转离；祂维护他们作为朋友，谴责迫害他们的人。但是，\Emph{只要}维护先知与消灭先知的马尔齐安之基督不一致，那么谴责迫害先知的人就正与创造主之基督相称，因为祂在一切事上成就了先知的预言。再者，以父辈的罪责备儿子，这更属于创造主，而不属于那连个人的罪行也不惩罚的神。但你会说：祂不能被视为维护先知，仅仅因为祂想确认犹太人对自己先知的不敬之罪。那么，在此情况下，不应将任何罪归咎于犹太人；他们反而值得赞扬和嘉许，因为他们虐待了那绝对善的神（\Emph{马尔齐安的}）在如此长的时间后才激动起来要消灭的人。然而，我想，到此时祂已不再是绝对善的神；祂在创造主那里寄居了相当长时间，已不再是纯粹而简单的伊壁鸠鲁的神。因为请看祂如何屈尊诅咒，证明自己会恼怒、会发怒！祂竟宣布了\Emph{祸哉}！但有人对此词的含义提出异议，仿佛它较少带有诅咒意，更多是告诫。然而，区别何在？因为即使告诫也少不了威胁的刺，尤其是当它被\Emph{祸哉}所加剧时。此外，告诫和威胁都将是那知道如何发怒之人的手段。因为没有人会以告诫或威胁禁止一件事，除非他会在做这件事时施加惩罚。没有人会施加惩罚，除非他能发怒。另一些人承认这个词意味着诅咒，但他们认为基督发出祸哉并非出于祂自己真实的情感，而是因为这祸哉来自创造主，祂想向他们展示创造主的严厉，以便更赞扬自己在其真福中的忍耐。就好像创造主不能同时作为善神和审判者，在祂两种属性的卓越中，使仁慈成为祝福的序言，将严厉置于诅咒的后续；从而在两方面完全展开祂的训导，既追求祝福，又防范诅咒。祂早已说过：\BibleQuote{看，我已将祝福和诅咒摆在你们面前。} \BibleRef{申 30:19} 这话实在是福音这种特质的预兆。此外，那为了使人相信自己的良善而恶意对比创造主之严厉的，是怎样的存在？以贬低他人为支撑的推荐是微不足道的。然而，通过这样展示创造主的严厉，祂实际上确认祂是可畏的。如果祂是可畏的，那么祂当然更值得顺服而非轻慢；如此，马尔齐安的基督开始有利于创造主的利益而教导。然后，\Emph{基于上述承认}，既然涉及富人的祸哉属于创造主，那么对富人生气的是创造主，而非基督；而基督认可富人的诱因——即骄傲、炫耀、爱世界和藐视天主，这些使他们配受创造主的祸哉。但是，谴责富人难道不是出于那刚刚赞赏穷人的同一位\Emph{天主}吗？没有人不谴责他所赞赏之物的反面。因此，如果创造主被归于对富人的祸哉，那么对穷人的祝福也应归于祂；如此，创造主的全部工程都归于基督。——如果马尔齐安的神被归于对穷人的祝福，那么对富人的诅咒也应归于他；如此，他将与创造主平等，既良善又审判；从而无需区分两个神；当这种区分被移除时，只有真理存留，即创造主是唯一的天主。既然\Quote{\Emph{祸哉}}是一个表示诅咒或某种异常严厉呼喊的词，且由基督对富人说出，我必须表明创造主也藐视富人，正如我已表明祂是穷人的维护者，以便证明基督在这件事上与创造主一致，即使祂曾使所罗门富足。\BibleRef{列上 3:5} 但\Emph{关于此人}，当有选择留给他时，他宁愿求那他知道能取悦天主的事——即智慧——他便赚得了那他不偏爱的财富。天主赋予人财富并非不合宜，因为借财富，富人得安慰和帮助；且借财富，许多公义和爱德之事得以施行。但财富也伴随严重过失；正是因为这些，福音中宣布了祸哉给富人。\BibleQuote{你们已得到了，}祂说，\BibleQuote{你们的安慰}；当然，是从他们的财富，从这些财富为他们买来的世间的虚荣和浮华中。因此，在申命纪中，梅瑟说：\BibleQuote{恐怕当你吃饱了，建造了好房屋，你的牛羊金银增多时，你的心就高傲，忘记了上主你的天主。} \BibleRef{申 8:12} 同样，当希则克雅王因财宝而骄傲，在从巴比伦来的使者面前夸耀财富而非天主时，\Emph{创造主}藉依撒意亚的口责备他说：\BibleQuote{看，日子将到，凡你家中所有的，你祖先所积蓄的，都要被带到巴比伦去。} \BibleRef{依 39:6} 同样，祂藉耶肋米亚说：\BibleQuote{不叫富人夸耀自己的财富；夸耀的，应因认识上主而夸耀。} \BibleRef{耶 9:23} 同样，祂藉依撒意亚斥责熙雍的女子，她们因炫耀和财富丰裕而骄傲，\BibleRef{依 3:16} 正如在另一处，祂警告骄傲和显贵者：\BibleQuote{阴府扩大了，张开了口，显赫的、伟大的、富有的（这就是基督的\InnerQuote{祸哉富人}）都要下去；人必屈服，}就是那以财富自高的；\Quote{强人必受辱，}就是那因财富而强壮的。\BibleRef{依 5:14} 关于这些人，祂又说：\BibleQuote{看，万军的上主将挫败浮夸者连同他们的力量：高傲的必被砍倒，崇高的必倒在刀下。} \BibleRef{依 10:33} 这些是谁？就是富人。因为他们确实从财富中得到了他们的安慰、荣耀、尊贵和高位。在圣咏第四十八篇，祂也使我们不再关心这些，说：\BibleQuote{不要惊惶，若有人成为富有，若他的光荣增加；因为当他死时，什么也不能带走；他的光荣也不随之下降。}同样在圣咏第六十一篇：\BibleQuote{不要贪图财富；若它们向你炫耀，也不要倾心。}最后，这同一\Emph{祸哉}早已由亚毛斯对那也沉溺于享乐的富人宣布。\BibleQuote{祸哉他们，}他说，\BibleQuote{那些躺在象牙床上，舒身躺在榻上，吃羊群中的羔羊，牧牛群中的牛犊，随琴声歌唱，以为自己长久不衰，并非过客；饮酒上品，以最贵的油抹身的人。} \BibleRef{亚 6:1} 因此，即使我只能表明创造主劝人远离财富，而并未同时首先谴责富人——用基督所用的同样措辞——也没有人能怀疑，出自同一权威，基督的祸哉中附加了针对富人的警告，而这警告原先也是从那劝戒这些人物质之罪（即其财富）的同一源头而来。因为这样的警告是此类劝戒的必要后续。祂也宣布了祸哉给\BibleQuote{饱足的人，因为他们将要饥饿；给现在欢笑的人，因为他们将要哀恸。} \BibleRef{路 6:25} 与此对应的是，我们在上面看到创造主的祝福中出现的相反情形：\BibleQuote{看，我的仆人将得饱饫，你们却要饥饿}——因为你们曾饱足；\BibleQuote{看，我的仆人将喜乐，你们却要蒙羞} \BibleRef{依 65:13}——你们将哀恸，现在却欢笑。因为正如圣咏所写：\BibleQuote{流泪播种的，必含笑收割。}福音也是如此：\Quote{笑着（即喜乐）播种的，必含泪收割。}这些原则是创造主古时奠定的；基督只是将它们凸显出来，而并未改变它们。\BibleQuote{凡众人都夸赞你们的时候，你们就有祸了！因为他们的父辈也这样对待假先知。} \BibleRef{路 6:26} 创造主同样藉祂的先知依撒意亚，严厉责备那些追求人阿谀和称赞的人：\BibleQuote{我的百姓啊，称你们有福的人误导你们，扰乱你们所行的路。} \BibleRef{依 3:12} 在另一处，祂禁止完全信赖人，也禁止信赖人的称赞；如藉先知耶肋米亚：\BibleQuote{凡信赖世人的人，是可咒骂的。} \BibleRef{耶 17:5} 而在圣咏第一一七篇中说：\Quote{倚靠上主胜过倚靠世人；倚靠上主胜过依赖权贵。} 因此，凡人所追求的一切，直到他们的好话，都为创造主所憎恶。谴责给予假先知及其父辈的赞美和谄媚之词，与谴责他们对（真）先知的折磨和迫害，同样属于祂。正如先知所受的伤害不能归咎于他们自己的神，假先知所受的称赞也不会使任何别的神不悦，除了\Emph{真}先知的神。\par

% structure: p0037 source=h2 target=subsection reason=relative-heading-rank; base=h2

\subsection{第十六章 爱仇的诫命。它在旧约中造物主的圣经和基督的山中圣训里同样被教导。摩西的\OriginalTerm{同态复仇法}{Lex Talionis}与耶稣基督来为造物主宣扬和施行的仁慈和爱相一致，被令人钦佩地解释。各种爱德诫命被解释。}

\Quote{但是我告诉你们这听道的人}（此处显示出造物主那条古老的诫命：\Quote{向那些侧耳听你们的人说话}），\Quote{要爱你们的仇敌，为那恨你们的人祝福，为那毁谤你们的人祈祷。}造物主曾藉祂的先知依撒意亚把这些命令包含在一句诫命里：\Quote{你们要对那恨你们的人说：你们是我们的兄弟。}\BibleRef{依 66:5}因为，如果那些是我们的仇敌、恨我们、说我们坏话、毁谤我们的人要被称为我们的兄弟，那么祂教导我们将他们视为兄弟时，实际上就是命令我们祝福那些恨我们的人，为那些毁谤我们的人祈祷。诚然，基督显然教导一种新的忍耐，当祂实际上禁止了造物主所允许的报复（即要求\Quote{以眼还眼，以牙还牙}\BibleRef{出 21:24}），反而吩咐我们\Quote{有人打你这边脸，连那边脸也由他打；有人夺你的外衣，连内衣也由他拿去。}\BibleRef{路 6:29}无疑，这些是基督的补充，但完全符合造物主的教导。因此，必须立即确定这个问题：忍耐的戒律是否由造物主所命令？当祂通过匝加利亚命令说：\Quote{你们不要心里谋害\Emph{兄弟}}\BibleRef{匝 7:10}，祂没有明确包括\Emph{近人}；但随后在另一处祂说：\Quote{你们不要心里谋害\Emph{近人}。}\BibleRef{匝 8:17}祂既劝人忘记伤害，就更可能劝人忍耐地承受它。然而，当祂说：\Quote{报复是属于我的，我必报应}，祂便教导人要忍耐地等待报复的施行。因此，既然那同一（天主）一方面似乎要求\Quote{以牙还牙，以眼还眼}来报复伤害，另一方面却禁止不仅一切报复，甚至连报复的念头或回忆也不容许，那么这一点就越发清楚地显明给我们：祂要求\Quote{以眼还眼，以牙还牙}——并非为了允许通过报复来重复伤害（这实际上在禁止报复时已经禁止了），而是为了首先遏制伤害，即通过以报复或对等为惩罚而禁止伤害；这样，每个人考虑到允许施加第二次（报复性）伤害，就会避免犯下第一次（挑衅性）过错。因为祂知道，通过报复的前景来压制暴力，比通过（不确定的）报复的应许要容易得多。然而，鉴于人的本性和信仰，必须提供两种结果：使信天主的人期望从天主那里得到报复，而使没有信仰（约束）的人害怕规定报复的法律。法律这个难以理解的目的，基督作为安息日和律法以及圣父一切安排的主，既启示又使之明白，当祂命令\Quote{把另一边脸也转过来（给打的人）}时，为的是更有效地消灭一切因伤害而起的报复——法律曾想通过报复的方法来防止这种伤害，（而）启示显然已通过禁止记忆伤害并将报复交给天主来限制它。因此，无论基督引入什么（新规定），祂都不是反对法律，而是进一步推进它，丝毫没有损害造物主的诫命。所以，如果仔细考察忍耐被命令（并且达到如此圆满的程度）的缘由本身，就会发现它若不依靠那应许报复、自任审判者的造物主之命令，便不能成立。否则——如果这如此沉重的忍耐（即不以打还打；转过另一边脸；不仅不以骂还骂，反而祝福；不仅不保留外衣，反而连内衣也给人）是由一位无意帮助我的主加在我身上——那么，祂教我忍耐就毫无意义，因为祂没有给我诫命的赏报——我指的是这种忍耐的果实。祂本该允许我报复，如果祂自己不施行报复；如果祂不给我那份许可，那么祂自己就该施行报复；因为为纪律本身起见，伤害应当被报复。因为对报复的恐惧遏制了一切不义。但若不加区别地允许（不义）放纵，它就会肆虐——它会在其不受惩罚的安全中挖出人的双眼，打掉每颗牙齿。然而，这却是你那善良而单纯仁慈之神的原则——对忍耐行不义，向暴力敞开大门，使义人不得保护，恶人不受约束！\Quote{求你的，就给他}\BibleRef{路 6:30}——当然是给穷人，或者更确切地说，尤其是给穷人，尽管也给富足的人。但为使人不贫穷，你在《申命纪》中看到造物主为债权人规定了这项规定：\Quote{在你手中不可有穷人，好使上主你的天主在你所赐给你的地上赐福与你}——你是指债权人，因为多亏了他，那人才不贫穷。不仅如此。对不求的人，祂也命令给予。祂说：\Quote{不可有穷人\Emph{在你手中}；}换句话说，要留意，就你所能避免的，不可有；这一命令又进而更强烈地要求人给那求的人，如下文也说：\Quote{如果在你中间有一个穷人，是你的兄弟，你不可狠心，也不可对你穷兄弟袖手。你要向他伸手，照他所缺乏的借给他。}\BibleRef{申 15:7}贷款通常只借给那些求借的人。不过关于借贷的事，以后再谈。现在，如果有人想争辩说，造物主的诫命只限于对待兄弟，而基督的诫命则限于所有求的人，从而使后者成为一条新的不同的诫命，（我回答说）从那些显示造物主的法律在基督里重复的原则中，只能得出一个规则。因为基督命令向所有人所做的，与造物主为人的兄弟所规定的，并不是不同的。因为虽然向陌生人显示的爱德更大，但并不优于先前向近人所欠的。因为谁能对陌生人施予（基于认识的）爱呢？然而，爱德的第一步是向近人，第二步是向陌生人；那么第二步属于那也拥有第一步的主，而非那从未拥有第一步的主。因此，造物主遵循本性之律，首先教导对\Emph{近人}的仁慈，打算以后命令对\Emph{陌生人}的爱德；并且遵循祂治理之法的顺序，首先将爱德限于犹太人，以后扩展到全人类。因此，只要\Emph{祂治理}的奥秘仅限于以色列，祂便适当地命令只对兄弟施以怜悯；但当基督将\Quote{外邦人赐给祂为基业，将地极赐给祂为田产}时，便开始实现欧瑟亚所说的：\Quote{你们不是我的人民，那是我的人民；你们未蒙怜悯，那曾经蒙了怜悯}——即（犹太）民族。从此，基督将祂圣父的慈悲法律扩展到所有人，不将任何人排除在祂的慈悲之外，正如祂在邀请中不遗漏任何人。所以，无论祂教导的范围如何扩大，祂都在祂继承万民的基业中接受了这一切。\Quote{你们愿意人怎样待你们，也要怎样待人。}\BibleRef{路 6:31}这条命令无疑暗示了其对立面：\Quote{你们不愿意人怎样待你们，也不要怎样待人。}如果这是那位新的、先前未知、尚未充分宣告的神明的教导——祂没有事先给我任何教训，使我首先知道该为自己选择或拒绝什么，并照我所愿别人待我的去待别人，不做我不愿别人待我的事——那么这无非是祂留给我自己意见的随意之谈，没有给我任何适当的愿望或行动规则，使我能照自己所喜欢的待别人，或避免照自己所不喜欢的待别人。因为事实上，祂并未界定我该为自己及别人愿望什么或不愿望什么，好让我按照自己意志的法则塑造行为，并有能力不向别人做我愿意别人向我做的事——爱、服从、安慰、保护以及诸如此类的善举；同样，也不向别人做我不愿别人向我做的事——暴力、错待、侮辱、欺骗以及诸如此类的恶行。的确，那些未受天主教导的异教徒，就是根据意志和行为这种不协调的自由而行事。因为虽然善与恶各自为本性所知，但生命并非因此处于天主的纪律之下——唯独天主纪律最终在信仰中，如同在天主敬畏中，教导人意志与行为的适当自由。因此，马西翁的神，虽经特别启示，却无法发表上述诫命的任何摘要——这诫命此前是如此局限、模糊、黑暗，并且除凭我个人任意思想外无法轻易解释——因为祂在这样一条诫命的事上没有提供先前的辨别。但我的天主并非如此，因为祂始终处处命令保护、扶助、安慰穷人、孤儿和寡妇；因此依撒意亚说：\Quote{把你的饼分给饥饿的人，将无家可归的穷人领到你家中，见到赤身露体的，给他衣服穿。}\BibleRef{依 58:7}厄则克耳也如此描述义人：\Quote{把自己的饼给饥饿的人，用衣服遮盖赤身露体的人。}\BibleRef{则 18:7}那教训当时足以引导我去照我所愿别人待我的去待别人。因此，当祂发出这样的禁令：\Quote{不可杀人；不可奸淫；不可偷盗；不可作假见证}\BibleRef{出 20:13}——祂便教导我避免向别人做我不愿别人向我做的事；因此，福音中发展出来的这条诫命，将属于那唯一的主——祂自古就拟定了它，并赋予了它突出的特点，按照祂自己教导的决断安排了它，现在又按其重要性地把它简化为一个简明公式，因为（正如另一处所预言的）主——即基督——\Quote{要在地上说简短的话}。\par

% structure: p0039 source=h2 target=subsection reason=relative-heading-rank; base=h2

\subsection{第十七章 论借贷。禁止高利贷和取利的精神。律法是福音的预备，其规定也是如此；当前情形亦然。论报复。基督的教导全程证明祂是由造物主所派遣的。}

现在，关于借贷的问题，当祂问：\Quote{你们若借给那些你们希望收回的人，有什么可感谢的呢？}请将这与以下\BibleBook{厄则克耳}{}的话比较，其中祂论及前述的义人说：\Quote{他没有将钱放高利贷，也不取任何增益}——意指多余的利益，就是高利贷。第一步是要消除借贷所生的果实，使人更易于习惯本金本身的损失——若发生损失的话——因为他已学会失去其利息。现在，我们断言，这就是法律作为福音预备的功用。它致力于逐步塑造那些将要学习者的信德，以达至基督徒规范之圆满的光照，通过它所能提供的最佳规诫，灌输一种尚在踌躇表达中的仁爱。因为在上引的\Emph{\BibleBook{厄则克耳}{}段落中}，祂说：\Quote{你要归还借贷的抵押品}——当然是指给那无力偿还的人，因为，理所当然，祂不会规定将抵押品归还给有偿还能力的人。在\BibleBook{申命纪}{}中更清楚地规定：\Quote{你不可睡在他的抵押品上；日落时务必将外衣归还给他，使他在自己的外衣中睡觉。}\BibleRef{申 24:12} 更清楚的还有\Emph{前面一段}：\Quote{你要豁免邻居欠你的一切债务；对你的兄弟不可追讨，因为称之为上主你的天主的豁免年。}\BibleRef{申 15:2} 现在，当祂命令豁免那无力偿还者的债务时（因为祂甚至禁止向有能力偿还的人追讨，这是更强的论证），祂教导的岂不是我们应该借给那些我们无法收回的人吗？既然祂在借贷上施加了如此巨大的损失。\Quote{你们将成为天主的儿女。}还有什么比这更无耻的呢？他竟使我们成为他的\Emph{孩子}，他自己却因禁止婚姻而不允许我们生育子女？他怎能把一个他已经抹去的名字授予他的追随者？我尤其不能成为一个阉人的儿子，因为我有的圣父是宇宙所认领的那位伟大存在！因为宇宙的创造者不也是全人类的圣父吗？难道不比（\Person{马尔强}）那个被阉割的神祇更配称为圣父吗？那神祇并没有创造任何存在之物。即使创造者没有结合男女，没有允许任何生物生育子女，我在乐园之前、堕落之前、被逐之前、两人成为一体之前，就已经与祂有这层关系。当祂亲手塑造我，用祂的嘘气赐我活动时，我就第二次成为了祂的儿子。如今祂再次称我为祂的儿子，不是生我进入自然生命，而是进入属灵的生命。祂说：\Quote{因为祂恩待忘恩的和恶的。}\BibleRef{路 6:35} 做得好，\Person{马尔强}！你多么巧妙地收回了祂的雨水和阳光，使祂不显得像一位创造者！但这个从未被认识过的仁慈者是谁呢？他怎能是仁慈的？他之前从未表现出这样的仁慈，即借给我们阳光和雨水？他并不注定要从人类那里得到（应归于那创造者的）敬意；直到此刻，为回报祂在赐予元素上的慷慨，祂容忍人向偶像（而不是向祂自己）献上祂恩惠的应得回报。但\Emph{但天主}在属灵的祝福中确实是仁慈的。\Quote{上主的话语比蜂蜜和蜂房更甘甜。}祂曾斥责人们忘恩负义，而他们本应感恩——正是祂的阳光和雨水，连你，\Person{马尔强}，也享受过，却不知感恩！然而，你的神没有资格抱怨人的忘恩负义，因为他没有用任何方法使他们感恩。祂也教导怜悯：\Quote{你们要慈悲}，祂说，\Quote{如同你们的圣父对你们慈悲一样。}这条命令与\Quote{把你的饼分给饥饿的人；若他无家可归，就领他进你的家；若你见赤身露体的，就给他遮身}同属一辙；\BibleRef{依 58:7} 也与\Quote{应为孤儿伸冤，为寡妇辩护}相同；\BibleRef{依 1:17} 我在此认出那一位的古老教训，祂\Quote{喜爱仁慈胜过祭献}；\BibleRef{欧 6:6} 然而，如果现在是另一位存在者因自己的慈悲而教导慈悲，那么他为何在如此漫长的岁月里对我没有慈悲呢？\Quote{你们不要判断，你们也就不会被判断；不要定罪，你们也就不会被定罪；宽恕，你们也就被宽恕；给予，也就会给予你们：好的量器，压实的、满溢的，人们会倒在你们怀里。因为你们用什么量器量，也用什么量器量给你们。}\BibleRef{路 6:37} 在我看来，这段经文宣告了按功过施报的报应。但这报应由谁而来？若仅由人而来，那么他就只是教导一种人性的规范和报偿；我们就得在一切事上服从人；若由创造者而来，作为功过的审判者和报偿者，那么祂迫使我们服从祂，因为祂将报应放在祂手中，这报应将按每个人如何判断或定罪、释放或对待他的邻人而令人接受或可怕；若由（\Person{马尔强}的神）自己而来，那么他就会行使一种\Person{马尔强}所否认的审判职能。因此让马尔强派的人做出选择：难道抛弃他们师傅的训诫，与基督为人的利益或为创造者的利益而教导，不是同样自相矛盾吗？但\Quote{瞎子领瞎子，两人都掉进坑里。}\BibleRef{路 6:39} 有些人相信\Person{马尔强}。但\Quote{门徒不能高过师傅。}\BibleRef{路 6:40} \Person{阿佩莱斯}本该记住这一点——他是\Person{马尔强}的纠正者，尽管是他的门徒。异端者应先去掉自己眼中的梁木，然后当他怀疑基督徒眼中有微尘时，才能纠正他。正如好树不能结坏果子，同样真理也不能产生异端；正如坏树不能结好果子，同样异端也不能产生真理。因此，\Person{马尔强}从\Person{克尔东}的恶库中没有拿出任何好东西；\Person{阿佩莱斯}从\Person{马尔强}那里也没有。因为将基督一般性地用于人身上的比喻性话语应用在这些异端者身上，我们会得出更合适的解释，胜于按照\Person{马尔强}可悲的解释从中推断出两位神。我认为我有最充分的理由坚持我一贯的立场，即在任何地方都找不到基督启示了另一位天主。我奇怪\Person{马尔强}的手竟然唯独在这里感到麻木而不继续其篡改工作。但即使是强盗有时也会良心不安。没有不带着恐惧的恶行，因为没有不带着良心自责的恶行。因此，犹太人一直只认识祂，他们不知道祂以外的任何神；他们也只呼求他们所独知的那一位。既然如此，那么说\Quote{你们为什么称呼我：主啊，主啊？}的那位显然是谁呢？\BibleRef{路 6:46} 是那从未被呼求、因为从未被启示的那一位呢？还是那从起初就被视为上主、因为从太初就被认识的那一位——即犹太人的天主？再者，谁还会加上\Quote{却不实行我所说的话呢？}会是那当时才尽力教导他们的那一位吗？还是那从起初就藉法律和先知向他们传话的那一位？祂可以责备他们的不顺服，即使祂在其他任何时候都没有理由谴责他们。事实上，那当时归咎于他们古老顽梗的，正是那在基督来临之前就对他们说\Quote{这民族用嘴唇尊敬我，他们的心却远离我}的那一位。\BibleRef{依 29:13} 否则，一个新神、一个新基督、一个新而如此伟大宗教的启示者，竟然谴责那些他从未有机会试炼过的人为顽梗和不服从，这是多么荒谬啊！\par

% structure: p0041 source=h2 target=subsection reason=relative-heading-rank; base=h2

\subsection{第十八章 关于百夫长的信德、复活寡妇的独子、洗者若翰及其对基督的派遣、以及罪妇——从中提取所有证明基督与造物主的关系}

同样，当祂称赞百夫长的信德时，这是多么不可思议的事：祂竟承认自己\BibleQuote{连在以色列也没有见过这样大的信德}。\BibleRef{路 7:1} 但以色列的信德对祂毫无兴趣！但祂不可能从（基督在此所说的）事实中产生任何兴趣去认可和比较那尚属粗浅的，甚至可以说尚属虚无的事物。然而，祂为何不能使用另一个神的信德的例子呢？因为，如果祂这样做了，祂就会说在以色列从未有过这样的信德；但实际情况是，祂暗示祂本应在以色列找到如此大的信德，因为祂确实为此而来，祂实在是以色列的天主和基督，现在祂只是作为一个要执行和维护信德的人而谴责了它。如果祂真是它的反对者，祂会宁愿找到这样的信德，因为祂来是为了削弱和摧毁它，而不是认可它。祂也使寡妇的儿子从死里复活。\BibleRef{路 7:11} 这不是一个奇怪的奇迹。创造主的先知曾行过这样的奇迹；那么为什么祂的儿子不应该更行呢？现在，主基督显然没有引入另一个神来行如此重大的奇迹，以至于所有在场的人都归光荣于创造主，说：\BibleQuote{在我们中间兴起了一位大先知，天主眷顾了祂的子民。}\BibleRef{路 7:16} 哪个天主？当然是他们子民的天主，并且先知们来自祂。但如果他们光荣了创造主，而基督（听到他们并知道他们的意思）却没有纠正他们，即使在他们正是在这个死人复活中显明\Emph{祂的光荣}的伟大彰显中呼求创造主的行为上，那么无疑祂要么没有宣布另一位神，只允许祂在自己的恩惠行为与奇迹中受尊崇，要么祂怎能静静地任凭这些人长久停留在错误中，尤其是祂来正是为了医治他们的错误？但是若翰听到基督的奇迹时，如同听到一位外邦神的奇迹那样，感到冒犯。好，我首先要解释他冒犯的原因，以便更容易地消除我们异端者的丑闻。现在，既然万有权能的主自己，圣父的圣言和圣神，在地上行事和宣讲，那么那部分圣神以先知恩赐的形式通过若翰预备主的道路，就应该从若翰离开，当然再回到主那里，如同回到其包容一切的源头。因此，若翰现在成为一个普通人，只是众人之一，确实作为一个人而冒犯，但不是因为他期待或认为另一位基督教导或做什么新事，因为他甚至没有期待这样一位。没有人会对他不认识（因为他知道他不存在）且不期待或想不到的人产生怀疑。若翰十分肯定没有别的神，只有创造主，即使作为一个犹太人，尤其作为先知。他感到的任何怀疑显然是关于那位他确实知道存在但不清楚是否是那\Emph{基督}本身的那一位。因此，怀着这样的恐惧，若翰甚至问这个问题：\BibleQuote{你是那要来的那一位，或者我们期待另一位？}\BibleRef{路 7:20}——只是询问祂是否作为他所期待的那一位而来。\BibleQuote{你是那要来的那一位吗？}\Emph{即}你是那要来的吗？\BibleQuote{或者我们期待另一位？}\Emph{即}如果我们期待的那一位不是你，如果祂不是你，那么我们要等的那一位是不是另一位？因为他像所有人当时所想的那样，假设从相似的神迹证据来看，可能有一位先知同时被派来，而当时所期待的主自己与这位先知不同，并且比祂更优越。这就是若翰的困难。他怀疑是否所有人所期待的那一位真的来了，而且他们本应通过祂预言的工作认出祂，正如主派人告诉若翰，正是通过这些工作祂将被认出。\BibleRef{路 7:21} 现在，既然这些预言显然与创造主的基督有关——正如我们在对每一个的审查中所证明的——那么如果祂自称\Emph{不是}创造主的基督，却用那些祂赖以被接受为创造主之基督的证据来支持自己的说法，这足够荒谬。更荒谬的是，当祂不是若翰的基督时，祂却给予若翰见证，肯定他是一位先知，甚至祂的使者，将圣经应用在他身上：\BibleQuote{看，我派遣我的使者在你面前，他要在你前面预备你的道路。}祂仁慈地引用了先知话，以困惑的若翰所提到的选项的更高意义，以便通过肯定祂自己的先行者已经以若翰的身份来了，来熄灭隐藏在他问题中的怀疑：\BibleQuote{你是那要来的那一位，或者我们期待另一位？}现在，既然先行者已经完成使命，主的道路已经预备，祂现在应该被承认为先行者预备道路的那一位（基督）。\Emph{那位先行者}确实\BibleQuote{比一切妇人所生的都大}；\BibleRef{路 7:28} 但是尽管如此，在天国里最小的却不受他管辖；仿佛那个最小的人比若翰大的天国属于一个神，而比一切妇人所生的都大的若翰却属于另一个神。因为无论祂是由于卑微的地位而说到\Emph{任何}\BibleQuote{最小的人}，还是说到祂自己，因为被认为比若翰小——既然众人都跑到旷野跟随若翰，而不是跟随基督（\BibleQuote{你们出去到旷野要看什么呢？}\BibleRef{路 7:25}）——创造主同样有权声称若翰（比一切妇人所生的都大）和基督（或每一个\BibleQuote{最小的人\Emph{在天国里}}）都是祂自己的，这最小的人注定要在那个天国里比若翰更大，虽然同样属于创造主，并且他将比先知大得多，因为他不会因基督而跌倒，而当时若翰的伟大因\Emph{这一软弱}而减小。我们已经说过罪的赦免。那个\BibleQuote{一个有罪的女人}的行为，她用亲吻覆盖主的脚，用眼泪洗祂的脚，用头发擦干，用香膏抹上，\BibleRef{路 7:36} 证明了祂所接触的不是一个空洞的幻影，而是一个真正实在的身体，并且她作为罪人的忏悔按照创造主的意愿应得赦免，创造主习惯于喜爱仁慈胜过祭献。\BibleRef{欧 6:6} 但即使她忏悔的动机来自于她的信德，她听到通过忏悔因信德而成义的话：\BibleQuote{你的信德救了你}，这话是由祂说的，祂曾藉哈巴谷宣告：\BibleQuote{义人必因他的信德而生活。}\BibleRef{哈 2:4}\par

% structure: p0043 source=h2 target=subsection reason=relative-heading-rank; base=h2

\subsection{第19章 因比喻而跟随耶稣基督教导的虔诚富妇。马西昂派从基督被告知他的母亲和兄弟时的言论中衍生出的吹毛求疵。对基督看似拒绝他们的解释。}

一些富裕的妇女跟随了基督，\Quote{以她们的财产供给祂}，其中有王的家宰的妻子，这件事是预言的题目。主曾借着依撒意亚呼唤这些富有的妇女——\Quote{起来，安逸的妇女，听我的声音}——为要首先证明她们是门徒，然后成为助手和帮助者：\Quote{女儿们，怀着希望听我的话；要纪念这日子，在劳苦中怀着希望。}因为她们是在\Quote{劳苦}中跟随祂，并且\Quote{怀着希望}供给祂。关于\Emph{比喻}，只需指出这种语言曾由造物主同样明确地应许过，便足够了。但祂对百姓说话的那种直接方式——\Quote{你们听是听，却不明白}\BibleRef{依 6:9}——现在需要留意，因为它为基督提供了祂殷切教导中那种常见的形式：\Quote{有耳可听的，就听吧！}\BibleRef{路 8:8}。并非基督因受另一种精神驱使，就允许人听闻，而造物主却拒绝了；乃是因为劝勉紧随着警告而来。首先是\Quote{你们听是听，却不明白}；接着是\Quote{有耳可听的，就听吧！}因为他们虽有耳朵，却固执地拒绝聆听。然而，祂教导他们，需要的是心的耳朵；而造物主曾说过，他们不会用\Emph{这些}来听。因此，祂藉着祂的基督加上：\Quote{你们应当小心怎样听}\BibleRef{路 8:18}，并且不要听——当然是指心的聆听，而非耳朵。如果你将适当的意义赋予\Emph{造物主}的训勉，与那位用\Quote{你们应当小心怎样听}这句话唤醒百姓听道者的含义相称，那么这便等于对那些不愿听者的一种威胁。事实上，你们那最仁慈的神，不审判也不发怒，却是威吓的。这由紧接着的句子便可证明：\Quote{凡有的，还要给他；凡没有的，连他自以为有的，也要从他夺去。}\BibleRef{路 8:18}将给什么呢？信德、理解力甚至救恩的增加。将夺去什么呢？当然就是那将要给予的。这给予和剥夺是谁做的呢？如果剥夺是由造物主做的，那么给予也应由祂做。如果给予是由马西昂的神做的，那么剥夺也应由马西昂的神做。现在，无论祂为何原因威胁\Quote{剥夺}，这都不会是一位不知道如何威胁、因为不能发怒的神的工作。此外，我感到惊奇的是，他说\Quote{灯通常是不隐藏的}\BibleRef{路 8:16}，而他自己却将更大更必要的光隐藏了那么长时间；并且他应许\Quote{一切都要从隐藏中显露出来，成为显明的}\BibleRef{路 8:17}，而他自己却一直将他的神藏在暗处，等待（我想）直到马西昂出生。我们现在来讨论所有质疑主降生的人所极力主张的论据。他们说，祂作证自己不是由母亲所生，因为祂问：\Quote{谁是我的母亲？谁是我的兄弟？}\BibleRef{玛 12:48}异端者就是这样，要么凭猜测将简单明了的词曲解成他们想要的任何意思，要么按字面强行解释那些含有条件意义且不能简单解决的话，就如此处。我们方面首先回答说：如果祂没有母亲和兄弟，那么告诉祂祂的母亲和兄弟站在外面想见祂——这原本是不可能的。那报告消息的人必定先前或在当时就知道他们，那时他们想见祂或托人给祂传话。针对这第一个论点，对方通常给出这样的回答：但也许他们给祂传话是为了试探祂呢？然而，经上没有这样说；既然经文通常指明试探的事（\Quote{看，有一个法学士站起来，试探祂说}\BibleRef{路 10:25}；又如，当问及纳税时，法利塞人来到祂那里，\Quote{试探祂}\BibleRef{路 20:20}），因此，当它没有提到试探时，就不得作试探的解释。然而，\Emph{尽管我不接受这种解释}，我还是可以作为一个多余的驳斥问一问：所谓的试探有何理由？他们提到祂的母亲和兄弟，想要达到什么目的？如果是想查明祂是否由母亲所生——何时在这个问题上产生过疑问，以致他们要用这种方式试探来解决？谁能在看到祂是一个真实的人时，还怀疑祂是由母亲所生呢？——他们曾听到祂称自己为\Quote{人子}？——他们看到祂完全拥有人的品质，怀疑祂是神还是天主子？——他们更以为祂是一位先知，确实是一位伟大的，但仍然是作为人而生的。即使有必要通过调查祂的出生来这样试探祂，也肯定会有其他更好的证据，而不是通过提及那些亲人——尽管祂确实由母亲所生，但那时很可能没有这些亲人。请告诉我，母亲在每个情况下都与儿子同时活着吗？每个人生来都有兄弟吗？一个人难道不能没有父亲和姐妹（活着），甚至根本没有亲人吗？但有历史证明，当时在犹大地方，由森提乌斯·撒图尔尼努斯进行了一次\Emph{人口登记}，这或许可以满足他们对基督家族和世系的调查。因此，这种试探的方法毫无根据；那些\Quote{站在外面的人}确实是\Quote{祂的母亲和兄弟}。我们还需要探究祂在使用非字面话时的含义，当时祂说：\Quote{谁是我的母亲或我的兄弟？}似乎祂的话等于否认了自己的家庭和出生；但实际上，这是由于情况的绝对性和祂话语要按条件解释而产生的。祂义怒，因为如此亲近的人\Quote{站在外面}，而陌生人却在里面聆听祂的话语，更何况他们想叫祂离开手头的庄严工作。祂与其说是否认，不如说是不承认他们。因此，当先前的问题\Quote{谁是我的母亲？谁是我的兄弟？}之后，祂补充回答\Quote{他们就是听我的话而实行的人}，祂将血缘关系的名称转移给了其他人，祂认为这些人由于信德与祂更亲近。没有人将某物从没有该物的人那里转移。因此，如果祂使那些不是\Quote{祂的母亲和兄弟}的人成为\Quote{祂的母亲和兄弟}，祂又怎能否认那些实际拥有这些关系的人呢？当然只是基于他们的功过，而不是对近亲的任何否认；用祂自己的实际榜样教导他们：\Quote{谁爱父亲或母亲或兄弟胜过天主的话，就不配做祂的门徒。}\BibleRef{玛 10:37}此外，祂承认母亲和兄弟，从祂不愿承认他们的事实中反而更加清楚。祂接纳其他人，只是证实了那些因过犯而被祂拒绝的人与祂的关系，祂用其他人替代他们，不是作为更真实的亲属，而是作为更相称的。最后，如果祂宁愿选择那亲属所没有的信德，这也算不上什么大事。\par

% structure: p0045 source=h2 target=subsection reason=relative-heading-rank; base=h2

\subsection{第二十章  比较基督掌管风浪的能力与\Person{梅瑟}掌管\Place{红海}和\Place{约但河}之水的能力。基督掌管邪魔的能力——军旅的例子。血漏的治愈——论\Person{梅瑟}法律在这方面所规定的不洁。}

但\BibleQuote{这是什么样的人？他竟命令风和海水！} \BibleRef{路 8:25} 他当然是元素的新主人和所有者，既然创造主已被废黜，并被排除出其所有权！决非如此。但元素承认它们自己的制造者，正如它们习惯于服从祂的仆人一样。请细查\WorkTitle{出谷纪}，\Person{马西昂}；看\Person{梅瑟}的棍杖，它挥动着祂的命令指向红海，比犹太的所有湖泊都更广阔。海如何从深处裂开，然后固定成两座凝固的庞然大物，并因此从它们之间的间隔中，为人们开出一条干地穿越的道路；再一次，同一根棍杖震动，海恢复了力量，在其水流的汇聚中，埃及的军兵被吞没！连风都服务于这个结局！也请读，约但河如何像一把剑，阻止那迁徙的民族渡河；它的上游水如何停止，下游的水流完全停止流动，遵\Person{若苏厄}的命令，\BibleRef{苏 3:9} 当他的司祭开始过河时！对此你会说什么？如果那\Emph{上文所指的是你的基督}，他不会比创造主的仆人更有能力。但我本应满足于已经列举的例子而不加补充，如果祂这次海上的行程的预言不是在基督来临之前就已出现。实际上，一篇圣咏正是因这次过湖而应验。圣咏作者说：\BibleQuote{主在多水之上。}当祂驱散其波浪时，\Person{哈巴谷}的话应验了，他说：\BibleQuote{在祂的行程中分散诸水。}当祂一斥责，海就平静时，\Person{纳鸿}也得到证实：祂斥责海，使海干涸，\BibleRef{鸿 1:4} 还包括使海翻腾的风。你愿意我用什么证据来证明我的基督？是从创造主的实例，还是从预言？你认为祂被预言为一个军事的、武装的战士，而不是一个以比喻和寓意的意义，要用灵性武器进行灵性战争对抗灵性敌人，进行灵性战役的人：来吧，当你在一个人身上就发现一群自称\Emph{\OriginalTerm{军团}{Legion}}的魔鬼时，\BibleRef{路 8:30} 当然是由邪灵组成的，你就应该明白，基督也必须被理解为灵性敌人的消灭者，祂挥舞灵性武器，进行灵性斗争；而且正是祂，现在不得不对抗甚至一个魔鬼的军团。因此，圣咏可能确实论及这样的战争：\BibleQuote{主是强而有力的，主在战斗中是大能的。}因为祂与最后的敌人——死亡战斗，并通过十字架的胜利得胜。现在，军团见证耶稣是谁的儿子？\BibleRef{路 8:28} 毫无疑问，是那他们知道并畏惧其折磨和深渊的天主的儿子。他们似乎不可能直到这时还不知道那新近出现且陌生的神的能力在世上运行，因为创造主对其无知是极不可能的。因为如果祂曾一度不知道在自己之上还有另一位神，那么到这时祂至少已经发现有一位神在祂的天下工作。现在，他们的主所发现的事，到这时在祂整个家族内已众所周知，这家族处在同一个世界和同一穹苍之下，那位陌生的神就在那里居住和行动。因此，正如如果那位神存在，创造主和祂的造物必然知道祂一样，既然祂不存在，魔鬼们实在只知道他们自己天主的基督。他们不求问那位陌生的神，他们记得必须向创造主祈求——不被投入创造主的深渊。他们的请求终于得到了应允。凭什么？因为他们说了谎？因为他们宣告祂是一位无情天主的儿子？那帮助说谎者、支持毁谤者的神会是怎样的神呢？然而，\Emph{无需这样想}，因为既然他们没有说谎，既然他们承认深渊的天主也是他们的天主，那么祂也确实亲口肯定祂就是那些魔鬼所承认的那一位——耶稣，审判者和报复之天主的儿子。现在，看啊，在基督身上有创造主缺陷和软弱的迹象；因为我这一方有意将无知归于祂。请允许我在对抗异端者的努力中有些宽容。耶稣被一个患血漏的妇人触摸，\BibleRef{路 8:43} 祂不知道是谁。\BibleQuote{谁摸了我？}祂问道，当时祂的门徒在辩解。祂甚至坚持祂无知的断言：\BibleQuote{有人摸了我，}祂说，并提出一些证据：\BibleQuote{因为我感觉有能力从我身上出去了。}我们的异端者说什么？\Emph{基督}能知道那个人吗？祂为什么装作不知道？为什么？当然是为了挑战她的信德，考验她的敬畏。正如祂曾经好像无知地问\Person{亚当}：亚当，你在哪里？这样，你就看到创造主同样被谅解，而基督也效法创造主。但在这里，祂作为法律的敌对者行动；因此，既然法律禁止接触有血漏的妇人，\BibleRef{肋 15:19} 祂不仅希望这妇人触摸祂，而且还要医治她。那么，这里有一位不是出于本性仁慈，而是出于敌对的神！然而，如果我们发现这妇人的信德有此功劳，以至于祂对她说：\BibleQuote{你的信德救了你，}\BibleRef{路 8:48} 你算什么，竟在那件事中察觉对法律的敌意，而主亲自向我们表明那事是作为信德的奖赏而做的？但你是否会认为这妇人的信德在于她对法律获得的轻视？谁能设想，一个一直不认识任何神、尚未接受任何新法律的妇人，竟会暴力违反她至今所遵守的法律？的确，她冒险触犯是基于什么信德？相信哪位神？貌视谁？创造主？她的触摸至少是一种信德的行为。如果是基于对创造主的信德，那么她怎么可能违反祂的法律，既然她不知道任何别的神？无论她违法到什么程度，都源于并正比于她对创造主的信德。但这两件事怎能相容？她违反了法律，并且是出于信德而违反，而信德本应阻止她如此违反？让我告诉你，她的信德尤其是这样的：它使她相信她的天主喜欢仁爱胜过祭献；她确定她的天主在基督内工作；因此她触摸祂，不是只当作一个圣人，也不是当作一个先知——她知道他们因人性而可能被玷污——而是当作真天主，她认为祂不可能被任何不洁玷污。因此，她并非无理地为自己解释法律，意思是，那些易受沾染的东西会变污秽，但天主不然，她确知天主在基督内。但她也想起，法律所禁止的是那种每月的正常出血及生育时的出血，而不是由于健康失调所致的出血。然而，她的情况是长期的健康失调，她知道这需要天主仁慈的救助，而不是时间的\Emph{自然}缓解。因此，她显然可以被视为辨别了法律，而不是违背了法律。这将是那同时要赐予智慧的信德。\BibleQuote{如果你们不相信，}先知说，\BibleQuote{你们必不理解。}\BibleRef{依 7:9} 当基督赞赏这妇人的信德——它单纯地信靠创造主——祂通过回答她，\BibleRef{路 8:48} 声明祂自己就是她所信的、祂所赞许的那信德的神性对象。我也不能忽视这一事实：祂的衣服被触摸也证明了祂身体的真实性；因为当然是身体，而非幻影，被衣服所覆盖。这确实不是我们现在的论点；但这个评论对我们所讨论的问题有自然的关联。因为如果不是一个真实的身体，而只是一个虚幻的身体，它作为无实质之物肯定不会受玷污。那么，那因实体之空虚而不能受玷污的祂，又怎么可能渴望被触摸呢？作为法律的敌对者，祂的行为是欺骗性的，因为祂不可能受到实际污染。\par

% structure: p0047 source=h2 target=subsection reason=relative-heading-rank; base=h2

\subsection{第21章 基督与创造者的关联由旧约数事及圣路加叙述门徒使命相比较而显明。群众吃饼。圣伯多禄的宣认。以基督为耻。这种羞耻只能对真基督发生。马西昂派的妄论荒谬可笑。}

祂派遣门徒去宣讲天主的国。\BibleRef{路 9:1}这里祂说的是哪一位天主呢？祂禁止他们为旅途携带任何食物或衣物。除了那喂养乌鸦、装饰田野花朵的天主，谁会颁布这样的诫命呢？自古谁曾命令：\BibleQuote{牛在场上踹谷的时候，不可笼住它的嘴}，好让它从劳作中自由获取饲料，因为工人应得他的工资？\BibleRef{申 25:4}马尔西雍可以删去这些规条，但无所谓，只要其含义存留。但当祂命令他们，对那些拒绝接待的人，要跺掉脚上的尘土作为\Emph{见证}时，祂也吩咐这样做。无人会在无需法律判决的案件中作证；而凡是命令将不人道的行为交由证词审判的，确实是以审判者的姿态发出威胁。再者，基督所推荐的不是一位新神，这由所有人的意见清楚证实，因为有人向黑落德说耶稣是基督；有的说祂是若翰；有的说祂是厄里亚；还有的说祂是古代先知中的一位。\BibleRef{路 9:7}无论祂是哪一位，祂当然不会在复活后为了宣告另一位神而兴起。祂在荒野地方喂饱了群众；\BibleRef{路 9:10}要知道，这是依照旧约的方式。否则，如果同样的伟大没有出现，那么祂现在比造物主更卑微。因为\Emph{祂}，不是一天，而是四十年；不是靠低劣的面包和鱼，而是靠天上的玛纳；不是维持五千人，而是六十万人的生命。然而，祂的奇迹的伟大在于，祂愿意这微薄的食物不仅足够，而且甚至有余；在此祂遵循了古代的典范。因为同样，在厄里亚时代的饥荒中，匝尔法特寡妇那少得可怜的最后一把面，因先知的祝福而倍增，持续了整个饥荒时期。你们有列王纪上卷。\BibleRef{列上 17:7}如果你们也翻到列王纪下卷，会发现基督的这一切行为都曾由那位天主的人所行：他吩咐将给他的十个大麦饼分给众人；当他的仆人将人数众多与食物微少对比后回答：\BibleQuote{难道我要将这些摆在一百人面前吗？}他再说：\BibleQuote{给他们，他们必吃；因为上主这样说：他们必吃，且必有余，正如上主的话。}\BibleRef{列下 4:42}基督啊，即使在祢的新事中，祢仍是旧的！因此，当亲眼目睹这奇迹、并与古代先例相比较、且在其中发现预言日后将成之事的伯多禄，作为众人的代言人回答主的询问：\BibleQuote{你们说我是谁？}\BibleRef{路 9:20}他说：\BibleQuote{祢是基督，}他不可能不领悟祂就是那位基督，因为在圣经中他不认识别的，而此刻他正在祂奇妙的作为中瞻仰祂。连祂自己也证实了这一点，祂到此为止容忍了这个回答，甚至下令对此缄默。\BibleRef{路 9:21}因为如果伯多禄无法承认祂是别的，只能承认祂是造物主的\Emph{基督}，而祂却吩咐他们\BibleQuote{不要把这话告诉任何人}，那么祂显然不愿宣讲伯多禄所得出的结论。当然，你会说，因为伯多禄的结论是错误的，所以祂不愿散播谎言。然而祂为缄默指定了另一个理由，就是因为\BibleQuote{人子必须受许多苦，被长老、司祭长和经师们弃绝，并且被杀，但第三天必要复活。}\BibleRef{路 9:22}既然这些\Emph{苦难}确实是为造物主的基督所预言（我们将在适当的地方充分展示），那么祂将这些预言应用在自己身上，就证明了祂就是那被预言的一位。无论如何，即使这些事未被预言，祂强加缄默的理由也足以说明伯多禄没有错，那个理由就是祂必须受这些苦。祂说：\BibleQuote{凡想要保全自己生命的，必丧失性命；凡为我丧失自己生命的，必保全性命。}\BibleRef{路 9:24}说这话的确实是人子。那么，请与巴比伦王一起仔细看那燃烧的火窑，在那里你将发现一位\BibleQuote{相似人子者}（因为祂尚未真正是人子，尚未从人生），早在当时就指定了这样的结局。祂救了那三位弟兄的性命，\BibleRef{达 3:25}他们原是同意为天主而丧失性命；但祂消灭了加色丁人的性命，因为他们宁愿藉拜偶像保全性命。你所声称的新奇在哪里？那教义有这些古老的证据。然而，关于将要发生的殉道，以及他们将从天主获得赏报的一切预言都已应验。依撒意亚说：\BibleQuote{看，义人如何丧亡，没有人放在心上；虔诚人被收去，没有人注意。}\EditorialNote{参阅依撒意亚 57:1}什么时候比在祂的圣人受迫害时更常发生这样的事？这确实不是寻常之事，不是自然律的普通意外；而是那光辉的热忱，为信仰而战，在其中凡为天主丧失生命的，必保全生命，因此你在此再次认出那位审判者，祂以毁灭回报生命的邪恶所得，以救恩回报生命的善意丧失。然而，祂在此向我呈现的是一位忌邪的天主，以恶报恶。因为祂说：\BibleQuote{凡以我为耻的，我也将以他为耻。}\BibleRef{路 9:26}这样的羞耻只可能归于我的基督。祂整个一生如此蒙羞，以致为异端者打开了嘲弄之路，他们竭尽苦毒地攻击祂出生和成长的极度耻辱，以及祂肉身的卑贱。但是，你那基督如何能蒙受一种他不可能经历的羞耻呢？他从未在女人的子宫中凝成血肉之躯（即使是童贞女）；从未从人的种子生长（尽管仅按有形体的法则，从女人的体液）；在子宫成形之前从未被视为肉身；成形之后从未被称为\OriginalTerm{胎儿}{fœtus}；从未在十月怀胎的阵痛后分娩；从未在分娩的剧痛间与生产时通过身体排泄的不洁之物一起被放倒在地，立即以眼泪开始生命之光，并以那最初的创伤将孩子与母体分离；从未接受丰富的沐浴，也没有盐和蜜的敷抹；也未用襁褓作为裹尸布；之后也从未在母亲膝上自己的污秽中翻滚，吮吸她的乳房，长期为婴儿，渐渐为孩童，缓慢成长为成人。但他是由天上显现，立刻成人，立刻完美，立刻为基督；纯粹是灵、能力和神。但因为他不是真实的，因而是不可见的，所以他无需因十字架的咒诅而蒙羞，他因缺乏身体而逃避了真正的承受。因此，他绝不可能说：\BibleQuote{凡以我为耻的。}但我们的基督不能不如此宣告；祂被父\BibleQuote{稍微逊于天使}，\BibleQuote{是虫，不是人，是人所羞辱，为百姓所藐视}；因为祂愿意\BibleQuote{因祂的创伤我们得痊愈}\BibleRef{依 53:5}，使我们的救恩因祂的卑微而确立。祂理当为祂自己的受造人而谦卑自己，为祂自己而非别人的肖像和模样，好使人——既然他在向石头或木头下拜时不感到羞耻——能以同样的勇气在信仰中显出同样程度的无耻，不以基督为耻，从而为他的偶像崇拜的无耻而补偿天主。那么，马尔西雍，就该羞耻而言，哪一种更适合你的基督呢？显然，你本人因给了他一个虚构的存在而应当羞愧。\par

% structure: p0049 source=h2 target=subsection reason=relative-heading-rank; base=h2

\subsection{第二十二章。\Emph{显圣容}支持同一结论。\Person{马尔基翁}将\Person{梅瑟}和\Person{厄里亚}这两位造物主的杰出仆人与在荣耀中的\Person{基督}联系在一起，这很不一致。按\OriginalTerm{孟他努主义原则}{Montanist Principle}解释\Person{圣伯多禄}的无知}

你在这方面也应当深感羞愧，因为你竟容许他出现在僻静的山上，与\Person{梅瑟}和\Person{厄里亚}在一起，\BibleRef{路 9:28}，而他正是来毁灭他们的。这无疑是他希望人们理解的那来自天上的声音的含义：\BibleQuote{这是我的爱子，你们要听从他}\BibleRef{路 9:35}——这里强调的是\Emph{他}，即不再听从\Person{梅瑟}或\Person{厄里亚}。因此，单凭这声音就足够了，无需显示\Person{梅瑟}和\Person{厄里亚}；因为，明说他们应当听从谁，这必然禁止了听从其他任何人。或者，难道他的意思是，那些他没有显示的\Person{依撒意亚}和\Person{耶肋米亚}等人应当被听从，因为他禁止了那些他所显示的人？现在，即使他们的出现是必要的，他们也不应当被描绘成彼此交谈，这是亲密的标志；也不应当与他在光荣中相联，这表示尊重和恩宠；而应当被显示在某个泥沼中，作为他们毁灭的确凿记号，甚至是在造物主的黑暗之中——基督被派来驱散的黑暗——远远远离那将要使他们的言语和著作与他的福音分离的那一位的荣耀。那么，这就是他证明他们是外人的方式：甚至让他们留在自己的陪伴中！这就是他表明他们应该被舍弃的方式：他反而让他们与自己联合！这就是他毁灭他们的方式：他以自己的光荣照亮他们！他们自己的基督会怎么做？我想他大概会效仿异端的乖僻，正如\Person{马尔强}的基督必定做的那样，或者至少会让自己与任何其他人而不是他自己的先知在一起！但还有什么比以下事情更适合造物主的基督呢？——让他在他自己的预报者们的陪伴中显示自己？——让他与那些他曾显现于启示中的人一同被看见？——让他与那些曾谈论他的人交谈？——与那些他曾被他们称为光荣之主的人分享他的光荣？甚至与他那些主要的仆人在一起，其中一位曾是祂子民的塑造者，另一位后来是祂子民的改革者；一位是旧约的创始者，另一位是新约的完成者？因此\Person{伯多禄}确实做得对，当他认识到基督的同伴与基督有着不可分解的联系时，他提议说：\BibleQuote{我们在这里真好}（真好：这显然意味着与\Person{梅瑟}和\Person{厄里亚}同在的地方）；\BibleQuote{让我们搭三座帐棚，一座为你，一座为\Person{梅瑟}，一座为\Person{厄里亚}。但他不知道自己在说什么。}\BibleRef{路 9:33}怎么不知道？他的无知是单纯错误的结果吗？还是出于我们在新预言事业中所坚持的原则，即恩宠中含有神魂超拔或狂喜？因为当一个人被圣神所夺，尤其当他目睹天主的荣耀，或当天主藉他说话时，他必然失去知觉，因为他被天主的能力所笼罩——这一点在我们与属血肉思想的人之间存在争议。现在，证明\Person{伯多禄}的神魂超拔并不困难。因为他怎么可能知道\Person{梅瑟}和\Person{厄里亚}，除非是在圣神内？人们不可能拥有他们的形象、雕像或肖像，因为法律禁止这样做。如果不是在圣神内看见他们，他怎么能知道呢？因此，因为他是凭圣神说的，而不是凭他的自然感官，所以他不知道他所说的。但另一方面，如果他是这样无知，因为他错误地以为耶稣是他们的基督，那么显然\Person{伯多禄}，当先前被基督问及\BibleQuote{他们认为他是谁}时，他回答说\BibleQuote{你是基督}时，指的是造物主的基督；因为如果他当时知道他是属于敌对的神，他在这里就不会犯错误。但如果他在这里错误是因为他先前错误的观点，那么你可以肯定，直到那一天，基督没有启示任何新的神性，而且\Person{伯多禄}到那时为止没有犯错误，因为至今基督没有启示过这类事；因此，基督不应被视为属于任何其他神，而只属于造物主，他实际上在这里描述了造物主的整个安排。他从他的门徒中挑选了三位见证人，来见证即将到来的异象和声音。这正是造物主的方式。\BibleQuote{凭两三个见证人的口，每句话都得成立。}他退到山上。在这个地方的性质上，我看到了深意。因为造物主当初就在一座山上，以可见的荣耀和祂的声音塑造了祂的古选民。新约正应当在这等高处得到证实，正如旧约曾在此处订立；同样也有一片云彩遮盖，没有人会怀疑，这云彩是从造物主的空气中凝聚而成的。除非，他确实把自己的云彩带到了那里，因为他自己强行穿过了造物主的天空；或者，他所用的不过是造物主的一片不定形的云彩。因此，在此时（如同在往昔），云彩并没有沉默；而是有习惯的声音从天上传来，以及圣父对圣子的见证；正如在第一篇圣咏中祂曾说：\BibleQuote{你是我的儿子，我今日生了你。}藉\Person{依撒意亚}的口，祂也曾问及他说：\BibleQuote{你们中间谁敬畏天主？让他听他儿子的声音。}因此，当祂在此以这些话介绍他：\BibleQuote{这是我的（爱）子}，这一句当然被理解为\Quote{我所许诺的那一位}。因为如果祂曾许诺，然后后来说\Quote{这就是他}，那么，对于成就祂旨意的那一位来说，发出声音以证实祂先前所作的许诺是合宜的；但对于一个可能受到反驳\Quote{你岂有权利说\InnerQuote{这是我的儿子}？关于他，你事前没有给我们任何信息，正如你没有赐给我们关于你自己先前存在的启示一样}的那一位来说，是不合宜的。因此，\BibleQuote{听从祂}，就是听从那位从一开始（造物主）就宣布有资格以先知的名义被听从的那一位，因为祂必须被人民视为先知。\Person{梅瑟}说：\BibleQuote{上主你的天主要从你弟兄中给你兴起一位先知（当然是指按肉身血统）；你们要听从祂，如同听从我一样。}\BibleRef{申 18:15}\BibleQuote{凡不听从他的，必从民中剪除。}\BibleRef{申 18:19}同样，\Person{依撒意亚}也说：\BibleQuote{你们中间谁敬畏天主？让他听他儿子的声音。}\BibleRef{依 50:10}圣父自己将要推荐这声音。因为，他说，当祂说\BibleQuote{这是我的爱子，你们要听从祂}时，祂是在确立祂儿子的言语。因此，即使顺从的\BibleQuote{听从}从\Person{梅瑟}和\Person{厄里亚}转移到基督，这仍不是来自另一位神，或转向另一位基督；而是从造物主转向祂的基督，是由于旧约的逝去和新约的来临。\Person{依撒意亚}说：\BibleQuote{不是使者，也不是天使，而是祂自己要拯救他们；}因为正是祂自己在宣布并实现法律和先知。圣父将新门徒赐给了圣子，在\Person{梅瑟}和\Person{厄里亚}与祂一同被显示于祂荣耀的尊荣中之后，他们就被遣散了，因为他们已完全履行了他们的职责和职务，明确是为了向\Person{马尔强}确认一个事实：\Person{梅瑟}和\Person{厄里亚}甚至分享基督的荣耀。但我们在哈巴谷书里也看到了同一异象的整个结构，那里圣神以某些宗徒的口说：\BibleQuote{上主，我听见了你的声音，就害怕。}这声音是什么呢？岂不就是那天上的声音所说的话：\Quote{这是我的爱子，你们要听从祂}？\BibleQuote{我默观你的作为，就惊讶。}什么时候比\Person{伯多禄}看见祂的荣耀而不知自己说什么时更合适呢？\BibleQuote{在二者之间，你将被人认识}——即\Person{梅瑟}和\Person{厄里亚}。\Person{匝加利亚}也曾在两棵橄榄树和橄榄枝的形像下看见他们。因为他们就是他所说的：\BibleQuote{这是两位受傅者，侍立在全地上主旁边。}哈巴谷又说：\BibleQuote{他的荣耀遮蔽诸天（即被那云遮蔽），他的光辉如同光——甚至那光，他的衣服也因之发光。}如果我们提及给\Person{梅瑟}的许诺，我们会发现它在此处应验了。因为当\Person{梅瑟}渴望看见上主，说：\BibleQuote{我若在你眼前蒙恩，求你将你自己显示给我，使我能清楚地看见你，}他渴望看见的，是他将来要作为人所取的那种状态，而作为先知他知道这事将要发生。然而，关于天主的\Emph{面容}，他早已听说：\BibleQuote{没有人看见我还能存活。}\BibleQuote{你所说的这件事，我必为你做。}然后\Person{梅瑟}说：\BibleQuote{求你显示你的荣耀给我。}上主同样指向未来，回答说：\BibleQuote{我要在你面前经过，显我的荣耀}等等。最后祂说：\BibleQuote{然后你要看见我的背。}他想要看的不是腰或腿肚，而是将在末世显现的荣耀。祂曾许诺要这样面对面地向他显现，当时祂对\Person{亚郎}说：\BibleQuote{你们中间若有先知，我必要藉异象显示给他，藉异象与他说话；但对待\Person{梅瑟}却不是这样；我要与\Emph{他}面对面、明明地说话（就是说，以祂将要取的人形），而不以谜语。}\BibleRef{户 12:6}现在，虽然\Person{马尔强}否认此处\Person{梅瑟}被描绘成与上主说话，而只是站着，但既然他站着\BibleQuote{口对口}，他也必须站着\BibleQuote{面对面}与\Emph{他}——用他的话来说——不远于他，在他的荣耀中——更不用说，在他面前。带着这荣耀，他从基督那里得到光照而离开，正如他从前从造物主那里得到光照一样；正如从前使以色列子民的眼目眩惑，如今也击打那盲目了的\Person{马尔强}的眼目，他未能看出这个论证也如何反对他。\par

% structure: p0051 source=h2 target=subsection reason=relative-heading-rank; base=h2

\subsection{第二十三章 马尔基昂的基督不可能责斥无信的一代；真基督对于婴孩所显示的仁爱，也是另一基督所不能的；论基督在撒玛黎雅所遇见的三种不同性格的人及其教训}

我以以色列人自居。让\Person{马尔西昂}的基督上前喊说：\Quote{无信德的世代！我同你们在一起，直到几时呢？我容忍你们直到几时呢？}\BibleRef{路 9:41} 他立即要听我的这番反驳：\Quote{无论你是谁，陌生的客，先告诉我们你是谁，从谁而来，你对我们有什么权利。至今你所拥有的一切全属造物主。当然，你若来自祂，为祂行事，我们就接受你的责备。但若来自别的神，我要你告诉我们，你曾赐给我们什么属于你自己的，以致我们有责任相信，因为你竟以\InnerQuote{无信德}来责备我们，而你自己从未向我们启示。你何时开始与我们交往，竟抱怨迟延？你在何事上容忍过我们，竟夸耀你的忍耐？就像\Person{伊索}的驴，刚从井边来，就到处嘶叫。} 此外，我又假设自己是门徒，就是他曾痛骂的：\Quote{乖僻的世代！我同你们在一起，直到几时呢？我容忍你们直到几时呢？}对他的这一爆发，我当然可以最合理地用以下话语反斥他：\Quote{无论你是谁，陌生的客，先告诉我们你是谁，从谁而来，你对我们有什么权利。至今，我想，你属于造物主，所以我们跟随了你，在你身上认出一切属于祂的事物。现在，你若来自祂，我们就接受你的责备。但若你为另一个行事，请告诉我们，你曾赐给我们什么纯粹是你自己的，以致我们有责任相信，因为你竟以\InnerQuote{无信德}来责备我们，尽管至今你没有给我们任何凭据。你何时开始与我们争辩，竟指责我们迟延？你在何事上容忍过我们，竟夸耀你的忍耐？驴刚从\Person{伊索}的井边来，就已经在嘶叫了。}现在，谁不会如此驳斥责备的不公，如果他以为责备者属于那尚无权利抱怨者？除非连祂也不会痛骂他们，如果祂没有从前藉着法律和先知，以及大能作为和许多仁慈住在他们中间，并且一直经验到他们是\Quote{无信的}。 但是，看，基督抱起婴孩，教导众人应当像他们一样，如果想成为更大的。\BibleRef{路 9:47} 相反地，造物主却放出熊来攻击儿童，为报复被他们嘲笑的大先知\Person{厄里叟}。\BibleRef{列下 2:23} 这种对比十分无礼，因为它把截然不同的事物——无辜的幼年和已有分辨能力的童年——混为一谈；童年即便不亵渎，也能嘲弄。因此，既然天主是正义的天主，祂没有宽恕不虔敬的儿童，因为祂要求每个年龄阶段的尊荣，尤其当然来自少年。而天主又是良善的，祂如此爱护婴孩，以至于在\Place{埃及}祝福保护希伯来婴孩的收生婆，这些婴孩因法郎的命令而处于危险中。\BibleRef{出 2:15} 因此，基督与造物主共享这份慈爱。至于\Person{马尔西昂}的神，他是婚姻的敌人，怎么可能看似热爱孩童呢？孩童不过是婚姻的产物。恨种子的人也必憎恨果实。是，他应当被看作比\Place{埃及}王更残忍。因为法郎禁止婴孩被养大，而\Emph{他}却不允许他们出生，剥夺他们在母胎中十个月的存在。那么，把对孩童的慈爱归给那位为人类繁衍而祝福婚姻，并在此祝福中包括了婚姻果实本身（首先就是婴孩）的应许的那一位，是多么可信！造物主应\Person{厄里亚}的请求，在关于那个假先知（\Person{巴耳则步布}）的事上，从天上降下火击打。\BibleRef{列下 1:9} 我在此看到审判者的严厉。而我，相反地，看到基督严厉责备门徒，因为他们想要对\Place{撒玛黎雅}那个无名村庄施以同样的惩罚。\BibleRef{路 9:51} 异端者也可发现，基督的这种温良是由同一位最严厉的审判者所预许的。祂说：\Quote{他不争辩，也不在街上听到他的声音；压伤的芦苇，他不折断；将熄的灯心，他不吹灭。}\BibleRef{依 42:2} 具有这样的性情，他当然更不愿焚烧人。因为在当时，主也对\Person{厄里亚}说：\Quote{他不在火中，而在轻微细弱的风声中。}\BibleRef{列上 19:12} 然而，为何这位最仁爱慈祥的天主拒绝那自愿作祂不可分离同伴的人呢？\BibleRef{路 9:57} 如果那人因骄傲或伪善说\Quote{你不论往哪里去，我要跟随你}，那么，以审判的姿态责备骄傲或伪善的行为为可弃，祂便履行了审判者的职务。当然，祂所拒绝的人，祂判定其失去跟随救主的益处。因为正如祂召唤那未被拒绝或自愿被邀请的人进入救恩，同样祂将所拒绝的人交付灭亡。然而，当祂回答那以埋葬父亲为借口的人说：\Quote{任凭死人埋葬他们的死人，你去宣传天主的国吧}\BibleRef{路 9:59}，祂明确确认了造物主的两条法律——肋未记中关于司祭职务，禁止司祭甚至出席父母的葬礼。祂说：\Quote{司祭不可进入有死人的地方；为父亲也不可染上不洁}；以及户籍纪中关于（\OriginalTerm{纳齐尔}{Nazarite}）离俗誓愿的规定；因为在那里，那献身于天主的人，除其他事项外，被命令\Quote{不可接触任何尸体}，甚至父亲、母亲或兄弟的也不可。\BibleRef{户 6:6} 我想，祂是为\OriginalTerm{纳齐尔}{Nazarite}人和司祭职务而预备这个人，激励他宣传天主的国。否则，若不然，那没有任何法律诫命却命令儿子忽视父母葬礼的人，必被宣称为不虔敬。的确，在我们面前的第三个例子中，（基督）禁止那人\Quote{回头看}，那人原想先\Quote{向家人告别}，祂只是遵循造物主的规矩。因为祂曾阻止那些从\Place{索多玛}被救出的人回头观看。\BibleRef{创 19:17}\par

% structure: p0053 source=h2 target=subsection reason=relative-heading-rank; base=h2

\subsection{第二十四章 论七十门徒的派遣及基督对他们的训示。取自旧约的先例。假设\Person{马尔西昂}的基督能赐予践踏蛇蝎的权柄之荒谬。}

此外，除了十二宗徒外，他还拣选了另外七十位门徒。既然十二宗徒对应于\Place{厄林}十二股水泉的数目，那么七十位不也应与那地方的七十棵棕树数目相符吗？比较中的\OriginalTerm{对立论}{Antitheses}无论是什么，造成它们的主要是原因上的差异，而非权能上的差异。但如果一个人不留意原因的差异，他就很容易推断权能的差异。当以色列子民出埃及时，造物主使他们带着金银器皿的战利品，以及衣服和无酵面团的重担出来；而基督却命令祂的门徒甚至连行路的棍杖都不带。前者被赶入旷野，而后者却被派往城市。试想场合所呈现的差异，你就会明白，安排祂的子民在一种情况下因贫乏而行事，在另一种情况下因丰裕而行事，是同一位权能。当祂的子民可以从城市得到补给时，祂削减了他们的供应；正如当他们在旷野面对匮乏时，祂为他们积存了物资。祂甚至禁止他们带鞋。因为正是在祂的庇护下，百姓在旷野中许多年之久鞋子没有穿破。\BibleRef{申 29:5}祂说：\BibleQuote{你们在路上不要向任何人请安。}\BibleRef{路 10:4}基督真是先知毁灭者吗？因为祂的训令也是从他们领受的！当\Person{厄里叟}派遣他的仆人\Person{革哈齐}先行去复活\Person{叔能妇人}的儿子时，我认为他给了这样的指示：\Quote{束上你的腰，手里拿我的棍杖，去你的路；若遇见任何人，不要请安；若有人向你请安，也不要回答。}因为途中的祝福不就是人们相遇时的相互请安吗？主也同样命令：\BibleQuote{无论进了哪一家，先说：愿这家平安。}\BibleRef{路 10:5}在此祂效法了同样的先例。因为\Person{厄里叟}命令他的仆人遇见\Person{叔能妇人}时，同样请安；要对她说：\Quote{愿你丈夫平安，愿你的孩子平安。}这倒更像是我们的\OriginalTerm{对立论}{Antitheses}；它们将基督与造物主相比较，而非分割开。\BibleQuote{工人应当得工钱。}\BibleRef{路 10:7}谁能比审判者更合适宣判这样的判决呢？因为判定工人应当得工钱，本身就是一种司法行为。没有任何赏赐不包含审判的过程。造物主的律法在这一点上也给我们提供了印证，因为祂判定耕牛也像工人一样应得工钱：祂说：\BibleQuote{牛在场上踹谷的时候，不可笼住它的嘴。}\BibleRef{申 25:4}那么，谁比那位也怜悯牲畜的天主对人更好呢？基督宣告工人应得工钱时，实际上正是为造物主关于夺取埃及人金银器皿的诫命开脱。因为那些为埃及人建造房屋和城市的人，无疑是应得工钱的工人，他们并未被教导行欺诈，而只是去索取无法以其他方式从主人那里获得的报酬。祂以这种方式肯定了天主的国并非新事或未曾听闻，同时命令他们宣告它已临近。\BibleRef{路 10:9}只有曾经遥远的东西，才可以恰当地说它变近了。然而，如果某物在接近之前从未存在过，就永远不能说它已临近，因为它从未存在于远处。一切新而未知的事物也是突然的。因此，突然的事物只有在被宣告时才首次获得时间属性，因为它那时才第一次呈现其形式。此外，一件事物要么在未被宣告时始终迟延，要么自开始被宣告时才临近，两者必居其一。\par

祂也加添说，他们应对那些不愿接待他们的人说：\BibleQuote{然而你们应当知道，天主的国已经临近你们了。}\BibleRef{路 10:11} 如果祂不是以威吓的方式命令这事，那么这项命令便是最无用的。因为对他们来说，天国的临近有什么要紧呢？除非它的来临伴随着审判——甚至是为那些接受这宣告的人带来救恩。如果威胁可以不应验，那么在一个施威胁的天主身上，你们怎能同时有一位也执行的天主，并且在这两者中，有一位是审判者呢？同样，祂又命令将脚上的尘土跺下来，反对他们，作为证据——即那片土地上可能沾在凉鞋上的灰尘颗粒，更不用说与他们有任何别的来往。\BibleRef{路 10:11} 但如果他们的无礼和不友善不会从祂那里受到任何报复，那么祂预先提出一个证据，又有何目的呢？这证据必定预示着某些威胁。此外，当创造者也在\BibleBook{申命}{纪}中禁止接纳阿孟人和摩阿布人进入教会，因为当祂的子民从埃及出来时，他们不人道、不友善地欺骗性地扣留了他们的供应\BibleRef{申 23:3}，那么明显可以看出，这种交往的禁令是从祂（创造者）传给基督的。祂所使用的形式——\BibleQuote{轻视你们的，就是轻视我}\BibleRef{路 10:16} ——创造者也对梅瑟说过：\BibleQuote{他们不是埋怨你们，而是埋怨我。}\BibleRef{户 14:27} 梅瑟确实是一位宗徒，正如宗徒们是先知一样。两种职分的权柄必须同等划分，因为它们都来自同一位主，即宗徒和先知的天主。谁将赐予\BibleQuote{践踏蛇和蝎子的权能}\BibleRef{路 10:19}呢？是那位一切生物的主宰，还是那位连一只蜥蜴都不能支配的神呢？幸运的是，创造者曾藉\BibleBook{依撒意亚}{先知书}应许，甚至将这种权能赐给小孩子，使他们将手放在蛇穴和幼蛇的洞口，而毫不受伤。\BibleRef{依 11:8} 事实上，我们知道（没有曲解经文的字面意思，因为即使这些毒虫在有信德的地方也实际不能伤害人）在蝎子和蛇的形象下，预表了邪灵，它们的首领被描述为蛇、龙以及创造者权能下一切最显著的野兽的名字。这权能创造者首先赐给了祂的基督，正如\BibleBook{}{圣咏集}第九十篇对祂说：\BibleQuote{你要践踏蝰蛇和毒蛇，将狮子和龙踏在脚下。} 同样\BibleBook{依撒意亚}{先知书}说：\BibleQuote{在那一日，上主天主必将抽出祂神圣、伟大而有力的剑}（即祂的基督）\BibleQuote{攻击那龙，那巨大而弯曲的蛇；并在那一日将它杀死。} 但当同一位先知说：\BibleQuote{那条路将被称作洁净而神圣的路；不洁之物不得从那里经过，那里也没有不洁的路；但分散的人要从那里经过，他们不会迷路；那里没有狮子，也没有猛兽上去；在那里找不到它}，他指出了信德之路，藉此我们将到达天主那里；然后他向这信德之路应许了彻底摧毁和制服一切毒虫。最后，如果你读那段经文前面的话，你就会发现这应许的适当时期：\BibleQuote{你们软弱的手要坚强，你们颤抖的膝要稳定；那时瞎子的眼睛要睁开，聋子的耳朵要听见；那时瘸子要跳跃如鹿，哑巴的舌头要得以说话。} 因此，当祂宣告祂医治的益处时，也将蝎子和蛇放在祂的圣人脚下——那正是首先从圣父那里领受这权能，以便将它赐予他人，然后按预言次序彰显出来的那位。\par

% structure: p0056 source=h2 target=subsection reason=relative-heading-rank; base=h2

\subsection{第二十五章 基督感谢圣父将祂向智慧者隐藏的事启示给婴孩。这隐藏之事由创造者明智地成就。\BibleBook{路加}{福音}第十章中的其他论点被表明惟独创造者的基督才可能做到。}

谁若没有被显示为天的主的创造者，谁又会被祈求为天的主呢？因为祂说：\Quote{我感谢祢，父，并承认祢是天的主，因为祢将这些事对明智和通达的人隐藏起来，却启示给了婴孩。}\BibleRef{路 10:21}这些事是什么？是谁的事？被谁隐藏？又被谁启示？如果是由马西昂的天主隐藏和启示的，那就是极其不义的行为；因为他从未产生过任何东西，在其中可以隐藏什么——没有预言、没有比喻、没有异象、没有事物、言语或名字的证据，被比喻和象征或模糊的谜语所遮蔽，但他甚至隐藏了自己的伟大，而他却正藉着他的基督全力启示出来。现在，明智和通达的人做了什么错事，致使天主向他们隐藏，而他们的明智和通达却不足以认识祂？祂自己没有提供任何途径，藉着祂的作为的任何宣告或任何痕迹，使他们得以成为明智和通达。然而，如果他们甚至对一位不认识的天主有所亏欠，假设他现在终于被认识了，他们也不应该在他身上，即被引入为与造物主不同的那一位身上，发现一位嫉妒的天主。因此，既然他没有提供任何材料使他可以隐藏任何东西，也没有任何冒犯者使他可以向他们隐藏自己；再者，即使他有，他也不应该向他们隐藏自己，那么现在他就不会是自己先前没有隐藏的启示者；所以，除了在祂身上所有这些属性都一致地相遇者之外，没有人会是天的主或基督的父。因为祂藉着预言晦涩的预备工作来隐藏，其理解对信德是开放的（因为\Quote{如果你们不信，你们必不领悟}\BibleRef{依 7:9}）；并且祂有那些明智和通达的冒犯者，他们不愿寻求天主，尽管祂在如此众多和伟大的工作中可以被发现，\BibleRef{罗 1:20}或者他们鲁莽地论及祂，从而为异端者提供了他们的技艺；最后，祂是一位嫉妒的天主。因此，基督感谢天主所做的事，祂早已藉着依撒意亚宣布了：\Quote{我要消灭智慧人的智慧，并隐藏通达人的聪明。}同样在另一处，祂暗示祂既隐藏了，也将启示：\Quote{我要赐给他们隐藏的宝藏，并向他们揭示隐秘的珍宝。}又说：\Quote{谁又能驱散腹语者的预兆，和那些凭自己内心占卜的人的计谋？使智者转退，使他们的谋略变为愚蒙？}现在，如果祂指定祂的基督为外邦人的光明，说：\Quote{我已立祢作外邦人的光明}；如果我们理解这些话指的是\Emph{婴孩}\BibleRef{路 10:21}——即曾经在知识上是矮子，在通达上是婴孩，甚至在信德的卑微上也是婴孩——那么我们就可以更容易地理解，那位曾经隐藏\Emph{这些事}、并应许藉着基督启示它们的天主，就是现在向婴孩启示了这些事的天主。否则，如果是马西昂的天主揭示了造物主先前隐藏的事，那么他就是通过展示造物主的作为而做了造物主的工作。但你说，他是为了毁坏造物主而做的，为要驳斥他们。因此，他应该向那些造物主向他们隐藏了这些事的人，就是明智和通达的人，来驳斥他们。因为如果他所做的有善意，知识的恩赐应该归于那些造物主扣留了知识的人，而不是归于\Emph{婴孩}，因为造物主没有吝啬任何恩赐给婴孩。但毕竟，我想，迄今为止我们在基督内所发现的是对法律和先知的建造而非拆毁。祂说：\Quote{一切都被我父交给了我。}\BibleRef{路 10:22}如果祂是造物主的基督，万有都属于祂，你可以信祂；因为造物主没有将比祂小的圣子所创造的一切（即藉着祂的圣言所创造的）交付给祂。相反，如果祂是那个臭名昭著的外来者，那么\Quote{\Emph{一切}}被父交给他的是什么？是造物主的一切吗？那么父交给圣子的一切都是好的，因此造物主是好的，因为祂的一切\Quote{事}都是好的；而那位侵占了别人的好（领域）来交给他的儿子、从而教导抢夺别人财物的人，就不再是好的了。他必定是最虚伪的存在，除了靠攫取别人的财产来使他的儿子致富外，别无他法！或者，如果造物主的一切都没有被父交给他，他有什么权利要求自己拥有（权柄）统治人？或者，如果人已经被交给了他，而且只有人，那么人就不是\Quote{一切}。但圣经清楚地说明，\Emph{万有}已经转移给了圣子。然而，如果你将这\Quote{\Emph{一切}}解释为全人类，即\Emph{万民}，那么连\Emph{这些}交给圣子也在造物主的计划之内：\Quote{我将赐给你外邦人作你的基业，地极作你的产业。}事实上，如果他有一些自己的东西，可以连同造物主的\Emph{人}一起全部交给他的儿子，那么请显出这一切中的一件作为样品，使我能够相信；免得我对那位我未见其任何事物者的\Quote{一切}所持的不信，与我确信那些未见之事属于那位我见其万有者一样合理。但\Quote{除了圣子，没有人知道父是谁；除了父和圣子愿意启示的人，没有人知道子是谁。}\BibleRef{路 10:22}所以基督所传讲的是未知的天主！别的异端者也以这段经文自撑，反对说造物主是人所共知的，对以色列是因熟识，对外邦人是因本性。但祂自己又如何见证祂不为人所知呢？\Quote{以色列却不认识我，我的百姓也不明白我}\BibleRef{依 1:3}；对外邦人也是如此：\Quote{因为看啊，在列国中我无人。}因此祂视他们\Quote{如桶中的一滴}，而\Quote{熙雍祂留下如葡萄园中的了望台。}你看，这里岂不证实了先知的言语，他斥责人对天主的无知，这无知持续到人子的日子？正是为此，祂插入了这句话：父被祂所启示的人认识，因为祂正是那被父立为外邦人之光明的那一位，自然需要关于天主的启蒙，同时也为以色列人，甚至赋予他们更圆满的天主知识。因此，对于相信那敌对的天主，那些适合造物主的论证是无用的，因为只有不适合造物主的论证才能推进对祂的敌对者的信仰。如果你也看下一句话：\Quote{看见你们所看见的，那眼睛是有福的。我告诉你们，从前有许多先知想看你们所看见的，却没有看见}\BibleRef{路 10:23}，你会发现它源于以上的意义：确实没有人按应有的方式认识天主，因为连先知也没有看见在基督内正被看见的事。现在，如果祂不是我的基督，祂在这段经文中就不会提到先知。因为如果他们看不见一位他们不认识、且在他们之后很久才被启示的天主的事，那有什么可奇怪的呢？然而，那些当时看见别人自然看不见之事的人有什么福气呢？因为是他们从未预言过的事，他们没有获得看见；除非因为他们可以正当地看见属于他们天主的事，这些事他们甚至曾预言过，但同时却没有看见？然而，这将是其他人的福气，就是那些看见别人仅预言过的事的人。我们稍后会表明，不，我们已经表明，在基督内看见了那些曾被预言却对预言它们的先知隐藏的事，为要也向明智和通达的人隐藏。在真福音中，一位法学士来到主面前问：\Quote{我该做什么才能承受\Emph{永生}？}在异端福音中，只提到生命，没有\Emph{永生}的限定；所以这位法学士似乎只是就造物主在法律中应许延长的生命请教基督，而主因此按法律回答他说：\Quote{你要全心、全灵、全力爱主你的天主}\BibleRef{路 10:27}，因为问题只关乎\Emph{单纯}生命的条件。但法学士当然很清楚法律所指的生命如何获得，所以他的问题不可能与他自己惯于教导的规则之生命有关。但既然连死人在基督内都已复活，他自己被这些复活者的例子所激发，盼望永生，就不再仅仅注目（使他如此充满盼望的奇妙之事）而浪费时间。所以他请教基督关于获得永生的事。因此，主自己仍是同一位，没有引入新的诫命，除了一条首要关乎（人）整个救恩，包括现在和将来生命的诫命，便将法律的精髓放在他面前：他要尽心、尽力爱主他的天主。如果他所问和基督所答的只是关于延长的生命（这生命在造物主手中），而不是关于永生（这在马西昂的天主手中），那么他如何获得永生呢？当然不会和延长生命的方式相同。因为酬报的不同必然意味着服务的不同。因此，你的门徒马西昂不会因爱你的天主而获得永生，正如那爱造物主的人会获得延长生命一样。但若那应许延长生命的祂应受爱戴，那赐予永生的祂岂不更应受爱戴呢？因此，两种生命都在同一位主手中；因为两种生命都要遵循同样的纪律。造物主教导要爱的，祂也必然藉着基督来维持，因为这里的原则是：从祂那里先有较小的证据，比从祂那里没有先前较小的迹象而要求相信更高之事的那一位，更容易相信更大的事。那么，无论\Emph{永生}一词是否被我们增添，对我来说已经足够：那邀请人得到永生（而非延长生命）的基督，在被问及祂所摧毁的现世生命时，并没有选择劝勉那人追求祂所引入的永生。请问，假如那使人爱造物主的基督不属于造物主，造物主的基督会怎么做？我想祂会说：不要爱造物主！\par

% structure: p0058 source=h2 target=subsection reason=relative-heading-rank; base=h2

\subsection{第二十六章 从圣路加福音第十一章的其他证据表明基督来自造物主。天主经与基督的其他话语。哑吧邪魔及逐魔时基督的言论。人群中妇人的呼喊。}

當祂在某處向那上方之父祈禱時，\newline{}\BibleRef{路 11:1}\newline{}以傲慢無禮的目光仰視造物主的天，按祂粗暴殘忍的本性，祂本可用冰雹和雷電擊碎祂——正如祂後來安排祂在耶路撒冷被釘在十字架上——\newline{}祂的一個門徒前來對祂說：\Quote{師傅，請教導我們祈禱，如同若翰教導了他的門徒。}\newline{}這話說得好像他認為不同的神需要不同的祈禱！\newline{}提出這種猜測的人應首先證明基督曾宣報另一位神。\newline{}因為無人會想知道如何祈禱，除非他先學會向誰祈禱。\newline{}如果他已學會了，就請證明。\newline{}若你找不到任何證明，讓我告訴你：他是向造物主請求教導祈禱，若翰的門徒也曾向祂祈禱。\newline{}但既然若翰引入了某種新的祈禱秩序，這個門徒合理地推測他也應問基督：他們是否也必須（按他們師傅的某種特別規則）祈禱——不是向另一位神，而是以另一種方式。\newline{}因此，基督在賜予門徒對天主的認識之前，不會教導他祈禱。\newline{}所以祂實際教導的是向門徒已認識的那位祈禱。\newline{}簡言之，你可從祈禱的含義中發現所呼求的是哪位天主。\newline{}我能向誰說\Quote{父啊？}\newline{}\BibleRef{路 11:2}\newline{}是向那與創造我無關，我並非從祂獲得起源的那位？\newline{}還是向那藉創造和塑造我而成爲我父母的那位？\newline{}我能向誰求賜聖神？\newline{}向那不賜予世俗之靈的那位，還是向那\Quote{使祂的天使成爲神靈}，\Quote{祂的神在起初運行在水面上}的那位？\newline{}\BibleRef{創 1:2}\newline{}我願誰的國來臨？\newline{}是那位我從未聽說其爲榮耀之王的，還是那位連君王之心也在祂手中的？\newline{}誰將賜予我日用的食糧？\newline{}\BibleRef{路 11:3}\newline{}是那位連一粒小米也不爲我產出的，還是那位甚至從天上天天賜給祂的子民天使之糧的？\newline{}誰將寬免我的罪過？\newline{}\BibleRef{路 11:4}\newline{}是那位因拒絕審判而不保留罪過的，還是那位若不寬免就保留罪過直至審判的？\newline{}誰將不讓我們陷入誘惑？\newline{}是那位誘惑者從不敢顫慄的，還是那位從起初就預先定了誘惑天使的罪的？\newline{}若有人以這樣的措辭呼求另一位神而非造物主，他不是祈禱，而是褻瀆。\newline{}同樣，我應向誰求，才能獲得？\newline{}向誰尋找，才能找到？\newline{}向誰敲門，才能爲我開門？\newline{}\BibleRef{路 11:9}\newline{}誰能賜予求者，除非是那萬物所屬、我也是求者所屬的那位？\newline{}但在那位神面前我損失了什麼，以致我要向他尋找並找到？\newline{}若是智慧和謹慎，那是造物主隱藏的。\newline{}難道我要向他求這些嗎？\newline{}若是健康和生命，它們在造物主手中。\newline{}任何事物都不可在其他地方尋找和找到，只應在它被隱藏以便顯露的地方。\newline{}同樣，我只會敲那扇我曾由此獲得特權的門。\newline{}總之，若獲得、找到、進入是求者、尋者、叩門者勞苦和懇切的果實，\newline{}你就明白這些義務和結果是由造物主規定的。\newline{}至於你那位最卓越的神，他自稱白白幫助非他創造的人，就不能給他施加任何勞苦或懇切。\newline{}因爲他若不自發地賜予不求者、讓不尋者找到、爲不叩門者開門，就不再是最卓越的神了。\newline{}相反，造物主能藉基督宣告這些義務和賞報，以便人——因犯罪得罪了他的天主——在試煉中勞苦，\newline{}並藉著祈禱的恆心而獲得、藉著尋找而找到、藉著敲門而進入。\newline{}因此，前面的比喻描述那夜間去求餅的人，\newline{}是作爲朋友而非陌生人，並使他敲朋友而非陌生人的家門。\newline{}即使人得罪了，他相較於馬爾西翁的神，更是造物主的朋友。\newline{}因此，他敲的是他有權進入之門，是他找到的大門，是他知道擁有餅的那位的門，\newline{}那位現在與祂的兒女同寢，祂曾願意他們誕生。\newline{}即使敲門已晚，仍然是造物主的時間。\newline{}祂擁有整個世代及其終結，故最晚的時刻也屬於祂。\newline{}至於那位新神，沒有人能在晚間敲他的門，因爲他幾乎還未見到晨光。\newline{}正是造物主曾向外邦人關上門，那時猶太人敲門，\newline{}祂就起來給予，即使不是作爲朋友給人，也不是作爲陌生人，而是如祂所說：\Quote{因他的喋喋不休。}\newline{}\BibleRef{路 11:8}\newline{}然而，那位新神在其短暫（出現）的時間內不可能允許任何人喋喋不休。\newline{}因此，你稱之爲造物主的那位，也請承認祂是\Quote{父}。\newline{}正是祂知道祂的兒女需要什麼。\newline{}因爲當他們求餅時，祂從天上賜下瑪納；\newline{}當他們要肉時，祂賜給他們充足的鵪鶉——\newline{}沒有以蛇代魚，也沒有以蠍子代蛋。\newline{}\BibleRef{路 11:11}\newline{}祂既有善也有惡在自己權能中，因此不將惡代替善賜予，乃是相稱的。\newline{}相反，馬爾西翁的神沒有蠍子，就不能拒絕賜予他沒有的東西；\newline{}只有祂——擁有蠍子卻不賜予——才能如此。\newline{}同樣，祂將賜予聖神，祂也命令不潔之神。\newline{}當祂驅逐\Quote{啞巴魔鬼}\BibleRef{路 11:14}（藉此治療應驗了依撒意亞\BibleRef{依 29:18}），\newline{}並被指控藉貝耳則步驅魔時，祂說：\Quote{我若藉貝耳則步驅魔，你們的子弟藉誰驅魔呢？}\BibleRef{路 11:19}\newline{}透過這樣的問題，祂豈不是說祂以同樣的能力——即藉造物主的能力——驅逐邪靈，如同他們的子弟所做的那樣？\newline{}因爲你若以爲意思是\Quote{我若藉貝耳則步⋯⋯你們的子弟呢？}——\newline{}好像祂指責他們擁有貝耳則步的能力——\newline{}你立刻被前面的話所駁斥：\Quote{撒殫不能自相紛爭。}\newline{}\BibleRef{路 11:18}\newline{}因此，即使是他們也不是藉貝耳則步驅魔，而是（如我們所說）藉造物主的能力；\newline{}爲使這一點被理解，祂補充說：\Quote{但我若藉天主的手指驅魔，天主的國豈不已臨近你們了？}\newline{}\BibleRef{路 11:20}\newline{}因爲站在法老面前反對梅瑟的術士稱造物主的能力爲\Quote{天主的手指}。\newline{}\BibleRef{出 8:19}\newline{}那是天主的手指，因爲它是一個記號，表明軟弱之物卻充滿力量。\newline{}基督也表明了這點，當祂回憶（而非抹殺）那些真正屬於祂的古代奇蹟時，\newline{}祂說應理解天主的能力是那手指，而非其他天主的手指，\newline{}正是在那位天主之下，這手指獲得了此稱號。\newline{}因此，祂的國已臨近他們，祂的能力被稱爲祂的\Quote{手指}。\newline{}所以，祂將\Quote{武裝的壯士}的比喻——\newline{}\Quote{一位更強壯者勝過了他}——\newline{}與魔鬼的首領（祂已稱之爲貝耳則步和撒殫）聯繫起來，\newline{}\BibleRef{路 11:21}\newline{}表明是那魔鬼被天主的手指所勝，而非造物主被另一位神所制伏。\newline{}此外，若整個世界都留給祂，馬爾西翁的神可能顯得像\Quote{更強於祂}而勝過了祂，\newline{}祂的國怎會仍然屹立，有其疆界、法律和職能呢？\newline{}若非因祂的法律，連馬爾西翁派也不斷死亡——藉著分解歸於塵土——\newline{}並常常被甚至一隻蠍子提醒：造物主絕未被勝過。\newline{}\Quote{人群中有一個婦人高聲說：\InnerQuote{懷過祢的胎及祢所吮吸過的乳房是有福的！}\newline{}但主說：\InnerQuote{是的，更有福的是聽天主的話並遵守的人。}}\newline{}\BibleRef{路 11:27}\newline{}祂先前曾以幾乎相同的話拒絕了祂的母親或兄弟，而寧可那些聽從天主的人。\newline{}然而，祂的母親當時並不在場。\newline{}因此，在以前的事件中，祂沒有否認祂按出生是她的兒子。\newline{}第二次聽到這問候時，祂第二次如先前一樣將\Quote{福氣}從祂母親的子宮和乳頭轉移到門徒身上，\newline{}然而，除非祂在她身上有一位真實的母親，否則祂無法轉移福氣。\par

% structure: p0060 source=h2 target=subsection reason=relative-heading-rank; base=h2

\subsection{第二十七章 基督斥责法利塞人寻求征兆。祂谴责他们喜爱外表炫耀胜过内在圣洁。圣经充满类似告诫。祂来自创造者的使命的证明。}

我宁愿在其他地方驳斥马西昂派在造物主身上找出的缺点。在此只要指出这些缺点也存在于基督身上就足够了。请看他是多么不均衡、不一致和反复无常！他教导一件事却做另一件事；他命令\BibleQuote{凡求你的，就给他}；而他自己却拒绝给那些\BibleQuote{求看征兆}的人。\BibleRef{路 11:29} 在漫长的世代中，他向人隐藏自己的光，却又说灯不可藏起，而应放在灯台上，使光照亮众人。\BibleRef{路 11:33} 他禁止\Emph{再度}咒骂，当然更禁止多次咒骂；然而他却将\Emph{祸哉}加于法利塞人和法学士身上。有谁比他的基督更像我的天主呢？我们经常已经确凿地指出：如果他宣告另一位神，他就不可能被斥为法律的破坏者。因此，甚至邀请他吃饭的法利塞人，在他面前也显得有些惊奇，因为他没有按照法律先洗手再入席——既然他宣告的是法律的天主。耶稣也向他解释了法律，告诉他说：他们\BibleQuote{洗净杯盘外面，里面却充满劫夺与邪恶}。\Emph{他说这话}，是为了表明：在清理器皿中，要理解为在天主面前的人的洁净，因为这位法利塞人在他面前所处理的是关于一个人，而不是关于一个未洗净的器皿。所以他说：\BibleQuote{你们洗净杯盘的外面}（即肉体），\BibleQuote{但你们不洁净自己的里面}\BibleRef{路 11:39}（即灵魂）；并加上：\BibleQuote{那造外面的（即肉体），不也造了里面的吗？}（即灵魂）——借此他明确宣告：同一个天主掌管人外在和内在的洁净，二者都在他权下，他喜爱仁爱胜于人洗手，甚至胜过祭献。因为他接着命令：\BibleQuote{你们把所拥有的施舍给穷人，那么一切对你们都是洁净的。}\BibleRef{路 11:41} 纵然另一位神可以命令仁爱，但在他被认识之前，他不可能这样做。此外，在此段经文中明显可见，他们受谴责并不是关于他们的神，而是关于他在教导中的某一点，即他比喻性地命令他们洁净器皿，而实际上是要有慈悲的心肠。同样，他责备他们奉献小茐的十分之一，却\BibleQuote{疏忽了款待和爱天主}。\BibleRef{路 11:42} 这是出于哪一位天主的召叫和爱，而不是那一位他们依据其什一奉献法律而奉献芸香和薄荷的天主？因为责备的全部要点在于：他们当然关心他服务中的小事，却忽视了更重要的责任，即他命令的：\BibleQuote{你当全心、全灵、全意爱上主你的天主，是他将你从埃及领出。}\BibleRef{申 6:5} 此外，时间还远远不够，不容许基督要求如此过早——甚至如此令人反感——的爱，去爱一位新的、最近的、甚至几乎尚未显现的神。当他再次责备那些抢坐首席并贪图公开致意荣耀的人时，他只是追随造物主的做法，他称这类野心家为\BibleQuote{索多玛的首领}\BibleRef{依 1:10}，并禁止我们\Quote{依赖王侯}，并说那信赖人的人是全然不幸的。但是，凡是追求高位的人，因他愿以他人的殷勤为荣，既然他禁止这种以寄望和信赖于人的殷勤，他也就同时谴责了所有好高鹜远的人。他也痛斥法学士本身，因为他们\BibleQuote{把沉重难担的重担加在人身上，自己却连一根指头也不肯动}；\BibleRef{路 11:46} 但这并非好像他因厌恶法律而嘲笑法律的重担。因为他如何会厌恶法律呢？他曾如此热切地责备他们疏忽法律更重要的事——施舍、款待和爱天主。的确，他不仅重视这些大事，甚至芸香的十分之一和杯盘的洁净。然而，他倒更愿意认为他们情有可原，因为他们无法承担不能背负的重担。那么他所责备的重担是什么？无非是他们自动添加的，以人的诫命当作道理教导人；为了私利，将房屋连接房屋，以致夺取邻人的房屋；欺哄人民，喜爱礼物，追求酬报，抢夺穷人公正审判的权利，使寡妇成为掠物，使孤儿成为赃物。关于这些人依撒意亚也说：\BibleQuote{祸哉，你们这些在耶路撒冷强横的人！}\BibleRef{依 28:14} 又说：\Quote{那要求你们的将统治你们。}还有谁比法学士更这样做呢？现在，如果他们得罪了基督，他们是因为属于他而得罪他。他不会对外来法律的教师进行打击。可是为什么\Quote{祸哉}被宣布给他们，因为他们\BibleQuote{建造先知们的坟墓，而先知正是被他们的祖先杀害的}？\BibleRef{路 11:47} 他们倒应受称赞，因为通过这种敬虔的行为，他们似乎表明自己不认同祖先的行为。这不是因为（基督）嫉恨这种态度吗？正如马西昂派所谴责的，他常常将父亲的罪过追罚到子孙直到三四代？这些法学士究竟有什么\BibleQuote{钥匙}\BibleRef{路 11:52}，除了法律的解释？他们自己既不进入对法律的领悟，因为不信（因为\Quote{除非你们相信，你们不会领悟}）；也不允许别人进入，因为他们宁愿以人的诫命当作道理教导人。所以，当他责备那些自己不进入并阻挡别人进入的人时，他应被视为法律的贬低者还是支持者？如果是贬低者，那些阻碍法律的人应该高兴；如果是支持者，他就不再是法律的敌人。但他发出所有这些诅咒，是为了玷污造物主为残酷的存在，得罪他的人注定要承受\Quote{祸哉}。谁不更加害怕触怒一位残酷的存在，因而背离对他的忠诚？因此，他越是描述造物主为可畏惧的，就越热切教导应当事奉他。这就是造物主基督所应当行的。\par

% structure: p0062 source=h2 target=subsection reason=relative-heading-rank; base=h2

\subsection{第二十八章 旧约中的例子：巴郎、梅瑟与希则克雅，显示基督的教导与行为完全符合造物主的旨意与计划。}

因此，法利塞人的虚伪令祂不悦是公义的，他们用嘴唇爱天主，心却远离祂。祂对门徒们说：\Quote{你们要谨慎，防备法利塞人的酵母，就是虚伪。}这不是对创造者的宣告。圣子恨恶那些拒绝服从圣父的人；祂也不愿门徒对祂（注意：不是对另一位神，因为对那一位神，这样的虚伪本可允许）表现出那种祂要门徒提防的态度。祂所禁止的是法利塞人的榜样。基督如今禁止门徒冒犯的，正是法利塞人所得罪的那一位。既然祂谴责了他们的虚伪——这种虚伪掩盖了内心的秘密，用表面的行为遮蔽了不信的奥秘，因为他们（拿着知识的钥匙）自己既不进去，也不允许别人进去——祂于是加上说：\Quote{没有掩盖的事，将来不被揭露的；也没有隐藏的事，将来不被知道的。}\BibleRef{路 12:2}这是为了不让任何人以为祂是在揭示和承认一位迄今未知且隐藏的神。当祂又论及他们的怨言和嘲讽，说祂\Quote{这个人只是靠贝耳则步驱魔}时，祂的意思是所有这些指控都将暴露在光天化日之下，并因福音的宣讲而显在人前。祂随后转向门徒，说了这些话：\Quote{我告诉你们，我的朋友，你们不要害怕那些只能杀害肉身，此后不能再做什么的。}\BibleRef{路 12:4}但他们将发现依撒意亚早已说过：\Quote{看，义人怎样被除去，没有人放在心上。}\BibleRef{依 57:1}\Quote{但我指示你们当怕的是谁：当怕那杀害以后，又有权柄投入地狱的}（当然，是指创造者）；\Quote{是的，我告诉你们，当怕祂。}\BibleRef{路 12:5}此时，只要祂禁止冒犯那祂命令人敬畏的，并且祂命令人尊敬那祂禁止冒犯的，而发出这些命令的祂正属于那为此恐惧、无冒犯、尊敬所同一位的\Emph{天主}，这便足以满足我的目的。但我也可以从以下的话得出这结论：\Quote{我告诉你们，凡在人面前承认我的，人子也要在天主面前承认他。}\BibleRef{路 12:8}那些承认基督的人将被人杀害，但被他们处死后便再无可受。因此，这些人就是祂以上所警告不要只害怕被杀的那些，祂发出这警告，是为了加上一句关于必须承认祂的补充：\Quote{凡在人面前否认我的，将在天主面前被否认。}\BibleRef{路 12:9}——当然是被那原本会承认他、若他只承认\Emph{天主}的那一位所否认。那承认承认者的，正是那也将否认否认祂自己者的同一天主。再者，如果承认者在惨死后将无所畏惧，那么否认者在自然死亡后一切都会变得可畏。因此，既然死后必须畏惧的事，即地狱的惩罚，属于创造者，那么否认者也属于创造者。然而，正如否认者一样，承认者若否认天主，显然必将受天主之苦，尽管在人处死他之后，他再无苦可受。所以基督属于创造者，因为祂表明所有否认祂的人都当畏惧创造者的地狱。在阻止\Emph{门徒}否认自己之后，祂又加上一句关于避免亵渎的劝诫：\Quote{凡说话干犯人子的，可得赦免；但亵渎圣神的，必不得赦免。}\BibleRef{路 12:10}既然罪的赦免与保留都显示一位审判的天主，那么不可亵渎的圣神便属于那不会赦免亵渎的那一位；正如在前文中，那不可否认的圣子属于那在杀害之后也将投入地狱的那一位一样。如今，既然基督使亵渎远离创造者，我不知道祂的对手怎能出现。否则，如果祂用这些话对那不会赦免亵渎并将杀至地狱者的严厉投下了一团指责的黑云，那么那敌对之神的灵便可以被亵渎而不受罚，他的基督也可以被否认；事实上，敬拜他和藐视他之间并无区别；既然蔑视没有惩罚，敬拜也没有奖赏，人无需期待。当\Quote{被带到官长面前}受审时，祂禁止他们\Quote{思虑怎样回答}；\Quote{因为，}祂说，\Quote{圣神要在那时辰教导你们应当说什么。}\BibleRef{路 12:11}如果这样的训令来自创造者，那么这诫命只能属于那位曾有先例的一位。在《户籍纪》中，先知巴郎被巴拉克王派去咒骂以色列（那时王正与以色列开战），他立时充满圣神。他原要来宣布咒骂，却说了圣神在那时默感他的祝福；他曾预先对王的使者，然后对王本人声明，他只能说天主放在他口中的话。\EditorialNote{户籍纪第二十二至二十四章}这新基督的新教义，创造者的仆人在很久以前就已开创！但请看梅瑟与基督的榜样有多么明显的区别。梅瑟自愿介入争吵的弟兄之间，责备那犯罪者：\Quote{你为什么打你的同伴？}然而，他被那人拒绝：\Quote{谁立了你作我们的首领和审判官？}\BibleRef{出 2:13}相反，基督被一个人请求调停他和他弟兄关于分家产的争吵时，拒绝了帮助，尽管是这么一件正直的事。那么，我的梅瑟比你的基督更好，因为他致力于弟兄的和平，并阻止了他们的不义。但\Emph{当然基督的情况必须不同}，因为他是那全然良善、不审判之神的那基督。\Quote{谁立了我作你们的审判官呢？}祂说。\BibleRef{路 12:13}再也找不到别的推托之词，只能用那邪恶之人和不敬弟兄者拒绝正直与虔诚的辩护者的话！总之，祂认可了这托词（尽管不好）通过使用它；也认可了那行为（尽管不好）通过拒绝调和弟兄。或者，祂岂非是要表达对梅瑟被那样的话拒绝的愤怒？因此，祂岂非不愿意在类似争闹弟兄的情况下，用那严厉的话的回忆来使他们困惑？显然是的。因为祂自己曾在梅瑟中临在，听到了那样的拒绝——祂本身就是创造者的圣神。我想我们已在另一处充分证明，财富的荣耀被我们的天主所谴责，祂\Quote{从高位上推下有权势者，从粪土中举扬贫穷者。}因此，出自祂的将是富人的比喻，那人因田产丰收而自夸，天主对他说：\Quote{糊涂人哪，今夜就要索回你的灵魂，你所预备的，将归谁呢？}\BibleRef{路 12:16}同样，\Person{希则克雅}王在向巴比伦使者夸耀自己的宝物和珍贵物品之后，也听到依撒意亚关于他王国悲惨结局的宣告。\EditorialNote{依撒意亚第三十九章}\par

% structure: p0064 source=h2 target=subsection reason=relative-heading-rank; base=h2

\subsection{从先知而来的类比，用以说明基督在\BibleBook{圣路加}{福音}本章其余部分中的教导。基督更严厉的属性，在其审判职分中，表明他来自创造者。附带责备了\Person{马西昂}的独身主义，以及他对福音经文的篡改。}

谁不愿意我们不该为生命的食粮、身体衣裳忧虑（\BibleRef{路 12:22}），不正是在那早已为人预备了这些的那一位，并因此将之分施给我们时，禁止我们对此一切忧虑，以为这是对他慷慨的侮辱？——那一位使\Quote{生命}本身本性成为\Quote{胜过食物}的状态，使\Quote{身体}的质料形成\Quote{胜过衣裳}；祂的\Quote{乌鸦也不种也不收，没有仓房，尚且被}祂\Quote{养活}；祂的\Quote{百合花和草也不劳苦，也不纺织，尚且被}祂\Quote{穿戴}；并且祂的\Quote{\Person{所罗门}极荣华的时候，尚且不如}那卑微的花那样\Quote{装扮}。此外，再没有比以下更突兀的：一位天主在分施祂的恩惠，而另一位却命令我们不要思考（如此仁慈的）分施——并且还意图贬抑（他的慷慨）。究竟他是否因轻视造物主而不愿这些小事被思念，因为这些事连乌鸦和百合花都不劳苦，因为，也许，由于它们的无价值反而自动到手，稍后会看得更清楚。同时，他怎责备他们为\Quote{小信德的人}呢？\BibleRef{路 12:28} 什么信德？他是指那信德，即他们尚不能完全显示在一个几乎尚未启示、他们仅尽所能学习的天主之上的信德？还是指那信德，即他们因此确切理由亏欠造物主的，因为他们相信祂自愿供应人类这些需要，因而不思考这些？当他添加说：\Quote{这一切都是世上外邦人所寻求的}\BibleRef{路 12:30}，甚至因他们不信天主是造物主和万物的赐予者，既然他不愿他们像这些外邦人，他便责备他们缺乏对同一天主的信德，他注意到外邦人对此完全缺乏信德。当他进一步说：\Quote{你们的圣父知道你们需要这些}\BibleRef{路 12:30}，我须先问：基督这里要理解为哪一位父？若他指向他们自己的造物主，他也肯定祂是善的，知道祂的儿女需要什么；但若他指另一位神，他如何知道食物衣裳是人必需的，既然他未曾为此预备？因为若他知道那需要，他必已预备。然而，若他知道人需要什么，却未能供给，他在这失败中或怀恶意或孱弱。但当他承认这些为人是必需时，他实际上肯定它们是\Quote{善的}。因为无邪恶之事是必需的。这样，他不会再是造物主作品和纵容的贬抑者，好使我能在此完成我上面延迟给予的回答。再者，若是另一位神预见了人的需要并供应，那么\Person{马西昂}的基督自己怎会应许这些？\BibleRef{路 12:31} 他慷他人之慨吗？他说：\Quote{你们寻求天主的国，这一切都要加给你们了}——当然，由他自己。但若\Quote{由他自己}，他是怎样的存在，竟能赐予别物？若\Quote{由造物主}，万物都属于祂，那么那应许属他之物的是谁？若这些是\Quote{加给}那国，它们应被置于第二等；第二等属于那亦拥有第一等者；食物和衣裳属于那国度也属于祂。因此，整个应许、其比喻的完整实在、其类比之完美均等属于造物主；因为这些除了指向那每一点都有相似关系者外，不指向任何他者。我们是仆人，因为我们有主在我们的天主内。我们应当\Quote{束上腰}\BibleRef{路 12:35}，即，摆脱困惑忙碌生活的羁绊；\Quote{点着灯}\BibleRef{路 12:35}，即，心灵被信德点燃，以真理的作为发光。这样\Quote{等候我们的主}\BibleRef{路 12:36}，即基督。从何处\Quote{回来}？若\Quote{从婚筵}，祂是造物主的基督，因为婚筵属于祂。若祂不属于造物主，连\Person{马西昂}自己也不会赴席，即使被邀请，因为他在他的神中发现一位憎恶婚床者。因此比喻在主身上失当，若祂不是一位婚筵与之相宜的存在。在下一个比喻中，他也犯下显著错误，当他将\Quote{窃贼}的角色赋予造物主——\Quote{家主若知道什么时辰有贼来，必不容许房屋被挖透}\BibleRef{路 12:39}。造物主怎能以任何方式带有窃贼面貌，祂是全人类的主？无人偷窃或抢劫自己的财物，但更似那攫取他人之物、使人背离其主者。再者，当他向我们指明魔鬼是\Quote{窃贼}，在起初世界时，人若知道他的时辰，就永不会被他侵入，他警诫我们\Quote{要预备}，因为\Quote{人子来的时辰你们不知道}\BibleRef{路 11:40}——并非祂自己是那贼，而是作为那审判不预备自己、不防贼之人的审判者。既然祂是人子，我视祂为审判者，并在审判者中我认出造物主。那么若在此处他以\Quote{人子}之称号显示造物主的基督，好给我们一些关于那来临时辰未知之贼的预兆，你仍已在上文定规：无人是己物之贼；此外，我们的原则也完整不变——即祂极力强调造物主为可敬畏者，祂就越属于造物主，行造物主之工。因此，当\Person{伯多禄}问祂说：这比喻是为\Quote{他们}说的，还是\Quote{也为众人}\BibleRef{路 12:41}，祂便为他们，也为所有应在教会中掌权者，设立了管家的比喻。那在他主人不在时善待同仆的管家，主人回来要立他管理一切所有的；但那行事相反的，要被腰斩，安排在不信者中，当主人回来之日，在他不期待的时候、不知不觉的时候\BibleRef{路 12:41}——即那人子，造物主的基督，并非贼，而是审判者。因此，在此段落中，祂或将主呈现为审判者，教训我们祂的性格，或为纯粹良善之神；若是后者，他现在也肯定审判属性，尽管异端者拒绝承认。因为有人试图缓和此意，将其应用于他的神——仿佛单纯地腰斩那人，将他安排在不信者中，是宁静温和之举，基于一种观念，以为他未被传唤（到审判者前），只回到他自己的状态！仿佛这过程本身不包含审判行为！何其愚拙！那被腰斩者的结局是什么？岂不就是丧失救恩，因为他们的分离是从那些将要获得救恩者？那不信者的境况又是什么？岂不就是定罪？否则，若这些被腰斩者和不信者将无所受苦，另一方面，被接纳者和信者也就无所获得。然而，若被接纳者和信者将获得救恩，被拒绝者和不信者必然遭受相反结局，即救恩的丧失。如今这是一个审判，而那在我们面前呈现它的属于造物主。除了报应之神，我还能从祂将\Quote{用鞭子责打仆人，或重或轻}并向他们追讨所交托之物理解为谁？服从谁对我合适，除非是那位酬报者？你的基督宣告：\Quote{我来把火投在地上。}\BibleRef{路 12:49} 那位最温和的存在，没有地狱的主，不久前曾阻止他的门徒要求火降在那无礼的村庄。而祂却用火暴焚烧了\Place{索多玛}和\Place{哈摩辣}。关于祂，圣咏作者咏唱：\Quote{火在祂面前行，烧灭祂四围的敌人。} 祂通过\Person{欧瑟亚}发出威胁：\Quote{我必降火在犹大的城邑}\BibleRef{欧 8:14}；通过\Person{依撒意亚}：\Quote{怒火在我中烧起。} 祂不能撒谎。若这不是那从燃烧荆棘中发声的那位，那无论你坚持理解为何种火都无关紧要。即使它是喻像之火，然而从他从我元素中取譬喻来表达自己的意思，祂属于我，因祂使用我的。火的比喻必属于拥有其实际的那位。但祂将亲自解释祂所提到之火的品质，当祂接着说：\Quote{你们以为我来是给地上带来和平吗？我告诉你们，不是，而是分裂。}\BibleRef{路 12:51} 经上写道\Quote{刀剑}，但\Person{马西昂}修改了词，仿佛\Quote{分裂}不是\Quote{刀剑}之工。因此，那拒绝赐予和平的，也意图毁灭之火。正如争战，也如焚烧。正如刀剑，也如火焰。两者对其主皆不适宜。最后祂说：\Quote{父亲将分裂敌对儿子，儿子敌对父亲；母亲敌对女儿，女儿敌对母亲；婆婆敌对儿媳，儿媳敌对婆婆。}\BibleRef{路 12:53} 既然这亲属间的争战正是先知号角所预言的话，恐怕\Person{米该亚}\BibleRef{米 7:6}必定已为\Person{马西昂}的基督预言之！为此祂称他们为\Quote{伪善者}，因他们能\Quote{分辨天地气象，却不能分辨这个时机}\BibleRef{路 12:56}，此时祂当然应被认出，成全一切关于他们预言的，并教训他们。但那时谁能认识一位其存在无证据可证之人的时候呢？祂也公正地斥责他们\Quote{连自己也不能判断什么是正义的}\BibleRef{路 12:57}。古时祂通过\Person{匝加利亚}命令：\Quote{执行真理与和平的审判}\BibleRef{匝 8:16}；通过\Person{耶肋米亚}：\Quote{执行公道和正义}\BibleRef{耶 22:3}；通过\Person{依撒意亚}：\Quote{审判孤儿，为寡妇辩护}\BibleRef{依 1:17}，把这当作索勒赫葡萄园的过错，因为\Quote{祂指望结好葡萄，反倒结了野葡萄}（压迫的呼喊）。那曾教训他们遵行祂命令的天主，现在要求他们自发而行。那播种诫命的，现在催迫丰收。但何等荒谬，祂现在命令他们公义审判，却消灭公义审判的天主！因为那把人下在监里，不还清\Quote{最后一个小钱}不放他出来的审判者\BibleRef{路 12:58}，他们以造物主身份论之，以贬抑祂。然而，我认为有必要用同样回答来应对这吹毛求疵。因为每当造物主的严厉被展示在我们面前，基督（被证明）正属于祂，祂以恐惧迫使人服从。\par

% structure: p0066 source=h2 target=subsection reason=relative-heading-rank; base=h2

\subsection{第三十章。芥子与酵母的比喻。转向家主关门后随之而来的庄严排除。这审判性的排除将由基督执行，由此显明祂拥有造物主的属性。}

当再次就安息日所行的医治提出疑问时，祂是如何讨论的？\BibleQuote{你们中谁在安息日不把自己的驴或牛从槽里解开，牵去饮水呢？}\BibleRef{路 13:15} 因此，当祂按法律所规定的条件行事时，祂是肯定了法律，而不是违背法律；法律命令除了为任何生物可行之事以外，不可做工；\Emph{若为任何一人}，何况是为\Emph{人的}生命呢？在比喻的事上，我处处要求一致性是可以允许的。祂说：\Quote{天主的国好像一粒芥子，一个人拿去种在自己的园子里。}那\Emph{人}应被理解为谁？无疑是基督，因为（尽管是马西昂的）祂被称为\Quote{人子}。祂从圣父领受了天国的种子，即福音之言，并把它撒在自己的园子里——当然是在世界中——例如在今人身上。然而，经上说\Quote{\Emph{在他的园子里}}，但世界和人都不属于祂，而属于造物主；因此，那撒种在自己土地上的就是造物主。反之，如果为回避这一陷阱，他们宁愿将\Emph{那人的}位格从基督转移到任何接受天国种子并撒在自己心园的人身上，即使这个含义也除了造物主之外不适合任何别的。因为，如果天国属于最宽容的神，它紧接着就被一个炽热的审判所跟随，其严厉导致哀号，这怎么可能呢？关于接下来的比喻，我实在担心它可能预示了敌对之神的国！因为祂将之比作，不是造物主更熟悉的无酵饼，而是\Emph{酵母}。\BibleRef{路 13:20} 对于寻求论点的人而言，这是一个极好的猜想。然而，我必须从自己这一面，用一个钟爱的概念驱散另一个，并主张甚至\Emph{酵母}也适合造物主的国，因为它之后是\Emph{烤炉}，或者如果你愿意，是地狱的火炉。祂已经多少次显明自己是一位审判者，并在审判者中显明造物主？事实上，祂多少次拒绝，并在拒绝中定罪？在目前的经文中，祂说：\BibleQuote{一旦家主起来，}\BibleRef{路 13:25} 但除了依撒意亚所说的\BibleQuote{当祂起来震撼大地时？}\BibleRef{依 2:19} 之外，还有什么意思呢？\BibleQuote{并关上了门，}当然是将恶人关在外面；当他们敲门时，祂要回答：\BibleQuote{我不认识你们是从哪里来的；}当他们述说如何在祂面前\BibleQuote{吃过喝过}时，祂还要对他们说：\BibleQuote{你们这些作恶的人，离开我去吧！在那里要有哀号和切齿。}\BibleRef{路 13:25} 但在哪里？毫无疑问在外面，当他们被祂关在外面并关上门之后。因此，那为了惩罚而排除的人将施以惩罚，当他们看见义人进入天主的国，而自己却被拘留在外。由谁拘留在外？如果是造物主，那么谁会在\Emph{里面}迎接义人进入天国？是良善神。那么，造物主在做什么呢？被祂的对手关在外面的人，祂却拘留他们受罚，而实际上，如果那些人必须落入祂手中，为更大大激怒祂的对手，祂本应仁慈地接纳他们。因此，要么恶人会\Emph{被造物主}违反排除者的意志而拘留，在这种情况下，排除者低于造物主，不情愿地服从祂；要么，如果这个过程是按祂的意志进行的，那么祂自己就依法决定了其执行；那么，这位造物主恶名的始作俑者，将不会证明自己比造物主好一分一毫。如今，如果这些观念与理性不相容——一个被认为惩罚，另一个释放——那么审判和国度将属于同一权能；既然两者都属于同一者，执行审判的只能是造物主的\Emph{基督}。\par

% structure: p0068 source=h2 target=subsection reason=relative-heading-rank; base=h2

\subsection{第三十一章。基督按照依撒意亚邀请穷人的劝诫。大宴席的比喻是造物主自己仁慈与恩宠安排的图景。邀请被拒与旧约引文相对应。马西昂的基督无法满足这比喻所指示的条件。马西昂派解释的荒谬。}

祂吩咐何人應當被邀請赴午宴或晚宴？\BibleRef{路 14:12} 正是指出祂曾藉依撒意亞所指明的那樣：\BibleQuote{你把你的糧食分給饑餓的人，並且把那些乞丐——甚至無家可歸的人——領到你的家裏來，}\BibleRef{依 58:7} 因爲他們無疑\Quote{不能報答}你的人道行爲。既然基督禁止現今期望報答，而應許\Quote{在復活時}報答，這正是造物主的計劃，祂不喜愛那些貪愛禮物、追求報酬的人。也請思想，那個發出邀請之人的譬喻更適合哪一位神：\BibleQuote{有一個人設擺了盛大的晚宴，請了許多人。}\BibleRef{路 14:16} 晚宴的預備無疑是永生豐盛預備的預象。我首先指出，通常陌生人以及與邀請者無親屬關係的人不被邀請赴晚宴；而是家中和家族成員更常是受邀的貴賓。那麼，發出邀請本屬於造物主，因爲那些將被邀請的人也屬於祂——無論是作爲\Emph{人}，因出自亞當，還是作爲\Emph{猶太人}，因他們的祖先；而不屬於那位既無本性權利也無特權擁有他們的神。其次我要說，如果預備晚宴的那位發出邀請，那麼按此意義，晚宴是造物主的，祂派人去提醒客人。這些人先前確實曾被祖先邀請，但需要藉先知來勸誡。\Emph{當然這絕不是那位的宴席}，祂從未派過使者去提醒——在發出邀請之前從未做過任何事，而是突然親自降臨——只在開始邀請時才被認識；只在強迫人赴宴時才邀請；爲晚宴和邀請指定同一個時刻。但當被邀請時，他們推辭了。\BibleRef{路 14:18} 這很合理，如果邀請來自另一位神，因爲它如此突然；如果推辭不合理，那麼邀請就不是突然的。如果邀請不是突然的，那麼它必定是由造物主發出的——甚至由祂自古老時代，他們最終拒絕了祂的召叫。他們第一次拒絕是當他們對亞郎說：\BibleQuote{爲我們造神，在我們前面走}；\BibleRef{出 32:1} 後來又當\BibleQuote{他們耳朵聽見，卻不明白}\BibleRef{依 6:10} 天主的召叫。祂以最貼切此譬喻的方式藉耶肋米亞說：\BibleQuote{聽從我的聲音，我要作你們的天主，你們要作我的子民；並要遵行我所吩咐你們的一切道路。}\BibleRef{耶 7:23} 這是天主的邀請。\Quote{但是，}祂說，\BibleQuote{他們不聽從，也不側耳。}\BibleRef{耶 7:24} 這是百姓的拒絕。\BibleQuote{他們離去，各隨自己邪惡心思的意念而行。}\BibleRef{耶 11:8} \BibleQuote{我買了一塊田——我買了幾對牛——我娶了一個妻子。}\BibleRef{路 14:18} 祂仍催促他們：\Quote{我打發了我所有的僕人先知到你們那裏去，清早起來打發。}這裏指的是聖神，客人的提醒者。\BibleQuote{然而我的子民不聽從我，也不側耳，卻硬着頸項。}\BibleRef{耶 7:26} 這事被報告給家主。於是他被觸動（他被觸動是對的；因爲馬西昂否認他的神有情感，所以他必然是我的天主），並吩咐去邀請\BibleQuote{城裏大街小巷}的人。\BibleRef{路 14:21} 讓我們看這是否不是與祂藉耶肋米亞所說的話含義相同：\BibleQuote{我對以色列家豈是曠野，或是未耕之地呢？}\BibleRef{耶 2:31} 這就是說：\Quote{那麼我沒有可以呼召到我這裏來的人嗎？我沒有可以領他們來的地方嗎？}\BibleQuote{既然我的子民說，我們不再到你那裏去。}\BibleRef{耶 2:31} 因此祂派人去召叫別人，但來自同一座城。\BibleRef{路 14:23} 我的第三點是：儘管那個地方人很多，祂仍命令他們從大道和籬笆那裏聚集人。換句話說，我們現在是從外邦陌生人中被聚集；無疑帶着那種嫉妒的怨恨，正如祂在申命記中所表達的：\BibleQuote{我要掩面不顧他們，並要向他們顯明在末日將要發生的事（即別人將佔據他們的位置）；因爲他們是乖僻的一代，是沒有信德的兒女。他們以那不是神的東西惹動我的嫉妒，以他們的偶像觸動我的忿怒；我要以那不是子民的民惹動他們的嫉妒，以愚昧的國民觸動他們的怒氣}\BibleRef{申 32:20} ——甚至以我們，猶太人仍然懷着希望的人。但主說他們不會實現這希望；\BibleQuote{熙雍被撇下，如同葡萄園裏的茅舍，如同瓜田中的草棚}，\BibleRef{依 1:8} 因爲那民族拒絕了對基督的最後邀請。（現在，我問）在遍覽了造物主這整個安排和預言之後，其中有什麼能歸給那匆忙完成一切工作、既沒有造物主的進程也沒有與譬喻和諧安排的那位呢？或者，他的第一次邀請將是什麼，第二階段的提醒又是什麼？有些人起初肯定會拒絕；另一些人後來必須接受。但現在他來邀請兩類人混雜在一起，從城裏、從籬笆那裏，與譬喻的意圖相反。他現在不可能定那些從未受他邀請、他如此熱切接近的人爲輕視他的邀請。如果他預先定他們爲將拒絕他的召叫的人，那麼他也預先預言了外邦人將取代他們被揀選。當然他第二次來是爲了向外邦人宣講。但即使他打算再來，我想他不會是爲了再邀請客人，而是爲了給他們席位。與此同時，你將此次晚宴的邀請解釋爲屬靈飽足和喜樂的天上盛筵的邀請，必須記住，酒、油、糧食，甚至城市的屬地應許也被造物主用作屬靈事物的預象。\par

% structure: p0070 source=h2 target=subsection reason=relative-heading-rank; base=h2

\subsection{第三十二章　一种逻辑学家所谓的连锁论，证明失羊和失银币的比喻与马西昂的基督毫无适当关系}

谁寻找了迷失的羊和丢失的银币？\newline{}\BibleRef{路 15:1}\newline{}岂不是丢失者么？但谁是丢失者？岂不是曾经拥有它们的人么？那么，那是谁？岂不是它们所属的那一位么？既然\Emph{人}除了是造物主的财产以外，不可能是任何其他者的财产，那么那一位拥有他便也拥有那拥有者；那一位失去了他，因为曾经拥有他；那一位寻找了他，因为失去了他；那一位找到了他，因为寻找了他；那一位喜悦了，因为找到了他。因此，这两个比喻的用意与被比喻的那一位毫无关系，因为羊和银币都不属于他——也就是说，不属于\Emph{人}。因为他没有失去他，因为他未曾拥有他；他没有寻找他，因为他未曾失去他；他没有找到他，因为他未曾寻找他；他没有喜悦，因为他没有找到他。因此，为罪人的悔改——即失去的人被寻回——而喜悦，是属于那一位很久以前就宣告他宁愿罪人悔改而不死的那一位的属性。\par

% structure: p0072 source=h2 target=subsection reason=relative-heading-rank; base=h2

\subsection{第三十三章　驳斥马西昂派对天主和玛门的解释。先知们为基督反对贪婪和骄傲的辩护。洗者若翰是造物主新旧两约之间的纽带。基督这样说——但依撒意亚也早已这样说过。只有一位天主，即造物主，按自己的旨意改变了诸时代。没有新神参与这一改变。}

祂说不能事奉的两个主人，\BibleRef{路 16:13} 理由是当一人喜悦时，另一人必然不悦，祂自己清楚指出，当祂提到天主和玛门时。那么，如果你身边没有诠释者，你可以再次从祂自己学习祂对\Emph{玛门}的理解。因为当祂劝我们用世上的朋友帮助自己，以那管家的榜样为例，他被撤职后，减轻他主人的债务人的债务，以期他们用帮助回报他，祂说：\Quote{我告诉你们：要用不义的玛门结交朋友，}也就是说，用钱财，正如那管家所做。现在我们都知道，钱财是不义的煽动者，是全世界的主宰。因此，当祂看见法利塞人的贪财，对钱财进行奴性的崇拜，祂就向他们抛出这句话：\Quote{你们不能事奉天主而又事奉玛门。} \BibleRef{路 16:13} 那时，贪爱钱财的法利塞人听见祂用玛门指钱财，就嘲笑祂。不要有人以为玛门一词是指造物主，并且基督叫他们不事奉造物主。何等的\Emph{愚蠢}！反而应从中学习，基督指出的是唯一的天主。因为祂提到了两个主人：天主和玛门——造物主和钱财。你们确实不能事奉天主——当然就是他们似乎事奉的那位——而又事奉玛门，他们却更愿意献身于玛门。然而，如果祂自称是另一位\Emph{神}，那祂指出的就不是两个主人，而是三个。因为造物主是主人，而且肯定比玛门更是主人，更应受敬拜，因为祂才是我们真正的主人。那么，那称玛门为主人、并将它与天主并列的祂，怎会对那真正是这些的主人，即造物主，只字不提呢？或者，祂保持沉默，是否就承认可以事奉\Emph{祂}，因为祂只说不能同时事奉祂自己和玛门？因此，当祂定下天主是独一的立场时，因为如果祂自己是造物主的对手，祂肯定会提到造物主，所以实际上祂没有坚持自己是唯一的主、没有敌对的神，就等于提到了造物主。因此，这将阐明以下经文的意义：\Quote{你们在不义的玛门上不忠信，谁还把真实的财富委托给你们呢？} \BibleRef{路 16:11} \Quote{在不义的玛门上}，也就是说，在不义的财富上，而不是在造物主上；因为连马西昂也承认造物主是公义的：\Quote{如果你们在别人的财物上不忠信，谁还把属于我的交给你们呢？} 因为凡是不义的，都应当与天主的仆人无关。但造物主对法利塞人怎么可能是外人呢？因为祂本是犹太民族自己的天主。既然这些话，\Quote{谁会把更真实的财富托付给你们呢？}和\Quote{谁会把属于我的交给你们呢？}只适合造物主，而不适合玛门，祂就不可能以与造物主无关、有利于敌对神的方式说出它们。祂似乎只有在这种意义上说了它们，即当祂指出他们对造物主（而不是对玛门）不忠信时，祂在造物主（祂提到祂的方式）和敌对神之间作了一些区别——后者不会把自己的真理托付给对造物主不忠信的人。那么，如果祂不被明确地表达为有别于所讨论的对象，祂怎么能看起来属于另一位神呢？但当法利塞人\Quote{在人面前自称为义}，\BibleRef{路 16:15} 并把得赏报的希望寄托于人时，祂就按耶肋米亚先知的话责备他们：\Quote{信赖人的人，是可诅咒的。} \BibleRef{耶 17:5} 因为先知继续说：\Quote{但是上主知道你们的心，} 祂赞扬那自称像灯、\Quote{洞察肺腑和心肠}的天主的权能。\BibleRef{耶 20:12} 当祂用以下言语打击骄傲时：\Quote{在人前是崇高的，在天主前却是可憎的，} \BibleRef{路 16:15} 祂引用了依撒意亚的话：\Quote{因为万军上主的日子要临到一切骄傲自大的人，一切高傲狂妄的人，他们也必降为卑下。} 我现在能明白为何马西昂的神被隐藏了如此长久。我想，祂是在等从造物主那里学会这一切。祂一直学习到若翰的时候，然后立刻开始宣讲天主的国，说：\Quote{法律及先知直到若翰为止；从此天主的国就被传扬。} \BibleRef{路 16:16} 就好像我们也不承认在若翰身上，旧制度和新制度之间有一个界限，犹太教在此结束，基督教开始——然而，并不认为法律和先知的中止以及那包含天主之国的福音（即基督自己）的开始是因另一位神的能力。因为，虽然正如我们所展示的，造物主预言旧事会过去，新事会到来，但是，既然若翰被证明是那位要带来福音、宣讲天主之国的主的先驱和预备道路者，那么从若翰已经来了这一事实就得出结论：基督必须是那位要追随其先驱若翰的那一位。所以，若旧历程已终止，新历程已开始，以若翰为两者间的中介，就没有什么可惊奇的，因为这是按造物主的旨意发生的；因此，你从任何来源，即使是不寻常的来源，获得关于天主之国的更好证据，也比从那个妄想（即法律和先知在若翰结束，新事在他之后开始）要好。\Quote{因此，天地——以及法律和先知——更容易过去，而主的话一点一划也不能落空。} \Quote{因为，}正如依撒意亚所说：\Quote{我们天主的话永远常存。} \BibleRef{依 40:8} 因为即使\Emph{那时}，藉着依撒意亚，也是基督——造物主的圣言和圣神——预言性地描述若翰是\Quote{在旷野里呼喊者的声音，预备主的道路，} \BibleRef{依 40:3} 并且他将要来，从此终止法律和先知的历程，为的是成全而非废除；为了使基督能宣讲天主的国，祂特意加上保证：天地更容易过去，而祂的话不能落空；祂也证实了另一事实：祂关于若翰所说的话没有落空。\par

% structure: p0074 source=h2 target=subsection reason=relative-heading-rank; base=h2

\subsection{第三十四章：\Person{梅瑟}允许离婚而\Person{基督}禁止，加以解释。\Person{洗者若翰}与\Person{黑落德}。\Person{马西昂}试图在富翁与穷人于\Place{阴府}中的比喻中发现\WorkTitle{反题}，被驳斥。创造者的安排在两种状态中显明。}

但是\Person{基督}禁止离婚，说：\BibleQuote{谁休妻另娶，就是犯奸淫；谁娶被休的妇人，也是犯奸淫。}\BibleRef{路 16:18} 为了禁止离婚，祂使娶被休的妇人成为非法。然而，\Person{梅瑟}在\BibleBook{}{申命纪}中准许休妻：\BibleQuote{如果一人娶了妻，占有她后，在她身上发现了什么难堪的事，因而不喜悦她，便给她写了休书，交在她手中，叫她离开他的家。}\BibleRef{申 24:1} 因此，你看出法律与福音之间、\Person{梅瑟}与\Person{基督}之间有区别？确实有！但你也拒绝了那为同一真理和同一\Person{基督}作证的别的福音。在那里，禁止离婚的同时，祂给了我们关于这个特别问题的解决方法：祂说：\BibleQuote{\Person{梅瑟}为了你们的心硬，准许你们休妻；但起初并不是这样。}\BibleRef{玛 19:8}——这确实是因为那\BibleQuote{造了他们一男一女}\BibleRef{玛 19:4}的那位也说过：\BibleQuote{两人成为一体；所以天主所结合的，人不可拆散。}\BibleRef{玛 19:6} 现在，通过这个答复（对法利塞人），祂既批准了\Person{梅瑟}的规定（\Person{梅瑟}是祂自己的仆人），又恢复了造物主制度的原始目的，祂就是那造物主的\Person{基督}。然而，既然你要从你所接受的圣经中被驳倒，我就按你的立场来与你交锋，仿佛你的\Person{基督}是我的\Person{基督}一样。因此，当祂禁止离婚，同时又代表那使男女结合的圣父时，祂岂非更应为\Person{梅瑟}的法律开脱而不是废除它吗？但是，请注意，如果这个\Person{基督}在教导与\Person{梅瑟}和造物主相反的事时是你的\Person{基督}，那么，如果我能证明祂的教导并不与他们相反，按同一原则祂必是我的\Person{基督}。那么，我主张，在祂现在所作的关于离婚的禁止中，有一个条件；所假定的是，一个男人休妻的明确目的是为了另娶。祂的话是：\BibleQuote{凡休妻另娶的，就是犯奸淫；凡娶被休的妇人的，也是犯奸淫。}\BibleRef{路 16:18}——\Quote{休了}，也就是指那种不该休妻的原因，即为了另娶。因为娶一个被非法休弃的妇人的人，与娶一个未被休离的妇人的人，一样是奸夫。婚姻若没有正当解除，是永久的；因此，在婚姻未解除时再婚，就是犯奸淫。因此，既然祂的离婚禁令是有条件的，祂就没有绝对禁止；凡祂没有绝对禁止的，在某些场合，当祂给予禁令的原因不存在时，祂就准许。事实上，祂的教导并不与\Person{梅瑟}相反，\Person{梅瑟}的诫命祂部分地辩护，我不说是确认。然而，如果你否认离婚在任何方面为\Person{基督}所准许，那么你们这一方又怎能破坏婚姻，不结合男女，也不准那些在别处结合的人领受圣洗圣事和圣体圣事，除非他们同意抛弃他们婚姻的果实，连造物主自己也抛弃呢？那么，在你们的教派中，如果妻子犯奸淫，丈夫该怎么办？他要留住她吗？但你们自己的宗徒，你们知道，不准\Person{基督}的肢体与娼妓结合。\BibleRef{格前 6:15} 因此，当离婚理应发生时，它在\Person{基督}那里也有辩护者。因此，\Person{梅瑟}从此必须被视为由祂所确认，因为他在妻子发生不贞时，与\Person{基督}一样禁止离婚。因为在\BibleBook{玛窦}{福音}中祂说：\BibleQuote{凡休妻的，除非因奸淫的缘故，就是叫她犯奸淫。}\BibleRef{玛 5:32} 凡娶被丈夫休弃的妇人的，也被视为同样犯奸淫。 然而，除了因奸淫的缘故，造物主并不拆散祂自己所结合的；同一\Person{梅瑟}在另一处规定，凡侵犯了少女而与她结婚的，从此以后不得休妻。\BibleRef{申 22:28} 现在，如果强逼成婚之后必须永久，那么自愿的、出于同意的婚姻，岂不更应如此吗？这有先知的话为证：\BibleQuote{不可背弃你青年时的妻子。}\BibleRef{拉 2:15} 如此，你看到\Person{基督}在准许离婚和禁止离婚两方面，都自发地追随造物主的踪迹。无论你向哪边逃避，你都发现祂也在保护婚姻。当祂愿婚姻不可侵犯时，祂禁止离婚；当婚姻被不忠玷污时，祂准许离婚。当你拒绝结合那些连你的\Person{基督}都结合了的人时，你应当脸红；当你没有你的\Person{基督}要他们分离的正当理由而分离他们时，你更应当脸红。\Person{主}现在要表明\Person{主}从哪里得出了祂的这个决定，以及祂把它导向什么目的。这样就会更充分地显明，祂的目的不是通过任何突然发明的离婚提议来废除\Person{梅瑟}的条例；因为这提议不是突然提出的，而是扎根于前面提到的\Person{若翰}。因为\Person{若翰}责备\Person{黑落德}，因为他非法娶了他已故兄弟的妻子，而这位兄弟留下了一个女儿（这种结合法律只允许在兄弟无子而死的情况下，那时法律甚至规定这种婚姻，为的是借自己的兄弟，从他自己的妻子，为死去的丈夫立嗣），\BibleRef{申 25:5}结果是\Person{若翰}被投入监牢，最后被同一个\Person{黑落德}处死。\Person{主}既提到了\Person{若翰}，自然也提到了他的死，便以非法婚姻和奸淫的形式，猛烈斥责\Person{黑落德}，甚至宣布那娶了被休之妇的人也是奸夫。祂这样说，是为了更严厉地使\Person{黑落德}罪上加罪；\Person{黑落德}娶了他兄弟的妻子，这妻子不但因死亡，也因离婚而摆脱了丈夫；他这样做是出于情欲，而不是出于（叔娶寡嫂的）法律的规定——因为他的兄弟留下了一个女儿，以此而论，娶那寡妇本身就是不合法的；当先知向他宣告法律时，他便因此杀害了他。我对此案所作的评论，也将有助于我说明随后关于富人在地狱中受折磨、穷人在\Person{亚巴郎}怀中安息的比喻。\BibleRef{路 16:19} 因为这段经文，从字面看，是突然出现在我们面前的；但若考虑其含义和主旨，它自然与提到\Person{若翰}被凶恶地杀害，以及被\Person{若翰}因不敬虔婚姻而谴责的\Person{黑落德}相吻合。它大胆地勾画出他们两人的结局：\Person{黑落德}的\Quote{痛苦}和\Person{若翰}的\Quote{安慰}，为的是让\Person{黑落德}现在就听到那警告：\BibleQuote{他们自有\Person{梅瑟}及先知，听从他们好了。}\BibleRef{路 16:29} 然而，\Person{玛尔基翁}粗暴地将这段经文转向另一个目的，断定痛苦和安慰都是造物主为来世为那些遵守法律和先知的人所保留的报应；而他把天上的怀抱和港口归属于\Person{基督}和他自己的神。我们对此的回答是：那使他眼花缭乱的圣经本身，明确区分了穷人所在的\Person{亚巴郎}的怀抱和阴间的痛苦之所。\Quote{地狱}（我认为）是一回事，\Quote{\Person{亚巴郎}的怀抱}是另一回事。据说\Quote{一个深渊}分隔了这些区域，并阻碍从一方到另一方的通行。此外，富人不可能\BibleQuote{举目一望}\BibleRef{路 16:23}，而且是从远处，除非是从一个较高的地方，并且从所述的距离，穿过巨大的高低深广。因此，每个听说过极乐世界的人，都应该清楚，有一个确定的地方叫做\Person{亚巴郎}的怀抱，它是为接受\Person{亚巴郎}子孙的灵魂而设的，甚至是来自外邦人（因为他是\BibleQuote{万民之父}，这些万民必须归入他的家族），并且具有与他本人相信天主时同样的信德，没有法律的轭和割损的标记。因此，我称这区域为\Person{亚巴郎}的怀抱。虽然它不在天堂，却比地狱更高，被指定为义人的灵魂提供一段安息的时间，直到万物的圆满成全所有人的复活，并赐予他们\Quote{报酬的全部偿报}。这圆满届时将在天堂的恩许中显示出来，然而\Person{玛尔基翁}却将这些恩许归于他自己的神，仿佛造物主从未宣布过它们一样。然而，\Person{亚毛斯}告诉我们关于\Quote{那通天的层级}，是\Person{基督}所\Quote{建造}的——当然是为祂的子民。那里也有\Person{依撒意亚}所问的永久的居所：\Quote{谁能向你宣告那永久的居所，除非是那行走在正义中，谈论正直之道，憎恨不义和邪恶的那位（当然就是\Person{基督}）？} 现在，虽然这永久居所已经应许，而且通天的上升层级是由造物主建造的，祂还进一步应许\Person{亚巴郎}的后裔要像天上的星辰一样，这当然是基于天堂的恩许；那么，在不损害这恩许的前提下，为什么不可以认为\Person{亚巴郎}的怀抱是指信友灵魂的某个临时收容所，其中甚至现在就已描绘出未来的图像，并给出两种审判荣耀的某种预见呢？如果是这样，那么，你们这些异端者，在你们现世的生命中，就有了一个警告：\Person{梅瑟}和先知宣告只有一位天主，造物主，和祂唯一的\Person{基督}，而且永罚和永生的赏报都归于祂，那唯一的天主，祂使人死，也使人活。但是，\Person{玛尔基翁}说，我们的神从天上来的告诫命令我们不要听从\Person{梅瑟}和先知，而要听从\Person{基督}；听从祂\Emph{就是命令}。这确实是真的。因为宗徒们那时已经足够听了\Person{梅瑟}和先知，因为他们被\Person{梅瑟}和先知说服，跟随了\Person{基督}。因为就连\Person{伯多禄}也不能说：\BibleQuote{你是基督。}\BibleRef{路 9:20}除非他事先听到并相信了\Person{梅瑟}和先知，因为直到那时，只有他们宣布了\Person{基督}。事实上，他们的信德理应得到这从天上而来的声音的确认，这声音命令他们\Emph{听从祂}——那位他们认出的宣讲和平、宣告喜讯、应许永久居所、为他们建造通往天堂的层级的那位。然而，在地狱里，关于他们却说：\BibleQuote{他们自有\Person{梅瑟}及先知，听从他们好了！}——甚至那些不信他们、或至少没有真诚相信死后有惩罚来对待财富的傲慢和奢侈的炫耀的人，这些惩罚的确是由\Person{梅瑟}和先知宣布的，但却是那位从高位上推下王侯、从粪堆中提拔穷人的天主所定断的。因此，既然造物主在赏善罚恶两个方向宣布不同的判决是完全一致的，我们必须得出结论说，这里没有神祇的多样性，而只有我们面前实际事理的不同。\par

% structure: p0076 source=h2 target=subsection reason=relative-heading-rank; base=h2

\subsection{第三十五章. \Person{基督}的审判严厉性与造物主的温柔，为反驳\Person{马尔西昂}而主张。十位癞病人的治愈。旧约类比。天主的国在你们内；这教导与\Person{梅瑟}相似。\Person{基督}，匠人所弃的石头。\Person{基督}来临中的严厉迹象。祂不是\Person{马尔西昂}所想象的不能受苦者的证据。}

然后，祂转身对门徒说：\BibleQuote{祸哉，那引人跌倒的人！他若没有生下来，或把一块磨石拴在他的颈上，投入海中，比使他使这些小子中的一个跌倒，为他更好。} \BibleRef{路 17:1} 这就是说，使祂的一位门徒跌倒。那么，请判断祂如此严厉警告的是什么惩罚。因为不是陌生人要为对祂门徒的冒犯复仇。也要在那一位身上认出审判者，祂以与创造主昔日所显示的同等的温情表达对追随者的安全：\BibleQuote{凡触动你们的，就是触动我眼中的瞳仁。} \BibleRef{匝 2:8} 同样关怀出自同一存在。祂要责备一个犯罪的弟兄。\BibleRef{路 17:3} 若有人未尽斥责之责，他实际上犯了罪，或因仇恨希望弟兄继续犯罪，或因错误友谊而姑息，尽管肋未纪中有训令：\BibleQuote{你不可心里怀恨你的兄弟；应切实规劝你的邻人，免得因他而负罪债。} 祂如此教导并不奇怪，因为祂禁止你连发现兄弟的牲畜在途中迷路也不牵回；何况更应将迷途的兄弟带回正路。祂命令你饶恕你的弟兄，即使他得罪你\BibleQuote{七次}。\BibleRef{路 17:4} 但这当然是小事；因为与创造主那里有更大的\OriginalTerm{恩宠}{grace}，祂不设限地赦免，无限期地命令你\BibleQuote{不可对你的兄弟怀有任何恶意}，\BibleRef{肋 19:18} 不仅要给予那求的，甚至也要给予那不求的。因为祂的旨意不是要你原谅冒犯，而是要忘记它。\newline{}关于癞病人的律法具有深刻含义，涉及疾病本身的形式和大司铎的检查。我们将负责阐明这一含义。然而，玛尔定辛勤地以律法的严格性对我们提出反对，以图维护基督在此处也是律法的敌人——甚至在治愈十个癞病人时抢先于祂的诫命。祂简单地命令他们去给司铎察看；\BibleQuote{他们去的时候，就洁净了}，\BibleRef{路 17:11} 没有触摸，没有言语，凭祂静默的大能和简单的意愿。然而，基督既然曾被一次而永远地宣告为我们疾病和罪恶的医治者，并以祂的行为证明了这一点，祂还有什么必要忙于探究治愈的性质和细节？或者创造主在基督身上被召来审查律法？如果此治愈的任何部分以不同于律法的方式由祂完成，祂自己却仍然完美地做到了；因为主当然能凭自己或祂的儿子以一种方式，通过祂的仆人先知以另一种方式产生那些祂大能和权柄的证据，这些证据（因是祂自己的作为而超越荣耀和力量）理所当然地将祂仆人所作的工作远远抛在后面。但前文已经对此点说得够多了。\newline{}现在，虽然祂在前一章中说，\BibleQuote{在厄里叟先知的日子，以色列中有许多癞病人，他们中没有一个得洁净，只有叙利亚人纳阿曼}，但数字当然不能证明神祇的不同，以贬低创造主只治愈一人，而抬高治愈十人的那一位。因为谁能怀疑，由那曾治愈一人者而治愈多人，比由那从未治愈一人者治愈十人更容易？但祂宣告的主要目的是打击以色列的不信或骄傲，因为（虽然他们中有许多癞病人，并且先知并不缺乏）竟没有一个人被如此显著的例子所激动，而投奔那在祂先知中工作的天主。因此，既然祂自己是天主圣父真正的至高司铎，祂按照律法的隐藏意图检查他们，这意图表明基督是人类污秽的真正辨识者和清除者。然而，祂命令执行律法所明显要求的：\BibleQuote{去，把你们自己给司铎察看。} \BibleRef{路 17:14} 然而，如果祂要先洁净他们，这是为什么呢？难道是作为律法的藐视者，为了向他们证明，既然在路上已经治愈，律法对他们已毫无意义，甚至司铎也无关紧要？当然，如果有人以为基督有这样的想法，事情必然以最好的方式过去！但我们可以找到对这段经文更可信的解释：他们之所以被洁净，是因为他们顺从，当他们被命令去司铎那里时，他们去了；而且不能相信遵守律法的人竟能从一位毁灭律法的神那里获得治愈。那么，为什么祂没有给第一个回来的癞病人这样的命令？因为厄里叟在叙利亚人纳阿曼的事上没有这样做，但因此他并非更少是创造主的代理人？这回答足够了。但信徒知道有更深层的理由。因此，请考虑真实动机。奇迹发生在撒玛黎雅地区，其中一个癞病人也来自那个地方。\BibleRef{路 17:17} 然而，撒玛黎雅已从以色列叛离，带走了反叛的九个支派，由先知阿希雅使之疏远，雅洛贝罕将他们安置在撒玛黎雅。此外，撒玛黎雅人总是以他们祖先的山和井为乐。因此，在若望福音中，撒玛黎雅妇人在井边与主交谈时说：\BibleQuote{毫无疑问，祢是更大的}，等等；又说：\BibleQuote{我们的祖先在这座山上朝拜；你们却说，人应当朝拜的地方是在耶路撒冷。} 因此，那位曾说\BibleQuote{祸哉，那些信赖撒玛黎雅山的人}，\BibleRef{亚 6:1} 如今屈尊恢复那地区，特意要求那些人\BibleQuote{去把你们自己给司铎察看}，因为司铎只在那里有，即圣殿所在之处；将撒玛黎雅人置于犹太人之下，因为\BibleQuote{救恩出自犹太人}，\BibleRef{若 4:22} 无论是对以色列人或撒玛黎雅人。确实，应许的基督完全属于犹大支派，为叫人知道在耶路撒冷有司铎和圣殿；那里也是宗教的母胎和它的\Emph{活}泉源，而非它的\Emph{仅仅}\Quote{井}。因此，既然他们认出在耶路撒冷律法应被履行，祂治愈了他们，他们的救恩将来自信德\BibleRef{路 17:19} 而无须律法的仪式。因此，当祂惊讶于十个中只有一个为释放感谢天主的恩宠时，祂没有命令他按律法献祭，因为他已经献上感恩的贡品，当他\BibleQuote{光荣了天主}时；\BibleRef{路 17:15} 因为主愿意这样解释律法的要求。然而，撒玛黎雅人是向哪一位神感恩？因为迄今甚至没有一个以色列人听说过另一位神？除了那位曾通过基督治愈一切的神，还有谁？因此，对他说：\BibleQuote{你的信德救了你，}\BibleRef{路 17:19} 因为他发现他应向全能的天主献上真正的祭品——即感恩——在祂真正的圣殿里，在祂真正的至高司铎\Emph{耶稣}基督面前。\newline{}但法利塞人不可能像是向主询问关于敌对神的国度来临的事，因为从未有其他神被基督宣布过；祂也不会答复他们关于任何其他神的国度，除了他们习惯询问的那位。\BibleQuote{天主的国，不是显然可见的；人也不说：看哪，在这里！或：在那里！因为，看，天主的国就在你们中间。}\BibleRef{路 17:20} 现在，谁会不把\Quote{\Emph{在你们中间}}解释为\Emph{在你们手中，在你们能力之内}，如果你们听从并遵行天主的诫命？然而，如果天主的国在于祂的诫命，请把梅瑟放在你们眼前，按照我们的\OriginalTerm{对立论}{antitheses}，你们会发现同样的观点。\BibleQuote{这诫命不是高不可及，也不是离你们很远。不是在天上，以致你们说：\InnerQuote{谁为我们上到天上，给我们取来，使我们听了而遵行？}也不是在海外，以致你们说：\InnerQuote{谁为我们渡海，给我们取来，使我们听了而遵行？}但这话离你们很近，就在你们口中，在你们心里，在你们手中，好使你们遵行。}\BibleRef{申 30:11} 这意思是：\BibleQuote{天主的国不在这里或那里；因为，看，它就在你们中间。}\BibleRef{路 17:21} 如果异端者胆大妄为，争辩说主没有答复关于祂自己的国，而只是关于他们所询问的创造主的国，那么以下的话对他们不利。因为祂告诉他们：\BibleQuote{人子必须受许多苦，并被弃绝，}\BibleRef{路 17:25} 然后祂才来临，那时祂的国才真正显示。在此陈述中，祂表明祂的回答所考虑到的是祂自己的国，这国正等待祂自己的苦难和弃绝。然而，祂必须被弃绝，然后被承认、被提升、被光荣；祂借用了\Quote{弃绝}这个词本身，来自达味以\Emph{一块石头}的比喻庆贺祂的双重显现——第一次在弃绝中，第二次在尊荣中——的段落：\BibleQuote{匠人所弃的石头，已成了屋角的基石。这是上主所作的。} 现在，如果我们相信天主曾预言任何\Emph{基督}的贬抑甚或光荣，那么祂将祂的预言指向除祂那一位之外的其他任何基督，就是徒劳的，因为祂以\Emph{石头}、\Emph{磐石}和\Emph{山}的比喻预言了祂。然而，如果祂说到自己的来临，为什么祂将它比作诺厄和罗特的日子，\BibleRef{路 17:26} 那些日子是黑暗可怕的——而祂是一位温和、仁慈的神？为什么祂命令我们\BibleQuote{记住罗特的妻子，}\BibleRef{路 17:32} 她藐视创造主的命令，并因她的轻视而受罚，如果祂不是带着审判来惩罚违反祂诫命的行为？如果祂真的像创造主一样惩罚，如果祂是我的审判者，祂就不该从祂所毁灭的那一位那里引用例子来教训我，免得祂看起来像是我的导师。但如果祂甚至在这里也不是说到自己的来临，而是说到希伯来基督的来临，那么我们要继续期待祂会屈尊告诉我们一些关于祂自己来临的预言；同时我们将继续相信祂不是别的，正是祂在每一段经文中提醒我们的那一位。\newline{}\par

% structure: p0078 source=h2 target=subsection reason=relative-heading-rank; base=h2

\subsection{第三十六章 不义判官和寡妇、法利塞人和税吏的比喻。\Person{基督}对富有的官的答复、瞎子得治。祂的问候——\Person{达味}之子。这一切证明\Person{基督}与创造主的关系，\Person{玛尔强}的\WorkTitle{对立论}关于\Person{达味}和\Person{基督}的对比被驳斥。}

当祂劝勉人在祈祷中要恒心切祷时，祂给我们设了一个比喻，即那位因寡妇的恒切恳求而被迫听从她的判官。\BibleRef{路 18:1} 祂表明，我们必须以祈祷恳求的那位判官是天主，而非祂自己——如果祂本人不是判官的话。但祂接着又说：\Quote{天主必会为祂的选民伸冤。}\BibleRef{路 18:7} 因此，既然判者也将亲自作伸冤者，祂便证明了创造者就是那特别良善的天主，祂将祂描绘为那为日夜向祂呼号的选民伸冤者。然而，当祂将创造者的圣殿\Emph{呈现在我们眼前}，并描述两人在其中以不同心情祈祷——法利塞人傲慢，税吏谦卑——并表明他们如何分别回家，一个被弃绝，另一个成了义，\BibleRef{路 18:10} 那么，祂藉此教导我们祈祷的正当训练，无疑已确定必须向那位天主祈祷，因为人要从祂那里领受这种祈祷的训练——或谴责骄傲，或使谦卑成义。我从基督那里找不到任何属于创造者以外之神的圣殿、祈求者或判决（赞同或谴责）。祂命令我们以谦卑敬拜那高举谦卑者的祂，而非以骄傲敬拜那贬抑骄傲者的祂。祂向我显现了哪一位别的神，好让我向祂献上祈求？我该以什么敬礼形式、怀着什么希望去接近祂？我想，没有。因为祂教导我们的祈祷，如我们所证明的，只适合创造者。当然，如果祂不愿受祈祷，因为祂是至善、自发地善的天主，那又是另一回事！但这位良善的天主是谁？祂说：\Quote{除了一个以外，没有良善的。}\BibleRef{路 18:19} 这并非好像祂向我们表明两位神中有一位是至善的；而是祂明确断言只有一个良善的天主，祂是唯一的良善，因为祂是唯一的天主。现在，毫无疑问，祂就是那良善的天主，祂\Quote{降雨给义人，也给不义的人，叫日头照好人，也照歹人；}\BibleRef{玛 5:45} 甚至养育、滋养、帮助马西昂派的人！后来，\Quote{有一个人问祂说：\InnerQuote{善师，我该做什么才能承受永生？}}耶稣问他是否\OriginalTerm{知道}{know}（也就是说，换言之，是否\OriginalTerm{遵守}{keep}）创造者的诫命，以证明永生是通过创造者的规诫获得的。\BibleRef{路 18:18} 当他肯定自己从幼年起就遵守了所有主要诫命时，（耶稣）对他说：\Quote{你还缺少一件：变卖你一切所有的，分给穷人，就必有财宝在天上；你还要来跟从我。}\BibleRef{路 18:21} 好了，马西昂，以及你们这些与那异端者同受苦难、同怀仇恨的伙伴，对此你们敢说什么？基督是否撤销了上述诫命：\Quote{不可杀人，不可奸淫，不可偷盗，不可作假见证，孝敬你的父母？} 还是祂既遵守了它们，又补充了所欠缺的？然而，这关于赒济穷人的诫命，在律法和先知书中已广泛散布。这位好夸口的诫命遵守者因此被证实将钱财看得远比爱德为重。福音的真理于是丝毫无损：\Quote{我来不是要废除律法和先知，而是为成全。}\BibleRef{玛 5:17} 祂也消除了其他疑虑，当祂宣布天主的名和善的名属于同一位存在者，永生、天上的财宝以及祂自己也都由那位存在者掌管——祂的命令，祂既持守，又用自己的补充规诫予以增益。我们或许在\BibleBook{米该亚}{先知书}的以下段落中也发现祂说：\Quote{人啊，已指示你什么是善；上主向你要求什么，无非是行公义，喜爱仁慈，并准备跟随上主你的天主。} 基督就是那告诉我们何为\Emph{善}的人，即律法的知识。祂说：\Quote{你知道诫命。} \Quote{行公义}——\Quote{变卖你一切所有的；} \Quote{喜爱仁慈}——\Quote{分给穷人：} \Quote{并准备与天主同行}——祂说：\Quote{你还要来，跟从我。} 犹太民族从起初就如此仔细地划分为支派、宗族、家族和家室，以致无人能清楚自己的出身——甚至奥古斯都最近的户口普查，当时可能仍然存在。但马西昂的耶稣（虽然一个人既被看见是人，他的出生无疑不可怀疑），既未出生，自然不可能拥有任何公开的出身证明，而应被视为那一类无人知晓、无名之辈中的一员。那么，那盲人听见祂经过，为何喊叫：\Quote{耶稣，达味之子，可怜我吧？}\BibleRef{路 18:38} 除非他无不确定地认为祂是达味之子（换言之，属于达味的家族），通过祂的母亲和兄弟，那些人先前某个时候因公众的声望而为他所知。\Quote{在前头走的人就责备他，叫他不要出声。}\BibleRef{路 18:39} 这很恰当，因为他太吵闹了，并非因为他在达味之子的事上错了。否则，你必须向我证明，那些责备他的人知道耶稣不是达味之子，这样他们才可能有理由让盲人安静。但即使你能向我证明这一点，那盲人仍然更可能会假定他们无知，而不是主可能允许他人对自己发出不真实的呼声。然而，\Quote{主站住了。}\BibleRef{路 18:40} 是的，但并非为了证实错误，相反，祂反而彰显了创造者。祂岂能先除去这人的盲疾，好让他后来不再视祂为达味之子呢？然而，为使你们不致诽谤祂的忍耐，或指控祂虚伪，或否认祂是达味之子，祂非常明确地确认了盲人的呼声——既通过实际赐予的医治，又通过为他的信德作证：\Quote{你的信德救了你。}\BibleRef{路 18:42} 你希望这盲人的信德是什么？是耶稣出自那（异方之）神（马西昂的神），为颠覆创造者、废除律法和先知？是祂并非那从叶瑟之根发出的枝条、达味腰间的果子，以及复明盲人的那一位？但我认为，当时并没有像马西昂那样瞎眼的人，以致这样的看法能构成盲人的信德，并使他只凭耶稣——达味之子——的\Emph{名字}就信赖祂。祂自己知道这一切，也愿意别人知道，便赐予这人——他的信德已经拥有更好的视力，并且已拥有真光——外在的视力，好让我们也学习信德的准则，并同时找到它的赏报。凡愿看见耶稣是达味之子的，必须相信祂；通过童贞女的诞生。凡不愿如此相信的，将不会从祂那里听到\Quote{你的信德救了你}的问候。于是他仍将盲目，跌入一个个\OriginalTerm{反命题}{Antithesis}中，它们互相摧毁，正如\Quote{瞎子领瞎子，一同掉进坑里。} 因为（这里是马西昂的一个\OriginalTerm{反命题}{Antithesis}）：从前大卫在夺取熙雍时，被那些阻碍他进入（堡垒）的瞎子冒犯了（就此而言，我更应该说，他们是同样瞎眼之人的预表，这些人日后不肯承认基督是达味之子）；与此相反，基督却帮助了那瞎子，以此表明祂不是达味之子，并且祂性情不同，对瞎子仁慈，而大卫却命令杀死他们。如果真是这样，为什么马西昂声称那瞎子的信德如此一文不值？事实是，达味之子如此行事，以致那\OriginalTerm{反命题}{Antithesis}因其自身的荒谬而失去效力。那些冒犯大卫的人是瞎子，现在向达味之子恳求的人也患有同样的残疾。因此，达味之子因瞎子的某种满足而平息了怒气，使他复明，并赞许他的信德，这信德使他相信了真理：他必须用恳切的祈求来赢得达味之子的帮助。但毕竟，我怀疑冒犯大卫的是（老耶布斯人的）傲慢，而非他们的疾病。\par

% structure: p0080 source=h2 target=subsection reason=relative-heading-rank; base=h2

\subsection{第37章 基督与匝凯。马西昂所否认的身体的救恩。十个仆人领取十\OriginalTerm{米纳}{mina}的比喻。基督是审判者，要执行那严厉之人的旨意，即造物主。}

\Quote{救恩临到了}匝凯\Quote{的家}。（\newline{}\BibleRef{路 19:9}）为甚么呢？是因为他也相信基督是由马西昂来的吗？但盲人的呼喊仍在众人耳中回响：\Quote{耶稣，达味之子，可怜我吧。}并且\Quote{众百姓都归光荣于天主}——不是马西昂的，而是达味的天主。虽然匝凯可能是个外邦人，但他与犹太人交往，略知他们的圣经，而且，更甚于此，他在不知不觉中实践了依撒意亚的训言：\Quote{把你的粮食}先知说\Quote{分给饥饿的人，将流离失所的人领到你的家中。}（\newline{}\BibleRef{依 58:7}）他以最好的方式做了这事：接待主，留祂在家中。\Quote{见赤身露体的人时要给他衣服穿。}他以同样令人满意的方式承诺这样做：将他的一半财物用于一切慈善工作。同样地，他\Quote{解开邪恶的锁链，废除轭上的绳索，使受压迫者获得自由，折断所有的轭}（\newline{}\BibleRef{依 58:6}），当他说：\Quote{我若欺骗过任何人，就偿还他四倍。}（\newline{}\BibleRef{路 19:8}）因此主说：\Quote{今天救恩临到了这一家。}（\newline{}\BibleRef{路 19:9}）祂就这样作证：由先知所说的造物主的训言导向救恩。但当祂加上：\Quote{因为人子来，是为寻找并拯救迷失的人。}（\newline{}\BibleRef{路 19:10}）我目前的争论不在于\Emph{祂}是否来拯救那迷失的，\Emph{那曾经属于谁的}，以及祂所救来的人\Emph{从谁那里}失落了；而是我接近了一个不同的问题。无疑地，\Emph{人}是这里的主题。既然他由两部分组成：身体和灵魂，需要探究的是，人在哪一部分显得是迷失的？如果是在身体里，那么迷失的是他的身体，而不是他的灵魂。然而，那迷失的，人子拯救。因此，身体获得救恩。如果（另一方面）人是在灵魂里迷失的，那么救恩是为那迷失的灵魂预备的；而没有迷失的身体是安全的。如果（取唯一剩下的假设）人完全迷失，在两重本性中都迷失，那么必然得出救恩是为整个人预备的；这样，异端者的意见就被粉碎了，他们说没有身体的救恩。这证实了基督属于造物主，祂效法造物主许诺整个人的救恩。那十个仆人的比喻也证明祂是一位审判的天主——甚至是一位严格结算的天主，不仅给予尊荣，也拿走他所拥有的。否则，如果这里描述的是造物主，当作那\Quote{严厉的人}，他\Quote{拿走他没有存放的，收割他没有播种的}（\newline{}\BibleRef{路 19:22}），那么在这里教导我的正是祂，（无论祂是谁），那钱就属于祂，祂教导我有效地使用。\par

% structure: p0082 source=h2 target=subsection reason=relative-heading-rank; base=h2

\subsection{第38章 基督对法利塞人的驳斥。给凯撒和天主纳税。接着对撒杜塞人关于复活时婚姻的驳斥。这些证明祂不是马西昂的基督，而是造物主的基督。马西昂为了让他的第二神登场而进行的篡改，被揭露和驳斥。}

基督知道\Quote{若翰的圣洗是从哪里来的}。\BibleRef{路 20:4}那么祂为何向他们发问，好像不知道似的？祂知道法利塞人不会给祂答案；那么祂为何徒然发问？岂不是要凭他们自己的口，或他们自己的心，来审判他们？假若你把这些问题归咎于创造主的托词，或是归咎于祂与基督的比较；那么就想想，如果法利塞人回答了祂的问题，会发生什么事。假设他们的回答是若翰的圣洗是\Quote{从人来的}，他们就会立刻被石头砸死。\BibleRef{路 20:6}某个马西昂，与马西昂竞争，会站起来说：哦，至善的天主！祂的行事与创造主何其不同！知道人会从那里一头栽下去，祂却把他们放在悬崖边缘。因为人对创造主关于那棵树的律法就是这样处理的。但若翰的圣洗是\Quote{从天而来的}。\Quote{那么，}基督问道，\Quote{你们为什么不信他？}\BibleRef{路 20:5}因此，那位希望人们相信若翰、并打算责备他们因不信他的人，是属于若翰所施行圣事的那一位。无论如何，当祂真的用这样的报复来回应他们不愿说出自己的想法时，祂说：\Quote{我也不告诉你们我凭什么权柄做这些事}，\BibleRef{路 20:8}祂是以恶报恶！\Quote{凯撒的归凯撒，天主的归天主。}\BibleRef{路 20:25}什么是\Quote{属于天主的事}？就是像凯撒的\OriginalTerm{德纳}{denarius}那样的事——即祂的肖像和模样。因此，祂命令\Quote{归还给天主}、即创造主的，是\Emph{人}，因为人印有祂的肖像、模样、名字和本质。让马西昂的天主照看自己的钱币吧。基督吩咐人的印记的\OriginalTerm{德纳}{denarius}要还给祂的凯撒（我说的是祂的凯撒，而非异神之凯撒）。但真理必须承认，这个天主没有属于自己的\OriginalTerm{德纳}{denarius}！在每一个问题中，公正且恰当的规则是，答案的含义应当适应所提出的询问。但若给出一个与所提交的问题完全不同的答案，那简直是疯狂。愿天主禁止我们期待基督做出连对一个普通人都不合适的行为！否认复活的撒杜塞人，在讨论这个问题时，向主提出了一个关于某个妇女的法律案例：她按照法律规定，嫁给了先后死去的七兄弟。\BibleRef{路 20:27}因此问题是：在复活中，她应属于哪一个丈夫？请注意，这就是询问的要点，这就是争论的焦点和实质。基督必须对此给出直接的回答。祂无所畏惧，所以祂似乎既不应该回避他们的问题，也不应该借此机会间接提出祂在其他时候并不公开教导的话题。因此祂给出了回答：\Quote{今世之子也娶也嫁。}\BibleRef{路 20:34}你看这与案例多么切合。因为问题涉及来世，而祂要宣布在那里没有人嫁娶，祂先奠定原则：此处有死亡，也有婚姻。\Quote{但那堪得来世及从死人中复活的，也不娶也不嫁；因为他们不能再死，既然与天使相等，成为天主和复活之子。}\BibleRef{路 20:35}那么，如果答案的含义不能偏离所提出的问题，而问题通过这个答案的含义已完全被理解，那么主的回答不允许有任何其他解释，只能通过它来清晰理解问题。你已经看到了婚姻被允许的时间，以及它被认为不合适的时间，这些时间并非出于自身，而是由于关于复活的询问。你也同样看到了复活本身的确认，以及撒杜塞人所提出的全部问题——他们并未询问关于另一个天主，也未询问关于婚姻的正确律法。现在，如果你让基督回答那些并未向祂提交的问题，你实际上是在表明祂无法解决真正被咨询的问题，并且当然是被撒杜塞人的诡计所困。我现在将额外按照关于问题和答案的规则，处理那些具有一致性的论点。他们弄到了圣经的抄本，并以这种方式断章取义地阅读：\Quote{那些被那世之神堪当认可的}。他们把\Quote{那世}这个短语加到\Quote{神}这个词上，从而造出了另一个神作为\Quote{那世之神}；而这段经文本应这样阅读：\Quote{那些被天主堪当认可获得那世的人}，即将限定短语\Quote{今世}移到从句末尾，换句话说：\Quote{那些被天主堪当认可获得那世并从死里复活的人。}因为向基督提交的问题与\Emph{那世}的\Emph{神}无关，只与\Emph{那世}的\Emph{状态}有关。问题是：\Quote{在复活后的那世，这妇人是谁的妻子？}\BibleRef{路 20:33}他们如此颠覆了祂关于婚姻本质问题的回答，并将祂的话\Quote{今世之子娶也嫁也}解释为指创造主的人以及祂允许他们结婚；而他们自己——即被那世之神（即敌对之神）堪当认可复活的人——甚至在这里也不结婚，因为他们不是今世之子。但事实上，既然被咨询的是关于\Emph{那世}的婚姻，而非今世，祂只是宣告了问题所涉及之事的不存在。确实，那些抓住了祂声音、发音和表达方式的力量的人，除了与问题相关的内容外，没有发现其他含义。因此，经师们喊道：\Quote{师傅，祢说得好。}\BibleRef{路 20:39}因为祂肯定了复活，通过描述其形式来反驳撒杜塞人的意见。祂并未拒绝那些认为祂的回答有此含义之人的作证。然而，如果经师们认为基督是达味之子，（但达味）自己称祂为主，\BibleRef{路 20:41}这与基督有何关系？达味并未实际驳斥经师们的错误，但达味肯定了基督的尊荣，因为他更突出地声称祂是主而非子——这个属性几乎不适合创造主的毁灭者。但我们的解释何等一致！因为不久前被盲人呼为\Quote{达味之子}的那一位，\BibleRef{路 18:38}当时并未对此事发表评论，因为经师们不在场；而如今祂故意在他们面前提出这个问题，而且是\Emph{主动}提出，\BibleRef{路 20:41}以便表明祂自己——盲人按照经师的教义仅称之为达味之子——也是他的主。这样，祂既尊崇了那承认祂为达味之子的盲人的信德，同时又打击了经师们的传统，这传统阻止他们认识祂也是（达味的）主。凡是关于创造主之基督荣耀的事，除了创造主之基督自己，无人会如此维护和持守。\par

% structure: p0084 source=h2 target=subsection reason=relative-heading-rank; base=h2

\subsection{第三十九章　关于那些以基督之名而来的人。祂来临的可怕征兆。在\WorkTitle{旧约}和\WorkTitle{新约}中如此宏伟描述的祂，正是创造者的基督。这个证明借助无花果树和所有树木的比喻得以加强。预言的平行段落。}

论及祂名号的恰当性，已经看到这两个名号都适合于祂，祂是第一个既向人类宣告祂的\Emph{基督}，又赐予祂\Emph{耶稣}这一进一步名号的那一位。因此，\Person{马尔基昂}的基督的无耻将显而易见，当他说将有许多人奉他的名而来时，而这名号根本不归属于\Emph{他}，因为他并非造物主的基督和耶稣，这些名号真正属于造物主；尤其当他禁止接纳那些与他完全等同的骗子时（因为他同样以属于另一位的名号而来）——除非他的任务是警告门徒远离一个虚假冒用的名号，这些门徒（属于那一位）因祂的名号被正当地赐予祂，也拥有其真实性。但是当\BibleQuote{他们将不久前来临且说：\InnerQuote{我是基督}}（\BibleRef{路 21:8}），你们将接纳他们，因为你们已经接纳了一个与他们完全相似的人。然而，基督是奉自己的名号而来。那么，当祂亲自来临，祂正是这些名号的真正拥有者，即造物主的基督和耶稣时，你们将做什么？你们会拒绝祂吗？但这是何等不义、何等不公正、何等冒犯好天主：你们竟不接受那奉自己名号而来者，却接受了一个奉祂名号而来的人！现在，让我们看看祂归之于时代的征兆。我观察到：\BibleQuote{战争，邦国与邦国相争，民族与民族相争，瘟疫、饥荒、地震、可怖的景象，以及来自天上的巨大征兆}（\BibleRef{路 21:9}）——这一切都适合于一位严厉可畏的天主。现在，当祂接着说\BibleQuote{这一切必须发生}，祂表明自己是何种身分？是造物主的毁灭者，还是维护者？因为祂肯定这些祂的定命必须完全实现；但既然祂是好天主，祂宁愿挫败而非促成如此悲惨可怕的事件，除非这些是祂自己的（法令）。\BibleQuote{但在这一切之前}，祂预言迫害和苦难将临到他们身上，这些迫害和苦难确实\BibleQuote{要为他们成为见证}，并带来救恩。请听在\BibleBook{匝加利亚}{先知书}中所预言的：\BibleQuote{万军的主必保护他们；他们要吞灭他们，用投石器征服他们；他们必饮他们的血如酒，并充满碗正如祭台一样。在那日，主必拯救他们，即祂的百姓，如羊一般；因为他们像圣石一样滚动，}等等。为使你们不以为这些预言是指他们从与外邦人的众多战争中将要遭受的苦难，请考虑（受苦的）性质。在关于以合法兵器进行的战争的预言中，没有人会想到列出石头作为武器，石头更常见于民众暴乱和手无寸铁的骚动。没有人用碗来测量战争中流出的滔滔血流，也不仅限于在一个祭台上所流。没有人称那些手持兵器、以武力抵抗武力而倒下者为羊，而只称那些殉道者为羊，他们顺服地、耐心地献出自己于本分之地，而非自卫作战。简言之，如他所说：\BibleQuote{他们像圣石一样滚动}，并不像士兵战斗。他们是石头，甚至是基石，我们自身被建立在其上——正如圣保禄所说：\BibleQuote{被建造在宗徒的基石上}（\BibleRef{弗 2:20}），他们如\BibleQuote{被祝圣的石头}，被滚来滚去，暴露于众人的攻击之下。因此，在此段落中祂禁止人们\BibleQuote{预先思考如何回答}当被带到法庭时（\BibleRef{路 21:12}），正如祂曾向\Person{巴郎}提示他未曾想到的信息，\EditorialNote{民数纪第二十二至二十四章}，甚至与他所想相反；并应许给\Person{梅瑟}\BibleQuote{一张口}，当他以口舌迟钝为推辞时（\BibleRef{出 4:10}），以及那智慧，祂通过\Person{依撒意亚}显示为不可抗拒的：\BibleQuote{一人说：\InnerQuote{我是主的}，并要以雅各伯的名字称呼自己；另一人要以以色列的名字签名。}（\BibleRef{依 44:5}）那么，有什么辩护比一位殉道者所做的简单而公开的宣认更智慧、更不可抗拒？这殉道者\Quote{与天主相争胜}——这正是\Quote{以色列}的含义？现在，无怪乎祂禁止\Quote{预先谋划}，因为祂自己实际从圣父领受了适时发言的能力：\BibleQuote{主赐给我受教者的舌头，使我知道怎样用言语扶助疲倦的人；}（\BibleRef{依 50:4}）但\Person{马尔基昂}却向我们介绍一位不服从圣父的基督。来自最亲密朋友的迫害以及因祂名号而起的诽谤被预言（\BibleRef{路 21:16}），我不必再提。但祂说：\BibleQuote{凭着忍耐，你们要拯救自己。}关于这忍耐，圣咏说：\BibleQuote{义人的耐心忍耐将永不灭亡；}因为另一篇圣咏说：\BibleQuote{义人的死，（在主眼中）是宝贵的}——这无疑出自他们的耐心忍耐，因此\BibleBook{匝加利亚}{先知书}宣称：\BibleQuote{忍耐者将获冠冕。}为使你们不致大胆争辩说宗徒被犹太人迫害是因为他们宣告了另一位神，请记住：就连先知也遭受了犹太人同样的对待，他们并非任何其他神的传报者，而只是造物主的传报者。然后，在指明毁灭的时期，即\BibleQuote{当耶路撒冷开始被军队包围时}（\BibleRef{路 21:20}），祂描述了万物的终结的征兆：\BibleQuote{太阳、月亮和星辰中出现的异兆，地上民族在困惑中痛苦不堪——如大海咆哮——因为他们期待即将临于地上的灾祸。}（\BibleRef{路 21:25}）\par

那\Quote{连天上的万军也必将动摇}，\BibleRef{路 21:26} 你可以在岳厄尔中找到：\BibleQuote{我要在天上地下显示奇事——血、火和烟柱；太阳要变为黑暗，月亮要变为血，在主伟大而可怕的日子来到以前。} \BibleRef{岳 3:30} 在哈巴谷书中你也有这样的陈述：\BibleQuote{大地将为江河所裂；万民看见你，必痛苦；你以脚步分散众水；深渊发声，惊惶之高耸；日月停在原处；你的电光变为光明；你的盾牌如同闪电之光辉；你发怒时踏碎大地，在怒气中打粮压榨万民。} 因此，我认为，主和先知有关大地震动、诸元素和万国的言论是一致的。但主后来是怎么说的？\BibleQuote{那时，他们要看见人子带着极大的权能从天上降临。当这些事发生时，你们应当挺身昂首，因为你们的救赎近了。} \BibleRef{路 21:27} 这指的是那比喻本身所论及的国度时期。\BibleQuote{同样，当你们看见这些事发生时，也应知道天主的国近了。} \BibleRef{路 21:31} 这就是主的日子，人子从天上荣耀降临的日子，达尼尔曾写道：\BibleQuote{看，有一位像人子者，乘着天上的云彩而来，} \BibleRef{达 7:13} 等等。\BibleQuote{他便被授予了王权，} \BibleRef{达 7:14} 这（在比喻中）\Quote{他往远方去为自己取得王权}，留给仆人们银钱去经商生利——即万民的（普世国度），圣咏中圣父曾应许赐给他：\BibleQuote{你向我请求，我必将外邦人赐你为产业。} \BibleQuote{一切荣耀都要侍奉他；他的统治是永远的，不会被夺去，他的国度不会被毁灭，} \BibleRef{达 7:14} 因为在其中\Quote{人们不再死亡，也不嫁娶，就如天使一样。} \BibleRef{路 20:35} 关于人子同样的来临及其益处，我们在哈巴谷书中读到：\BibleQuote{你出来拯救你的百姓，拯救你的受膏者，} \BibleRef{哈 3:13} ——换句话说，就是那些在祂国度时期得救赎而挺身昂首的人。既然如此，一方面这些关于应许的描述相互一致，正如另一方面那些关于大灾难的描述也是一致——无论是在先知的预言中，还是在主的宣告中——你不可能在它们之间划出任何区别，仿佛灾难可归给创造者（那位可怕的天主），而这正是\OriginalTerm{善神}{good god}（马西昂的）所不应允许、更不应期望的；与此同时，应许则应归给善神，而创造者由于不知此神，无法预言它们。然而，如果创造者确实将这些应许作为自己的应许预言了，既然它们与基督的应许毫无区别，那么祂在恩赐的慷慨方面就与善神本身不相上下；显然，你的基督所应许的，不会比我的《人子》所应许的更多。（如果你仔细查考）这段福音经文的全部内容，从门徒的询问到无花果树的比喻，\BibleRef{路 21:33} 你会发现其意义在关联上处处与《人子》相合，以至于始终一致地将悲伤与喜乐、灾难与应许都归于祂；你不能在任何一个方面将它们与祂分开。因此，既然只有一个《人子》，祂的来临介于灾难与应许两种结局之间，那么必要的结果是，对万民的审判和圣徒的祈祷都归于那同一个《人子》。祂这样居于中间，共同关联两种结局，祂将在来临时通过审判万民来终止其中一个，同时也通过实现圣徒的祈祷来开始另一个：因此，如果你承认人子的来临是我的\Emph{基督}的来临，那么当你将祂显现之前施行的审判归于祂时，你也必然将与这审判同出的福分归于祂。另一方面，如果你坚持这是你的\Emph{基督}的来临，那么当你将祂来临所导致的福分归于他时，你就不得不也将祂显现之前降下的灾祸归于他。因为人子来临之前的灾祸和紧随其后的福分，两者都与那事件不可分割地联系在一起。那么，思考一下，你选择在《人子》的角色中安置哪一位基督，你就能把两种经世之工的执行归于谁。你要么将创造者变成最仁慈的天主，要么将你自己的神变成本性可怕的神！简言之，思考一下比喻中的画面：\BibleQuote{你们看看无花果树及各种树木；当它们结出果子时，你们就知道夏天近了。同样，当你们看见这些事发生时，也应知道天主的国近了。} \BibleRef{路 21:29} 如今，如果普通树木的结果是夏天临近的预兆，那么同样，世界的巨大冲突也预示着它们所先导的国度的来临。但每一个征兆都属于那事物所属者；每一事物都由其所属者指定其征兆。因此，如果这些苦难是国度的征兆，就像树木成熟是夏天的征兆一样，那么国度就属于创造者，因为苦难——国度的征兆——归于祂。既然仁爱的神早已断言这些事必须发生，尽管如此可怕恐怖，正如法律和先知所预言的，那么祂并没有废除法律和先知，当祂肯定其中所预言的必定应验时。祂进一步宣告：\BibleQuote{天地绝不能不等到一切都成就，就先行过去。} \BibleRef{路 21:33} 请问，这些事是什么？是创造者所造的吗？那么，万有将顺从地忍受其创造者经世的完成。然而，如果它们来自你那个卓越的神，我很怀疑天地会平静地容许它们创造者的敌人所决定之事的完成！如果创造者默默地顺从，那么祂就不是一位\Quote{忌邪的天主}。但天地可以过去，既然它们的主如此决定；但愿祂的话永远长存！依撒意亚也这样预言过。\BibleRef{依 40:8} 门徒们也应受警告：\Quote{恐怕他们因贪食醉酒和今生的挂虑而心神麻木，那日子如同网罗突然临到你们}——如果他们真的在世俗的富足和忙碌中忘记了天主。梅瑟的劝诫也是如此——所以，从那日子的\Quote{网罗}中拯救出来的，正是那位早在许久以前就向人发出同样劝诫的那一位。\BibleRef{申 8:12} 在耶路撒冷有些地方可以教导，有些地方则在城外可以退隐，\BibleRef{路 21:37}——\Quote{祂白天在圣殿里施教}；正如祂藉欧瑟亚所预言的：\Quote{在我的家中，他们找到了我，我在那里与他们说话。} \Quote{但夜间祂出去，到橄榄山上。} 因为匝加利亚曾指明：\Quote{在那一天，祂的脚要站在橄榄山上。} \BibleRef{匝 14:4} 那里也有适合听众的时辰。\Quote{清晨} \BibleRef{路 21:38} 他们必须到祂那里去，祂曾藉依撒意亚说：\BibleQuote{上主赐给我受教者的舌头}，又加上：\BibleQuote{祂为我安排了清晨，也赐给我耳朵聆听。} \BibleRef{依 50:4} 如果这是\Emph{毁灭}先知，那么\Emph{成全}他们又是什么呢？\par

% structure: p0087 source=h2 target=subsection reason=relative-heading-rank; base=h2

\subsection{第四十章　救主受难的步骤如何在预言中预定。逾越节。犹达斯的背信。主最后晚餐的建立。耶稣基督的身体和血驳斥了马西昂的幻影论谬误。}

同样，祂也知道祂必须受苦的准确时间，因为法律预表了祂的苦难。因此，在犹太人的所有节期中，祂选择了逾越节。\EditorialNote{路加福音22:1}梅瑟曾在这节期中宣告过一项神圣的奥秘：\Quote{这是主的逾越节。}\BibleRef{肋 23:5}那么，祂是多么恳切地显露出祂内心的倾向：\Quote{我渴望又渴望，在我受苦以前，同你们吃这一次逾越节晚餐。}\BibleRef{路 22:15}这位实际渴望遵守逾越节的，是何等一位法律的毁灭者！难道祂如此喜爱犹太人的羔羊吗？不正是因祂必须被\Quote{牵到宰杀之地，又像羊在剪毛的人前不出声，所以祂不开口}，\BibleRef{依 53:7}祂才如此深切地渴望完成祂自己赎罪之血的预表吗？祂本也可被任何外人出卖，但我发现连这里祂也实现了圣咏：\Quote{那吃过我饭的人，也举起脚后跟踢我。}祂甚至可以不需代价就被出卖。因为对于一位公开将自己献给民众、既能被武力轻易捕获也可被诡计擒拿的人来说，出卖者有何必要？这对于另一位基督也许完全适合，但对于一位正在实现预言的基督却不合适。因为经上记载：\Quote{他们为了银子出卖了义人。}\BibleRef{亚 2:6}那笔钱的数额和用途——犹达斯后悔后，从\Emph{从最初作为酬金的目的}收回，用来购买陶工的田地——如同玛窦福音所叙述的，已在耶肋米亚书中清楚预言：\Quote{他们拿了三十块银币，即被估价者的价钱，就用来买了陶工的那块田。}当祂如此恳切地表达要吃逾越节晚餐的愿望时，祂把这视为\Emph{祂自己的}筵席；因为天主渴望分享不属于自己的东西是不相称的。然后，祂拿起饼，递给祂的门徒们，使之成为祂自己的身体，说：\Quote{这是我的身体}，也就是说，我身体的形象。然而，除非先有真实的身体，否则不可能有形象。空无之物或幻影是不能有形象的。如果（如马西昂可能说的）祂假装饼是祂的身体，因为祂缺少真实的身体实质，那么推论就是祂必须把饼给了我们。这对于马西昂的幻影身体理论会很有帮助——饼竟然被钉在十字架上！但为什么称祂的身体为饼，而不称为别的可食之物（比如说甜瓜），马西昂一定是用甜瓜代替了心！祂不明白这基督身体的形象是多么古老，祂自己曾通过耶肋米亚说：\Quote{我好像一只温柔的小羊被牵去宰杀，我却不知道他们设计谋害我，说：\Emph{我们把这棵树放在祂的饼上。}}这当然指的是十字架放在祂的身体上。因此，祂一如既往地启示古代的预言，清楚地表明了祂用\Emph{饼}指什么，即把饼称为祂自己的身体。同样，当祂提到杯，并订立了\Quote{以祂的血}所立的\Emph{新}约时，\BibleRef{路 22:20}证实了祂身体的真实性。因为没有血属于不是\Emph{一个}肉身的身体。如果向我们呈现的任何身体不是肉身，那么它不是属血肉的，就不会有血。因此，从肉身的证据我们获得身体的证明，从血的证据获得肉身的证明。然而，为让你发现酒自古以来就被用作血的形象，请转向依撒意亚，他问：\Quote{那从厄东，从波责辣来，身穿染红衣服，衣着华美，力量强大，阔步前行的，是谁？你的衣服为何是红色，你的衣裳好像踏过酒醡者的衣裳？}先知的神视看到主仿佛已在前往受难途中，穿着祂的肉体本性；由于祂要在其中受苦，祂将祂肉体的流血状况比作染红的衣服，仿佛在踏压和破碎酒醡的过程中变红，工人从酒醡下来时被葡萄汁染红，好像沾染血迹的人。创世纪预言得更清楚，当（在犹大的祝福中，基督按肉体要出自犹大支派）它甚至在当时就以这位宗祖的形象描绘了基督，说：\Quote{他在酒中洗净了衣服，在葡萄汁中洗净了外氅。}\BibleRef{创 49:11}——在衣服和外氅中，预言指向了祂的肉体，在酒中指向了祂的血。如此，祂现今在酒中祝圣了祂的血，而祂从前（通过宗祖）用酒的形象描述祂的血。\par

% structure: p0089 source=h2 target=subsection reason=relative-heading-rank; base=h2

\subsection{第四十一章　对出卖者宣布的祸患是一个司法行为，否定了基督是马西昂所希望的那位。基督在公议会前的行为解释。基督甚至那时引导祂的审判官们思考关于祂使命的预言证据。这些人的道德责任被确认。}

\Quote{祸哉}，祂说，\Quote{那出卖人子的人！} \BibleRef{路 22:22} 如今可以肯定，在这\Emph{祸哉}中，必须理解为一位愤怒并激怒的主的咒骂和威胁，除非犹大在犯了如此大罪之后还能免受惩罚。若他本意是免受惩罚，那\Emph{\Quote{祸哉}}就是一句空话；若非如此，他当然要受到他所犯背信之罪所冒犯的那位的惩罚。如今，若祂明知那祂特意拣选作为自己同伴的人，却任凭他陷入如此大罪，你便不能再以亚当的案件来反对造物主，因为现在这话会反噬你自己的天主：要么祂是\Emph{无知的}，没有预见到阻止那将来的罪人；要么祂是\Emph{无能为力的}，即使无知也无法阻止他；要么祂是\Emph{不愿意的}，即使有预知和能力，因此应得恶意的污名，竟任凭自己所拣选的人在自己的罪中丧亡。因此，（乐意地）我劝你，在你的那位天主身上承认造物主，而不是违背你的意愿，将你那卓越的天主与祂相比。因为关于伯多禄的例子，祂也给了你证明祂是一位嫉妒的天主，当祂定意在宗徒狂妄地表达热心之后，让他彻底否认祂，而不是\Emph{阻止他的跌倒}。此外，先知们所预言的基督注定要被一个亲吻出卖，\BibleRef{路 22:47} 因为祂确实是那位人民的嘴唇尊崇者的儿子。\BibleRef{依 29:13} 当祂被带到公议会前，他们问祂是否是基督。\BibleRef{路 22:66} 犹太人还能问什么基督，除了他们自己的？那么，为什么祂甚至在那时刻不向他们宣告那位对立的基督呢？你回答：为了能够受苦。换句话说，为让这位最卓越的天主使人陷入罪恶，而这些人祂仍让他们保持无知。但即使祂告诉他们，祂仍要受苦。因为祂说：\Quote{我若告诉你们，你们也不会相信。} \BibleRef{路 22:67} 他们既然不信，就会继续坚持要祂死。而且，如果祂宣布自己是那位对立的天主所差遣的，因此是造物主的敌人，难道祂岂不更有可能要受苦吗？那么，祂在那关键时刻没有宣布自己是另一位\Emph{基督}，并不是为了受苦，而是因为他们想从祂口中逼供，即使祂给了他们，他们也不打算相信，而他们本应因祂所行的、应验他们圣经的工程而承认祂。因此，祂在那时刻保持不显明，显然是理所当然的，因为祂应得到自发的承认。但尽管如此，祂以庄严的姿态说：\Quote{从此，人子要坐在天主大能的右边。} \BibleRef{路 22:69} 因为祂是依据达尼尔先知的预言，向他们暗示祂是\Quote{人子}，\BibleRef{达 7:13} 并依据达味圣咏，暗示祂要\Quote{坐在天主的右边}。因此，在祂说了这话，并暗示要比较圣经之后，似乎有一线光明向他们显明祂要他们理解的祂是谁；因为他们说：\Quote{那么你是天主子吗？} \BibleRef{路 22:70} 除了他们独一知道的那位天主，还能是哪个天主？除了他们在圣咏中记念的那位曾对祂的儿子说\Quote{坐在我的右边}的天主，还能是哪个？然后祂回答：\Quote{你们说我是。} \BibleRef{路 22:70} 仿佛祂的意思是说：是\Emph{你们}说这话——不是我。但同时祂允许自己成为他们第二次问题中所说的一切。然而，你如何向我们证明，他们是以疑问的语气而非肯定地说了这句\Quote{\Emph{\OriginalTerm{那么你是天主子吗？}{Ergo tu filius Dei es}}}？正如（一方面）因为祂以间接的方式，通过圣经章节，向他们表明他们应将祂视为天主子，所以他们自己的话\Quote{那么你是天主子}应当以类似（间接）的意义来理解，等于说\Quote{你不愿明白地自称这话}，同样（另一方面）祂也回答他们\Quote{你们说我是}，其意义同样明确，甚至是肯定的；而祂的陈述如此完全地达到这个效果，以至于他们坚持接受祂陈述所指示的意义。\par

% structure: p0091 source=h2 target=subsection reason=relative-heading-rank; base=h2

\subsection{第42章 受难的其他细节与预言逐一比较。比拉多与黑落德。巴拉巴被偏爱胜过耶稣。十字架受刑的细节。地震与正午的黑暗。一切奇妙地预见于造物主的圣经中。基督交付灵魂并非\Person{马西昂}的\OriginalTerm{幻象论}{Docetic}观点的证据。在祂的安葬中存在着对此的反驳。}

因为当祂被带到\Person{比拉多}面前时，他们便以严重的指控来逼迫祂，即祂宣称自己是\Emph{基督君王}；这无疑是指祂是坐在天主右边的天主之子。然而，若他们不确定祂是否曾自称\Emph{天主之子}，他们本会用别的头衔来为难祂——如果祂不曾说出\BibleQuote{你说是的}这些话，以承认祂就是他们所说的那一位。同样，当比拉多问祂：\BibleQuote{你是基督（君王）吗？}祂便像从前（对犹太议会）那样回答说：\BibleQuote{你说是的}\BibleRef{路 23:3}，以免显得是因惧怕他的权力而给出更充分的回答。\Quote{主就这样受审。}祂也使祂的子民受审。主亲自与\Quote{民间的长老和首领}一同受审，正如\Person{依撒意亚}所预言的。然后祂应验了所有关于祂受难的记载。那时，\BibleQuote{外邦人咆哮，众民图谋虚妄；地上的君王齐集，首领一同聚首，反对上主，反对祂的受傅者。}\Emph{外邦人}是指比拉多和罗马人；\Emph{众民}是指以色列各支派；\Emph{君王}以\Person{黑落德}为代表，\Emph{首领}以司祭长为代表。确实，当祂被比拉多无偿地送到黑落德那里时，\BibleRef{路 23:7}\Person{欧瑟亚}的话便应验了，因为他曾预言基督说：\BibleQuote{他们要把他捆绑着，当作礼物献给君王。}黑落德看见耶稣，便\Quote{非常欢喜}，却未从祂听到一句话。\BibleRef{路 23:8}因为，\BibleQuote{他如同羔羊在剪毛者前缄口，默不作声}，\BibleRef{依 53:7}因为\Quote{上主赐给他一个受教的舌头，使他晓得如何以及何时应当说话}——甚至那\BibleQuote{贴在颚上的舌头}，如圣咏所说，因祂不言。然后，最邪恶的囚犯巴辣巴被释放，仿佛他是无辜的人；而最公义的基督被交付处死，仿佛他是杀人犯。\BibleRef{路 23:25}此外，有两个凶犯被钉在祂两旁，使祂被列于不法之徒之中。虽然祂的衣服无疑被士兵分开，一部分抽签分了，但\Person{马尔强}却把它们全部删去了（从他的福音中），因为他着眼那圣咏：\BibleQuote{他们分了我的外衣，为我的内衣拍阄。}你倒不如把十字架本身也除去！但即使如此，圣咏也没有对此沉默：\BibleQuote{他们刺穿了我的手和脚。}实际上，整个事件的细节都在那里读到：\BibleQuote{恶犬围困了我，恶党环绕了我。凡看见我的人都嗤笑我；他们张着嘴，摇着头，（说：）他倚靠上主，让上主救他吧。}现在（你篡改）祂衣服的见证有何用？你若把它当作你假基督的战利品，那整篇圣咏（仍弥补）基督的衣裳。但是，看，连元素都震动。因为它们的主正在受苦。然而，如果这所有的伤害是加在它们的敌人身上，天空会发光，太阳会更加耀眼，白昼会延长\BibleRef{苏 10:13}——高兴地注视马尔强的基督悬挂在他的刑架上！这些证据仍然适合我，即使它们不是预言的题材。\Person{依撒意亚}说：\BibleQuote{我要以黑暗遮盖天空。}\BibleRef{依 50:3}这就是那日子，关于这日子\Person{亚毛斯}也写道：\BibleQuote{在那一天，上主说，太阳要在正午落下，大地要在光天白日下变黑。}\BibleRef{亚 8:9}（正午）圣所的帐幔裂开了\BibleRef{路 23:45}，因革鲁宾的离去，\BibleRef{则 11:22}他们\BibleQuote{使熙雍女儿遗下，有如葡萄园中的草棚，有如瓜田中的茅庐。}\BibleRef{依 1:8}祂在\BibleBook{}{圣咏}第三十篇中怀着何等的坚定，努力将真正的基督呈现给我们！祂大声向天父呼喊：\BibleQuote{我将我的灵魂交在你手中。}这样，即使在死亡时，祂也可耗尽最后一口气以应验先知。说完这话，祂便断了气。\BibleRef{路 23:46}谁？是灵魂自己交出自己，还是肉体交出灵魂？但灵魂不能自己呼出。呼吸是一回事，被呼出是另一回事。如果灵魂被呼出，它必须由另一个呼出。然而，如果那里只有灵魂，就应当说是\Emph{离开}而非\Emph{断气}。但除了肉体，又有什么能呼出灵魂？肉体既在拥有灵魂时\Emph{呼吸}灵魂，又在失去灵魂时\Emph{呼出}灵魂。确实，如果那不是肉体（在十字架上），而是肉体的幻影（而幻影不过是灵魂，所以灵魂自己呼出自己，并随之离去），那么当作为幻影的灵魂离去时，幻影无疑也离去了：于是幻影和灵魂一起消失，无处可寻。因此，在\Quote{交出灵魂}之后，十字架上什么也没有留下，什么也没有悬挂；没有东西可向比拉多求取，没有东西可从十字架上卸下，没有东西可裹在殓布里，没有东西可安放在新坟墓中。然而，在那里的并不是\Emph{虚无}。那么有什么呢？如果有一个幻影基督还在那里。如果基督已经离去，祂也带走了幻影。留给异端者厚颜无耻的唯一托词，就是承认留在那里的是幻影的幻影！但如果\Person{若瑟}知道那是一个身体，并如此虔诚地处理它呢？就是那个\BibleQuote{并未同意}犹太人行恶的若瑟？\BibleRef{路 23:51}那个\BibleQuote{不随从恶人的计谋，不站罪人的道路，不坐亵慢者的座位}的有福之人。\par

% structure: p0093 source=h2 target=subsection reason=relative-heading-rank; base=h2

\subsection{第四十三章 结论。耶稣被证实为创造主的基督，依据\WorkTitle{路加福音}最后一章的事件。虔敬的妇女在坟墓旁。复活时的天使。基督复活后的多次显现。祂派遣宗徒到万民中。这一切都显示与全能圣父的智慧相符，正如先知所预言的。基督死后的身体非虚幻。马西昂在这一点上对\WorkTitle{福音}的操纵。}

埋葬主的那人理当在预言中被记念，并从此被称为\Quote{有福的}；因为预言并未省略那些天未亮就带着预备好的香料来到坟墓的妇女们的虔敬行为。\BibleRef{路 24:1}因为关于此事，欧瑟亚曾说过：\BibleQuote{他们必从清晨就寻求我的面，对我说：来，让我们归向上主，因为祂撕裂了我们，也必医治我们；祂打伤了我们，也必包扎我们；两天后祂必使我们复生，第三天祂必使我们起来。}因为谁能拒绝相信这些话常在这些妇女的思绪中反复，徘徊于她们似乎被主离弃的忧伤（她们当时感到被主打击）与复活本身的希望（她们正确地认为一切都将恢复）之间呢？但当\BibleQuote{她们找不到（主耶稣的）身体}\BibleRef{路 24:3}，\Quote{祂的坟墓从她们中间被挪去}，正如依撒意亚的预言。\BibleQuote{然而，有两位天使在那里显现}\BibleRef{路 24:4}因为天主的话正好要求这么多尊贵的伙伴，它通常规定\Quote{\Emph{两个见证人}}。此外，妇女们从坟墓回来，并离开这天使的显现，早被依撒意亚预见，他说：\BibleQuote{来吧，你们这些从异象回来的妇女们；}意即，\Quote{来吧，}去报告主的复活。然而，门徒们的不信如此固执，这是好的，为使我们至终能一致地坚持：耶稣向门徒启示自己正是先知们所预言的基督，而非其他。因为当其中两人正在走路时，主加入了他们，却没有显出是祂，并且祂假装不知道刚发生的事，\BibleRef{路 24:13}他们说：\BibleQuote{但我们原希望祂就是要救赎以色列的那一位}，\BibleRef{路 24:21}——意指他们自己的，即创造主的基督。祂远未向他们宣告自己是另一位基督！但他们不能认为祂是创造主的基督；而且，如果他们这样认为，祂也不能容忍这种关于祂自己的意见，除非祂真正是他们所认为的那位。否则，祂实际上就是错误的制造者，真理的欺骗者，与良善天主的本性相悖。但在祂复活之后，祂从未向他们启示为别的，而正是他们自己所表明一直认为祂是的。祂特意责备他们：\BibleQuote{无知的人哪！你们的心信得太迟钝了，不信祂对你们所说的一切。}\BibleRef{路 24:25}祂这样说，证明祂不属于敌对之神，而是属于同一天主。因为同样的话也由天使对妇女们说过：\BibleQuote{你们当记得祂还在加里肋亚时怎样对你们说过：人子必须被交付，被钉在十字架上，并在第三日复活。}\BibleRef{路 24:6}\BibleQuote{\Emph{必须}被交付；}为什么呢？岂不是因为这是创造主天主所记载的？因此祂责备他们，因为他们仅因祂的受难而跌倒，并且怀疑妇女们向他们报告的复活真相，由此（他们表明）不相信祂就是他们所认为的那一位。因此，为使他们如此相信祂，祂宣布自己正是他们所认为的——创造主的基督，以色列的救赎者。但关于祂身体的真实性，还有什么更明显的呢？当他们怀疑祂是否只是一个幻影——甚至认为祂是幻影时——祂对他们说：\BibleQuote{你们为什么惊慌？为什么心里起疑惑？看我的手脚，就知道是我自己；因为鬼神没有骨肉，如同你们看见我有的。}\BibleRef{路 24:37}马西昂不愿意从他的福音中删去一些甚至对他不利的陈述——我怀疑，他是刻意如此，以便从他未删去的经文中获得能力，当他本可以删去时，要么否认他删除了任何东西，要么为自己可能做出的删除辩护。但他只保留那些他能够同样通过曲解或删除来颠覆的经文。因此，在我们面前的这段经文中，他试图将\BibleQuote{鬼神没有骨，如同你们看见我有的}这句话倒置，意思变为\Quote{一个灵体，如同你们看到我这样，没有骨}；也就是说，灵体不具有骨的本性。但何必如此曲折的构造呢？祂本可以简单说：\Quote{鬼神没有骨，正如你们观察我没有一样？}再者，如果祂没有骨，为什么祂伸出手和脚给他们检查——这些由骨组成的肢体？为什么祂还加上\BibleQuote{知道是我自己}\BibleRef{路 24:39}，而他们先前已经知道祂是有形体的？否则，如果祂完全是一个幻影，为什么祂责备他们以为祂是幻影？但当他们仍不相信时，祂向他们要了一些食物，\BibleRef{路 24:41}，明确目的是向他们表明祂有牙齿。\par

现在，我敢相信，我们已经完成了我们的使命。我们表明\Person{耶稣基督}正是造物主的基督。我们的证据取自祂的教义、格言、情感、感受、奇迹、苦难，甚至复活——正如先知所预言的。直到最后，祂教导我们（祂使命的同一真理），当祂派遣宗徒去宣讲祂的福音\Quote{于万民之中；}因为祂如此应验了圣咏：\BibleQuote{他们的声音传遍天下，他们的言语直达地极。}\BibleRef{咏 19:4}。\Person{马尔基昂}，我怜悯你；你的劳作是徒然的。因为出现在你福音中的\Person{耶稣基督}正是我的那一位。\par

\SourceRecord{Translated by Peter Holmes.；Ante-Nicene Fathers；Vol. 3.；Edited by Alexander Roberts, James Donaldson, and A. Cleveland Coxe.；Buffalo, NY: Christian Literature Publishing Co.,；1885.；<http://www.newadvent.org/fathers/03124.htm>.}

% structure: p0001 source=h1 target=section reason=page-title

\section{\WorkTitle{驳斥马尔强，第五卷}}

\Emph{其中\Person{特土良}证明，就\Person{圣保禄}书信而言，他在前卷中关于\BibleBook{路加}{福音}所证明的，即它们非但不冲突，反而与\WorkTitle{旧约}著作完全一致，因而见证造物主是唯一的天主，\Person{主耶稣}是祂的基督。正如前几卷一样，\Person{特土良}以深刻的推理和许多对圣经的恰当比喻来支持他的论点。}\par

% structure: p0003 source=h2 target=subsection reason=relative-heading-rank; base=h2

\subsection{第一章 引言。\Person{宗徒保禄}本人并非新天主的宣讲者。他虽在其它宗徒之后蒙\Person{耶稣基督}召唤，但他的使命来自造物主。如何如此？论证，正如福音的情形一样，仅限于\Person{保禄}著作中\Person{马尔强}所接受的部分。}

除天主外，没有一物是没有起源的。既然起源在万物的状态中占据首位，那么若要对其状态获得清晰的认识，就必须在处理它们时优先考虑起源；因为你无法找到考察任何事物性质的方法，除非你确知它的存在，并且是在发现其起源之后。因此，既然我在我的小作品进程中到达这一点，我需要向马尔基昂询问他的宗徒的起源——我，在某种程度上是一个新门徒，不追随任何其他师傅；同时我什么都不相信，除非相信什么也不应轻信（而\Emph{那}，我还可以进一步说，就是没有查考其起源便相信的急遽相信）。简言之，我最有理由进行最仔细的探究，因为有人向我断言一个人是宗徒，而我在福音的宗徒名单中找不到他。确实，当我听说这个人是在主到达天上安息之后被选中的，我感到一种疏忽可以归咎于基督，因为祂以前不知道这个人对祂是必要的；并且祂认为必须以偶然相遇而非精心挑选的方式将此人加到宗徒团体中；可以说是出于必要，而非自愿选择，尽管宗徒职务的成员已经适当任命，并已被差派到各自的使命。因此，哦，本都的船主，如果你从未将任何违禁品或走私货物装上你的小船，如果你从未丢弃或擅自处理过货物，那么我毫不怀疑，在神圣的事上你更加谨慎和认真；所以如果你能告诉我们，你是在什么提货单下接受了宗徒保禄上船，谁给他打了标记，什么货主送了他，谁把他交给你，这样你就可以毫无疑虑地让他上岸，以免他属于那个人，而那个人可以通过出示他所有的宗徒著作来证明他的所有权。他自称是\BibleQuote{宗徒}——用他自己的话说——\BibleQuote{不是由于人，也不是藉着人，而是藉着耶稣基督。}（\BibleRef{迦 1:1}）当然，任何人都可以自我介绍；但他的介绍只有通过第二人的权威才有效。一人签署，另一人联署；一人盖印，另一人在公共记录中登记。没有人能同时是自己的提议者和附议者。此外，毫无疑问你读过：\BibleQuote{将有许多人来说：我是基督。}（\BibleRef{路 21:8}）如果有人能假装自己是基督，那么一个人更可以自称是基督的宗徒！但对我而言，我以门徒和探究者的身份出现；以便我既能以此驳斥你的信仰——你没有任何支持它的东西——又能羞辱你的无耻，你声称拥有却缺乏确立它们的手段。让有一个基督，让有一个宗徒，尽管是另一个神的；\Emph{但有什么关系呢？}既然他们只能从造物主的盟约中提取证明。因为早在很久以前，\BibleBook{}{创世纪}就向我许诺了宗徒保禄。因为在雅各伯对他子女所说的预象和预言祝福中，当他转向本雅明时，他喊道：\BibleQuote{本雅明如豺狼，早晨吞食猎物，晚间分给食物。}（\BibleRef{创 49:27}）他预见到保禄将从本雅明支派兴起，一只贪食的狼，早晨吞食猎物：换句话说，在他生命的早期，他将摧毁主的羊群，作为教会的迫害者；但在傍晚他将给予他们食物，这意味着在他晚年，他将作为外邦人的导师教育基督的羊栈。然后，在撒乌耳对待达味的行为中，首先表现为对他激烈的迫害，然后表现为懊悔和补偿，因为他从达味那里得到了以善报恶，我们无非看到了保禄在撒乌耳中的预表——他们同样属于同一支派——以及耶稣在达味中的预表，耶稣按照童贞女的族谱由达味下降。然而，如果你不赞同这些预象，至少\BibleBook{宗徒}{大事录}向我传达了保禄的这一历程，你不能拒绝接受。由此我证明，他从迫害者变成了\BibleQuote{宗徒，不是由于人，也不是藉着人}（\BibleRef{迦 1:1}）；由此我被引导相信\Emph{宗徒}本人；由此我找到理由驳斥你对他的辩护，并无所畏惧地承受你的嘲讽。\Quote{那么你否认宗徒保禄。}我并不诽谤我所辩护的人。我否认他，是为了迫使你证明他。我否认他，是为了让你相信他是我的。如果你尊重我们的信仰，你就应该承认构成它的细节。如果你向我们的信仰挑战，请告诉我们你的信仰的基础是什么。要么证明你所信仰的真理，要么在证明失败时，告诉我们你是如何相信的。否则，你的行为是什么？你相信与那位唯一能提供你所信仰之证据者对立。现在请从我的角度看待宗徒，就像你已经接受基督一样——宗徒被证明属于\Emph{我的}程度与基督相同。并且在这里，我们将在同一阵线上战斗，仅仅根据一个简单的规则挑战我们的对手：即一个据说不属于造物主——甚至表现出实际敌视造物主——的宗徒，可以被合理地认为不教导、不知道、不渴望任何有利于造物主的事，而他的首要原则就是极力提出另一个神，正如他极力使我们脱离造物主的法律一样。他绝不可能将人们从犹太教召出来，同时不向他们指明他所邀请他们信仰的神是谁；因为没有人能在不知道要过渡到谁的情况下离开对造物主的忠诚。因为要么基督已经揭示了另一个神——在这种情况下，宗徒的见证也会产生同样的效果，以免他不再被视为基督所启示之神的宗徒，并且因为他不应该被已经由基督启示过的宗徒所隐瞒——要么基督没有关于神的这种启示；\Emph{那么}宗徒就更需要启示一位现在无法被其他人认识的神，如果连宗徒都不启示他，那他无疑将完全没有任何信仰。我们以此作为首要原则，因为我们希望立即声明，我们将在宗徒的事上采用与基督的事上相同的方法，以证明他没有宣扬任何新神；也就是说，我们将从圣保禄本人的书信中提取证据。现在，我们在异端者的福音中发现的篡改形式已经让我们准备好期望他的书信也被他以同样的乖僻所截断——甚至就数量而言也是如此。\par

% structure: p0005 source=h2 target=subsection reason=relative-heading-rank; base=h2

\subsection{第二章 论\BibleBook{迦拉达}{书}。梅瑟法律规条的废除并非另一天主的证明。神圣的立法者，即创造主自己，就是废除者。宗徒在第一章的教导显示与旧约的教导一致。\BibleBook{}{宗徒大事录}的真实性驳斥\Person{马西昂}。本书与\BibleBook{}{保禄书信}一致。}

我们也认为对犹太主义最具决定性的一封书信，便是那位宗徒教导迦拉达人的书信。因为我们完全承认古代法律的废除，并且坚持这实际上出自创造者的安排——这一点在讨论过程中我们已经多次谈到，当指出革新是由我们天主的先知所预言的。现在，如果创造者确实许诺\Quote{旧事将要过去}，代之以将要出现的新事物，而基督在说\BibleQuote{法律及先知直到若翰为止}\BibleRef{路 16:16}时标记了分离的时期——以此使洗者若翰成为那终结的旧事与那开始的新事两种安排之间的界限——那么那位宗徒（他是在若翰之后启示的基督内）当然只能做一件事：否定\Quote{旧事}并肯定\Quote{新事}，并借此促进创造者的信仰，毕竟创造者曾预言旧事要过去。因此，法律的废除与福音的建立都有助于我的论点，即使在这封书信中，二者都关乎迦拉达人那种狂妄的假设，即他们以为信仰基督（当然是创造者的基督）是必要的，但不必废除法律，因为他们仍然觉得法律被其作者自己废除不可思议。此外，如果他们从宗徒那里听说过任何其他神，他们难道不会立刻自己得出结论，必须放弃他们所离开的那位神的法律，以便跟随另一位吗？因为一个人开始跟随新神后，他需要多久才会学会遵守新的纪律呢？然而，既然在福音中宣告的正是那位在法律中早已为人熟知的神，所改变的只是安排，那么唯一需要讨论的问题是，创造者的法律是否应当在创造者的基督里被福音排除？取消这一点，争论就无从谈起。现在，如果取消这一点，他们都会自己知道，当然必须因相信另一位神而放弃对创造者的所有服从，那么宗徒就没有必要如此热切地教导他们那些他们自己的信德自会暗示的事。因此，这一书信的全部主旨仅在于向我们显示，法律的取代来自创造者的任命——这一点我们仍需牢记。既然他也没有提及任何其他神（他不可能找到更合适的时机来这样做，尤其是当他旨在陈述废除法律的理由时——特别是新神的规范本可提供一个特别好的、最充分的理由），那么他写\BibleQuote{我奇怪你们这么快就离开了那召你们进入祂恩宠的那一位，而归向\Emph{另一个}福音}\BibleRef{迦 1:6}的意思就很清楚了——他指的是\Quote{另一个}在于所规定的行为，而不在于崇拜；\Quote{另一个}在于所教导的纪律，而不在于神性；因为基督福音的职责是召叫人们从法律进入恩宠，而不是从创造者到另一个神。因为没有人诱使他们背弃创造者，以致仅仅当他们重新回到创造者时，就似乎\Quote{被移到了另一个福音}。当他再加上\BibleQuote{并不是另一个}\BibleRef{迦 1:7}这句话时，他确认了他所维护的福音是创造者的。因为创造者自己应许了福音，当祂通过依撒意亚说：\Quote{你这位给熙雍报喜讯的，请登上高山；你这位给耶路撒冷报福音的，请竭力高呼。}另外，关于宗徒本人，祂说：\BibleQuote{那传报和平福音、传报喜讯的脚踪是多么美丽！}\BibleRef{依 52:7}甚至向外邦人宣讲福音，因为祂也说：\Quote{外邦人将要信赖祂的名字；}即基督的名字，祂对他说：\BibleQuote{我已立你作外邦人的光明。}\BibleRef{依 42:6}然而，你们会坚持认为当时宗徒所传的是一个新神的福音。这样就有两个神和两个福音；而宗徒说\BibleQuote{并没有另一个福音}\BibleRef{迦 1:7}就犯了大错，因为（按假设）还有一个；所以他本可以更好地辩护自己的福音，通过证明这一点，而不是坚持只有一个。但也许，为了避免这个困难，你会说他随后立即加上\BibleQuote{即使从天上来的天使，给你们宣讲另一个福音，当受诅咒}\BibleRef{迦 1:8}，因为他知道创造者要带来一个福音！但这样就使你更加纠缠不清。因为你现在陷入了网罗。断言有两个福音，不是一个已经否认有另一个的人会做的事。他的意思很清楚，因为他首先提到了自己（在诅咒中）：\BibleQuote{但是，即使是我们，或是从天上来的天使，给你们宣讲另一个福音。}\BibleRef{迦 1:8}他这是举例。即使他自己也不可宣讲另一个福音，那么天使也不行。他这样提到\Quote{天使}，是为了表明人更不该被相信，因为连天使或宗徒也不该被相信；而不是他要将天使应用于创造者的福音。然后他简要地谈到自己从迫害者转变为宗徒——以此证实\WorkTitle{宗徒大事录}，那本书中正是这封书信的主题，即某些人插话说应当受割损，并应遵守\Person{梅瑟}法律；以及宗徒们被咨询后，借着圣神的权威决定\BibleQuote{不要把重轭放在他们的颈上，连他们的祖先也不能负荷}。既然\WorkTitle{宗徒大事录}与\Person{保禄}如此一致，你们拒绝它们的原因就显而易见了。因为它们宣告的神除了创造者外没有别的，并证明基督属于创造者，而非任何别的神；而圣神的应许被证实实现了，除了在\WorkTitle{宗徒大事录}中再无别的文献。现在，不大可能的是，一方面这些记载与宗徒一致，描述他的行程符合他自己的陈述；另一方面，它们又在宣布创造者基督的神性时与他相矛盾——好像保禄没有接受宗徒们的宣讲，当他从他们那里接受\Quote{不要教导法律}的禁令时。\par

% structure: p0007 source=h2 target=subsection reason=relative-heading-rank; base=h2

\subsection{第三章　\Person{圣保禄}完全符合\Person{圣伯多禄}与其他受割礼的宗徒的立场。他对\Person{圣伯多禄}的责备得到解释，并从\Person{马西昂}的误用中救回。本封书信对犹太主义者的强烈抗议。然而其教导显示出与法律和先知的相符。\Person{马西昂}对\Person{圣保禄}著作的篡改受到谴责}

但关于伯多禄和其他宗徒的态度，他告诉我们，\BibleQuote{十四年后，他上耶路撒冷去}（\BibleRef{迦 2:1}），为与他们商议他所持守的福音准则，免得他那些年奔跑，以及仍在奔跑，都是枉然的——当然，如果他所传的福音不及他们的方式，就会如此。他多么渴望得到那些人的认可和支持，而你总想将他们理解为与犹太教结盟！实际上，当他说\BibleQuote{连希腊人弟铎也没有被强迫受割损}（\BibleRef{迦 2:3}），他首次向我们表明，割损是维护法律的问题，当时这问题已被那些他称为\Quote{潜入的假弟兄}（\BibleRef{迦 2:4}）所搅动。这些人只坚持继续遵守法律，无疑仍真诚地相信创造主。他们在教导中歪曲了福音，但并非以篡改圣经的方式抹去创造主的基督，而是通过保留\Emph{旧制度}，不排除创造主的法律。因此他说：\BibleQuote{因为有潜入的假弟兄，他们偷偷进来，要窥探我们在基督内所享有的自由，好使我们作奴隶；我们一时也没有顺服他们，臣服于他们。}（\BibleRef{迦 2:4}）我们只需注意清楚的意义和事情的缘由，圣经的曲解就会显明。当他首先说\BibleQuote{与我同去的弟铎，虽是希腊人，也没有被\Emph{强迫}受割损}（\BibleRef{迦 2:3}），然后加上\BibleQuote{却因着潜入的假弟兄}，他让我们洞察到他以一种完全相反的方式行动的理由，向我们表明他为什么做了那件事——若没有那件事促使他如此行，他既不会做也不会向我们显示。但我要你告诉我们：如果那些假弟兄没有潜入来窥探他们的自由，他们是否会对所要求的顺服让步？我想不会。他们之所以让步（部分让步），是因为有些人的软弱信德需要顾及。因为那些对遵守法律仍犹豫不决的初步信仰，理应得到这种通融的对待，甚至宗徒自己也怀疑他可能白白奔跑，并且仍在奔跑（\BibleRef{迦 2:2}）。同样，那些窥探基督徒自由的假弟兄必须受到挫败，他们想迫使自由服从于他们自己的犹太教，直到保禄发现他的劳苦是否徒然，直到那些在他之前作宗徒的人与他行右手相交之礼，直到他按照他们的安排开始向外邦人传教的职务。因此，他做出了一些暂时的必要让步；这就是为什么他让弟茂德受割损（\BibleRef{宗 16:3}），并带纳齐尔人进入圣殿（\BibleRef{宗 21:23}），这些事在《宗徒大事录》中已有记载。这些事的真实性可从与宗徒自己的声明一致推断出来，他说\Quote{对犹太人，我就成为犹太人，为赢得犹太人；对在法律下的人，我就成为在法律下的人}——此处也是如此，对那些潜入的人——\Quote{最后，对一切人，我成了万有，为能救一切人。}既然情形需要这样的解释，没有人会否认保禄所宣讲的天主及基督，正是他始终在排除其法律的那一位，尽管限于时期，他允许了它，但若他宣布一个新神，他早就该彻底废除了。因此，伯多禄、雅各伯和若望正确地向保禄伸出右手相交之礼，并同意如此分工：保禄向外邦人，而他们向受割损的人（\BibleRef{迦 2:9}）。他们同意\BibleQuote{记念穷人}（\BibleRef{迦 2:10}）也完全符合创造主的法律，该法律照顾穷人和匮乏者，正如我们在评论你们的福音时所指出的。因此，问题显然只涉及法律，同时，什么样的法律部分便于遵守也很明显。然而，保禄责备伯多禄没有按福音的真理正直地行走。无疑他责备他；但仅是因为他在\Quote{饮食}上的不一致，他随着不同的人而变化，\Quote{怕那些受割损的人}（\BibleRef{迦 2:12}），而不是由于有关另一位神的错误见解。因为若这样的问题出现，其他人也会被那连伯多禄在小事上也不放过的人\Quote{当面抵抗}。但马西昂派希望相信什么呢？另外，必须允许宗徒继续他自己的声明，他说\BibleQuote{人不是因法律的作为成义，而是因信德}（\BibleRef{迦 2:16}）：信德，当然是在同一位天主的信德，法律也属于祂。因为若果真有另一位神，神的不同自然会产生这样的分离，他就不会费力将信德从法律中分离。因此，他正确地拒绝\BibleQuote{重建他所拆毁的（法律结构）}。确实，法律必须被拆毁，从若翰\BibleQuote{在旷野里喊说：预备主的道路，修直祂的途径}（\BibleRef{路 3:4}）的时候起，以使山谷、小山和山岗都得填满和铲平，弯曲和崎岖的道路变为正直平坦——也就是说，使法律的困难变为福音的容易。\par

因为他记起那时刻已到，就是圣咏所说的：\BibleQuote{让我们挣断他们的绳索，从我们身上摆脱他们的轭。}自那时起，\BibleQuote{列邦喧嚷，万民图谋虚妄的事；}\BibleQuote{地上的君王起来，掌权者聚集一起，敌对主和祂的基督。}为的是从此以后，人因信德的自由而成义，而非因法律的奴役，因为\BibleQuote{义人必因他的信德而生活。}\BibleRef{哈 2:4}现在，虽然\Person{哈巴谷}先知首先说了这话，但你在此见到宗徒印证先知，正如\Person{基督}所做的一样。因此，义人将赖以生活的信德，其对象必是同一的天主，法律也属于祂；遵行法律，无人能成义。既然在法律中同样有诅咒，在信德中也有祝福，你就看到这两种情形都由造物主陈明：祂说：\BibleQuote{看，}祂说，\BibleQuote{我已在你们面前设置了祝福和诅咒。}\BibleRef{申 11:26}你不能因事情合一就确立作者的多样性；因为多样性本身就是由同一作者提出的。然而，为什么\BibleQuote{基督为我们成了诅咒，}\BibleRef{迦 3:13}宗徒自己以相当有助于我们这边的方式宣告，这是造物主安排的结果。但这决不意味着，因为造物主古时说过：\BibleQuote{凡被挂在树上的，都是可诅咒的}，\Person{基督}就属于另一位神，并因此甚至在法律中就已受诅咒。而且，造物主如何能预先诅咒一位祂所不认识的呢？然而，为何造物主将自己的儿子交付给祂自己的诅咒，而不是将祂置于你那位神的咒诅之下——而且还是为了人，而人是与那位神不相干的？这样岂不更合适吗？现在，如果造物主关于祂儿子的安排在你看来是残忍的，那么在你自己的神那里也同样残忍；反之，如果在你神那里是合理的，那么在我这里同样合理——甚至更加合理。因为那位曾将祝福和诅咒都摆在人面前的天主，通过\Person{基督}的诅咒为人预备了祝福，这比那位据你所说从未宣告过两者中任何一位的神这样做，更可信。\BibleQuote{因此，我们领受了圣神的恩许，}正如宗徒所说，\BibleQuote{藉着信德，}就是那义人赖以生活的信德，按造物主的旨意。那么，我的意思是：那位预像了信德之恩宠的天主，才是信德的对象。但当他还加上：\BibleQuote{因为你们都是信德的子女，}\BibleRef{迦 3:26}就显明异端者的劳作所删去的是\Person{亚巴郎}名字的提及；因为藉着信德，宗徒宣告我们是\BibleQuote{\Emph{亚巴郎的子孙}}，并在提及他之后，也明确称我们为\BibleQuote{信德的子女}。但我们如何是信德的子女？若不是\Person{亚巴郎}的信德，又是谁的？因为\BibleQuote{\Person{亚巴郎}信了天主，天主就以此算为他的义德；}\BibleRef{迦 3:6}也因为他因此当得起被称为\BibleQuote{万民之父}。而我们，在信天主的事上更像他，因此像\Person{亚巴郎}一样被成义，也因此获得生命——因为义人因信德而生活——所以就发生了这样的事：正如他在前面称我们为\BibleQuote{亚巴郎的子孙}，因为他在信德上是我们的（共同）父亲，同样在这里他也称我们为\BibleQuote{信德的子女}，因为是由于他的信德，才应许\Person{亚巴郎}要成为（万）民之父。至于他使信德脱离割损礼这一事实本身，难道不正是借此使我们成为\Person{亚巴郎}的子女吗？\Person{亚巴郎}在肉身受割损之前就已经相信了。简言之，对两位神中一位的信德，不可能使我们得以进入另一位神的安排，以致那信德将义德归给信祂的人，并藉着祂使义人生活，又宣告外邦人因信德成为祂的子女。这种安排完全属于那位天主，藉着祂的安排，通过这同一个\Person{亚巴郎}的召唤早已显明，这一点由自然意义确凿地证明。\par

% structure: p0010 source=h2 target=subsection reason=relative-heading-rank; base=h2

\subsection{第四章。\Person{马西昂}篡改\Person{圣保禄}文本的另一实例。时期的圆满，由宗徒宣告，由先知预述。\Person{梅瑟}礼制由造物主亲自废除。\Person{马西昂}关于\Person{亚巴郎}名字的诡计。造物主藉祂的\Person{基督}，是\Person{圣保禄}所宣告的恩宠与自由的泉源。\Person{马西昂}的\OriginalTerm{幻象论}{Docetism}被驳斥。}

\Quote{但，}他说，\Quote{我按人的方式来说：当我们还是孩童时，我们被置于世界的基本要素的奴役之下。}然而，这并非\Quote{按人的方式}说的，因为这里没有比喻，而是字面的真理。因为（就本节的后半部分而言），哪一个孩童（指外邦人作为孩童的意义上）不是被奴役于世界的基本要素，并视之为神明呢？至于前半部分，确实\Emph{有}一个比喻（如宗徒所写的），因为他在说了\Quote{我按人的方式来说}之后，又加上：\Quote{即使是人的盟约，也没有人废除或增添。}因为他藉着人的盟约的永久性这一比喻，来辩护天主的盟约。\Quote{恩许是向亚巴郎和他的后裔说的。他没有说\InnerQuote{后裔们}，如同指许多人；而是如同指一个，\InnerQuote{你的后裔}，就是基督。}\BibleRef{迦 3:16}马尔强擦掉的部分，我何必多费笔墨呢？其实从他保留的部分，更能有效地驳倒他。\Quote{但时期一满，天主就派遣了自己的圣子来}\BibleRef{迦 4:4}——这位天主当然是那整个时间序列的主，这个序列构成了一个时代；祂也指定了作为时间\Quote{\Emph{标记}}的太阳、月亮、星座和星辰；祂还预定并预言了祂圣子的启示将延迟到时期的终结。\Quote{在\Emph{最后的日子}，上主（殿）的山必要彰显；}\Quote{在\Emph{最后的日子}，我要将我的圣神倾注在一切血肉上}，如岳厄尔所说。唯有祂才有资格耐心等待时期的圆满，因为时期终结与开始同属祂。至于那个无所事事的神祇，既无作为，也无预言，因而也无时间可言，\Emph{他}究竟做了什么来促成时期的圆满，或耐心等待其完成呢？如果什么也没做，那么他只能卑微地等待造物主的时间，这是何等无能的状态！但祂派遣祂的圣子是为了什么目的？\Quote{为赎回那在法律之下的人}\BibleRef{迦 4:5}，换言之，为使\Quote{弯曲的路径变直，崎岖的道路变为平坦}，如依撒意亚所说\BibleRef{依 40:4}——为使旧事过去，新秩序开始，即\Quote{新的法律出自熙雍，上主的话来自耶路撒冷}\BibleRef{依 2:3}，并且\Quote{使我们获得义子的地位}\BibleRef{迦 4:5}，也就是外邦人，他们原本不是儿子。因为祂要做\Quote{外邦人的光明}，\Quote{外邦人将寄望于祂的名。}\BibleRef{依 42:4}因此，为使我们有把握自己是天主的儿女，\Quote{祂派遣了祂的圣神到我们心中，呼喊说：\InnerQuote{阿爸，父啊！}}\BibleRef{迦 4:6}因为\Quote{在最后的日子}，祂说，\Quote{我要将我的圣神倾注在一切血肉上。}\par

那么，这恩宠从谁而来，岂非从那宣告此恩许者而来？谁是我们的圣父，岂非那同时也是我们的创造者？因此，在如此丰盛（的恩宠）之后，他们不该再回到\Quote{软弱贫乏的\OriginalTerm{元素}{elements}}。\BibleRef{迦 4:9} 然而，罗马人惯于将学习的初基称为\OriginalTerm{元素}{elements}。因此，他并非试图通过贬低世俗的元素，使人们远离他们的神；尽管他刚说过：\Quote{然而，那时你们事奉那些本来不是神的}，\BibleRef{迦 4:8} 他指责的是那种认为元素就是神的物理或自然迷信的错误；但他并未在指责中针对那些元素的天主。他本人清楚地告诉我们，他所说的\Quote{\OriginalTerm{元素}{elements}}是什么意思，即法律的初阶：\BibleQuote{你们遵守日子、月份、节期、年份}\BibleRef{迦 4:10}——我想是指安息日、\Quote{预备日}、斋戒和\Quote{大节期}。就连这些的停止，如同割损一样，都是由创造者的谕令所规定的；他曾通过依撒意亚说：\BibleQuote{你们的月朔、安息日和大节期，我无法忍受；你们的斋戒、庆节和典礼，我的灵魂憎恶}\BibleRef{依 1:13}；也通过亚毛斯说：\BibleQuote{我憎恨、厌恶你们的节期，也不闻你们庄重集会的香气}\BibleRef{亚 5:21}；又通过欧瑟亚说：\BibleQuote{我要止息她的一切欢乐、节期、安息日、月朔和她的一切庄重集会}\BibleRef{欧 2:11}。他所设立的制度，他自己就毁掉了吗？是的，与其让别的任何一位来毁掉。或者，若另一位毁掉了它们，他不过是帮助了创造者的目的，移除了连他自己也曾谴责的东西。但这里不是讨论为什么创造者废除自己的法律的问题。向我们而言，足以证明他有意如此废除，因此可以断言，宗徒并未决定任何有损创造者的事，因为废除本身出自创造者。但是，正如窃贼的情况，赃物有时会沿途掉落，成为抓获他们的线索；在我看来，这事也发生在\Person{马尔强}身上：他保留了书信中最后提到\Person{亚巴郎}名字的地方（尽管没有哪段经文更需要他删除），甚至部分改动经文。\BibleQuote{因为经上记载，亚巴郎有两个儿子，一个出于婢女，一个出于自由妇；但那出于婢女的，是按肉身生的；那出于自由妇的，是藉恩许生的：这些都是\OriginalTerm{寓意化}{allegorized}了的}（也就是说，它们预示了超越\OriginalTerm{字面历史}{literal history}的意义）；\BibleQuote{因为这两个是两种盟约}，或是两种（神圣计划的）显现，正如我们发现这个词的解释，\BibleQuote{一种出自\Place{西乃山}}，与犹太人的会堂相关，按照法律，\BibleQuote{那产生奴役的}——\BibleQuote{另一种产生}（自由，被提升）超越一切率领者、掌权者、宰制者，以及一切可称呼的名衔，不仅今世，也来世，\BibleQuote{她是我们众人的母亲}，我们在其中拥有（基督）神圣教会的恩许；因此他在结论中补充说：\BibleQuote{所以，弟兄们，我们不是婢女的儿女，而是自由妇的儿女}。在这段话中，他无疑表明，基督徒有高贵的出生，正如寓意的奥秘所暗示的，是从亚巴郎那出于自由妇的儿子而生；而来自婢女的儿子则产生了犹太主义的法律奴役。因此，两种分施都出自同一位天主，我们发现，二者早已由他预先描绘。当他谈到\BibleQuote{基督使我们获得自由的那种自由}时，\BibleRef{迦 5:1} 这个短语本身岂不表明那解放者就是那曾是主人的那位？因为\Person{加尔巴}本人从不解放不属于自己的奴隶，即使他即将使自由人重获自由。因此，自由必将由那一位赐予，法律的奴役权曾听命于他。这是非常恰当的。那些已获得自由的人不应再\BibleQuote{被奴役的轭缠住}\BibleRef{迦 5:1}——即法律的轭；如今那圣咏的预言已经应验：\BibleQuote{让我们挣开他们的捆绑，摆脱他们的绳索，因为地上的统治者聚集起来，敌对主和他的受傅者}。因此，所有从奴役之轭下被释放的人，他都恳切地要求他们消灭奴役的标记——即割损，依据先知的预言。他记起耶肋米亚曾说过：\BibleQuote{割损你们的心包皮}\BibleRef{耶 4:4}；同样，梅瑟也命令：\BibleQuote{割损你们坚硬的心}\BibleRef{申 10:16}——而不是\Emph{字面}的肉体。那么，如果他是新神的使者，要排除割损，为什么他说\BibleQuote{在基督内，割损不算什么，不割损也不算什么呢？}\BibleRef{迦 5:6} 因为如果他受命于那敌视割损的神，他就有责任优待他所废除之事的对立原则。\par

此外，既然割损与不割损都归于同一天主，两者在基督内都失去了效力，这是由于信德的卓越——就是关于那信德，经上曾记载：\BibleQuote{外邦人将因祂的名而信赖吗？}\BibleRef{依 42:4}——就是那信德，他说，是\BibleQuote{藉爱德而运作的}\BibleRef{迦 5:6}。藉此他说表明，那创造主是那恩宠的源头。因为无论他所说的爱是由于天主的爱，或是那爱近人的爱——无论哪种情形，所指的都是创造主的恩宠：因为正是祂以这些话命令了前者：\BibleQuote{你要全心、全灵、全力爱天主你的天主；}\BibleRef{申 6:5}并在另一处也命令了后者：\BibleQuote{你要爱你的邻人如同你自己。}\BibleRef{肋 19:18}\BibleQuote{但那扰乱你们的，必要承受审判。}\BibleRef{迦 5:10}来自哪一位天主？来自（马西昂的）至善之神吗？但他并不执行审判。来自创造主吗？但祂也不会定罪那维持割损的人。如今，如果除了创造主以外没有别神被断定执行审判，那么必然只有那位决定了法律终止的，才能定罪法律的捍卫者；而且，如果他也在法律当保持常存的部分中肯定法律，那又如何呢？他说：\BibleQuote{因为全部法律在你们身上由这一句话得以成全：\InnerQuote{你要爱你的邻人如同你自己。}}\BibleRef{迦 5:14}如果，他执意认为通过\Quote{\Emph{得以成全}}这词句，是暗示法律不再\Emph{需要}被遵守，那么他当然不再意指我应当爱邻人如同自己，因为这条诫命必须随同法律一同废止。但不！我们必须继续永久遵守这条诫命。因此，创造主的法律获得了那敌对神的认可，他实际上赐予它的并非草率驱逐的判决，而是精简接纳的恩惠；其全部精义都集中在这一条诫命上了！但这法律的浓缩，事实上只有制定它的那位才可能做到。因此，当他说：\BibleQuote{你们要彼此分担重担，这样就成全了基督的法律}\BibleRef{迦 6:2}，既然除非人爱邻人如同自己，否则这事不能成就，那么显然\BibleQuote{你要爱你的邻人如同你自己}这条诫命——它事实上是\BibleQuote{彼此分担重担}这命令的基础——实在是\BibleQuote{基督的法律}，尽管从字面看是创造主的法律。所以，基督是创造主的基督，正如基督的法律是创造主的法律。\BibleQuote{你们不要被欺骗，天主是轻慢不得的。}\BibleRef{迦 6:7}但马西昂的神却\Emph{可以}被轻慢，因为他不知道如何发怒，或如何报仇。\BibleQuote{因为人种什么，也就收什么。}\BibleRef{迦 6:7}那么，正是那赏报和审判的天主以此警告。\BibleQuote{让我们行善不可厌倦}\BibleRef{迦 6:9}，并\BibleQuote{一有机会，就要行善}\BibleRef{迦 6:10}。现在，如果你否认创造主曾下达行善的诫命，那么诫命的差异就可能支持神祇的不同。但是，如果祂也宣告赏报，那么死亡和生命的收成都必须来自同一位天主。但\BibleQuote{到了时候，我们必要收获}\BibleRef{迦 6:9}，因为在\BibleBook{}{训道篇}上说：\BibleQuote{凡事都有定时}\BibleRef{训 3:17}。此外，\Quote{世界已被钉在十字架上与我}，我身为创造主的仆人——\Quote{世界}，我说，但不是造世界的天主——\Quote{我也与世界一同被钉}\BibleRef{迦 6:14}，不与造世界的天主。宗徒所说的\Emph{世界}，在此指点按照世俗原则的生活与交往；正是在弃绝这些时，我们与它们彼此钉死、彼此杀害。他称他们为\Quote{迫害基督的人}。但当他补充说，他\Quote{在身上带着基督的\Emph{伤痕}}——因为伤痕当然是身体的偶性——他就因此表达了真理：基督的肉身不是虚假的，而是真实和实在的，他将那伤痕表现为在他身上所受的。\par

% structure: p0014 source=h2 target=subsection reason=relative-heading-rank; base=h2

\subsection{第五章 \BibleBook{格林多}{前书}。保禄的恩宠与平安问安被证明是反马西昂的。基督的十字架是创造主所预定的。马西昂只是以不敬地将福音与创造主分离，而延续了基督十字架的绊脚石和愚妄。法律与福音在软弱、愚拙和卑贱之事上的类比。}

我的初步评论在前一封书信中使我未能处理其标题，因为我相信会有另一个机会来考虑此事，它在每封书信中都反复出现，且形式相同。那么，关键就是：宗徒为他所写信的人指定的不是通常的\Emph{健康}，而是\Quote{恩宠与平安}。\BibleRef{格前 1:3}我确实不问，一位摧毁犹太主义的人与犹太人仍在使用的一个公式有何关系。因为直到今日，他们还用\Quote{平安}的问候彼此致意，并且从前在他们的圣经中他们也这样做。但我从他的做法清楚地理解到，他确认了造物主的宣告：\Quote{那传布喜讯、宣扬和平福音者的脚步是多么美丽啊！}\BibleRef{依 52:7}因为传报\Emph{福音}——即天主的\Quote{恩宠}——的使者很清楚，连同\Quote{平安}也要被宣扬。现在，当他宣告这些祝福是\Quote{从天主圣父和主耶稣基督而来}时，\BibleRef{格前 1:3}他使用了双方共有的称号，这些称号也适应我们信仰的奥迹；我认为，如果不考虑它们各自更合适的属性，就不可能准确地确定哪个天主被宣告为圣父和主耶稣。首先，我断言，除了人和宇宙的创造者与维持者之外，没有其他可以被称为\Emph{圣父}和\Emph{主}；其次，圣父也因其大能而拥有主的称号，圣子也通过圣父接受同样的称号；然后，\Quote{恩宠与平安}不仅属于发布它们的那位，也属于被冒犯的那位。因为\Emph{恩宠}只有在冒犯之后才存在，\Emph{平安}也只有在战争之后才存在。现在，子民（以色列人）因违反祂的法律而犯罪，整个人类因忽视自然义务也犯罪，他们都曾对造物主犯罪和反叛。然而，\Person{马西昂}的天主不可能被冒犯，因为他既不为任何人所知，又不可能被激怒。那么，对于一个从未被冒犯的天主，能有什么\Emph{恩宠}呢？对于一个从未经历过反叛的天主，能有什么\Emph{平安}呢？他说：\Quote{十字架的道理，为丧亡的人是愚妄，但为我们得救的人，却是天主的德能和天主的智慧。}\BibleRef{格前 1:18}然后，为让我们知道这从何而来，他补充说：\Quote{因为经上记载：\InnerQuote{我要摧毁智者的智慧，废除明达者的明达。}}现在，既然这些话是造物主的，并且他把与十字架道理相关的事视为愚妄，那么十字架，以及因十字架而有的基督，都将属于造物主，因为十字架的事件是由祂预言的。但如果造物主作为敌人夺走了他们的智慧，以使基督的十字架被视为愚妄——基督被视为祂的对手——那么造物主在发表预言时，怎么能对一位不属于祂的、祂毫不知情的基督的十字架有任何预言呢？但是，再问：在一个如此善良、如此慈悲的主的体系里，有些人获得救恩，因为他们相信十字架是天主的德能和智慧；而另一些人却遭受丧亡，因为他们认为基督的十字架是愚妄——（这怎么可能，我再说一遍）——除非这是出自造物主的安排，祂惩罚了\Emph{以色列}子民和人类，因为他们对祂犯下了某种大罪，剥夺了他们的智慧和明达？接下来的话将证实这个建议，当他问：\Quote{天主岂不是使这世界的智慧变成了愚妄吗？}\BibleRef{格前 1:20}并补充原因：\Quote{因为世人没有凭自己的智慧认识天主，天主就乐意用愚妄的宣讲来拯救那些相信的人。}\BibleRef{格前 1:21}但首先说一下\Quote{\Emph{世界}}这个词；因为在这方面，异端者特别费尽心机地表明，\Emph{世界}指的是\Emph{世界的主宰}。然而，我们理解这个词是指世上任何一个人，按照人类语言的简单习惯，即用容器指代内容物。\Quote{马戏团在喊叫，}\Quote{论坛在发言，}\Quote{大教堂在低语，}都是常见的表达，意思是这些地方的人们如此行事。既然世上的人（而不是天主）凭自己的智慧没有认识天主——他们本应认识祂（犹太人通过圣经知识，全人类通过认识天主的工程）——因此，那位在祂的智慧中未被认识的天主，决意用祂的愚妄来击打人的知识，拯救所有相信所宣讲的十字架的愚妄的人。\Quote{因为犹太人要求神迹，}——他们本应对天主已有定见——\Quote{希腊人寻求智慧，}\BibleRef{格前 1:22}——他们依赖自己的智慧，而不是天主的。然而，如果所宣讲的是一位新天主，犹太人要求神迹才能相信，或者希腊人追寻他们更愿意接受的智慧，这又有什么罪过呢？因此，临到犹太人和希腊人身上的报应本身证明了天主既是忌邪的天主又是审判者，因为祂以愤怒和审判的报应使世界的智慧变为愚妄。既然这些原因在于那位赐给我们所用圣经的那一位，那么宗徒在论述造物主时——祂是犹太人和外邦人尚未认识的——无疑是要教导我们，那位（在基督内）将要被认识的天主就是造物主。他宣告基督是\Quote{犹太人的绊脚石}\BibleRef{格前 1:23}，这明确指向造物主关于祂的预言，当祂通过\Person{依撒意亚}说：\Quote{看，我在熙雍安放了一块绊脚石，一块使人跌倒的磐石。}\BibleRef{依 8:14}这磐石或石头就是基督。\BibleRef{依 28:16}这块绊脚石\Person{马西昂}仍然保留。那么，什么是\Quote{天主的愚妄比人更智慧}呢？不就是基督的十字架和死亡吗？什么是\Quote{天主的软弱比人更刚强}\BibleRef{格前 1:25}呢？不就是天主的诞生和降生成人吗？然而，如果基督不是由童贞女所生，不是由人类血肉构成，因而没有真正经历死亡和十字架，那么祂身上就没有任何愚妄或软弱；那么，\Quote{天主偏选了世上愚妄的，为羞辱那有智慧的；}\Quote{天主偏选了世上软弱的，为羞辱那坚强的；}\Quote{天主偏选了世上卑贱的}和微小的\Quote{，以及那些一无所有的（即那些实际上不存在的事物），为废除那些存在的（即那些实际存在的事物），}\BibleRef{格前 1:27}就不再真实了。因为在天主的安排中，没有任何事物是卑微、不尊贵和可轻视的。这些只出现在人的安排中。造物主的\WorkTitle{旧约}本身，无疑也可能被指控为愚妄、软弱、不体面、卑贱和可轻视。还有什么比天主要求血祭和芬芳的全燔祭更愚妄、更软弱？还有什么比清洗器皿和床铺更软弱？还有什么比发红的皮肤病变色更不体面？\BibleRef{肋 13:2}还有什么比以眼还眼的律法更卑贱？还有什么比食物和饮料的禁忌更可轻视？我相信，整个\WorkTitle{旧约}都受到异端者的嘲笑。因为天主偏选了世上愚妄的，以羞辱它的智慧。\Person{马西昂}的天主没有这样的训练，因为他没有效法（造物主）以对立面来混淆对立面的过程，以致\Quote{没有人可以夸耀；但正如经上记载：夸耀的，应当因主而夸耀。}在哪一位主呢？当然是在那位颁布这条诫命的主。除非，果然，造物主命令我们因\Person{马西昂}的天主而夸耀。\par

% structure: p0016 source=h2 target=subsection reason=relative-heading-rank; base=h2

\subsection{智慧、伟大与大能的天主之道。天主的自我隐藏与后来的启示。对于马西昂的天主，这样的隐藏与显现是不可能的。天主的预定。对于先前未知的天主（如马西昂的），不可能有目的的先存系统。钉死基督的世间权势。圣保禄作为智慧的建筑师，与预言联合。宗徒的诸多训导与旧约教训平行。}

因此，通过这些陈述，他向我们指明他所意指的天主，因为他说：\BibleQuote{我们在成全的人中讲论天主的智慧。}\BibleRef{格前 2:6}正是这位天主，使智慧人的智慧变成愚拙，使明达人的聪明归于无有，使世上的智慧成为愚妄，拣选了世上愚拙的事物，安排它们以达成救恩。他说，这智慧一度隐藏在愚拙、软弱和无荣耀的事物中；也一度潜藏在预像、寓意和谜语般的类型中；但后来要在基督内被启示出来，基督被安置为\BibleQuote{作为外邦人的光明}\BibleRef{依 42:6}，这是创造者通过依撒意亚的口所应许的，祂要发现\BibleQuote{眼睛未曾见过的隐藏宝藏。}那么，那个天主若从未做过隐藏的遮盖，却曾隐藏任何事物，这本身就是完全难以置信的想法。如果它存在，它对自己的隐藏是不可能的——更不必说它的任何宗教规条了。相反地，创造者本身正如祂的法令一样为人熟知。我们知道，这些法令在以色列公开设立；但它们被隐含的意义所遮蔽，天主的智慧就隐藏在其中，直到时候来到，要在\Quote{成全的人}中逐渐显明，而这些是\BibleQuote{天主在万世之前所预定的计划}\BibleRef{格前 2:7}。但这些万世若不是属于创造者的，又是谁的？因为万世由时间组成，时间由日、月、年构成；而且日、月、年由太阳、月亮和星宿来度量，这些是祂为此目的而规定的（因为祂说：\BibleQuote{它们要作为月份和年岁的标记}），由此可见万世属于创造者，并且凡在万世之前所预定的一切，不能说是属于任何其他存在，只能属于那位同样宣称万世为其己有的存在。否则，让马西昂证明万世属于他的天主。那么他也必须为他的天主要求世界本身；因为万世是在世界内计算的，世界仿佛是时间的容器，以及它们的标记或顺序。但他没有这样的论证可以给我们看。因此我回到原题，问他这个问题：为什么（他的天主）在创造者的万世之前就预定我们的荣耀？如果他在时间之初就启示了它，我能够理解他是在万世之前预定了它。但是当他几乎在创造者的所有万世将近结束时才这样做，他在万世之前的预定（而不是在万世之内）是徒劳的，因为他无意在万世几乎跑完其行程之前启示他的目的。因为他在计划目的时如此急切，在启示这些目的时却如此迟缓，这完全是自相矛盾的。\par

然而，在创造主那里，这两条路线完全相容——既在万世之前预定，又在末世启示出来，因为祂所预定和启示的，也在中间的时间借着预象、象征和寓言的预先服务而宣告出来。但因为（宗徒）在我们荣耀的事上补充说：\BibleQuote{今世掌权者没有一个知道这智慧；因为若是知道了，他们决不会把荣耀的主钉在十字架上。}\BibleRef{格前 2:8} 异端者就认为，今世掌权者钉死了主（即敌对之神的基督），是为了让这一 \Emph{打击} 甚至反作用到创造主自己身上。但是，任何从我们已经说过的话中看出我们的荣耀必须被视为出自创造主的人，就已经得出结论：既然创造主按自己隐秘的旨意定下了这荣耀，它理所当然不为创造主所有的掌权者和权势所知，因为仆人不允许知道主人的计划，何况堕落的天使和罪恶的领袖本身——魔鬼；因为我会主张，\Emph{这些}（后者）由于他们的堕落，对创造主任何的安排更加陌生。但我甚至可以不再将今世的掌权者和权势解释为属于创造主，因为宗徒将无知归咎于\Emph{他们}，然而按照我们的福音，连魔鬼在诱惑中也认出了\Person{耶稣}（\BibleRef{玛 4:1}），而且根据双方（马西昂派和我们）都共有的记载，邪魔知道\Person{耶稣}是天主的圣者，知道\Person{耶稣}是祂的名字，知道祂来是要毁灭他们（\BibleRef{路 4:34}）。而且关于武装的壮士，被一位比他更强壮者战胜并夺取其财物的比喻，\Person{马西昂}也承认是指创造主：因此，创造主不可能再对荣耀的主无知了，既然祂被他战胜；祂也不能钉死那位祂无法对付的。因此，在我看来，必然的推论就是，我们必须相信创造主的掌权者和权势确实是明知地钉死了在祂基督内的荣耀的主，带着那种最堕落的奴隶甚至不惜杀害主人的绝望和过分的恶意。因为在我的福音中记载着：\BibleQuote{撒殚进入了\Person{犹达斯}。}\BibleRef{路 22:3} 但是，根据\Person{马西昂}，宗徒在所考虑的段落（\BibleRef{格前 2:8}）中不允许将关于荣耀之主的无知归咎于创造主的权势；因为他确实主张，这些并不对应于\Quote{今世掌权者}。但（宗徒）显然不是指属灵的掌权者；所以他是指世俗的掌权者，即那些属于选民的掌权者（虽然是天主经纶中的首领，但当然不是世上的外邦人），以及他们的统治者，\Person{黑落德}王，甚至\Person{比辣多}，以及由他所代表的那时世上最大的罗马权势，当时正由他掌管。这样，对方的主张就被推倒了，我们自己的证据也就建立起来了。但你仍然坚持我们的荣耀来自你的神，那荣耀也隐藏在你神那里。那么，为什么你的神使用宗徒也引用的同一部圣经？你的神与先知的话有什么相干？\BibleQuote{谁曾知晓主的心意，或谁曾作过祂的参谋？}\BibleRef{依 40:13}\Person{依撒意亚}如此说。他（\Person{保禄}）又怎么与来自我们天主的例证有关呢？因为当（宗徒）自称\BibleQuote{一个聪明的工头}\BibleRef{格前 3:10}时，我们发现创造主通过\Person{依撒意亚}用同样的称号指称那位描绘神圣规范的教育者：\BibleQuote{我要从犹大中除掉\Emph{巧匠}，}等等。难道不就是\Person{保禄}自己在那里被预言，注定\Quote{要从犹大中除掉}——也就是从犹太主义中——以建立基督教，好\BibleQuote{安放那唯一的根基，就是基督}\BibleRef{格前 3:11}吗？关于这工作，创造主也通过同一位先知说：\BibleQuote{看，我要在\Place{熙雍}安放一块基石，一块宝贵的、尊贵的石头；凡信赖它的，必不致蒙羞。}\BibleRef{依 28:16} 除非是，天主宣称\Emph{自己}是建造地上工程者，以便祂不给予任何关于祂基督的征兆，基督注定要为相信祂的人成为根基，每个人都可以随意在其上建造或纯正或无价值的教义的上层建筑；因为当一个人的工程要受火试验（或）当他因火得赏报时，是创造主的职责；因为火是用来考验你们在祂所立的根基——即祂基督的根基——之上所建造的工程。\BibleQuote{你们不知道你们是天主的宫殿，天主的圣神住在你们内吗？}\BibleRef{格前 3:16} 现在，既然人是创造主的财产、作品和肖像与模样，他的肉体由祂用尘土塑造，他的灵魂由祂的\OriginalTerm{气息}{afflatus}所造，那么由此推出，如果我们不是创造主的宫殿，\Person{马西昂}的神就完全住在属于别人的宫殿里。但\BibleQuote{若有人毁坏天主的宫殿，天主必要毁坏他}——当然，是由那宫殿的天主。\BibleRef{格前 3:17} 如果你以报应者相威胁，你就是以创造主来威胁我们。\BibleQuote{你们必须成为愚拙，好成为有智慧。}\BibleRef{格前 3:18} 为什么？\BibleQuote{因为这世界的智慧在天主面前是愚妄。}\BibleRef{格前 3:19} 在哪个天主面前？即使古经迄今没有提供任何支持我们观点的证据，但在（宗徒）在这里附加的内容中出现了一个极好的见证：\BibleQuote{因为经上记载：祂以诡计俘获智者；又说：主知道智者的思想是虚妄的。} 因为大体上我们可以确定：他不可能引用他受命要摧毁的那位天主的权威，因为他不会为祂教导。\Quote{因此，}他说，\BibleQuote{谁也不可拿人来夸耀；}\BibleRef{格前 3:21} 这训诫符合创造主的教导：\BibleQuote{信赖世人的人，是可咒骂的；}\BibleRef{耶 17:5} 又说：\BibleQuote{信赖主，胜过依赖世人；}关于（在掌权者中）夸耀，也说了同样的话。\par

% structure: p0019 source=h2 target=subsection reason=relative-heading-rank; base=h2

\subsection{第七章。圣保禄的措辞常常受犹太圣经启发。基督我们的逾越节——这句话将我们引入古代恩赐时期的核心。基督真实的肉体。已婚和未婚状态。\Quote{时期短促}的意义。在劝勉和教义中，宗徒完全按照旧约天主的意念和旨意教导。创造者废除了对肉类和饮品的禁令。}

\BibleQuote{祂要亲自照明暗中的隐密，} \BibleRef{格前 4:5} 甚至藉着基督；因为祂曾应许基督成为光，\BibleRef{依 42:6}，并且祂宣告自己是一盏灯，\Quote{探寻人心和肺腑。} 由祂，\BibleQuote{每个人也要获得称赞}，\BibleRef{格前 4:5}，因为从祂，如同从审判者，也发出称赞的反面。但至少在此，你说他解释世界就是世界之神，当他说：\BibleQuote{我们成了一台戏，给世界、天使和人们观看。} \BibleRef{格前 4:9} 因为如果他用世界指世上的人们，他就不会后来特别提到\BibleQuote{\Emph{人们}}。然而，为阻止你们使用这样的论点，圣神预见地解释了这段经文的含义，说：\BibleQuote{我们成了一台戏，给世界观看，} 即\BibleQuote{给在此服务的\Emph{天使，}} 和\BibleQuote{给受他们服侍的\Emph{人们。}} 当然，我们宗徒的高贵勇气（更不用说圣神了）在写信给他在福音中所生的儿女时，不敢自由地谈论世界之神；因为他似乎不可能对它说任何话，除非以直截了当的方式！我完全承认，按照创造者的法律，\BibleRef{肋 18:8}，这人是个罪犯，\BibleQuote{他娶了他父亲的妻子。} \BibleRef{格前 5:1} 无疑，他遵循自然法和公开法的原则。然而，当他判那人\BibleQuote{交与撒殚，} \BibleRef{格前 5:5} 他成为报应之天主的传令官。他也说\BibleQuote{使肉体毁灭，好使灵魂在主的日子得救，} \BibleRef{格前 5:5} 这并不要紧，因为在肉体毁灭和灵魂得救两方面，都有祂的审判过程；当他命令\BibleQuote{将那恶人从他们中间除去，} \BibleRef{格前 5:13} 他只是提及创造者常常宣判的刑罚。\BibleQuote{你们应把旧酵母除净，好使你们成为新和的面团，正如你们原是无酵的。} \BibleRef{格前 5:7} 因此，在创造者的规条中，无酵饼是我们（基督徒）的预像。\BibleQuote{因为我们的逾越节羔羊基督，已被祭杀，作了牺牲。} \BibleRef{格前 5:7} 但是，基督为何是我们的逾越节，如果逾越节不是基督的预像，在救赎之血的相似和羔羊（即基督）的相似中？\EditorialNote{出谷纪 12} 为什么（宗徒）用创造者庄严礼仪的象征来装扮我们和基督，除非它们与我们有关？再者，当他警告我们提防奸淫时，他启示了肉身的复活。他说：\BibleQuote{身体，} 他说，\BibleQuote{不是为淫乱，而是为主；主也是为身体。} \BibleRef{格前 6:13} 正如圣殿是为天主，天主也是为圣殿。所以，圣殿将与它的神一同消亡，它的神也与圣殿一同消亡。那么，你们看，\BibleQuote{那使主复活的，也要使我们复活。} \BibleRef{格前 6:14} 祂要在身体内使我们复活，因为身体是为主，主也是为身体。适当地，他加上问题：\BibleQuote{你们不知道你们的身体是基督的肢体吗？} \BibleRef{格前 6:15} 异端者要说什么？这些基督的肢体不会复活，因为它们不再是我们的？\BibleQuote{因为，} 他说，\BibleQuote{你们是用高价买来的。} \BibleRef{格前 6:20} 代价！肯定没有支付任何代价，因为基督是个幻影，祂没有任何可以支付我们身体的实质！但实际上，基督有赎买我们的资本；既然祂以高价赎买了我们这些身体，我们必须不行奸淫（因为它们现在是基督的肢体，不再是我们自己的），祂当然会确保那些祂以如此代价使之成为自己的人的得救。现在，我们如何在我们的身体内光荣天主、高举天主，\BibleRef{格前 6:20} 而这身体注定要消亡？现在我们必须面对婚姻的主题，玛尔基昂比宗徒更节制，却禁止婚姻。因为宗徒虽然更偏爱独身的恩宠，\BibleRef{格前 7:7} 却允许缔结婚姻并享受其乐，且劝告继续婚姻而非解除婚姻。\BibleRef{格前 7:27} 基督明确禁止离婚，梅瑟无疑准许离婚。\par

现在，当玛尔基昂完全禁止信友的一切肉身交往（关于他的慕道者，我们暂且不提），并且规定在婚前解除所有婚约时，他跟随的是谁的教导，是梅瑟的还是基督的？然而，连基督，当祂在此命令\BibleQuote{妻子不可离开丈夫，若是离开了，就该独身不嫁，或与丈夫和好}\BibleRef{格前 7:10}时，既允许离婚——祂从未绝对禁止过——又确认了婚姻（的神圣性），首先禁止离婚；若已分离，则希望藉和好重续婚姻纽带。但（宗徒）为节欲提出了什么理由？因为\BibleQuote{时期短促了}\BibleRef{格前 7:29}。我几乎以为那是因为在基督内有另一位神！然而，那位发出这短暂时期的，也将赐予适合此短暂的事物。没有人会为属别人的时期做准备。你贬低你的神，玛尔基昂啊，当你使他完全受限于创造主的时期。显然，当（宗徒）规定婚姻应\BibleQuote{在主内}\BibleRef{格前 7:39}，即基督徒不应与外教人通婚时，他维护了创造主的法律，创造主处处禁止与外人通婚。但当他说道\BibleQuote{虽有称为神的，不论在天上或在地上}\BibleRef{格前 8:5}时，他的话意思很清楚——并非真的有神，而是有些被称为神的，其实并非真神。他引入关于祭偶像之肉的讨论时，首先声明关于偶像本身：\BibleQuote{我们知道偶像在世界上算不了什么}\BibleRef{格前 8:4}。然而，玛尔基昂并不说创造主不是神；所以很难认为宗徒将创造主列在那些被称为神而其实不是的行列；因为即使他们是神，\BibleQuote{为我们只有一个天主，就是圣父}\BibleRef{格前 8:6}。那么，万物从哪里来给我们，不正是从那位拥有万有者而来吗？请问，这些是什么呢？你在书信的前面部分看到了：\BibleQuote{一切全是你们的，无论是保禄，或是阿颇罗，或是刻法，或是世界，或是生命，或是死亡，或是现世，或是来世}\BibleRef{格前 3:21}。那么，他使创造主成为万有的天主，世界、生命和死亡都出自祂，这些不可能属于另一位神。因此，从祂那里，在\Quote{\Emph{一切}}\BibleRef{格前 3:23}之中，也来了基督。当他教导每人应当以自己的劳力为生\BibleRef{格前 9:13}时，他以大量的例子开始——士兵、牧人和农夫\BibleRef{格前 9:7}。但他需要神的权威。然而，援引他所要摧毁的创造主的权威有什么用呢？这样做是徒劳的，因为他的神没有这样的权威！（宗徒）说：\BibleQuote{不可笼住牛嘴，免得它吃所踏出的谷粒}，并补充说：\BibleQuote{天主岂关心牛吗？}是的，关心牛，是为了人的缘故！因为他说：\BibleQuote{这是为我们写的}\BibleRef{格前 11:10}。这样他表明法律对我们有预表意义，并且它支持那些靠福音生活的人。他也表明，传福音的人因此是由那拥有法律、为他们准备的神所派遣的，因为他说：\BibleQuote{这是为我们写的}。然而，他拒绝使用法律赋予他的这项权力，因为他宁愿毫无约束地工作。他以此为荣，不容任何人剥夺他的这种光荣\BibleRef{格前 9:15}——当然不是要摧毁法律，因为他证明别人可以使用它。因为，看，玛尔基昂在盲目中，绊倒在那磐石上，我们的祖先曾在旷野中喝过它的水。既然\BibleQuote{那磐石就是基督}\BibleRef{格前 10:4}，当然，它是属于创造主的，人民也属于祂。但为何要借用外来神给予的神圣记号的形象呢？难道是为了教导这一真理：古事预表了要从它们中产生出来的基督？因为，正要简要回顾（以色列）人民所遭遇的事，他一开始就说：\BibleQuote{这些事都是我们的鉴戒}\BibleRef{格前 10:6}。现在，告诉我，这些鉴戒是创造主给予属于敌对神的人吗？或者是一位神从另一位神那里借用了鉴戒，而且是从敌对神那里？他惊恐地把我从那位他转移我忠诚的神那里拉向他自己。他的对手会使我更倾向于他吗？如果我现在犯下与以色列人民同样的罪，我将承受同样的惩罚吗？还是不会\BibleRef{格前 10:7}？但如果惩罚不同，他向我提出我不会遭受的恐怖是多么徒劳！再者，我将从谁那里承受它们？如果是从创造主，那么祂施加什么恶事呢？而且，既然祂是嫉妒的天主，祂怎么可能惩罚得罪祂对手的人，而不是鼓励他呢？然而，如果是从另一位神——但他不知道如何惩罚。所以，宗徒的整个声明缺乏合理的基础，如果它不是指创造主的管束。但事实上，宗徒的结论与开头一致：\BibleQuote{这些事发生在他们身上，是作为预像，并记载下来，是为了劝戒我们，这世代结局已来到我们身上了}\BibleRef{格前 10:11}。何等的创造主！已经如此预见，并关心警告属于另一位神的基督徒！每当遇到已经处理过的类似异议，我就略过；其他一些我简要处理。对另一位神的一大论证是允许吃各种肉类，违反法律\BibleRef{格前 10:25}。仿佛我们自己不承认法律的重担已被废除——但由那制定法律者废除，祂也应许了新的事物秩序。因此，那位禁止某些肉食的，也恢复了它们的使用，正如祂起初允许的那样。然而，如果某位陌生的神来临要摧毁我们的神，他的首要禁令必定是，他自己的追随者不应依赖他对手的资源来维持生命。\par

% structure: p0022 source=h2 target=subsection reason=relative-heading-rank; base=h2

\subsection{第八章　人是造物主的肖像，基督是人的元首。神恩。依撒意亚论及七重之神。宗徒与先知相比。斥责玛尔基昂要显示他的神明有何类似先知所预言的神恩。}

\BibleQuote{男人的头是基督。} \BibleRef{格前 11:3} 这基督是谁，倘若祂不是人的创造者？他在这里用\Emph{头}喻指\Emph{权柄}；\Quote{权柄}只能归于\Quote{创造者}。他实在是哪一个人的头呢？无疑是那随后他提及之人的：\BibleQuote{男人不应蒙头，因为他是天主的肖像。} \BibleRef{格前 11:7} 既然他是创造者的肖像（因为\Emph{祂}，当注视祂的圣言基督、那将要成为人的，说：\BibleQuote{让我们照我们的肖像，按我们的模样造人。} \BibleRef{创 1:26}），我怎能另有别的头，而非祂——祂的肖像我乃是？因为我若是创造者的肖像，在我内就没有另一头的空间。但是为什么\BibleQuote{女人应当在头上带有权柄，为了天使的缘故？} \BibleRef{格前 11:10} 若是因\BibleQuote{女人是为了男人而被造的} \BibleRef{格前 11:9}，并从男人取出，符合创造者的旨意，那么宗徒也就这样维护了那位天主的规定，他从该天主的制度解释其规定的理由。他补充说：\BibleQuote{为了天使的缘故。} \BibleRef{格前 11:10} 什么天使？换句话说，谁的天使？如果他是指创造者的堕落天使，那他的意思极为恰当。那曾成为他们陷阱的面容，理应带有某种谦卑外表和暗淡美丽的标记。然而，如果指的是对立之神的天使，他们有何畏惧？因为连马西昂的门徒（不提他的天使）也不垂涎女人。我们之前屡次指出，宗徒把异端列为邪恶 \BibleRef{格前 11:18}，属于\Quote{肉性之事}，并且他视那些躲避异端如避恶事的人为可敬。同样，在论及福音时，我们从饼和杯的圣事证明主身体的真实性与血的真实性，以对抗马西昂的幻影；而在我的整部作品中几乎一致主张，所有关于审判属性的提及都确凿地指向创造者，即一位审判的天主。现在，关于\Quote{神恩} \BibleRef{格前 12:1} 的主题，我必须指出，这些也是创造者藉基督应许的；我认为我们可以从此得出非常正确的结论：恩赐的赏赐不是另一位神的工作，而正是那位被证明给予应许的天主。以下是依撒意亚的预言：\BibleQuote{由叶瑟的树干将生出一根枝条，从他的根上将发出一朵花；上主的神将安息在祂身上。} 随后他列举了同一神的各种恩赐：\BibleQuote{智慧和聪敏的神，超见和刚毅的神，知识和孝爱的神。上主的神要以敬畏上主充满祂。} \BibleRef{依 11:1} 藉着\Emph{花}这一形象，他显示基督将从叶瑟树干所生的枝条中兴起，换句话说，从叶瑟之子达味后裔的童贞女而生。在这基督身上，圣神的整个\OriginalTerm{本体}{substantia}必须安息，但这并非意指某种后来才加给祂的东西，因为祂在降生成人之前始终是圣神；因此你不能由此论证这预言指向那位基督（仅仅作为达味后裔的人），他将获得他的天主之神。反之，（先知说）从那时起，当（真正的基督）以肉身显现，如同从叶瑟之根升起的\Emph{预言之花}时，整个恩宠之神的运作必须安息在祂身上，这运作对犹太人而言将停止并终结。事实本身证明了这一点；因为此后创造者的神不再在\Emph{他们}中间吹拂。从犹大被夺走了\Quote{智者、精巧工匠、谋士和先知}，以便应验\BibleQuote{法律和先知直到若翰为止。} \BibleRef{路 16:16} 现在请听他是如何宣告，藉基督自己，当祂返回天上时，这些神恩将被赐下：\BibleQuote{祂升上高处，} 即升入天堂；\BibleQuote{祂掳掠了俘虏，} 意指死亡或人的奴役；\BibleQuote{祂将恩赐赐给了人的儿子们，} 即那些恩惠，我们称之为\Emph{神恩}。他特意说\Quote{\Emph{人的儿子们}}，而非泛指的人；如此向我们指明那些真正被称为人子的人，蒙拣选的人，宗徒。\Quote{因为，} 他说，\Quote{我藉着福音生了你们；} \BibleRef{格前 4:15} 和\Quote{你们是我的孩子，我为你们再受产痛。} \BibleRef{迦 4:19} 那藉约厄尔之言所赐的圣神应许，如今完全应验了：\BibleQuote{在末日，我要将我的神倾注在一切有血肉的人身上，你们的儿女要说预言；连我的仆人和我的婢女，我也要倾注我的神。} 既然创造者应许在末后的日子赐下祂的圣神之恩；既然基督在这些末后的日子已显现为神恩的分配者（如宗徒所说：\BibleQuote{及至时期一满，天主就派遣了自己的儿子来，} \BibleRef{迦 4:4} 又说：\BibleQuote{弟兄们，我对你们说：时期是短促的。}），那么显然与此末后的预言相联，圣神的恩赐属于祂——那预言者的基督。现在请比较圣神的特定恩宠，如宗徒所描述、依撒意亚先知所应许的。他说：\BibleQuote{有人因圣神领受了智慧的言语；} 我们立刻看出这就是依撒意亚所称的\Quote{智慧之神}。\BibleQuote{另有人领受知识的言语；} 这就是\Quote{（先知的）聪敏和超见之神}。\BibleQuote{另有人藉同一圣神领受信德；} 这就是\Quote{孝爱和敬畏上主之神}。\BibleQuote{另有人领受治病的恩赐，另有人行奇迹的恩赐；} 这就是\Quote{刚毅之神}。\BibleQuote{另有人作先知预言，另有人辨别神类，另有人能说各种语言，另有人能解释语言；} 这就是\Quote{知识之神}。请看宗徒与先知在分施同一圣神以及解释祂的特定恩宠上如何一致。这一点我也可自信地说：那将我们身体通过其众多不同肢体的合一比作圣神各种恩赐的结合的人，也表明只有一个主，掌管人的身体与圣神。这圣神（按宗徒的展示）并非意指这些恩赐的服事应在身体内，也没有把它们放在人的身体内）；而关于爱德超越这一切恩赐的主题，祂甚至教导宗徒，这是首要的诫命，正如基督所显示的：\BibleQuote{你当全心、全灵、全力、全意爱上主，你的天主，并爱近人如你自己。} \BibleRef{路 10:27} 当他提及\Quote{\Emph{在法律中记载着}}，创造者要用异语和异唇说话，他固然以此提及证实了舌语的恩赐，却不可被认为因引用创造者的预言而确认这恩赐来自另一位神。\BibleRef{格前 14:21} 完全同样，当他命令妇女在教会内沉默，不可为学习而发言 \BibleRef{格前 14:34}（尽管他已藉着使说预言的妇女蒙头表明她们也有说预言的权利），他引法律作为女人应服从的根据。现在，让我一劳永逸地说，他对这法律本应无其他认识，除了废除它。但为使我们现在离开神恩的主题，事实本身足以证明我们中谁在妄称这些恩赐属于他的神，以及它们是否可能与我们对立，即使创造者为那尚未启示的基督（只预定给犹太人）应许了这些恩赐，要在祂的时期、在祂的基督、在祂的子民中运作。那么，让马西昂展示他的神的各种恩赐：一些先知，不是凭人的感觉说话，而是凭天主之神说话，既能预言未来，又能揭示心中的秘密；\BibleRef{格前 14:25} 让他提出一篇圣咏、一个异象、一篇祈祷 \BibleRef{格前 14:26}——但必须出于圣神，在出神状态，即神魂超拔中，每当有舌语的解释临到他；也请向我证明，他那群特别神圣的姊妹中，任何一个多嘴的女人曾说过预言。而今，这些（神恩的）标记从我方毫无困难地显现，它们也与创造者的规则、分施和教导相符；因此毫无疑问，基督、圣神和宗徒各自属于我的天主。这就是我坦白的声明，给任何愿意询问的人。\par

% structure: p0024 source=h2 target=subsection reason=relative-heading-rank; base=h2

\subsection{第九章。复活的道理。身体将复活。基督审判者的特质。犹太人对预言的曲解被揭露和驳斥。默西亚圣咏得到辩护。犹太人和理性主义者对此点的解释相似。耶稣——而非希则克雅或撒罗满——是这些圣咏中预言的对象。唯独他是旧约和新约的基督。}

同时，马西翁派不会展示任何这类东西；此时他不敢说哪一方对一位尚未启示的基督拥有更大的权利。正如我的基督是应被期待的，祂从起初就被预告了，因此他的基督根本不存在，因为没有从起初被\Emph{预告}。我们的信仰更好，相信一位未来的基督，胜过异端者的信仰，他们根本没有可相信的对象。关于死者的复活，\BibleRef{格前 15:12} 让我们首先探究当时有些人如何否认它。无疑是以如今依然被否认的同样方式，因为肉体的复活从来都有人否认。但许多智者声称灵魂具有神圣本性，确信其不朽的命运，甚至民众也因大胆存想灵魂不死而敬拜死者。至于我们的身体，显然它们要么立即被火或野兽毁灭，即使被最小心保存，也终因时间而腐朽。因此，当宗徒驳斥那些否认肉体复活的人时，他确实在捍卫他们所否认的正是那个东西，即身体的复活。整个答案都包含在这里。其余都是多余的。如今在这称为死者复活的一点上，必须准确保持词语的恰当力量。\Quote{死者}一词仅仅表达失去生命原理、藉以活动的存在。身体是失去生命并因此成为死者的东西。因此，死者一词只适用于身体。此外，既然复活属于死者，而死者的称谓只适用于身体，所以只有身体才具有与之相关的复活。同样，\Quote{复活}一词或\Quote{再次升起}，只涵盖那跌倒的。\Quote{升起}确实可以用于从未跌倒、但一直躺卧的东西。但\Quote{再次\Emph{升起}}只可以用于那跌倒的；因为正是由于跌倒后再次升起，它才被称为已经\Emph{复}起。因为音节\OriginalTerm{\Quote{复}}{RE}总蕴涵重复（或再次发生）。因此我们说，身体因死亡而倒在地上，正如事实本身所表明，符合天主的律法。因为曾对身体说：（\BibleQuote{直到你归回土地，因为从那里你被取来；因为）你是灰土，还要归于灰土。}因此，那来自土地的必归于土地。凡归于土地的，就是跌倒的；凡跌倒的，就能再次升起。\BibleQuote{因为死既因一人而来，复活也因一人而来。}\BibleRef{格前 15:21} 在此，\Quote{人}一词，指由身体实体组成的，正如我们已多次指出的，向我呈现了基督的身体。但如果我们都在基督内这样被复活，如同我们在亚当内死去，那么必然的结果是，我们在基督内作为身体实体被复活，因为我们在亚当内作为身体实体死去。这相似确实不完全，除非我们在基督内的复活与我们在亚当内的死亡在实体上一致。但在此处，（宗徒）插入了一句关于基督的说明，这既然与我们当前的讨论有关，不可忽视。因为身体的复活将得到更好的证明，根据我能否成功表明基督属于那位被认为在其救恩计划中提供了肉体复活的天主。当他说：\BibleQuote{因为祂必须为王，直到把一切仇敌放在祂脚下}，我们立刻从这话看出，他谈论的是一位复仇的天主，因此就是向基督作如下应许的那位：\BibleQuote{坐在我的右边，直到我使你的仇敌作你的脚凳。上主从\Place{熙雍}伸出你能力的权杖，祂要在你的仇敌中与你一同统治。}我必须主张那些\Place{以色列}人试图剥夺我们的圣经，并表明它们支持我的观点。他们声称这篇圣咏是纪念\Person{希则克雅}的颂歌，因为\Quote{他上到上主的殿}，天主就转回并除去他的敌人。因此（他们还认为），那些话，\BibleQuote{在晓星之前，我从子宫生了祢}，适用于\Person{希则克雅}和\Person{希则克雅}的生产。我们方面则发布了\WorkTitle{福音书}（其可信性我们还须感谢他们已在一个如此重大的主题上给予了某种确认）；这些福音宣告主在夜间出生，以便应验\Quote{在晓星之前}，这从那颗星尤其明显，也从天使的见证，他在夜间向牧羊人宣布基督那时已经出生，还有出生地点，因为临近夜间人们才到达（东方的）\Quote{客栈}。也许，基督在夜间出生也有一个奥秘目的，因为祂注定要在无知黑暗的阴影中成为真理之光。再者，天主不会说\Quote{我生了祢}，除非是对祂的真子。因为虽然祂论及全体人民（\Place{以色列}）说：\BibleQuote{我养育了子女}，\BibleRef{依 1:2} 但祂没有加上\Quote{从子宫}。那么，祂为何要多此一举加上\Quote{从子宫}这句话（好像有任何人对某人从子宫出生有疑问），除非圣神特意要这些话被理解为论及基督？\BibleQuote{我从子宫生了祢}，也就是说，仅从子宫，没有男人的种子，使肉身的条件成为从子宫出来。这里（在圣咏中）加上的\BibleQuote{祢是永远的司铎}，涉及（基督）自己。\Person{希则克雅}不是司铎；即使他是，也不会是永远的司铎。\BibleQuote{按品位}，祂说，\Quote{属\Person{默基瑟德}}。\Person{希则克雅}与\Person{默基瑟德}有何关系？他是至高天主的司铎，且未受割损，祝福了受割损的\Person{亚巴郎}，并从\Person{亚巴郎}接受了什一奉献。然而对于基督，\Quote{默基瑟德的品位}非常合适；因为基督是天主正确合法的\OriginalTerm{大司祭}{High Priest}。祂是未割损司祭职的大司祭，即使在那时，为了外邦人而被立为这样的，外邦人将更充分地接纳祂，虽然在祂末次来临时，祂也将以接纳和祝福宠爱受割损的人，即\Person{亚巴郎}的后裔，他们不久将承认祂。那么，还有另一篇圣咏，以这些话开头：\BibleQuote{天主，赐给君王你的审判}，即赐给将作为君王来的基督，\BibleQuote{将你的公义赐给君王的儿子}，即赐给基督的子民；因为祂的儿子就是在祂内重生的人。但这里会有人说这篇圣咏指向\Person{撒罗满}。然而，圣咏中唯独适用于基督的部分，岂不足以教导我们其余部分也涉及基督，而非\Person{撒罗满}吗？\BibleQuote{祂要降临，像雨落在羊毛上，像滴露落在土地上，}描述祂从天上降下成为肉身的温柔和隐蔽。然而\Person{撒罗满}，如果真有降下的话，也不是像阵雨那样下来，因为他不是从天上降下。但我要提出更字面的要点。圣咏作者说：\BibleQuote{祂要统治，从海到海，从大河直到地极。}这惟独赐给了基督；而\Person{撒罗满}只统治了中等规模的\Place{犹大}王国。\BibleQuote{是的，众王都要俯伏在祂面前。}他们这样敬拜谁呢，除了基督？\BibleQuote{万民都要事奉祂。}众人向谁这样致敬，除了基督？\BibleQuote{祂的名要存留永远。}谁的名有这永久的名声，除了基督的？\BibleQuote{祂的名要比太阳长久，}因为天主的圣言，即基督，要比太阳长久。\BibleQuote{万民要在祂内蒙福。}在\Person{撒罗满}内没有国家蒙福；在基督内每个国家都蒙福。如果圣咏证明祂甚至就是天主呢？\BibleQuote{他们称祂有福。}（凭什么？）因为上主\Place{以色列}的天主是有福的，惟有祂行奇事。\BibleQuote{祂荣耀的名也是有福的，愿全地充满祂的荣耀。}相反，\Person{撒罗满}（我敢断言）甚至失去了祂从天主而来的荣耀，被祂对女人的爱引诱而陷入偶像崇拜。因此，在这篇圣咏中部出现的话\BibleQuote{祂的仇敌要舔尘土}（当然，既然是，借用宗徒的话，\Quote{被放在祂脚下}），将有助于我所着眼的目标，当我引用这篇圣咏并坚持我对它含义的看法时——即，我欲表明祂王国的荣耀和仇敌的屈服，符合造物主自己的计划，以得出这个结论：除祂以外，无人可被视为造物主的基督。\par

% structure: p0026 source=h2 target=subsection reason=relative-heading-rank; base=h2

\subsection{第十章 继续论身体复活的道理。死人怎样复活？带着什么样的身体而来？这些问题的解答是为了维护复活身体的真实性，以驳斥\Person{马尔西昂}。\Person{基督}作为第二亚当，与第一个人的创造者相联系。让我们佩带天上的肖像。依照先知们论死亡的胜利。\Person{欧瑟亚}与\Person{圣保禄}的比较。}

现在让我们回到复活的话题。我们在另一部著作中确实已足够注意为复活辩护，以对抗各种异端。但即便如此，我们也不应在此欠缺（对教义的辩护），考虑到那些不熟悉那篇小论文的人。 \Quote{什么？}他问，\Quote{他们为死者受圣洗，若死者不复活，他们将做什么呢？} \BibleRef{格前 15:29} 现在，不必在意那种做法（无论它可能是什么）。\OriginalTerm{二月的洗涤礼}{Februarian lustrations}也许可以回答他（同样有效），通过为死者祈祷。因此，不要以为宗徒在此暗示某个新的\Emph{神作为}这种（为死者受圣洗）的作者和倡导者。他提及此事的唯一目的，是为了更坚定地坚持身体的复活，因为那些徒然为死者受圣洗的人，正是出于对\Emph{这样一种}复活的信仰才采用这种做法的。宗徒在另一处定了\Quote{只有一个圣洗}。\BibleRef{弗 4:5} 因此，\Quote{为死者受圣洗}事实上意味着为身体受圣洗；因为，正如我们所指出的，是\Emph{身体}变为\Emph{死者}。那么，那些为身体受圣洗的人，若身体不再复活，他们将做什么呢？因此，我们站稳了（我们说），宗徒接下来讨论的问题同样与身体有关。\Quote{有人要说：\InnerQuote{死人怎样复活呢？他们带着什么身体来呢？}} \BibleRef{格前 15:35} 在确立了曾被否认的复活教义之后，自然地要讨论复活中的身体是何种情形，对此无人知晓。在这一点上，我们有其他对手需要应对。因为马西昂丝毫不承认肉身的复活，他只应许灵魂的救恩；因此他提出的问题不是关于身体的\Emph{种类}，而是关于身体的\Emph{本体}。然而，即使从宗徒针对那些问\InnerQuote{死人怎样复活？他们带着什么身体来？}的人所阐述的关于身体性质的话中，他也是最明显地被驳倒的。因为他既然讨论了身体的\Emph{性质}，当然\OriginalTerm{事实上}{ipso facto}就在论证中宣称是\Emph{身体}将要复活。的确，既然他举出例子\InnerQuote{麦粒或别的种子，天主给它一个身体，如他愿意；} \BibleRef{格前 15:37} 他又说，\InnerQuote{每种种子都有它自己的身体；} \BibleRef{格前 15:38} 因此，\InnerQuote{人的肉体是一种，牲畜的肉是另一种，鸟的肉又是一种；还有天上的形体和地上的形体；太阳的光荣是一种，月亮的光荣是另一种，星辰的光荣又是另一种} \BibleRef{格前 15:39} ——难道他不是在暗示将有肉体或身体的复活吗？他用肉体的和物质的样本来说明这一点。难道他不是也在保证，复活将由那位天主完成，一切（作为他例子的受造物）都出自祂吗？他说，\InnerQuote{死人的复活也是这样。} \BibleRef{格前 15:42} 怎样呢？正如种子，被播下一个身体，却长出一个身体。他把身体播下称为在土里的解体，\InnerQuote{因为它播下的是可腐朽的，}（但\InnerQuote{复活起来的是）光荣和权能。} \BibleRef{格前 15:42} 现在，正如种子那样，在这里：谁安排了身体解体过程，谁就拥有使身体复活的工程。然而，如果你从你所交付解体的复活中除去身体，那结局中的多样性又会变成什么呢？同样，\InnerQuote{它虽然被播下是属魂的身体，却被复活为属神的身体。} \BibleRef{格前 15:44} 现在，虽然生命的自然原则和圣神各有其专属的身体，以致\InnerQuote{属魂的身体}可合理地理解为灵魂，\InnerQuote{属神的身体}可理解为圣神，但这并不是理由让人以为宗徒说灵魂要在复活中成为圣神，而是说\Emph{身体}（它由于与灵魂同生，并藉灵魂维持生命，所以可称为生灵的（或\Emph{自然的}）身体）将要成为\Emph{属神的}，因为它藉着圣神复活进入永生。简言之，既然不是灵魂，而是肉体\InnerQuote{被播下在腐败中，}当它在土里化为腐朽时，由此推出（在如此解体之后）灵魂不再是自然的身体，而是那曾是自然身体的肉体（是未来变化的主体），因为从自然的身体变成属神的身体，正如他下文所说：\InnerQuote{那属神的不是在先的。} \BibleRef{格前 15:46} 因为正是为此，他先前论到基督自己说：\InnerQuote{第一个人亚当成了有生命的生灵，最后的亚当成了赋予生命的圣神。} \BibleRef{格前 15:45} 然而，我们的异端者，由于愚昧过度，不愿让这话保持原样，把\InnerQuote{最后的亚当}改为\InnerQuote{最后的主}；因为他当然害怕，如果他允许主是最后的（或第二个）亚当，我们就会争辩说，基督既是第二个亚当，就必然属于那也拥有第一个亚当的天主。但篡改是明显的。因为为什么有第一个亚当，除非也有第二个亚当？因为事物除非各各相似，并且在名称、本质或起源上具有同一性，否则不能归为一类。现在，即使在个别不同的事物中，也必须有第一个和最后一个，但它们必须有一位作者。然而，如果作者是不同的，他本人确实可以被称为\InnerQuote{最后的}。但他所引入的事物是第一个，而只有在本性上像这第一个的事物才能是最后一个。然而，当它不是同一作者的作品时，它在本性上就不像第一个。同样地，（异端者）也将被\Emph{人}这个词驳倒：\InnerQuote{第一个人是出于地，属于土；第二个人是主，来自天上。} \BibleRef{格前 15:47} 既然第一个人是\Emph{人}，除非第二个也是\Emph{人}，否则怎能有第二个？不然，如果第二个是\InnerQuote{主}，难道第一个也是\InnerQuote{主}吗？然而，对我而言，这已足够：因为他在他的福音中承认人子既是基督，也是人；所以他不能（在此处）否认祂是（宗徒所说的）\InnerQuote{亚当}和\InnerQuote{人}。随后的话也将使他难以承受。因为宗徒说：\InnerQuote{那属土的怎样，属土的人也怎样}——当然，又是\Emph{人}；\InnerQuote{那属天的怎样，属天的人也怎样。}因为祂不可能将任何不是\Emph{人}的属天存在与属土的\Emph{人}相对立；祂的目的正是通过共同使用这个名称，更准确地区分他们的状态和期望。因为就他们目前的状态和未来的期望，祂称他们为属土的人和属天的人，仍然保留他们名称的平等，根据他们（在最终状态上）被算在亚当内或在基督内。因此，在劝勉他们珍惜天主的希望时，他说：\InnerQuote{我们既然带有了属土之人的形像，也要带有了属天之人的形像。}——这句话不是关于复活生命的状态，而是关于现世的准则。他说\Emph{让我们戴上}，是一条诫命；而不是\Emph{我们将要戴上}，作为一个应许——希望我们照他本人所行的去行，并在肉体的行为上脱下属土之人的样式，即旧人。接下来是什么话呢？\InnerQuote{弟兄们，我说：血肉不能继承天主的国。} \BibleRef{格前 15:50} 他指的是血肉的行为，这些行为在他写给迦拉达人的书信中使人丧失天主的国。 \BibleRef{迦 5:19} 在其他地方他也惯于用自然状态代替其中所做的行为，例如他说：\InnerQuote{凡属于血肉的人，不能中悦天主。} \BibleRef{罗 8:8} 然而，除非我们在这肉身内，何时才能中悦天主呢？我想，没有任何其他时间人可以工作。但是，如果我们甚至在肉身中自然地生活，却避免肉身的行为，那么我们就不是在肉身内；因为虽然我们并未离开肉身的实体，但我们却与肉身的罪无关。既然在\Emph{肉身}一词中，我们被命令脱去的不是实体，而是肉身的行为，因此，在同样这个词的使用中，天主的国被否认给肉身的行为，而不是给肉身的实体。因为那在其中行恶的实体并不被定罪，只有其中所行的恶才被定罪。下毒是罪行，但盛毒的杯子却不是有罪的。同样，\Emph{身体}是肉身行为的器皿，而里面的灵魂调配了邪恶行为的毒药。那么，灵魂，作为肉身行为的真正作者，如何在身体中所行的事被补赎后得以进入天主的国，而身体，不过是灵魂的服役工具，却必须留在定罪中呢？杯子该受惩罚，而下毒者却逃脱吗？我们当然不是为肉身要求天主的国：我们所做的一切，只是断言它的实体会有复活，作为进入天国之门。但复活是一回事，天国是另一回事。复活在先，然后是天国。因此，我们说，肉身复活，但改变后获得天国。因为\InnerQuote{死人必要复活成为不朽的，}甚至那些在身体腐朽时曾是可朽坏的；\InnerQuote{我们都要改变，在一刹那，一眨眼之间。 \BibleRef{格前 15:52} 因为这可朽坏的}——宗徒说话时似乎指着自己的肉身——\InnerQuote{必须穿上不朽的，这必死的必须穿上不死的，} \BibleRef{格前 15:53} 确实，为了使之成为天国之国的合适实体。\InnerQuote{因为我们必将如同天使。}这将是我们肉身的完美改变——只能在复活之后。反之，如果没有肉身，它怎会穿上不朽和不死呢？既然因改变而成为别的，它将获得天主的国，不再是（旧的）血肉，而是天主将赐给它的身体。因此，宗徒正确地宣告：\InnerQuote{血肉不能继承天主的国；} \BibleRef{格前 15:50} 因为他把这（尊荣）归之于复活后随之而来的改变状态。因此，那时将要应验那由创造主所写的话：\InnerQuote{死亡，你的胜利在哪里？}——或你的争斗？\InnerQuote{死亡，你的刺在哪里？} \BibleRef{格前 15:55} ——我说，由创造主所写，因为祂藉祂的先知写了这些话——那宣告了将要在天国中应验的话的那一位，将属于这恩赐，也就是天国。而且，他没有告诉我们应当向别的神献上\Quote{感谢}，因为祂使我们甚至战胜了死亡，而只应献给那一位，从祂那里他得到了这向必死的仇敌欢庆和胜利挑战的言辞。\par

% structure: p0028 source=h2 target=subsection reason=relative-heading-rank; base=h2

\subsection{第十一章 \BibleBook{格林多}{后书}。造物主是仁慈的圣父——在旧约中以及在基督内均如此显示。新约的新颖。顽固瞎眼的面纱遮住以色列——按玛尔定的原则无可指责。犹太人因拒绝造物主的基督而有罪。撒殚，这世界的神。瓦器中的宝藏——反对玛尔定的解释。造物主与这些器皿——即我们的身体——的关系。}

如果由于人类错误的过失，\Emph{天主}一词成了一个普通名词（因为世上被认为和相信有\BibleQuote{许多的神}\BibleRef{格前 8:5}），然而\BibleQuote{那有福的天主}（就是\BibleQuote{我们主耶稣基督的父}\BibleRef{格后 1:3}）将被理解为不是别的神，而是造物主；祂祝福了祂所造的一切，如你在创世纪中所见\BibleRef{创 1:22}，并且祂自己\Quote{受万有祝福}，如\Person{达尼尔}所说。现在，如果\Emph{父}的称号可归于\Person{马尔基昂}那不结果实的神，岂不更应归于造物主？没有别的神更适合这称号，祂也是\BibleQuote{慈悲的父}\BibleRef{格后 1:3}，且在先知书中被描述为\Quote{富有同情、宽仁大方、广施慈恩}。在\Person{约纳}书中，你看到祂向祈祷的尼尼微人所显的慈恩\BibleRef{纳 3:8}。祂对\Person{希则克雅}的眼泪是何等坚定！祂对\Person{纳波特}之血的事上，当\Person{阿哈布}哀求时，祂是何等易于宽恕！祂在\Person{达味}认罪时何其迅速赦免\BibleRef{撒下 12:13}——的确，祂宁愿罪人悔改而不愿其死亡，当然是由于祂慈恩的属性\BibleRef{则 33:11}。现在，如果\Person{马尔基昂}的神曾显示或宣告过任何这类事，我就承认他是\Quote{慈悲的父}。然而，既然他只在启示之时才将这称号归于他的神，仿佛他只在开始解放人类时才成为慈悲的父，那么我们这边也采用他所谓启示的那个确切时间；但这是为了否认他！因此，他没有权利将任何属性归于他的神，因为他正是通过这种归属性来宣扬他的神；除非事先证明他的神确实存在，才允许他赋予属性。所赋予的属性只是偶然；但偶然要以对事物本身的陈述为前提，尤其当另一个人宣称那属性属于一个尚未被证明存在的神时。我们对他的存在的否认将更加强硬，因为作为证据的属性已属于那已被启示的天主。所以，\Quote{新约}将只属于那应许它的那一位——即使不是\Quote{它的文字}，也是\Quote{它的精神}\BibleRef{格后 3:6}；而它的\Emph{新}就在这里。事实上，那曾在石版上刻写其文字的，正是那关于其精神曾说过\BibleQuote{我要将我的圣神倾注在一切血肉之人身上}的那一位\BibleRef{岳 2:28}。即使\Quote{文字叫人死，精神却叫人活}\BibleRef{格后 3:6}，两者都属于那说\BibleQuote{我使人死，也使人活；我击伤，也医治}的那一位\BibleRef{申 32:39}。我们已经确立了造物主对审判与仁慈这双重特征的拥有权——通过法律\Quote{在文字中杀死}，通过福音\Quote{在精神中使人活}。现在，这些属性尽管不同，却不能造成两位神；因为它们在旧约的预先安排中已经集中于一位。他暗指\Person{梅瑟}的面纱，带着它\Quote{以色列子民不能定睛看他的面容}\BibleRef{格后 3:7}。他这样做是为了维护新约荣耀的优越性，因为新约的荣耀是永恒的，胜过那\Quote{将要消逝}的旧约\BibleRef{格后 3:7}，这一事实支持了我的信念——它高举福音于法律之上；你必须注意它是否不止于此。因为唯有先前有之事物，才谈得上超越。但接着他说：\BibleQuote{但他们的心思蒙蔽了}——指的是世界的心；当然不是造物主的心思，而是世上众人的心思。关于以色列，他说：\BibleQuote{直到今日，那同样的面纱还留在他们心上}\BibleRef{格后 3:15}，表明\Person{梅瑟}脸上的面纱是那至今仍蒙蔽这民族之心的面纱的表象；因为现在他们心中看不见\Person{梅瑟}，正如当时他们眼中看不见他。但如果造物主的基督（\Person{梅瑟}所预言的）还没有来，\Person{保禄}与那仍使犹太人看不见\Person{梅瑟}的面纱有何相干？如果\Person{梅瑟}关于基督的预言（犹太人本该通过他而相信）尚未应验，犹太人的心如何能被描述为仍然蒙蔽和遮盖？如果犹太人不理解他们自己天主的神秘宣告，那异邦基督的宗徒凭什么抱怨，除非这面纱指的是那使他们看不见\Person{梅瑟}的基督的盲目？然后，随后的话：\BibleQuote{但何时转向主，这帕子就会被除去}\BibleRef{格后 3:16}，正确指向犹太人，\Person{梅瑟}的面纱覆盖着他们的目光，以致当他们转向基督信仰时，就会理解\Person{梅瑟}如何谈论基督。但别位神的基督怎能除去造物主的面纱？造物主不可能掩盖未知之神的神秘——因为那是未知之神的未知奥秘。所以他说：\BibleQuote{我们如今以揭开的面容}（意指心中的坦率，这在犹太人中被面纱遮盖）\BibleQuote{观看基督，就变成同样的形象，从那种荣耀}（\Person{梅瑟}因主的荣耀而变容）\BibleQuote{到另一种荣耀}\BibleRef{格后 3:18}。通过如此描述那从与天主会面中照耀\Person{梅瑟}面容的荣耀，以及那因百姓的软弱而遮掩同一荣耀的面纱，并进而加上在基督身上圣神的启示与荣耀——即如他所说，\Quote{藉着主的圣神}——他证明整个\Person{梅瑟}体系是基督的预像，犹太人对此无知，但我们基督徒认识他。我们很清楚，有些段落因阅读方式或标点而模棱两可。\Person{马尔基昂}采用了后者，将随后的段落读作\BibleQuote{在这世界的神中}\BibleRef{格后 4:4}，仿佛描述造物主为这世界的神，以便藉此暗示有另一位神掌管另一世界。然而，我们说这段落在\Quote{神}后加逗号，意思为：\Quote{天主在这世界蒙蔽了不信者的眼睛}。\Quote{在其中}指犹太的不信者，其中一些人直到现在福音仍被\Person{梅瑟}的面纱遮掩。正是这些人，天主管威胁他们\Quote{嘴唇亲近我，心却远离我}\BibleRef{依 29:13}，用这些怒言：\Quote{你们听是要听见，却不明白；看是要看见，却不识透}；以及\Quote{你们若不相信，必不得理解}；又说：\Quote{我要夺去智慧人的智慧，废除精明人的精明}。当然，这些话祂不是因为他们隐藏了未知之神的福音而说的。无论如何，如果有一位这世界的神，他蒙蔽这世界不信者的心，因为他们没有自主地认识他的基督，这基督本应从圣经中被认识。我满足于我的优势，愿意不再过多注意这个标点歧义，以免给我的对手任何优势；事实上，我本可以完全省略讨论。一个更简单的答案唾手可得：将\Quote{这世界的神}解释为魔鬼，他曾如先知所述说：\BibleQuote{我要与至高者相似；我要高举我的宝座在云上}\BibleRef{依 14:14}。的确，这世界的整个迷信都落在他的手中，所以他有效地蒙蔽不信者的心，尤其是背教者\Person{马尔基昂}的心。他没有注意到这句话的这一点对他多么不利：\BibleQuote{因为那吩咐光从黑暗里照耀出来的天主，已经照在我们心里，为把（祂荣耀的）知识之光（赐给）在（耶稣）基督的面容上}\BibleRef{格后 4:6}。现在，是谁说了\BibleQuote{要有光}\BibleRef{创 1:3}？又是谁对基督论及照亮世界说：\BibleQuote{我已立你作外邦人的光}\BibleRef{依 49:6}——即那些\BibleQuote{坐在黑暗中死荫里的人}\BibleRef{路 1:79}？除了祂，还有谁呢？圣神在圣咏中预见未来，回答说：\BibleQuote{上主，你面容的光已显示在我们身上}\BibleRef{咏 4:7}。这里上主的面容就是基督。因此，宗徒上面说：\BibleQuote{基督是天主的像}\BibleRef{格后 4:4}。既然基督是那说\BibleQuote{要有光}的造物主的面容，那么基督、宗徒、福音、面纱、\Person{梅瑟}——以至整个救恩计划——都属于那创造这世界的神，根据上述句子的见证，当然不属于那从未说\BibleQuote{要有光}的神。我在此略过关于另一封书信的讨论，我们认为是写给厄弗所人的，而异端认为是写给劳狄刻雅人的。他在其中告诉他们要记住，当他们还是外邦人时，他们\BibleQuote{没有基督，与以色列团体隔绝，无来往，无盟约，无希望，也无天主\BibleRef{弗 2:12}，即使在他自己的世界中}，即这世界的造物主。既然他说外邦人没有天主，而他们的神是魔鬼，不是造物主，那么很清楚，他应被理解为这世界的主，外邦人把他当作自己的神——而不是造物主，他们不认识他。但怎么会这样：\BibleQuote{我们在这瓦器中的宝贝}\BibleRef{格后 4:7}不应被视为属于那拥有瓦器的神？既然天主的荣耀在于如此大的宝贝藏在瓦器中，而这些瓦器是造物主所造，那么荣耀属于造物主；不仅如此，这些属于祂的瓦器如此散发天主大能的卓越，那大能本身也必属于祂！的确，所有这些都托付给这些\Quote{瓦器}，正是为了彰显祂的卓越。从此，那竞争的神便无权宣称这荣耀，因而也无权宣称那大能。相反，耻辱和软弱将归于他，因为他从未沾手的瓦器却接受了全部的卓越！那么，如果正是在这些瓦器中，他告诉我们我们必须忍受如此大的苦难\BibleRef{格后 4:8}，并在其中携带天主本身的死，那么\Person{马尔基昂}的神真是忘恩负义、不公义，如果他不打算在复活中恢复我们这同一实体——其中如此多的事为了效忠于他而忍受，在其中基督的死被携带，在其中他大能的卓越也被珍藏\BibleRef{格后 4:10}。因为他着重声明：\BibleQuote{为叫耶稣的生命也彰显在我们身上}\BibleRef{格后 4:10}，以对比前面所说的，他的死在我们身上被携带。现在，他在这里说的是\Emph{基督的什么生命}？是我们现在活着的生命？那么，为什么在接下来的话中，他劝戒我们不要看那看得见的、暂时的，而要看那看不见的、永远的\BibleRef{格后 4:16}？换句话说，不是现在的，而是将来的？但如果他指的是基督将来的生命，暗示它将在我们身上彰显\BibleRef{格后 4:11}，那么他就清楚预言了肉身的复活\BibleRef{格后 4:14}。他还说\BibleQuote{我们外面的人朽坏}\BibleRef{格后 4:16}，意思不是死后永远的灭亡，而是由劳苦和苦难所致，提及此他之前说：\BibleQuote{为此，我们不颓丧}\BibleRef{格后 4:16}。现在，当他添加关于\Quote{里面的人}也说\BibleQuote{日渐更新}时，他在这里表明了两种结局——身体因考验的磨损而消逝，灵魂因默想应许而更新。\par

% structure: p0030 source=h2 target=subsection reason=relative-heading-rank; base=h2

\subsection{第十二章。\newline{}天上的永恒居所。\Person{戴尔都良}对宗徒关于克服死亡恐惧的安慰教导的优美阐释——这种恐惧在反基督教的迫害下极易产生。基督的审判台——反马西昂的观念。乐园。基督的审判特性与异端对祂的观点不符；宗徒的严厉表明祂是造物主基督的合格宣讲者。}

论到我们这地上居所的房子，当他说：\BibleQuote{我们在天上有一个永久的居所，不是人手所造的} \BibleRef{格后 5:1}，他绝无意暗示，因为它是由造物主的手所建造，就必在死后永远朽坏消散。他讨论这个主题，是为了提供安慰，以对抗对死亡和对这朽坏本身的恐惧，这从下文更明显可见，当他补充道：\BibleQuote{我们在这地上的身体的帐幕里叹息，切望穿上那来自天上的住所，但愿我们脱下之后，不致被发现赤身露体} \BibleRef{格后 5:2}，换句话说，将重新获得我们所失去的，即我们的身体。他又说：\BibleQuote{我们在这帐幕里叹息，并非不愿意脱下，而是想要穿上} \BibleRef{格后 5:4}。他在这里明确说出他在前书中所轻描淡写的话，他写道：死人要复活成为不朽的（指那些已经经历死亡的人），\BibleQuote{我们也要改变}（指天主发现尚在肉身中的人）\BibleRef{格前 15:52}。\Emph{那些}人要复活成为不朽的，因为他们将重新获得他们的身体——且是一个更新的身体，由此而来的不朽；而\Emph{这些}人也要在最后一刻的危机中，由于他们的瞬间死亡，在遭遇敌基督的迫害时，经历改变，在此获得的不是脱去身体，而是\Quote{穿上}那来自天上的外衣。因此，当这些人要在他们（改变了的）身体上穿上这天上衣服时，死人也同样要恢复他们的身体，并在其上穿上一种外衣，即天上的不朽；正是为此，他说：\BibleQuote{这必朽坏的必须穿上不朽的，这必死的必须穿上不死的} \BibleRef{格前 15:53}。一类人是在恢复身体时穿上这（天上的）衣服；另一类人则是作为外衣穿上，虽然他们几乎不会失去身体（由于改变的突然性）。因此，他并非毫无理由地将他们描述为\BibleQuote{确实不愿意脱下}，而是\BibleQuote{愿意穿上} \BibleRef{格后 5:4}，换句话说，他们愿意不经历死亡，而被生命吞没，\Quote{好使这必死的（身体）被生命吞没}，通过在其改变状态的外衣中被从死亡中解救出来。这就是为什么他向我们表明，如果我们被死亡所袭，我们不应悲伤，反倒更好，并告诉我们，我们甚至从天主领受了\OriginalTerm{圣神的凭质}{earnest of His Spirit} \BibleRef{格后 5:5}（由此仿佛被保证拥有\OriginalTerm{穿上}{the clothing upon}，这是我们的希望所在），并且\BibleQuote{我们既然在肉身内，就是与主远离} \BibleRef{格后 5:6}；此外，我们因此应当宁愿\BibleQuote{离开肉身与主同住} \BibleRef{格后 5:8}，从而准备好甚至欣然面对死亡。正是基于此，他告诉我们：\BibleQuote{我们众人必须在基督的审判台前显露出来，为使各人按他在肉身内所行的，或善或恶，领取相应的报应} \BibleRef{格后 5:10}。然而，既然那时要按照人的功过来报应，谁还能与天主算账呢？但他既提到审判台，又提到善行与恶行的区分，就给我们呈现了一位将要宣布两种判决的审判官\BibleRef{格后 5:10}，从而肯定了所有的人都必须以身体出现在法庭前。因为除非针对身体，否则不可能宣判身体所做的事。天主若不在成就功过的那同一状况中惩罚或奖赏某人，便是不公义。\BibleQuote{所以谁若在基督内，他就是新造的人；旧的已经过去；看，一切都成了新的} \BibleRef{格后 5:17}，这样依撒意亚的预言就应验了\BibleRef{依 43:19}。当他（在后面一段中）命令我们\BibleQuote{从一切肉身和血液的污秽中洁净自己}（因为这种实体不能进入天主的国\BibleRef{格前 15:50}）；当他再次\BibleQuote{将教会如同贞洁的童女许配给基督} \BibleRef{格后 11:2}，确实是将配偶许配给配偶，一个形象不能与它所比较的事物的真实本性相反的东西结合和比较。因此，当他称\BibleQuote{假宗徒、欺骗的工人，冒充成他的样子} \BibleRef{格后 11:13}，当然是指他们的伪善，他指责他们有行为混乱的罪，而非虚假教义。因此，这对立是行为上的，而非神明上的。如果\BibleQuote{连撒旦自己也装作光明的天使} \BibleRef{格后 11:14}，这样的断言不可用来损害造物主。造物主不是天使，而是天主。如果他不打算称撒旦为\BibleQuote{天使}——我们和\Person{马西昂}都知道他正是天使——那么就必须说撒旦被变成光明的神明，而不是光明的天使。\WorkTitle{论乐园}是我们一篇论著的标题，其中讨论了该主题所容许的一切。在此，关于此事，我只是感到奇怪，一个在地上没有任何\OriginalTerm{分施}{dispensation}的神，怎么可能拥有一个属于自己的乐园——或许只是像乞丐一样利用造物主的乐园，正如他利用造物主的世界一样？然而，关于人从地上被提升到天上，造物主在\Person{厄里亚}身上给了我们一个实例\BibleRef{列下 2:11}。但更令我惊讶的是（\Person{马西昂}接下来假设的），一个如此良善、仁慈、如此厌恶打击和残暴的天主，竟然指使天使撒旦——还不是他自己的，而是造物主的——\BibleQuote{击打}宗徒\BibleRef{格后 12:7}，然后在他三次恳求释放他时拒绝了他的请求！因此，看来马西昂的神是在模仿造物主的行为，造物主是骄傲者的敌人，甚至\Quote{从高位上推下有权势者}。那么，他就是那位赐给撒旦权力攻击\Person{约伯}的天主，好使他的\Quote{能力在软弱中显得完全}吗？为什么这位谴责迦拉达人的使徒\BibleRef{迦 1:6}仍然保留着法律的公式：\BibleQuote{凭两三个人的口供，句句都可定准} \BibleRef{格后 13:1}？又为何他威胁罪人说\BibleQuote{他必不宽容}他们\BibleRef{格后 13:2}——而他却是最温柔天主的传道者？他甚至宣称：\BibleQuote{主赐给了他权柄，可以严厉对待他们} \BibleRef{格后 13:10}。如今，异端者啊，当他的宗徒让自己如此可畏时，你（付出代价地）否认你的神是可畏惧的对象吧！\par

% structure: p0032 source=h2 target=subsection reason=relative-heading-rank; base=h2

\subsection{第十三章 \BibleBook{罗马}{书}。\Person{圣保禄}不由自主地使用表明天主公义的措辞，即使当他颂扬福音的仁慈时。马西翁尤其恶劣地肢解了这封书信。然而本书作者基于共同立场辩论：最终的审判将依照福音；因信德而成义者被劝勉与天主和好；旧约与新约的治理出于同一只手。}

既然我的小作品即将结束，我对仍在出现的问题只能简要处理，而那些频繁出现的问题则必须搁置。我仍然遗憾必须就法律进行争辩——尽管我已多次证明法律被（福音）取代并不能为另一位神提供论据，因为法律已经在基督内被预言，并且在造物主自己的计划中已为\Emph{祂的}基督所预定。（但我必须回到那个讨论）只要（宗徒引导我，因为）这封书信本身看起来非常像是废除了法律。然而，我们此前已多次指出，宗徒宣称天主是一位审判者；在审判者中隐含了伸冤者；在伸冤者中隐含了造物主。于是，在他所说的这段话中：\BibleQuote{我不以（基督的）福音为耻，因为它是天主的能力，为拯救一切相信的人，先是犹太人，后是希腊人；因为天主的正义在这福音中，从信德到信德被启示出来，} \BibleRef{罗 1:16} 他无疑将福音和救恩都归于那一位——（按照我们异端者自己的区分）——我所称为\Emph{正义的天主}，而非\Emph{善良的天主}的。正是祂将（人）从对法律的信赖转移到对福音的信赖——也就是说，祂自己的法律和祂自己的福音。此外，当他宣告\BibleQuote{天主的忿怒从天上被启示，要攻击一切不虔敬和不正义的人，他们以不正义压制真理} \BibleRef{罗 1:18} 时，（我问）这是什么天主的忿怒？当然是造物主的。因此，真理也将属于祂，忿怒也属于祂，这忿怒必须被启示出来，为真理伸冤。同样，当他加上\BibleQuote{我们确信天主的审判是依照真理} \BibleRef{罗 2:2} 时，他既为那忿怒辩护——这审判为真理而来——同时又提供了另一个证明，即真理出自同一天主，因为他见证祂的审判，从而证实了祂的忿怒。\Emph{马西昂的主张}则完全是另一回事，即造物主在愤怒中报复那被不正义压抑的敌对的神的真理。但是马西昂在这封书信中尤其造成了多大的漏洞，他随意删除整段经文，从我们完整无缺的抄本中便可清楚看出。对我而言，只需接受他认为适合保留下来而未删除的内容作为真理的证据就够了，这些内容同时也奇怪地显明了他的疏忽和盲目。那么，如果天主将审判人的隐秘事——无论是那些在法律内犯罪的人，还是那些没有法律而犯罪的人（因为不认识法律的人，却因本性遵行法律所包含的事）\BibleRef{罗 2:12} ——那么，将要审判的天主必定是那法律和那本性都属于祂的一位，这本性是为不认识法律之人的准则。但祂如何进行这审判？\BibleQuote{照我的福音，藉着（\Emph{耶稣}）基督}，（\Emph{宗徒}）说，\BibleRef{罗 2:16}。因此，福音和基督必然属于祂，法律和本性也属于祂，这两者将由福音和基督来辩护——甚至在那天主的审判中，正如他先前所说，这审判是依照真理 \BibleRef{罗 2:2}。因此，那要辩护真理的忿怒，只能由忿怒的天主从天上启示 \BibleRef{罗 1:18}；因此，这句话与先前宣告审判属于造物主的那句话完全一致，绝不能归于另一位神，他既不是审判者，也不会发怒。它只符合那一位，祂的属性包括我所说的审判和忿怒，并且必然也属于那执行这些属性的媒介，即福音和基督。因此他严厉谴责法律的违犯者，他们教导人不可偷窃，自己却偷窃 \BibleRef{罗 2:21}。（他发出这谴责）完全是为了尊崇天主的法律，并非要指责造物主本人曾命令 \BibleRef{出 3:22} 对埃及人施行欺诈以获取他们的金银，而同时祂却禁止人偷窃，——他们（无耻地）惯于在其他细节上也这样指控祂。那么，我们是否要认为宗徒因害怕而不敢公开诽谤天主，而他却不犹豫地使人离开祂？然而，他在谴责犹太人时，竟引用先知的话来指责他们：\BibleQuote{因你们，天主的名在（外邦人中）被亵渎} \BibleRef{罗 2:24}。但这多么荒谬，他自己竟亵渎那位因亵渎祂而被他斥责为作恶的人！他更喜欢内心的割损，胜过在肉体上忽略割损。现在，法律的目的是使割损成为内心的，而非肉体的；在圣神内，而非在文字上 \BibleRef{罗 2:29}。因为这是耶肋米亚所推荐的割损：\BibleQuote{你们要（为上主）割损（你们的心，除掉）心上的包皮} \BibleRef{耶 4:4}；甚至梅瑟也说：\BibleQuote{因此，要割损你们心的顽硬} ——那割损内心的\Emph{圣神}将出自那位也制定了割损肉体的\Emph{文字}的那一位；而\BibleQuote{内心做犹太人的}将是同一位天主的子民，正如\BibleQuote{外表做犹太人的} \BibleRef{罗 2:28} 也是同一位天主的子民；因为宗徒宁愿根本不提犹太人，除非他是犹太人的天主的仆人。曾经是法律；现在却是\BibleQuote{天主的正义，因（耶稣）基督的信德} \BibleRef{罗 3:21}。这区别是什么意思？你的神是否在侍奉造物主的安排，为祂和祂的法律提供时间？那\Emph{现在}是否在属于\Emph{那时}的那一位手中？那么，法律当然属于祂，现在天主的正义也属于祂。这是时期的区别，而非神的区别。他命令那些因信基督而非因法律成义的人，与天主和好。与哪一位天主？我们从未在任何时期成为祂的敌人的那一位？还是我们反抗过的那一位，无论是在祂的成文法还是自然法方面？现在，既然和好只能与曾经有过战争的那一位达成，我们既将由祂成义，基督也必将属于祂，我们因信德在基督内成义，也唯有藉着祂，天主的敌人才得以和好。\BibleQuote{此外，}他说，\BibleQuote{法律侵入，为叫过犯增多} \BibleRef{罗 5:20}。这是为什么呢？\BibleQuote{为的是，}他说，\BibleQuote{（罪恶在哪里增多），恩宠在那里也格外增多} \BibleRef{罗 5:20}。谁的恩宠，若不是来自那法律也出自祂的那位天主？除非，实际上，造物主插入祂的法律，仅仅是为了给敌对的神的恩宠——一位与祂自己为敌的神（我几乎要说，一位祂不认识的神）——提供一些工作，好\BibleQuote{使罪恶在祂自己的时期如何掌权致死，恩宠也照样藉着正义掌权，藉着耶稣基督获得永生} \BibleRef{罗 5:21}，祂自己的对手！为此（我想，正是如此）造物主的法律已\BibleQuote{把众人都圈在罪恶之中} \BibleRef{迦 3:22}，并带来\BibleQuote{整个世界在（天主面前）都成为有罪}，并\BibleQuote{堵住了众口} \BibleRef{罗 3:19}，以致无人能藉着它夸耀，以使恩宠得以保持，为光荣基督，不是造物主的基督，而是马西昂的基督！我在此可以提前对基督的实体作一评论，因为即将出现一个问题。因为他说\BibleQuote{我们已死于法律}。有人可能会争辩说，基督的身体确实是身体，但不完全是血肉。现在，无论实体是什么，既然他提到\BibleQuote{基督的身体}，并且他紧接着说这身体已\BibleQuote{从死人中复活} \BibleRef{罗 7:4}，那么这身体只能理解为血肉的身体，法律正是针对这身体被称为死亡的法律。但是，请看，他为法律作证，并以罪为理由为它辩解：\BibleQuote{那么，我们该说什么呢？法律是罪吗？绝对不可！} \BibleRef{罗 7:7} 呸，马西昂。\BibleQuote{绝对不可！}（请看）宗徒如何回避对法律的一切指控。然而，除了通过法律，我不知何为罪。但由此我们得到对法律多么高的赞扬，因为通过法律，罪恶的隐藏显现被揭示出来！因此，不是法律使我迷惑，而是\BibleQuote{罪恶，藉着诫命找着机会} \BibleRef{罗 7:8}。那么，你（哦，马西昂）为何将祂的宗徒连法律本身都不敢归咎的事，归咎于法律的天主？不仅如此，他还加上一个高潮：\BibleQuote{法律是圣的，它的诫命是正义而良善的} \BibleRef{罗 7:13}。现在，如果他如此尊敬造物主的法律，我不知道他怎能摧毁造物主本身。谁能作出区分，说有两位神，一位是正义的，另一位是良善的，而祂应该被相信既是正义的又是良善的，因为祂的诫命既是\Quote{\Emph{正义的}和\Emph{良善的}}？接着，当他肯定法律是\BibleQuote{属神的} \BibleRef{罗 7:14}，他因此暗示法律是预言的，是预表性的。现在，即使从这一事实，我也必须断定，基督被法律预言，但是以预表的方式，因此确实并非所有犹太人都能认出祂。\par

% structure: p0034 source=h2 target=subsection reason=relative-heading-rank; base=h2

\subsection{第十四章。在基督降生成人中所显示的天主德能。\Person{圣保禄}短语的意义。罪恶肉身的相似。其中无幻象论。我们真实身体的复活。\Person{马西昂}的删削在书信中造成的巨大鸿沟。当\Person{宗徒}斥责犹太人对天主的恶行时；既然那神是创造者，这事实恰好证明\Person{圣保禄}的神就是创造者。书信末尾\Person{马西昂}所容忍的规诫，显明与创造者的圣经完全一致。}

如果圣父\BibleQuote{派遣了祂的圣子以罪恶肉身的模样}，\BibleRef{罗 8:3}就不该因此说祂所似乎具有的肉身只是一个幻影。\newline{}因为他在前一节将罪归于肉身，并指出它是\BibleQuote{居住在他肢体中的罪恶法律}，并\BibleQuote{与心神的法律交战}。\newline{}因此，（他是否想说）圣子被派遣以罪恶肉身的模样，是为了用一个类似的实体，即一个带有肉身性的实体，来救赎这罪恶的肉身，这实体与罪恶的肉身相似，但它本身却无罪。\newline{}这正是天主性完美的能力，在一个与人性相似的本性中实现（人的）救恩。\newline{}因为若天主的圣神医治了肉身，那不算是大事；但当那一个肉身，正是罪恶实体的复制品——本身也是肉身——只是没有罪，（实现这医治，那么无疑就是一件大事）。\newline{}因此，这\Emph{模样}所指涉的是罪恶性的特质，而非实体的任何虚假性。\newline{}因为他若有意将\BibleQuote{模样}如此加于实体从而否认其实在性，他就不会添加\BibleQuote{有罪的}这一属性；在那种情况下，他只会用\BibleQuote{肉身}一词，而省略\BibleQuote{有罪的}。\newline{}但既然他将两者放在一起，说了\BibleQuote{有罪的肉身}（或\BibleQuote{罪恶的肉身}），他就既肯定了实体，即肉身，又将这\Emph{模样}指向实体的过失，即它的罪。\newline{}但即使假设模样是实体的属性，那实体的真相也不会因此被否认。\newline{}那么为什么\Emph{称}真实的\Emph{实体为相似}呢？因为它确实是真实的，只是不属于与我们同等状态的种子；但它仍然是真实的，因为它的本性并非真的与我们不同。\newline{}再者，在相反的事物中没有相似。\newline{}因此，肉身的相似不会被称为\Emph{灵}，因为肉身不能与灵有任何相似；但如果它似乎是它实际所不是的，则它会被称为\Emph{幻影}。\newline{}然而，它被称为\Emph{模样}，因为它是它看似所是的。\newline{}现在它是（它看似所是的），因为它与另一事物（与之相比的事物）处于同等地位。\newline{}但一个幻影，仅仅如此而无他，不是模样。\newline{}然而，宗徒本人在这里\Emph{帮助了我们；因为}，当他解释在何种意义上他不愿我们\BibleQuote{生活在肉身中}，尽管在肉身中——甚至是不生活在肉身的行为中——他表明当他写下\BibleQuote{血肉不能继承天主的国}，\BibleRef{格前 15:50}时，他并非意在谴责（肉身的那）实体，而是其行为；因为我们仍在肉身中时，这些行为是可能不犯的，因此它们应被恰当地归咎于肉身的实体，而是其行为。\newline{}同样，如果\BibleQuote{身体确实因罪而死}（由此我们看出，所指的并非灵魂的死，而是身体的死），\BibleQuote{但灵因义而成为生命}，\BibleRef{罗 8:10} \Emph{那么这生命}归于那因罪而招致死亡的，即\Emph{正如我们刚才所看到的}，身体。\newline{}现在\Emph{身体}只被归还给失去它的人；因此死者的复活意味着他们身体的复活。\newline{}他于是接着说：\BibleQuote{那从死者中复活了基督的，也必使你们有死的身体复活。}\BibleRef{罗 8:11}在这些话中，他既肯定了肉身的复活（没有肉身就不可正当地称为身体，也不能适当地视为有死的），又证明了基督的身体性实体；因为我们将来的有死的身体将按祂复活的方式而复活；而祂复活的方式无非是在身体内。\newline{}我在这里要跳过一段很宽的经文被删去的鸿沟；然而，我落到宗徒为以色列作证的地方：\BibleQuote{他们对天主怀有热忱}——当然是他们自己的天主——\BibleQuote{但不是按着真知。因为，}他说，\BibleQuote{他们不认识天主的正义，而要建立自己的正义，就没有服从天主的正义；因为基督是法律的终结，使一切相信的人都成义。}\BibleRef{罗 10:2}于此，我们将面对异端者的一个论点，即犹太人不知道那更高的天主，因为他们反对他而建立自己的正义——即他们法律的正义——不接受基督，法律的终结（或完成者）。\newline{}但既然如此，他如何能证明他们对他们自己的天主的熱忱，如果他指责他们的无知不是关于同一个天主？\newline{}他们确实对天主有热忱，但这不是明智的热忱：事实上，他们不认识祂，因为他们不认识祂通过基督的安排，基督将使法律圆满；他们就这样坚持自己的正义来反对祂。\newline{}但造物主自己也证明他们对祂的无知：\BibleQuote{以色列不认识我；我的子民不明白我；}\BibleRef{依 1:3}至于他们偏爱建立自己的正义，（造物主又描述他们）\BibleQuote{将人的诫命当作教义教导人；}此外，\BibleQuote{聚集起来反对上主和祂的基督}——当然是由于对祂的无知。\newline{}现在，没有任何关于另一个天主的事可以用到造物主身上；否则宗徒责备犹太人对于他们一无所知的天主的无知，就不会是公正的。\newline{}他们如果只是坚持他们自己的天主的正义来反对他们不认识的天主，那么他们的罪在哪里呢？\newline{}但他呼喊说：\BibleQuote{啊，天主的富裕和智慧何其深奥！祂的决断何等不可探究！}\BibleRef{罗 11:33}这情感的迸发从何而来？无疑是由于他之前反复思考的圣经的回忆，以及他上面所阐述的与来自法律的基督信仰有关的奥秘的默想。\newline{}如果马西昂在其删改中有目的，为什么他的宗徒发出这样的惊叹，因为他的天主没有任何财富可供他默想？他是如此贫穷匮乏，以至于他未曾创造任何东西，未曾预言任何东西——简言之，他什么也没有拥有；因为他是进入另一个天主的世界而下降的。\newline{}事实上，造物主的资源和财富，曾经被隐藏的，如今被揭示。\newline{}因为祂曾这样应许：\BibleQuote{我要赐给他们隐藏的宝藏，并将人们未见过的事物开启给他们。}\BibleRef{依 45:3}因此，就有了这惊叹\BibleQuote{啊，天主的富裕和智慧何其深奥！}因为祂的宝藏正在开启。\newline{}这就是依撒意亚所说的话，以及（宗徒自己）随后引用的同一先知的话的要旨：\BibleQuote{谁知道主的心意？谁做过祂的参谋？谁先给了祂，而得到偿还呢？}\newline{}现在，（马西昂，）既然你从圣经中删除了这么多，为什么你保留了这些话，好像它们也不是造物主的话？\newline{}但来吧，让我们无误地看看你的新天主的规诫：\BibleQuote{恶要厌恶，善要持守。}\BibleRef{罗 12:9}那么，这规诫与造物主的教训不同吗？\BibleQuote{从你们中间除掉恶，远离它，并要行善。}\newline{}\Emph{然后又说}：\BibleQuote{要以兄弟之爱彼此亲爱。}\BibleRef{罗 12:10}这不正等同于：\BibleQuote{你应爱你的邻人如你自己}？\BibleRef{肋 19:18}（再者，你的宗徒说：）\BibleQuote{在希望中喜乐；}\BibleRef{罗 12:12}即，对天主的希望。\newline{}\Emph{造物主的圣咏作者这样说}：\BibleQuote{依赖上主，胜过依赖王侯。}\newline{}\BibleQuote{在患难中忍耐。}\BibleRef{罗 12:12}你在圣咏中看到：\BibleQuote{愿上主在患难之日应允你。}\newline{}\BibleQuote{祝福，不可咒骂，}\BibleRef{罗 12:12}（你的宗徒说。）但你还能找到比那位创造了万物并祝福了它们的那一位更好的老师吗？\newline{}\BibleQuote{不可心高气傲，要俯就卑微的人。不要自以为聪明。}\BibleRef{罗 12:16}因为依撒意亞曾对这样的心态宣告祸哉。\BibleRef{依 5:21}\newline{}\BibleQuote{不要以恶报恶。}\BibleRef{罗 12:17}（与此相似的是造物主的规诫：）\BibleQuote{不可心里恨你的兄弟。}\BibleRef{肋 19:17}（又说：）\BibleQuote{你们不要为自己复仇；}\BibleRef{罗 12:19}\Emph{因为经上记载}，\BibleQuote{复仇在我，我必报应，主说。}\newline{}\BibleQuote{要尽力与众人和睦。}\BibleRef{罗 12:18}因此，法律的报复并不允许对伤害进行报应；它反而通过害怕报复来抑制任何这样的企图。\newline{}因此，他非常恰当地将造物主的全部教训总结为这一条规诫：\BibleQuote{你应爱你的邻人如你自己。}\BibleRef{罗 13:9}如今，如果这是出自法律本身的综合总结，我茫然不知谁是法律的天主。我担心（毕竟）它一定是马西昂的天主。\newline{}如果基督的福音也在这同一规诫中得以实现，但不是造物主的基督，那么我们继续争论基督是否曾说过\BibleQuote{我来不是为废除法律，而是为成全}又有什么用呢？\BibleRef{玛 5:17}那（我们的）本都人徒劳地努力否认这句话。\newline{}如果福音没有成全法律，那么我只能说，法律成全了福音。\newline{}但好在在后面一节中他用\BibleQuote{基督的审判台}来威胁我们——当然是审判者和报应者，因此是造物主的基督。\newline{}这位\Emph{造物主}，尽管他可能极力宣扬另一位天主，但若他如此将祂作为可畏惧的对象提出，他肯定向我们呈现为一位应受事奉的存在。\par

% structure: p0036 source=h2 target=subsection reason=relative-heading-rank; base=h2

\subsection{第十五章  \WorkTitle{致得撒洛尼人前书}。短篇书信涵意尖锐而极有价值。\Person{圣保禄}指责犹太人先杀害先知，后杀害\Person{基督}。这证明\Person{基督}与先知同属一位天主。自然律——实即造物主的诫律——与\Person{基督}的福音都命令贞洁。复活在旧约中已由\Person{基督}预备好。人的复合本性。}

我不会不乐于留心那些较短的书信，即使篇幅简短，其中仍有许多锐利之处。犹太人杀害了他们的先知。\BibleRef{得前 2:15} 我可以问：这跟那位敌对之神的宗徒有何关系？他本人如此和蔼，几乎连自己子民的过失也不谴责；而且，他本人也在消灭那些他正在摧毁的先知中有一份。以色列在杀害那些他也已弃绝的先知时对他做了什么伤害？因为他率先对他们作出了敌意的判决。然而，\Emph{以色列}得罪了他们自己的\Emph{天主}。他向那受损的\Emph{天主}指责他们的不义；他绝不是受损之\Emph{神的对头}。否则，他就不会用这些话加重他们杀害主的罪责：\BibleQuote{他们杀害了主耶稣和他们的先知}，虽然\Emph{他们的}这个词是异端者添加的。然而，他们杀害那宣告新神的基督，又杀害他们自己神的先知，这有什么特别尖刻之处？然而，他们杀害主及其仆人，被作为递进的事例提出。那么，如果他们杀害的是一位神的基督和另一位神的先知，他必定会将这两种不敬之罪同等并列，而不是以递进的方式提及；但它们不能相提并论；因此，递进之所以可能，只是因为罪行实际上是在两种情况下对同一位主犯下的。因此，基督和先知属于同一位主。那\BibleQuote{我们的圣化}，他宣称是\BibleQuote{天主的旨意}，你可以从他禁止的相反行为中得知。我们应该\BibleQuote{远避邪淫}，而不是远避婚姻；每个人都应该\BibleQuote{晓得怎样以圣洁守住自己的器皿}。\BibleRef{得前 4:3} 如何做？\BibleQuote{不像那些不认识天主的外邦人纵情于贪欲}。\BibleRef{得前 4:5} 然而，贪欲并不归于外邦人中的婚姻，而是归于放荡、不自然和巨大的罪。自然律反对奢侈、粗俗和不洁；它不禁止婚姻的结合，而是禁止贪欲；它借着婚姻的尊荣来保护我们的器皿。我要这样处理这段经文，以维持另一种更高圣洁的优越性，将节欲和童贞置于婚姻之上，但绝不禁止后者。因为我的敌意是针对那些想要毁灭婚姻之神的人，而非那些追随贞洁的人。他说那些\BibleQuote{直到主来还活着的人}连同\BibleQuote{在基督内死了的人}将\BibleQuote{先复活}，并被\BibleQuote{提到云中，到空中迎接主}。\BibleRef{得前 4:15} 我发现，正是因为预见了这一切，天上的智慧才曾赞叹\BibleQuote{那在上面的耶路撒冷}，\BibleRef{迦 4:26} 并曾借依撒意亚的口说：\BibleQuote{那些如云飞来，如鸽子回到巢窝的，是谁呢？}\BibleRef{依 60:8} 既然基督为我们预备了这升天的途径，\Emph{祂}必是亚毛斯所说的那一位：\BibleQuote{祂建造了自己的苍穹上达于天}，\BibleRef{亚 9:6} 为祂自己和祂的子民。那么，除了那位我听见应许这事的天主，我还能从谁那里期待这一切的实现呢？他禁止我们\BibleQuote{消灭圣神}，\BibleQuote{轻看先知之恩}？\BibleRef{得前 5:19} 当然不是造物主的神恩，也不是造物主的先知之恩——马尔西昂会这样回答。因为他已经消灭并轻看了他所摧毁的东西，无法禁止他所轻看的。因此，马尔西昂现在必须在他的教会中展示他那不可消灭的神恩（即他那神的灵）和不可轻看的先知之恩。既然他按自己所愿作出了展示，让他知道我们将以圣神和先知之恩的恩宠与权能的准则来检验它——即预言未来、揭示心中的秘密和解释奥秘。当他未能产生并证明任何这样的准则时，我们将反过来展示造物主的圣神和先知之恩，它们按照祂的旨意发出预言。这样，就清楚看出宗徒所说的话，正是那些要在他天主的教会中发生的事；只要祂存在，祂的圣神也照样运行，祂的应许也照样被宣告。来吧，你们这些否认肉身救恩的人，每当这种场合明确提到\Emph{身体}时，便将此词解释为任何东西，唯独不是肉身实体，那么，宗徒是如何给予所有（我们的官能）各自独特的名称，并将它们全包含在一个祈祷中，愿我们的\BibleQuote{精神、灵魂和身体得蒙保全，直到我们的主耶稣基督的来临}呢？他在这里提出了灵魂和身体作为两个各自不同的事物。因为灵魂虽有其独特性质的身体，正如精神也有，但既然灵魂和身体都被明确命名，灵魂有其特有的名称，不需要共同称谓\Quote{身体}。这称谓留给了\Quote{肉身}，它在（这段经文中）没有专有名称，因此必然使用共同称谓。的确，在\Emph{精神}和\Emph{灵魂}之后，我在人身上看不到其他可以用\Quote{身体}一词指代的实体，除了\Quote{肉身}。因此，每当\Quote{身体}一词未被明确提及时，我理解它就是指肉身。在当前段落中更是如此，因为肉身明确被称为\Quote{身体}。\par

% structure: p0038 source=h2 target=subsection reason=relative-heading-rank; base=h2

\subsection{第十六章　\BibleBook{得撒洛尼人}{后书}。\Person{马尔西昂}荒谬的删节；其意图昭然若揭。对外邦人和犹太人的最后审判不能由\Person{马尔西昂}的基督执行。那罪人是什么？\Person{马尔西昂}观点的不一致。假基督。最后背道中的大事件都在造物主的天主眷顾和旨意之内，万物从一开始就属于祂。\Person{保禄}的规诫与造物主的规诫相似。}

我们有时必须回到某些主题，以肯定与之相关的\Emph{真理}。我们在此重申：既然宗徒宣告主是美善与灾祸的施予者，那么祂必是造物主，或是（马尔西昂不愿承认的）一位像造物主者——\BibleQuote{以祂为公义，要将患难报应那加患难给我们的人，却使我们这受患难的人得平安，那时主耶稣将从天上，同祂大能的天使，在\Emph{那有火焰的火}中显现。}\BibleRef{得后 1:6} 然而，那异端者却删去了\Emph{那有火焰的火}，无疑是为了抹去其中我们天主的一切痕迹。但这删除的愚昧显而易见。因为宗徒宣告主将到来，\BibleQuote{要报应那不认识天主、不听从我们主耶稣基督福音的人，}他说，\BibleQuote{他们要受刑罚，就是永远沉沦，离开主的面和祂权能的荣光}\BibleRef{得后 1:8}——由此可知，既然祂来施罚，就必须要求\Quote{那有火焰的火}。因此，基于这一考虑，尽管马尔西昂反对，我们也必须断定：基督属于一位点燃（报应）之火的天主，因而属于造物主，因为祂向那些不认识主的人，即外邦人，施行报应。因为他分别提到了\BibleQuote{不听从\Emph{我们主耶稣基督}福音的人}\BibleRef{得后 1:8}，无论是基督徒中的罪人还是犹太人中的罪人。然而，对外邦人——他们很可能从未听过福音——施加惩罚，并非那本性不可知、仅在福音中启示、因此不能为所有人认识的天主的职分。然而，造物主甚至可由自然（之光）认识，因为可从祂的化工领悟祂，因而成为更广泛认知的对象。因此，惩罚那些不认识天主的人属于祂，因为无人应不认识祂。在（宗徒的）措辞中，\BibleQuote{离开主的面和祂权能的荣光}\BibleRef{得后 1:9}，他采用了依撒意亚的话，后者特意使同一主\Quote{兴起，猛烈震动大地}。然而，那罪人，那丧亡之子，必须首先在主到来之前被显露出来；\BibleQuote{他敌对且高举自己在一切称为神或受崇拜者以上，以致坐在天主的殿中，自充为神}\BibleRef{得后 2:3}。按我们的观点，他确实是\OriginalTerm{假基督}{Antichrist}；正如古先知和新先知，\Emph{尤其}是宗徒若望所教导的，他说：\BibleQuote{已经有许多假先知来到了世界上}\BibleRef{若一 4:1}，这些假基督的前驱，否认基督曾降生成人，也不承认耶稣（是基督），即在造物主天主内。然而，按马尔西昂的观点，实在难以知道祂是否（终究）不是造物主的基督；因为按他的说法，祂尚未到来。但无论两者中的哪一个，我想知道为何他\BibleQuote{用各种异能、征兆和欺骗的奇迹}\BibleRef{得后 2:9}而来？\BibleQuote{因为，}他说，\BibleQuote{他们没有接受爱真理的心，为得救；为此，天主就给他们一种错误的能力，使他们相信谎谬，叫一切不信真理、反喜爱不义的人，都被定罪}\BibleRef{得后 2:10}。因此，如果他是假基督（如我们所持），并按照造物主的旨意而来，那么必是造物主天主差遣他，使那些不信真理的人固守他们的错误，以使它们得救；真理和救恩同样属于祂，祂借差遣错误作为替代来报应（对真理的蔑视）——即造物主，正是那愤怒的属性属于祂，祂用谎言欺骗那些不被真理吸引的人。然而，如果祂不是假基督，如我们所料，那么祂就是造物主的基督，如马尔西昂所愿。既然如此，祂怎能唆使造物主的基督为祂的真理报仇呢？但若他最终同意我们，此处所指的是假基督，那么我也必须同样追问，他为何认为撒殚——造物主的一位天使——对他的目的是必要的？再者，假基督为何要被祂击杀，而他正受造物主差遣执行激起人对不真理的热爱之职能？简言之，无可争辩的是：那使者、真理和救恩都属于祂，愤怒、嫉妒和\BibleQuote{发送那强有力的迷惑}\BibleRef{得后 2:11}给轻视、嘲笑以及不认识祂的人，也属于祂；因此，连马尔西昂现在也必须退一步，向我们承认他的神是\Quote{\OriginalTerm{一个嫉妒的天主}{a jealous god}}。（既然这是无可置疑的立场，我请问）哪一位天主更有权发怒？我想，是从万有之始就赐给人——作为认识祂的首要见证——自然界的种种化工、仁慈的眷顾、灾祸和（神性）征兆，但尽管有这一切证据仍未被认识的那一位？还是那仅仅在一份福音抄本中（甚至没有可靠权威）被一次展示，实际上毫不隐瞒地宣告另一位天主的他？如今，那有权施行报应的，也唯一有权拥有引起报应的事物，即福音；（换言之，）真理和（与之相伴的）救恩。那命令：\BibleQuote{若有人不肯工作，就不应当吃饭}\BibleRef{得后 3:10}，与那颁布禁令者的诫命严格一致：\BibleQuote{牛在场上踹谷的时候，不可笼住它的嘴}\BibleRef{申 25:4}。\par

% structure: p0040 source=h2 target=subsection reason=relative-heading-rank; base=h2

\subsection{第十七章 \WorkTitle{致劳狄刻雅人书}。正确称呼是\WorkTitle{致厄弗所人书}。在基督内从创造之初总括万有。这里没有马西昂的基督的位置。本书信与\WorkTitle{旧约}章节之间的众多平行对应。空中权势的首领和今世的神——是谁？创造与再造是同一天主的工程。基督如何使法律过时。马西昂的涂改是徒然的。宗徒们与先知们一样来自造物主。}

我们由教会的真确传承得知，这封书信是写给厄弗所人的，而非劳狄刻雅人。然而，马尔吉昂却极欲给它加上一个（劳狄刻雅人的）新标题，仿佛他在考究这一点上极为精确。但标题有何要紧呢？因为宗徒写给某一教会时，实际上就是写给所有教会。故凡他所写信的对象是谁，总可以肯定：他宣称，在基督内的天主，就是那一切预言的都与之一致的天主。现在，关于那\Quote{好主意，就是天主在祂旨意的奥秘中所预定的，要在时期圆满的分期——若我可依此希腊词的准确含义这样说——使天上和地上的万有，总归于基督}的一切事，最合宜地属于哪一位天主呢？\BibleRef{弗 1:9} 岂非那万物从起初便属祂，连起初本身也属祂；从祂产生时期和时期圆满的分期，藉此分期直到最初的一切都在基督内总归的那一位？另一位神有什么\Quote{起初}呢？也就是说，既然他没有任何工程可显示，从他能产生什么呢？若无起初，怎能有\Quote{时期}？若无时期，怎能有\Quote{时期圆满}？若无圆满，又怎能有\Quote{分期}？实在，他在地上曾做过什么，以致任何长久的、待完成的时期的分期可算在他账上，为使万有——连天上的——都在基督内完成？我们绝不可能设想，在天上曾有任何事由另一位神完成，而所有人公认天下万事都是由那一位完成的。倘若，从起初的这一切事除归给造物主外，不可能归给任何别的神，那么谁会相信一位异样的神将这些事总归于一位异样的基督，而非由它们自己的作者在自己的基督内完成？反之，若这些事属于造物主，它们就必然与另一位神分开；若分开，就与他敌对。但敌对者怎能被总归于那正是要消灭它们的一位？再则，当宗徒说：\Quote{我们为颂扬祂的光荣，那首先在基督内怀着希望的人}\BibleRef{弗 1:12}，这些话宣布的是怎样的基督？谁能在天主来临之前首先——即预先——信赖天主，除非是犹太人，因为基督从起初就预先被宣布给他们？那被预先宣告的，也被预先信赖。因此宗徒将这话归于自己，即归于犹太人，以便随后与异民区分：（他接着说）\Quote{在你们听见真理之言，就是你们得救的福音后，也在祂内怀着希望；你们相信了，就受了所应许的圣神的印证。}\BibleRef{弗 1:13} 这是什么应许？那是通过岳厄尔所作的应许：\Quote{在末日，我要将我的神赐在一切有血肉的人身上}，\BibleRef{岳 2:28} 即赐给万民。因此，圣神和福音必在那被预先信赖（因为预先宣告）的基督内被发现。再者，\Quote{光荣的父}\BibleRef{弗 1:17} 就是那在圣咏中，当祂的基督升天时被歌颂为\Quote{光荣的君王}的一位：\Quote{谁是这位光荣的君王？万军的上主，祂就是光荣的君王。}（咏 24:8-10）从祂那里也求得\Quote{智慧和启示的神}，\BibleRef{弗 1:17} 依撒意亚（依 11:2）曾陈述那分赐恩宠之神的七重。祂也赐予\Quote{心眼的光照}，\BibleRef{弗 1:18} 祂曾用光丰富了我们的肉眼；此外，百姓的眼瞎令祂不悦：\Quote{谁眼瞎，不是我的仆人吗？……是的，天主的仆人成了眼瞎的。}在祂的恩赐中也有\Quote{圣徒中所应得的光荣产业的丰裕}，\BibleRef{弗 1:18} 这产业是祂在召叫外邦人时所应许的：\Quote{你向我请求，我必将万民赐给你作产业。}（咏 2:8）是祂\Quote{在基督身上施展了那强大能力，使祂从死者中复活，并叫祂坐在自己右边，将万有放在祂脚下}，\BibleRef{弗 1:19} 正是那说：\Quote{你坐在我右边，等我使你的仇敌作你的脚凳}的一位。（咏 110:1）因为在另一处，圣神论及圣子对圣父说：\Quote{你使万有屈服在祂脚下。}（咏 8:6）如今，若由所有在造物主内发现的这些事，竟然推论出另一位神和另一位基督，那么让我们去寻索造物主吧。我想，我们找到祂了，当祂谈到这样一些人：\Quote{你们因过犯和罪恶原是死的；那时你们生活在过犯中，随从这世界的潮流，顺服空中权势的首领，就是现今在悖逆之子身上运行的灵}时。\BibleRef{弗 2:1} 但马尔吉昂不可在此将\Quote{世界}解释为这世界的神。因为受造物不与造物主相似；被造之物不与造物者相似；世界不与天主相似。此外，那万代权势的首领，不该被认为称为空中权势的首领；因为那高于更高权势的主，不会从较低权势得名，虽然这些也可归于祂。再者，祂看来也不可能是那不信的煽动者，这不信祂自己必须在犹太人和外邦人那里忍受。因此我们可简单断定，这些称号不适合造物主。另有别者更适合它们——宗徒很清楚那是谁。那么是谁？无疑就是那自祂占有自己的\Quote{空气}以来，就兴起\Quote{悖逆之子}反对造物主本身的一位；正如先知使祂说：\Quote{我要将我的宝座安置在星辰之上；……我要升到云层之上；我要与至高者相似。}（依 14:13-14）这必是指魔鬼，在另一处（因为那里他们要有宗徒的意思），我们会在\Quote{这世界的神}这个称呼中认出他来。因为他以自己神性的虚假伪装充满了整个世界。确实，若他不存在，我们或许可能将这些描述用于造物主。但宗徒也曾生活在犹太教中；当他插入论及（那段时期的）过犯时说：\Quote{从前我们在过犯中也都在其中生活过}，\BibleRef{弗 2:3} 他不可被理解为指向造物主是有罪之人的主和这空气的君王；而是指在他自己的犹太教中，他本是悖逆之子的一员，有魔鬼作他的煽动者——当他迫害教会和创造主的基督时。因此他说：\Quote{我们也都是可怒之子}，但\Quote{本性的}。\BibleRef{弗 2:3} 但让异端者不要争辩说，因为造物主称犹太人为\Quote{子女}，所以造物主是忿怒的主。因为当宗徒说：\Quote{我们本性的可怒之子}时，既然犹太人不是\Quote{本性的}造物主的子女，而是由祖先蒙拣选的，他（必然）将他们是可怒之子归于本性，而非归于造物主，最后加上：\Quote{也和别人一样}，\BibleRef{弗 2:3} 那些人当然不是天主的子女。显然，罪、肉体的私欲、不信和忿怒，都归于全人类共同的本性，但魔鬼引领那本性步入歧途，这本性他已用植入的罪种感染了。他说：\Quote{我们是祂的化工，在基督耶稣内受造的。}\BibleRef{弗 2:10} 作（如工匠）与创造是不同的事。但他将两者归于同一位。人是造物主的化工。因此，那起初造人者，也在基督内创造了他。就本性的实体而言，他\Quote{造}了他；就恩宠的工作而言，他\Quote{创造}了他。也看这些话之后接着说的：\Quote{所以你们应当记起：从前你们按肉体是外邦人，被那些称为受割损的人——在肉体上受手工割损者——称为未受割损的；那时你们没有基督，与以色列社团隔绝，对恩许的盟约是外方人，没有希望，在世界上没有天主。}\BibleRef{弗 2:11} 现在，这些外邦人没有哪一位天主和没有哪一位基督？当然是那以色列社团、盟约和恩许所属的那位。\Quote{但现今在基督耶稣内，你们从前远离的人，因着基督的血，成为亲近的了。}\BibleRef{弗 2:13} 从前他们远离谁？远离他上面所说的（特恩），即远离造物主的基督、以色列社团、盟约、恩许的盼望、甚至天主自己。既是这样，外邦人如今在基督内便亲近这些恩典了，从前他们本是远离的。但若我们在基督内如此接近以色列社团（它包括神圣造物主的宗教）、盟约和恩许，甚至到他们自己的天主面前，那就十分荒谬（假设）另一位神的基督将我们从远处带到与造物主如此接近。宗徒想起曾预言外邦人从遥远的疏离中被召叫，正如这样的话：\Quote{那远离我的人都来到我的正义。}因为造物主的正义和祂的和平一样，都是在基督内宣布的，正如我们多次指出的。因此他说：\Quote{祂是我们的和平，使两者合而为一}，\BibleRef{弗 2:14} 即犹太民族和外邦世界。那近的和那远的，如今因\Quote{拆毁了中间阻隔的墙壁}，即他们的\Quote{仇恨}，\Quote{在祂的肉身上}合而为一了。\BibleRef{弗 2:15} 但马尔吉昂删除了\Quote{祂的}这个代词，为要使仇恨指向肉身，仿佛（宗徒）说的是血肉的仇恨，而非那与基督争敌的仇恨。这样你就如我已在别处说过的，展示了澎都人的愚拙，而非玛鲁基人的灵巧；因为你在此否认他有\Quote{血肉}，而在上一节你却允许他有\Quote{血}！然而，既然祂以自己的一切诫命废除了法律，甚至由自己成全了法律（因为\Quote{不可奸淫}是多余的，当祂说：\Quote{凡注视妇女而有意贪恋她的，已犯奸淫}；\Quote{不可杀人}也是多余的，当祂说：\Quote{不可毁谤你的邻人}），所以不可能使那如此促进法律的人成为法律的仇敌。\Quote{为将两者在自己内造成一个新人}，因为那\Quote{造}的也是\Quote{创造}的同一位（正如我们在上面已发现的：\Quote{因为我们本是祂的化工，在基督耶稣内受造的}），\BibleRef{弗 2:10} \Quote{成就和平}——真正的新人、真正的人，而非幻像——却是新的、由天主圣神生于童贞女，\Quote{为将双方和好于天主}，\BibleRef{弗 2:15} 即那两民族（犹太人和外邦人）都得罪了的天主，\Quote{藉着十字架，在祂一个身体内消灭了仇忿}。\BibleRef{弗 2:16} 由此我们也从这段经文发现，在基督内有一个能忍受十字架的血肉身体。\Quote{因此祂来，向你们远离的人，也向那近的人传布了和平}；我们双方都得以\Quote{进到父面前}，如今\Quote{不再是外方人和旅客，而是与圣徒同国，是天主家里的人}\BibleRef{弗 2:19}——即我们上面已显示曾经是疏远、远离的那位——\Quote{被建筑在宗徒的基础上}，\BibleRef{弗 2:20}（宗徒又加上）\Quote{和先知的基础上}；然而异端者删除了这些字，忘记主在教会内不仅设立了宗徒，也设立了先知。他无疑害怕我们的建筑在基督内是建立在古先知的基础上，因为宗徒本人从未停止处处以先知的话建造我们。他从哪里学会称基督为\Quote{屋角的基石}，\BibleRef{弗 2:20} 岂不是从圣咏中给他的形象：\Quote{匠人所弃的石头，已成了屋角的基石}？（咏 118:22）\par

% structure: p0042 source=h2 target=subsection reason=relative-heading-rank; base=h2

\subsection{第十八章 \Person{马西昂}另一愚蠢删改的揭露。宗徒某些源自旧约语言的比喻性表述。将此书信诸多段落与\WorkTitle{梅瑟五书}、\WorkTitle{圣咏集}及\WorkTitle{先知书}中的训诫和记载相互对照。它们都教导我们造物主的旨意和计划。}

既然我们的异端者如此喜爱他的修剪刀，当他亲手删去音节时我并不惊讶，因为整页整页的内容通常是他施以涂抹功夫的对象。宗徒对自己声明说，\Quote{我原是众圣人中最微小的，竟蒙得了这恩宠}，得以光照众人，使他们知道\Quote{这奥秘的共融，是从永远在创造万有的天主内隐藏的}。\BibleRef{弗 3:8} 异端者删去了介词\Emph{在...内}，使句子变成：\InnerQuote{（这奥秘的共融）是从永远在创造万有的天主面前隐藏的。} 然而，这种篡改是极其荒谬的。因为宗徒接着推论（从他的陈述）：\Quote{为使天上的率领者和掌权者，借着教会，得知天主各种的智慧。}\BibleRef{弗 3:10} 他指的是\Emph{谁的}率领者和掌权者？如果是创造者的，那么像祂这样的天主怎么会愿意将祂的智慧显示给率领者和掌权者，而对自己却不显示呢？因为显然，没有任何率领者能在没有他们主权之主的情况下理解任何事情。或者，如果宗徒在这段里没有提及天主，是因为祂（作为他们的首领）自己也被列在这些（率领者）中间，那么他就已经清楚地说过，这奥秘是从创造万有的那一位的率领者和掌权者面前隐藏的，而把祂包括在他们中间。但如果他说这是向他们隐藏的，就必须理解为对祂是明显的。因此，奥秘并非\Emph{从天主}隐藏，而是\Emph{在天主内}，即创造万有的天主内，向祂的率领者和掌权者隐藏的。因为\Quote{谁曾知晓上主的心意，或谁曾作过祂的参谋？}\BibleRef{依 40:13} 落入这个陷阱后，异端者很可能更改了经文，企图说明\Emph{他的}神愿意向他的率领者和掌权者显明他自己奥秘的共融，而创造万有的天主对此一无所知。但他硬要提出\Emph{创造者}的这种无知有什么用呢？创造者与至高神是陌生的，且远远游离于祂，连他自己的仆人们对他都一无所知？然而，未来对创造者却是清楚知晓的。那么，为何那必须在祂的天穹下、在祂的大地上被启示之事，祂却不知呢？由此，我们先前所确立的便得到了证实。因为既然创造者迟早必定会知道至高神那隐藏的奥秘，即使假定正确的读法是（如\Person{Marcion}所主张的）\InnerQuote{向创造万有的天主隐藏}，那么，他就应该这样表达结论：\InnerQuote{为使天主各种的智慧得以向祂显明，然后向天主的率领者和掌权者显明，无论祂是谁，创造者注定要与他们共享这知识。} 如此，当按照其本真含义来读时，这段中的涂改之迹就十分明显了。而我，现在想与你就宗徒的寓意表达进行讨论。那新神能在先知书中找到什么（适合他自己的）修辞呢？宗徒说：\InnerQuote{祂掳掠了俘虏。} 用什么武器？在什么战役中？由于哪个国家的沦亡？哪个城市的倾覆？这征服者把哪些妇女、哪些孩童、哪些王子投入锁链？因为当达味歌颂基督\InnerQuote{腰间佩剑束带}，或以撒意亚说\InnerQuote{夺去撒玛黎雅的掠物和大马士革的权势}\BibleRef{依 8:4}时，你便认为祂是肉眼可见的真实战士。那么现在要知道，祂的武装和战争是属灵的，因为你们已经发现那俘虏是属灵的，为使你们进一步知道\Emph{这}也属于祂，正是因为宗徒从同一批先知书中引用了俘虏的提及，同样也从他那里得到了他的规劝：他说：\InnerQuote{你们应戒绝谎言，各人与近人说真话。}\BibleRef{弗 4:25} 又，用圣咏表达其意的同样字句，他说：\InnerQuote{你们应发怒，却不要犯罪；}\BibleRef{弗 4:26} \InnerQuote{不可让太阳在你们的愤怒中落下。}\BibleRef{弗 4:26} \InnerQuote{不要与黑暗无益的作为有分；}\BibleRef{弗 5:11} 因为（在圣咏中写道：）\InnerQuote{与圣洁的人，你将成为圣洁；与乖僻的人，你将成为乖僻；}并且：\InnerQuote{你应从你们中间除去邪恶。} 又说：\InnerQuote{从他们中间出来，不要触摸不洁之物；你们这些抬上主器皿的人，要自洁。} （宗徒接着说：）\InnerQuote{不要醉酒，醉酒使人放荡，}\BibleRef{弗 5:18} ——这一规诫是由先知书中的一段所暗示的，那里谴责了那些引诱献身者（纳齐尔人）醉酒的人：\InnerQuote{你们给我圣者酒喝。}\BibleRef{亚 2:12} 这饮酒的禁令也给了大司祭亚郎和他的儿子们，\InnerQuote{当他们进入圣所时。}\BibleRef{肋 10:9} 命令\InnerQuote{用圣咏和赞歌向上主歌唱}\BibleRef{弗 5:19}，这适合出于那知道\InnerQuote{那些用鼓和瑟喝酒}\BibleRef{依 5:11}的人受天主责备的人。现在，当我发现这些规诫属于哪一位天主，无论在其萌芽还是发展，我就毫不费力地知道宗徒也属于谁。但他宣称\InnerQuote{妻子应当服从自己的丈夫}：他给出什么理由呢？他说：\InnerQuote{因为丈夫是妻子的头。}\BibleRef{弗 5:23} 请告诉我，\Person{Marcion}，你的神是在创造者的工作上建立他法律的权威吗？然而这还算小事；因为他实际上从同一源头引出了他的基督和教会的情况；因为他说：\InnerQuote{如同基督是教会的头；}\BibleRef{弗 5:23} 又同样说：\InnerQuote{谁爱自己的妻子，就是爱自己的身体，正如基督爱了教会一样。} 你看，你的基督和你的教会是如何与创造者的工作相比较的。以教会的名义，肉身获得了何等的尊荣！宗徒说：\InnerQuote{从来没有人恨恶自己的肉身}（当然，只有\Person{Marcion}例外）\InnerQuote{却保养顾惜它，如同主对教会一样。}\BibleRef{弗 5:29} 但你是唯一恨恶自己肉身的人，因为你剥夺了它的复活。你也应恨恶教会才是合理的，因为基督基于同样的原则爱了教会。是的，基督爱了肉身如同爱了教会。因为没有人会爱他妻子的画像而不加以照料、尊崇和加冕。相似者分享了真实的特权尊荣。现在我将尽力从我的角度证明，同一的天主是人和基督的天主、女人和教会的天主、肉身和神的天主，借着宗徒的帮助，他应用了创造者的诫命，甚至加上注解：\InnerQuote{因此，人要离开父亲和母亲，（要与妻子结合），二人成为一体。这奥秘是伟大的。}\BibleRef{弗 5:31} 顺便说，对我来说已经足够：创造者的工作在宗徒的估价中是伟大的奥秘，尽管它们被异端者视为如此卑贱。他说：\InnerQuote{但我是指基督和教会而言的。}\BibleRef{弗 5:32} 他说这话是为了解释奥秘，而不是为了破坏它。他向我们表明，这奥秘是由那也是奥秘的作者所预表的。现在\Person{Marcion}的意见是什么？创造者不可能为一位未知的神提供预象，或者，若已知，则为一位与祂自己敌对的神。事实上，至高神不应从较低者借用任何东西；他宁可消灭他。\InnerQuote{作子女的，要服从你们的父母。}\BibleRef{弗 6:1} 现在，虽然\Person{Marcion}删去了（下一句）\InnerQuote{这是第一条带有应许的诫命}，但法律仍明说：\InnerQuote{应孝敬你的父亲和你的母亲。}\BibleRef{出 20:12} 又说，（宗徒写道：）\InnerQuote{作父母的，你们要带着主的敬畏和劝勉养育儿女。}\BibleRef{弗 6:4} 因为你们已经听见古时如何对他们说：\InnerQuote{你们应将这事传于子孙，子孙再传于他们的子孙。}\BibleRef{出 10:2} 当纪律只有一个时，两位天主对我有何用？如果必须有两个，我打算跟随那首先教导这功课的一位。但既然我们的斗争是针对\InnerQuote{今世的统治者}\BibleRef{弗 6:12}，那么必然有一大群创造者之神！因为我为什么不应该在这里坚持这一点：如果他意指只有创造者是他之前提到的所有权势所归属的那一位，他就应该只提\Emph{一位}\InnerQuote{今世的统治者}。再者，当他在前一节命令我们\InnerQuote{穿上天主的全副武装，为能抵抗魔鬼的阴谋}\BibleRef{弗 6:11}时，他岂不表明他在魔鬼之名后提到的一切事物确实属于魔鬼——\Emph{率领者和掌权者，以及这黑暗世界的统治者}\BibleRef{弗 6:12}——我们也把这些归于魔鬼的权下吗？否则，如果\InnerQuote{魔鬼}意指创造者，那么在创造者的安排中谁是魔鬼？既然有两位神，也必须有两个魔鬼，以及许多的权势和今世的统治者吗？但是创造者怎能同时是魔鬼又是神呢？魔鬼并非同时既是神又是魔鬼。因为要么他们俩都是神，如果俩都是魔鬼；否则，那位是神的就不是魔鬼，正如那魔鬼也不是神。我确实想知道，是通过什么曲解，\Emph{魔鬼}一词竟然适用于创造者。也许他歪曲了至高神的某个目的——就像祂自己从天使长那里所经历的，那天使长确实为此而撒谎。因为祂禁止（我们的原祖）尝那不幸之树的果子，并非担心他们会成为神；祂的禁令是为了防止他们在违命后死亡。但\InnerQuote{属灵的邪恶}\BibleRef{弗 6:12}并非意指创造者，因为宗徒额外描述说\InnerQuote{在属天的地方}；因为宗徒很清楚，当天使被人的女儿诱入罪恶时，属灵的邪恶曾在属天的地方活动。但是，当宗徒在\InnerQuote{传扬福音的奥秘，我为这奥秘作了带锁链的使者}时，因着宣讲的自由，向教会显示了如此坚定和坦白的言语——实际上还请求（厄弗所人）向天主祈祷，使这\InnerQuote{放胆开口}能继续给他——的时候，他怎么会为了贬低创造者而求助于含混的描述和不知什么晦涩的谜语呢？\BibleRef{弗 6:19}\par

% structure: p0044 source=h2 target=subsection reason=relative-heading-rank; base=h2

\subsection{第十九章　\BibleBook{哥罗森}{书}．时间是真理与异端的准则．准则的应用．不可见天主的肖像释义．我们的基督在造物主古代安排中的先存．基督的圆满所含的一切．\Person{马西昂}的天主的伊壁鸠鲁特性．与之对立的天主教真理．法律之于基督，如同影子之于实体。}

我习惯于在我的针对一切异端的辩驳中，将\Emph{时间}的见证作为我简明的（真理）准则；主张其中的\Emph{优先性}是我们的规则，并断言\Emph{后来}是一切异端的特征。如今这一点也要由宗徒来证明，当他这样说：\BibleQuote{因为那为你们存留在天上的希望，你们以前在福音真理之言中听过了；这福音也传到你们那里，如同传到普世一样。}\BibleRef{哥 1:5}如今，如果遍传各处的是我们的福音，而非任何异端福音，更不是马西昂的——后者仅始于安多尼在位时期——那么我们的福音便是宗徒们的福音。但即使马西昂的福音成功充满全世界，它也无权享有宗徒性的特征。因为显然，这一品质只属于那首先充满世界的福音；换言之，属于那位自古便如此预言其传扬的天主的福音：\BibleQuote{他们的声音传遍全地，他们的话语传到地极。}他称基督为\BibleQuote{不可见的天主的肖像。}\BibleRef{哥 1:15}我们同样说基督的父是不可见的，因为我们知道，在古代（每当有人以天主的名义显现于人间时）是圣子作为（圣父）本身的肖像被看见。然而，不可认为他是在区别一个可见的天主和一个不可见的天主；因为早在他写这些话之前，我们就发现对我们的天主有如此描述：\BibleQuote{没有人能见主而存活。}\BibleRef{出 33:20}如果基督不是\BibleQuote{一切受造物中的首生者}，正如那\BibleQuote{天主圣言，万物藉祂而造，没有祂什么也不能造；}\BibleRef{若 1:3}如果\BibleQuote{万物}不是\BibleQuote{在祂内受造，无论是在天上或地上，可见的或不可见的，无论是有位的、主治的、宰制的、掌权的；}如果\BibleQuote{万物}不是\BibleQuote{藉祂而造并为祂而造}（因为这些真理马西昂不应承认关乎祂），那么宗徒就不可能如此肯定地断言：\BibleQuote{祂在万有之先。}因为祂怎能是在万有之先，如果祂不是在\Emph{万有}之先？祂又怎能是在万有之先，如果祂不是\BibleQuote{一切受造物中的首生者}——如果祂不是创造主的圣言？那在万有之后显现的，又怎能被证明是在万有之先？谁又能说祂先前存在，若找不到任何祂存在过的证据？此外，又如何能\BibleQuote{（圣父）乐意使一切的圆满居在祂内？}\BibleRef{哥 1:19}因为，首先，那不被马西昂从其中删除的成分所构成的圆满是什么——甚至那些\BibleQuote{在基督内受造的，无论是在天上或地上}，无论是天使还是人？那不由可见和不可见之物所成的圆满是什么？那不由有位、主治、宰制、掌权的所成的圆满是什么？反之，如果我们的假宗徒和犹太化福音传播者从他们自己的库存中引入了这一切，而马西昂将它们应用于构成他自己天主的圆满（这个假设尽管荒谬，却单独能为他辩护）；因为，在任何其他假设下，创造主的对手和毁灭者怎么会愿意让祂的圆满住在自己的基督内？再者，祂\BibleQuote{藉祂自己使万有与自己和好，藉祂十字架的血成就了和平，}\BibleRef{哥 1:20}除了与那曾被这一切所触怒、他们曾因过犯而悖逆、但最终又回归的那位之外，还能与谁呢？他们可能\Emph{被安抚}于一个陌生的神，但他们不可能\Emph{与任何别的神和好}，除了他们自己的神。因此，我们这些\BibleQuote{从前因恶行而在心意上与祂疏远、为敌的人}\BibleRef{哥 1:21}祂使我们与创造主和好——我们曾冒犯祂，因崇奉受造物而有损于创造主。然而，正如祂在别处所说，\BibleRef{弗 1:23}教会是基督的身体，同样在这里（宗徒）宣称他\BibleQuote{在自己肉身上，为基督的身体，即教会，补充基督苦难所欠缺的。}\BibleRef{哥 1:24}但你不可因此认为，每当提到祂的身体，这个词都只是一个隐喻，而非指真实的肉体。因为他在上面说了我们\BibleQuote{藉祂的身体受死而得以和好；}\BibleRef{哥 1:22}当然，意思是祂在那因肉体而可能性死亡的身体中死了：（因此他补充说，）不是\Emph{借着教会}（\OriginalTerm{借着教会}{per ecclesiam}），而是明确地\Emph{为了教会}（\OriginalTerm{为了教会}{proper ecclesiam}），以身体换取身体——以一个血肉之躯换取一个属灵的身体。再者，当他警告他们要\BibleQuote{谨慎诡辩和哲学}，认为它是\BibleQuote{虚妄的欺骗}，是\BibleQuote{照世间的蒙学}（不是指天地间俗世的构造，而是指世俗的学识，和\BibleQuote{人的传统}，在言辞和哲学上诡谲），\BibleRef{哥 2:8}若要指出在这句话（宗徒的话）中，一切异端如何因其所依靠的诡辩资源与哲学原理而被定罪，那将是冗长的，且是另一部著作的合宜主题。但（一次性地）让马西昂知道，他信条的首要术语出自伊壁鸠鲁学派，意谓主是愚钝和漠不关心的；因此他拒绝说祂是可畏的对象。此外，他从斯多葛学派的门廊引出了\Emph{物质}，并将其置于与神圣创造主同等的地位。他还否认肉身的复活——这是没有一个哲学学派共同持守的真理。但我们的（公教）真理与这异端者的诡计相距何其遥远——它畏惧激起天主的义怒，坚信祂从无中创造了万有，应许我们那（已死的）同一肉身的坟墓复活，并毫不羞愧地坚持基督是由童贞女的子宫所生！对此，哲学家、异端者，甚至外教人都嘲笑讥讽。因为\BibleQuote{天主拣选了世上的愚拙，好叫智慧人羞愧}\BibleRef{格前 1:27}——无疑就是那位天主，关于祂自己的这个安排，祂早已威胁要\BibleQuote{毁灭智慧人的智慧。}感谢这与哲学的诡辩和虚妄欺骗如此对立的真理的纯朴，我们绝不可能对这类邪僻观点有任何喜好。那么，如果天主\BibleQuote{使我们与基督一同复活，赦免了我们的过犯，}\BibleRef{哥 2:13}我们便不能设想，罪是由一位长期以来一直未被认识、因而罪不可能犯下的祂赦免的。现在告诉我，马西昂，你对宗徒这话有何看法，他说：\BibleQuote{不要让任何人在饮食上，或关于节期、月朔或安息日的事上审判你们，这些东西是将来事物的阴影，但实体是属于基督的。}\BibleRef{哥 2:16}我们现在不是要讨论律法，除了（指出）宗徒在这里清楚地教导律法如何被废除，即从阴影到实体——也就是从预表到现实，即基督。因此，阴影属于谁，实体也属于谁；换言之，律法属于祂，基督也属于祂。如果你将律法和基督分开，将这一个归给一位神，那另一个归给另一位，那就如同试图将阴影与它所投射的身体分开。显然，基督与律法有关联，正如实体与它的阴影有关联一样。但是当他责备那些以天使显现为权威而说人必须禁戒某些食物的人——\BibleQuote{不可摸，不可尝}——出于一种自贬的谦卑，同时\BibleQuote{随着肉体的心思徒然自大，不紧握元首}，（宗徒）这些话并非攻击律法或梅瑟，仿佛他是在迷信天使的暗示下颁布了各种饮食的禁令。因为梅瑟显然是从天主接受了律法。因此，当他提到他们\BibleQuote{随从人的命令和教理}\BibleRef{哥 2:22}时，他指的是那些\BibleQuote{不紧握元首}之人的行为，这元首就是万有聚集于祂的那一位；因为万有都回忆到基督，并集中围绕在祂作为它们的源始原则——甚至包括那些本质上无涉的饮食。他其余的一切规诫，正如我们在处理另一封信中出现的这些规诫时已充分指出的，都出自创造主；祂预言\BibleQuote{旧事要过去}，并将\BibleQuote{更新一切}，命令人\BibleQuote{为自己开垦荒地}，并借此教导他们甚至在那时就要脱去旧人，穿上新人。\par

% structure: p0046 source=h2 target=subsection reason=relative-heading-rank; base=h2

\subsection{第二十章\WorkTitle{斐理伯书}。传扬基督者之间的分歧，并不证明基督不止一位。圣\Person{保禄}的措辞——仆人的形体、人的模样与形态——并不支持\OriginalTerm{幻象论}{Docetism}。从本书所提的某些对比中，不能推导出\OriginalTerm{对立说}{Antithesis}（如\Person{马尔西翁}所声称的）存在于犹太教的天主与福音的天主之间。与\WorkTitle{创世纪}一处经文的平行对照。身体的复活及其改变。}

当（宗徒）提到那些宣讲福音者的种种动机，有些人\BibleQuote{因他的锁链而更有胆量，更勇于宣讲圣言}，另一些人\BibleQuote{有的出于嫉妒和纷争宣讲基督，有的出于善意}，许多人也是\BibleQuote{出于爱德}，某些人\BibleQuote{出于争竞}，还有些人\BibleQuote{出于敌视他}。\BibleRef{斐 1:14}他无疑有良机指责他们所宣讲的教义分歧，似乎正是这分歧导致了他们态度上的巨大不同。但他将这些态度揭示为分歧的唯一原因，避免责备信仰的正统奥秘，并断言尽管如此，只有一个基督和祂的一位天主，无论人们宣讲祂的动机如何。因此他说，对我来说无关紧要\BibleQuote{无论基督是被假意还是真诚宣讲}。\BibleRef{斐 1:18}因为只有一位基督被宣布，无论他们是在\Quote{假意}或\Quote{真诚}的信德中。因为他将\Emph{真理}一词用于他们宣讲的信实，而非规则本身的正确性，因为确实只有一个规则；而宣讲者的行为各不相同：有些人是真实的，即纯一的，另一些则因过多的学识而变得矫揉造作。既然如此，显然他们所宣讲的基督正是先知们一贯的主题。现在，如果宗徒引入的是一个完全不同的基督，那么新事物就会产生（信仰上的）分歧。因为尽管有新颖的教导，仍不会缺少一些人将所宣讲的福音解释为造物主的基督，因为现今各处的大多数人都与我们观点一致，而非站在异端一边。所以宗徒不会在这样的段落中不去评论和谴责分歧。然而，既然没有对分歧的责备，就没有新颖性的证据。当然，马西昂派认为他们在以下段落中得到了宗徒关于基督实体的支持—祂内只有肉身的幻象。因为关于基督，他说：\BibleQuote{祂虽具有天主的形体，却没有以自己与天主同等为抢夺的，反而空虚自己，取了奴仆的\Emph{形体}}，\Emph{不是实体}，\BibleQuote{成了人的\Emph{相似}}，\Emph{不是人}，\BibleQuote{以人的\Emph{样式}被找到}，\BibleRef{斐 2:6}\Emph{不是祂的实体}，即祂的肉身；仿佛实体不会同时具有\Emph{形体}和\Emph{相似}和\Emph{样式}。幸好我们在另一段落中（宗徒）称基督为\BibleQuote{不可见的天主的形象}。\BibleRef{哥 1:15}难道不是同样有力地推出来：基督不是真实的天主，因为宗徒将祂置于天主的\Emph{形象}中，如果（如马西昂所坚持）祂不是真实的人，因为祂取了人的\Emph{形体}或\Emph{形象}？因为在这两种情况下，真实的实体都被排除，如果\Emph{形象}（或\Quote{样式}）和\Emph{相似}和\Emph{形体}被归于幻象。但是，由于祂作为圣父之子，在祂的样式和形象中真是天主，已经通过这一结论的力量被确定为真实的人，作为人子，\BibleQuote{以人的样式被找到}和形象\BibleQuote{的一个人}。因为当他这样提出祂在人的样式\BibleQuote{\Emph{被找到}}时，他实际上\Emph{事实上}肯定了祂确实是人。因为凡\Emph{被找到}的，显然具有存在。因此，正如祂因大能而被找到是天主，同样地，祂因肉身而被找到是人，因为宗徒不可能宣告祂\BibleQuote{成了服从以至于死}，\BibleRef{斐 2:8}如果祂不是由可朽的实体构成。更加明显的是从宗徒附加的话：\BibleQuote{甚至死在十字架上}。\BibleRef{斐 2:8}因为他很难将此视为人类痛苦的高潮，以赞扬祂服从的美德，如果他知晓这一切都是幻象的虚幻过程，它更逃避十字架而非经历十字架，并且在痛苦中未显示任何美德，只有幻象。但是，\BibleQuote{他先前以为有利的事}，他在前几节中列出的—\BibleQuote{依恃肉身}，割损的记号，\BibleQuote{希伯来人中的希伯来人}，源自\BibleQuote{本雅明支派}，法利塞人的尊严荣誉—他现在认为只是\BibleQuote{损失}；\BibleRef{斐 3:7}换句话说，他所弃绝的不是犹太人的天主，而是他们愚蠢的顽劣。这些也是\BibleQuote{他以视为粪土，为获得基督的卓越知识}\BibleRef{斐 3:8}（但绝不是为了拒绝造物主天主）；\BibleQuote{他不要那出自法律的自己的义德，而要那借着祂，}即基督，\BibleQuote{来自天主的义德}。\BibleRef{斐 3:9}那么，你说，根据这种区分，法律并非出自基督的天主。真够微妙！但还有更微妙之处。因为当（宗徒）说：\BibleQuote{不是出自法律的义德，而是借着祂的义德}，他不会对法律所属的那位以外的任何一位使用\Emph{借着祂}这个短语。\BibleQuote{我们的家乡}，他说，\BibleQuote{是在天上}。\BibleRef{斐 3:20}我在此认出造物主对亚巴郎的古老应许：\BibleQuote{我要使你的后裔如天上的繁星一样多}。\BibleRef{创 22:17}因此\BibleQuote{星与星的光辉各有不同}。\BibleRef{格前 15:41}再者，如果基督从天降临时\BibleQuote{将改变我们卑贱的身体，使之相似祂荣耀的身体}，那么我们的这个身体必要复活，它现在处于痛苦中的卑贱状态，并按可朽定律落入地里。但如果它没有真实存在，怎能被改变？然而，如果这只是对那些在造物主的来临时将被发现仍在肉身中，且必须被改变的人说的，\BibleRef{格前 15:51}那就先复活的人怎么办？他们没有可被改变的实体。但他说（别处）：\BibleQuote{我们将和他们在云上一同被提到空中，迎接主}。\BibleRef{得前 4:16}那么，如果我们只和他们一起被提，我们当然也一同被改变。\par

% structure: p0048 source=h2 target=subsection reason=relative-heading-rank; base=h2

\subsection{第21章 \WorkTitle{致费肋孟书}。此书未被割裂。\Person{马尔基昂}接受此书却拒绝另外三封致个人的书信，前后矛盾。结论。\Person{戴尔都良}为其反对\Person{马尔基昂}之作的对称结构与有意安排申辩。}

此书信因其简短而得以免受\Person{马尔基昂}的篡改之手。然而，令我惊讶的是，他接受了这封只写给一个人的书信（收入其\WorkTitle{宗徒集}），却拒绝了那两封\WorkTitle{弟茂德书}和那一封\WorkTitle{弟铎书}——这些书信全都论及教会的纪律。我想，他的目的是要将其窜改活动甚至延伸到（\Person{圣保禄}）书信的数目上。读者啊，现在请你记住，我们在此已从宗徒的著作中援引了证据，以支持我们先前必须处理的主题，并且我们现在已结束我们推迟到此（部分）工作的讨论。（我请求你赐予这一恩惠，）使你不认为在\Emph{此处}的任何重复都是多余的，因为我们只是履行了先前对你的承诺；也不要对\Emph{彼处}的任何延迟心存疑虑，我们只是陈述了（论证的）要点。若你细心审视全部作品，你会凭着你公正的判断，免除我们在此处过于冗繁或彼处信心不足之责。\par

\SourceRecord{Translated by Peter Holmes.；Ante-Nicene Fathers；Vol. 3.；Edited by Alexander Roberts, James Donaldson, and A. Cleveland Coxe.；Buffalo, NY: Christian Literature Publishing Co.,；1885.；<http://www.newadvent.org/fathers/03125.htm>.}
